% Generated mechanically from the frozen sites.tex target.
% Do not edit this generated file; edit the bound locale source instead.

% OK, start here.
%

\setcounter{chapter}{6}
\chapter{位点与层}



\phantomsection
\label{sites-section-phantom}


\section{引言}
\label{sites-section-introduction}

\noindent
Grothendieck 引入位点这一概念，是为了研究概形的 \'etale 拓扑中的层。
这一概念的基本参考文献或许是 \cite{SGA4}。我们采用的位点概念与
\cite{SGA4} 中的不同；我们所谓的位点，在 \cite[Expos\'e II, D\'efinition 1.3]{SGA4}
中称为赋有预拓扑的范畴。之所以这样处理，是因为
在代数几何中，使用一个给定的覆盖类往往很方便，例如定义概形的某个性质何时
在给定拓扑中是局部性质时即是如此；见 Descent，第
\ref{descent-section-descending-properties} 节。本文的论述将紧随
\cite{ArtinTopologies}。我们不使用宇宙。












\section{预层}
\label{sites-section-presheaves}

\noindent
设 $\mathcal{C}$ 为范畴。一个{\it 集合预层}是从 $\mathcal{C}$ 到
$\textit{Sets}$ 的反变函子 $\mathcal{F}$（见 Categories，备注
\ref{categories-remark-functor-into-sets}）。因此，对 $\mathcal{C}$ 的每个
对象 $U$，都有一个集合 $\mathcal{F}(U)$。这个集合的元素称为
$\mathcal{F}$ 在 $U$ 上的{\it 截面}。对每个态射 $f : V \to U$，映射
$\mathcal{F}(f) : \mathcal{F}(U) \to \mathcal{F}(V)$ 称为{\it 限制映射}，
并且常记作 $f^\ast : \mathcal{F}(U) \to \mathcal{F}(V)$。也可以说，
$f^*(s)$ 是 $s$ 沿 $f$ 的{\it 拉回}。函子性意味着
$g^* f^* (s) = (f \circ g)^*(s)$。有时我们采用记号
$s|_V := f^\ast(s)$。这一记号与拓扑学中的函数限制概念相容，因为若
$W \to V \to U$ 是 $\mathcal{C}$ 中的态射，而 $s$ 是 $\mathcal{F}$ 在
$U$ 上的截面，则由 $\mathcal{F}$ 的函子性可知
$s|_W = (s|_V)|_W$。当然，我们必须谨慎，因为态射 $V \to U$ 很可能不止一个，
而相应的拉回映射肯定未必相等。

\begin{definition}
\label{sites-definition-presheaves-sets}
$\mathcal{C}$ 上的一个{\it 集合预层}，是从 $\mathcal{C}$ 到
$\textit{Sets}$ 的反变函子。{\it 预层态射}是函子之间的变换。集合预层的
范畴记作 $\textit{PSh}(\mathcal{C})$。
\end{definition}

\noindent
注意，对 $\mathcal{C}$ 的任意对象 $U$，点函子 $h_U$（见 Categories，例
\ref{categories-example-hom-functor}）都是一个预层。这些预层称为
{\it 可表预层}。它们具有如下良好性质：对任意预层 $\mathcal{F}$，都有
\begin{equation}
\label{sites-equation-map-representable-into-presheaf}
\Mor_{\textit{PSh}(\mathcal{C})}(h_U, \mathcal{F})
=
\mathcal{F}(U).
\end{equation}
这就是 Yoneda 引理（Categories，引理 \ref{categories-lemma-yoneda}）。

\medskip\noindent
类似地，我们可以定义阿贝尔群预层、环预层等概念。更一般地，还可以定义取值于
某个范畴的预层。

\begin{definition}
\label{sites-definition-presheaf}
设 $\mathcal{C}$、$\mathcal{A}$ 为范畴。$\mathcal{C}$ 上一个取值于
$\mathcal{A}$ 的{\it 预层} $\mathcal{F}$，是从 $\mathcal{C}$ 到
$\mathcal{A}$ 的反变函子，即
$\mathcal{F} : \mathcal{C}^{opp} \to \mathcal{A}$。
$\mathcal{C}$ 上取值于 $\mathcal{A}$ 的预层 $\mathcal{F} \to \mathcal{G}$
的一个{\it 态射}，是从函子 $\mathcal{F}$ 到 $\mathcal{G}$ 的一个变换。
\end{definition}

\noindent
这些分别构成 $\mathcal{C}$ 上取值于 $\mathcal{A}$ 的预层范畴的对象和态射。

\begin{remark}
\label{sites-remark-big-presheaves}
如前所述，我们可以考虑取值于 Categories，备注
\ref{categories-remark-big-categories} 所列任意一个“大全畴”的预层范畴。
这些同样会是“大全畴”；随着论述推进，它们也会列入上述备注中。
\end{remark}






















\section{预层的单射与满射}
\label{sites-section-injective-surjective}

\begin{definition}
\label{sites-definition-presheaves-injective-surjective}
设 $\mathcal{C}$ 为范畴，并设 $\varphi : \mathcal{F}
\to \mathcal{G}$ 是集合预层的映射。
\begin{enumerate}
\item 若对 $\mathcal{C}$ 的每个对象 $U$，映射
$\varphi_U : \mathcal{F}(U) \to \mathcal{G}(U)$ 都是单射，则称
$\varphi$ 是{\it 单射的}。
\item 若对 $\mathcal{C}$ 的每个对象 $U$，映射
$\varphi_U : \mathcal{F}(U) \to \mathcal{G}(U)$ 都是满射，则称
$\varphi$ 是{\it 满射的}。
\end{enumerate}
\end{definition}

\begin{lemma}
\label{sites-lemma-mono-epi}
上面定义的单射（相应地，满射）恰好是
$\textit{PSh}(\mathcal{C})$ 中的单态射（相应地，满态射）。一个映射是同构，
当且仅当它既是单射又是满射。
\end{lemma}

\begin{proof}
我们先证明 $\varphi : \mathcal{F} \to \mathcal{G}$ 是单射，当且仅当它是
$\textit{PSh}(\mathcal{C})$ 中的单态射。“仅当”方向是直接的，故只需证明
“当”方向。假设 $\varphi$ 是单态射。取
$U \in \Ob(\mathcal{C})$；我们需要证明 $\varphi_U$ 是单射。于是取满足
$\varphi_U (a) = \varphi_U (b)$ 的
$a, b \in \mathcal{F}(U)$；需要验证 $a = b$。在同构
(\ref{sites-equation-map-representable-into-presheaf}) 下，
$\mathcal{F}(U)$ 的元素 $a$ 和 $b$ 对应于两个自然变换
$a', b' \in \Mor_{\textit{PSh}(\mathcal{C})}(h_U, \mathcal{F})$。
类似地，在相应的同构
$\Mor_{\textit{PSh}(\mathcal{C})}(h_U, \mathcal{G})
= \mathcal{G}(U)$ 下，$\mathcal{G}(U)$ 中两个相等的元素
$\varphi_U (a)$ 和 $\varphi_U (b)$ 对应于两个自然变换
$\varphi \circ a', \varphi \circ b'
\in \Mor_{\textit{PSh}(\mathcal{C})}(h_U, \mathcal{G})$，
因此这两个自然变换也必相等。故
$\varphi \circ a' = \varphi \circ b'$；由于 $\varphi$ 是单态射，便有
$a' = b'$，进而 $a = b$。这就证明了 (1)。

\medskip\noindent
下面证明 $\varphi : \mathcal{F} \to \mathcal{G}$ 是满射，当且仅当它是
$\textit{PSh}(\mathcal{C})$ 中的满态射。“仅当”方向是直接的，故只需证明
“当”方向。假设 $\varphi$ 是满态射。

\medskip\noindent
对 $\textit{Sets}$ 范畴中的任意两个态射 $f : A \to B$ 和
$g : A \to C$，用 $\text{inl}_{f,g}$ 和 $\text{inr}_{f,g}$ 表示从
$B$ 和 $C$ 到 $B \coprod_A C$ 的两个典范映射。（这里，推出在
$\textit{Sets}$ 中计算。）

\medskip\noindent
现在定义 $\mathcal{C}$ 上的集合预层 $\mathcal{H}$：对每个
$U \in \Ob(\mathcal{C})$，令
$\mathcal{H}(U) = \mathcal{G}(U) \coprod_{\mathcal{F}(U)} \mathcal{G}(U)$
（其中推出在 $\textit{Sets}$ 中计算，并由映射
$\varphi_U : \mathcal{F}(U) \to \mathcal{G}(U)$ 诱导）；它在态射上的作用
以显然的方式定义（利用推出的函子性）。于是存在两个自然变换
$i_1 : \mathcal{G} \to \mathcal{H}$ 和
$i_2 : \mathcal{G} \to \mathcal{H}$，它们在对象
$U \in \Ob(\mathcal{C})$ 处的分量分别由映射
$\text{inl}_{\varphi_U, \varphi_U}$ 和
$\text{inr}_{\varphi_U, \varphi_U}$ 给出。推出的定义表明
$i_1 \circ \varphi = i_2 \circ \varphi$，因而 $i_1 = i_2$（因为
$\varphi$ 是满态射）。所以对每个 $U \in \Ob(\mathcal{C})$，都有
$\text{inl}_{\varphi_U, \varphi_U}
= \text{inr}_{\varphi_U, \varphi_U}$。于是 $\varphi_U$ 必为满射（一个简单的
组合论论证表明，若 $f : A \to B$ 是 $\textit{Sets}$ 中的态射，则
$\text{inl}_{f,f} = \text{inr}_{f,f}$ 当且仅当 $f$ 是满射）。换言之，
$\varphi$ 是满射，(2) 得证。

\medskip\noindent
最后证明 $\varphi : \mathcal{F} \to \mathcal{G}$ 既是单射又是满射，
当且仅当它是 $\textit{PSh}(\mathcal{C})$ 中的同构。这一次，“当”方向是
直接的。为证明“仅当”方向，只须注意：若 $\varphi$ 既单又满，则对每个
$U \in \Ob(\mathcal{C})$，映射 $\varphi_U$ 都可逆；对所有 $U$，这些映射的逆
可组合成自然变换 $\mathcal{G} \to \mathcal{F}$，它正是 $\varphi$ 的逆。
\end{proof}
\begin{definition}
\label{sites-definition-sub-presheaf}
若对每个对象 $U \in \Ob(\mathcal{C})$，集合 $\mathcal{F}(U)$ 都是
$\mathcal{G}(U)$ 的子集，并且这一关系与限制映射相容，则称 $\mathcal{F}$
是 $\mathcal{G}$ 的一个{\it 子预层}。
\end{definition}

\noindent
换言之，包含映射 $\mathcal{F}(U) \to \mathcal{G}(U)$ 可以粘合成一个
（单射的）预层态射 $\mathcal{F} \to \mathcal{G}$。

\begin{lemma}
\label{sites-lemma-image}
设 $\mathcal{C}$ 为范畴。假设
$\varphi : \mathcal{F} \to \mathcal{G}$ 是 $\mathcal{C}$ 上集合预层的态射。
则存在唯一的子预层 $\mathcal{G}' \subset \mathcal{G}$，使得 $\varphi$ 分解为
$\mathcal{F} \to \mathcal{G}' \to \mathcal{G}$，并且第一个映射是满射。
\end{lemma}

\begin{proof}
为证明存在性，只须对每个 $U \in \Ob(C)$ 令
% P01-ERRATA: P01-E0057 -- the frozen source omits \mathcal around C.
$\mathcal{G}'(U) = \varphi_U \left(\mathcal{F}(U)\right)$
（并从 $\mathcal{G}$ 继承其在态射上的作用），再证明这定义了
$\mathcal{G}$ 的一个子预层，且 $\varphi$ 分解为
$\mathcal{F} \to \mathcal{G}' \to \mathcal{G}$，其中第一个映射是满射。
唯一性是直接的。
\end{proof}

\begin{definition}
\label{sites-definition-image}
采用引理 \ref{sites-lemma-image} 的记号。我们称 $\mathcal{G}'$ 为
{\it $\varphi$ 的像}。
\end{definition}















\section{预层的极限与余极限}
\label{sites-section-limits-colimits-PSh}

\noindent
设 $\mathcal{C}$ 为范畴。范畴 $\textit{PSh}(\mathcal{C})$ 中存在极限和
余极限。此外，对任意 $U \in \Ob(\mathcal{C})$，函子
$$
\textit{PSh}(\mathcal{C})
\longrightarrow
\textit{Sets}, \quad
\mathcal{F}
\longmapsto
\mathcal{F}(U)
$$
与极限和余极限交换。或许证明这些陈述的最简便方法如下。给定一个图表
$
\mathcal{F} :
\mathcal{I}
\to
\textit{PSh}(\mathcal{C})
$，定义预层
$$
\mathcal{F}_{\lim} :
U
\longmapsto
\lim_{i \in \mathcal{I}} \mathcal{F}_i(U)
\text{  且  }
\mathcal{F}_{\colim} :
U
\longmapsto
\colim_{i \in \mathcal{I}} \mathcal{F}_i(U)
$$
显然存在投影映射 $\mathcal{F}_{\lim} \to \mathcal{F}_i$ 和典范映射
$\mathcal{F}_i \to \mathcal{F}_{\colim}$。这些映射满足 Categories，定义
\ref{categories-definition-limit}（相应地，Categories，定义
\ref{categories-definition-colimit}）中极限（相应地，余极限）映射的要求。
事实上，它们显然在 $\mathcal{F}$ 上构成一个锥，相应地构成一个余锥。
此外，若 $(\mathcal{G}, q_i : \mathcal{G} \to \mathcal{F}_i)$ 是另一组
这样的数据（如极限定义中那样），则对每个 $U$，我们得到一族满足适当函子性
条件的映射 $\mathcal{G}(U) \to \mathcal{F}_i(U)$，因而得到唯一映射
$\mathcal{G}(U) \to \mathcal{F}_{\lim}(U)$。容易验证，随着 $U$ 变化，
这些映射彼此相容，并来自所需的映射
$\mathcal{G} \to \mathcal{F}_{\lim}$。余极限情形的论证类似。





















\section{预层范畴的函子性}
\label{sites-section-functoriality-PSh}

\noindent
设 $u : \mathcal{C} \to \mathcal{D}$ 为范畴之间的函子。在这种情况下，记
$$
u^p :
\textit{PSh}(\mathcal{D})
\longrightarrow
\textit{PSh}(\mathcal{C})
$$
为这样的函子：它把 $\mathcal{D}$ 上的预层 $\mathcal{G}$ 送到预层
$u^p\mathcal{G} = \mathcal{G} \circ u$。注意，由上一节可知，这个函子与所有
极限和余极限交换。

\medskip\noindent
对 $V \in \Ob(\mathcal{D})$，用 $\mathcal{I}^u_V$ 表示如下范畴：
\begin{equation}
\label{sites-equation-colim-category}
\begin{matrix}
\Ob(\mathcal{I}^u_V)
&
=
&
\{
(U, \phi)
\mid
U \in \Ob(\mathcal{C}),
\phi : V \to u(U)
\}
\\
\Mor_{\mathcal{I}^u_V}((U, \phi), (U', \phi'))
&
=
&
\{
f : U \to U' \text{ 于 }\mathcal{C}
\mid
u(f) \circ \phi = \phi'
\}
\end{matrix}
\end{equation}
有时我们省略记号中的上标 ${}^u$，只写 $\mathcal{I}_V$。我们将利用这些范畴
定义函子 $u^p$ 的左伴随。在此之前，先证明几个技术性引理。

\begin{lemma}
\label{sites-lemma-almost-directed}
设 $u : \mathcal{C} \to \mathcal{D}$ 为范畴之间的函子。假设
$\mathcal{C}$ 有纤维积和等化子，且 $u$ 与它们交换。那么范畴
$(\mathcal{I}_V)^{opp}$ 满足 Categories，引理
\ref{categories-lemma-split-into-directed} 的假设。
\end{lemma}

\begin{proof}
需要验证两个条件。

\medskip\noindent
首先，假设给定三个对象
$\phi : V \to u(U)$、$\phi' : V \to u(U')$ 和
$\phi'' : V \to u(U'')$，以及态射 $a : U' \to U$、
$b : U'' \to U$，使得
$u(a) \circ \phi' = \phi$ 且 $u(b) \circ \phi'' = \phi$。
我们需要证明存在另一个对象 $\phi''' : V \to u(U''')$ 以及态射
$c : U''' \to U'$ 和 $d : U''' \to U''$，使得
$u(c) \circ \phi''' = \phi'$、$u(d) \circ \phi''' = \phi''$ 且
$a \circ c = b \circ d$。取 $U''' = U' \times_U U''$，并以 $c$ 和 $d$
为投影态射。由于 $u$ 与纤维积交换，这一构造有效；验证从略。

\medskip\noindent
其次，假设给定两个对象 $\phi : V \to u(U)$ 和
$\phi' : V \to u(U')$，以及态射
$a, b : (U, \phi) \to (U', \phi')$。我们需要找到一个等化 $a$ 和 $b$
的态射 $c : (U'', \phi'') \to (U, \phi)$。令
$c : U'' \to U$ 为 $\mathcal{C}$ 中 $a$ 和 $b$ 的等化子。由于 $u$ 与
等化子交换，并且
$u(a) \circ \phi = u(b) \circ \phi = \phi'$，我们得到一个态射
$\phi'' : V \to u(U'')$。
\end{proof}

\begin{lemma}
\label{sites-lemma-directed}
设 $u : \mathcal{C} \to \mathcal{D}$ 为范畴之间的函子。假设
\begin{enumerate}
\item 范畴 $\mathcal{C}$ 有终对象 $X$，且 $u(X)$ 是 $\mathcal{D}$ 的终对象；
\item 范畴 $\mathcal{C}$ 有纤维积，且 $u$ 与纤维积交换。
\end{enumerate}
那么指标范畴 $(\mathcal{I}^u_V)^{opp}$ 是滤过的（见 Categories，定义
\ref{categories-definition-directed}）。
\end{lemma}

\begin{proof}
这些假设蕴含引理 \ref{sites-lemma-almost-directed} 的假设成立（见 Categories，
第 \ref{categories-section-finite-limits} 节中的讨论）。由 Categories，引理
\ref{categories-lemma-split-into-directed}，可知 $\mathcal{I}_V$ 是若干定向范畴
的不交并（可能为空）。因此只需证明 $\mathcal{I}_V$ 连通。

\medskip\noindent
首先证明 $\mathcal{I}_V$ 非空。事实上，令 $X$ 为由假设存在的
$\mathcal{C}$ 的终对象。又因假设 $u(X)$ 是 $\mathcal{D}$ 的终对象，便有
态射 $V \to u(X)$。它给出 $\mathcal{I}_V$ 的一个对象。

\medskip\noindent
其次证明 $\mathcal{I}_V$ 连通。设
$\phi_1 : V \to u(U_1)$ 和 $\phi_2 : V \to u(U_2)$ 属于
$\Ob(\mathcal{I}_V)$。由假设，$U_1\times U_2$ 存在，且
$u(U_1\times U_2) = u(U_1)\times u(U_2)$。考虑由积的泛性质中
$(\phi_1, \phi_2)$ 所对应的态射
$\phi : V \to u(U_1\times U_2)$。显然，对象
$\phi : V \to u(U_1\times U_2)$ 同时映到
$\phi_1 : V \to u(U_1)$ 和 $\phi_2 : V \to u(U_2)$。
\end{proof}

\noindent
给定 $\mathcal{D}$ 中的 $g : V' \to V$，在对象上令
$\overline{g}(U, \phi) = (U, \phi \circ g)$，便得到函子
$\overline{g} : \mathcal{I}_V \to \mathcal{I}_{V'}$。给定
$\mathcal{C}$ 上的预层 $\mathcal{F}$，得到函子
$$
\mathcal{F}_V :
\mathcal{I}_V^{opp}
\longrightarrow
\textit{Sets}, \quad
(U, \phi)
\longmapsto
\mathcal{F}(U).
$$
换言之，$\mathcal{F}_V$ 是 $\mathcal{I}_V$ 上的集合预层。注意，
$\mathcal{F}_{V'} \circ \overline{g} = \mathcal{F}_V$。定义
$$
u_p\mathcal{F}(V) =
\colim_{\mathcal{I}_V^{opp}} \mathcal{F}_V
$$
由余极限的性质，对每个 $(U, \phi) \in \Ob(\mathcal{I}_V)$，得到典范映射
$\mathcal{F}(U)\xrightarrow{c(\phi)}u_p\mathcal{F}(V)$。对上述
$g : V' \to V$，由 Categories，引理
\ref{categories-lemma-functorial-colimit}，存在一个与
$\mathcal{F}_{V'} \circ \overline{g} = \mathcal{F}_V$ 相容的典范限制映射
$g^* : u_p\mathcal{F}(V) \to u_p\mathcal{F}(V')$。它是使得对所有
$(U, \phi) \in \Ob(\mathcal{I}_V)$，图表
$$
\xymatrix{
\mathcal{F}(U) \ar[r]^{c(\phi)} \ar[d]_{\text{id}}
&
u_p\mathcal{F}(V) \ar[d]^{g^*}
\\
\mathcal{F}(U) \ar[r]^{c(\phi \circ g)}
&
u_p\mathcal{F}(V')
}
$$
交换的唯一映射。这些映射的唯一性表明，我们由此得到一个预层。这个预层记作
$u_p\mathcal{F}$。

\begin{lemma}
\label{sites-lemma-recover}
存在一个与限制映射（$\mathcal{F}$ 上和 $u_p\mathcal{F}$ 上的限制映射）相容的
典范映射 $\mathcal{F}(U) \to u_p\mathcal{F}(u(U))$。
\end{lemma}

\begin{proof}
这正是上面引入的映射 $c(\text{id}_{u(U)})$。
\end{proof}

\noindent
注意，预层的任意映射 $\mathcal{F} \to \mathcal{F}'$ 都给出函子之间的相容
映射族 $\mathcal{F}_V \to \mathcal{F}'_V$，从而给出预层映射
$u_p\mathcal{F} \to u_p\mathcal{F}'$。换言之，我们定义了函子
$$
u_p :
\textit{PSh}(\mathcal{C})
\longrightarrow
\textit{PSh}(\mathcal{D})
$$

\begin{lemma}
\label{sites-lemma-adjoints-u}
函子 $u_p$ 是函子 $u^p$ 的左伴随。换言之，公式
$$
\Mor_{\textit{PSh}(\mathcal{C})}(\mathcal{F}, u^p\mathcal{G})
=
\Mor_{\textit{PSh}(\mathcal{D})}(u_p\mathcal{F}, \mathcal{G})
$$
对 $\mathcal{F}$ 和 $\mathcal{G}$ 双函子自然地成立。
\end{lemma}

\begin{proof}
设 $\mathcal{G}$ 是 $\mathcal{D}$ 上的预层，$\mathcal{F}$ 是
$\mathcal{C}$ 上的预层。我们将通过构造两个方向的映射来证明所示公式；留给
读者验证它们互为逆映射。

\medskip\noindent
给定映射 $\alpha : u_p \mathcal{F} \to \mathcal{G}$，得到
$u^p\alpha : u^p u_p \mathcal{F} \to u^p \mathcal{G}$。引理
\ref{sites-lemma-recover} 表明存在映射
$\mathcal{F} \to u^p u_p \mathcal{F}$。两者复合即给出所需映射。（这一构造的
好处是，它对眼前的一切显然都具有函子性。）

\medskip\noindent
反过来，给定映射 $\beta : \mathcal{F} \to u^p\mathcal{G}$，得到映射
$u_p\beta : u_p\mathcal{F} \to u_p u^p\mathcal{G}$。取
$V \in \Ob(\mathcal{D})$。我们断言，$\mathcal{I}_V$ 上的函子
$u^p\mathcal{G}_V$ 到取常值 $\mathcal{G}(V)$ 的常值函子有一个典范映射。
事实上，对 $\mathcal{I}_V$ 的每个对象 $(U, \phi)$，
$u^p\mathcal{G}_V$ 在该对象上的值是 $\mathcal{G}(u(U))$，它通过
$\mathcal{G}(\phi) = \phi^*$ 映到 $\mathcal{G}(V)$。由于
$\mathcal{G}$ 本身是函子，这构成函子之间的变换。由此得到映射
$u_p u^p \mathcal{G}(V) \to \mathcal{G}(V)$。另一个直接验证表明，该映射
对 $V$ 具有函子性，从而给出预层映射
$u_p u^p \mathcal{G} \to \mathcal{G}$。复合
$u_p\mathcal{F} \to u_p u^p\mathcal{G} \to \mathcal{G}$
即为所需映射。
\end{proof}

\begin{remark}
\label{sites-remark-functoriality-presheaves-values}
假设 $\mathcal{A}$ 是这样一个范畴：任意图表
$\mathcal{I}_Y \to \mathcal{A}$ 在 $\mathcal{A}$ 中都有余极限。那么显然，
在取值于 $\mathcal{A}$ 的预层范畴上，也存在与上面完全相同方式定义的函子
$u^p$ 和 $u_p$。此外，伴随对 $u^p$ 与 $u_p$ 的伴随性在此情形下仍然成立。
\end{remark}

\begin{lemma}
\label{sites-lemma-pullback-representable-presheaf}
设 $u : \mathcal{C} \to \mathcal{D}$ 为范畴之间的函子。对
$\mathcal{C}$ 的任意对象 $U$，都有 $u_ph_U = h_{u(U)}$。
\end{lemma}

\begin{proof}
由 $u_p$ 与 $u^p$ 的伴随性，有
$$
\Mor_{\textit{PSh}(\mathcal{D})}(u_ph_U, \mathcal{G})
=
\Mor_{\textit{PSh}(\mathcal{C})}(h_U, u^p\mathcal{G})
=
u^p\mathcal{G}(U) =
\mathcal{G}(u(U))
$$
因而由 Yoneda 引理可知，作为预层，$u_ph_U = h_{u(U)}$。
\end{proof}
















\section{位点}
\label{sites-section-sites-definitions}

\noindent
我们的位点概念使用下面这类结构。

\begin{definition}
\label{sites-definition-family-morphisms-fixed-target}
设 $\mathcal{C}$ 为范畴，见 Conventions，第
\ref{conventions-section-categories} 节。$\mathcal{C}$ 中一个
{\it 具有固定靶的态射族}，由一个对象 $U \in \Ob(\mathcal{C})$、一个集合
$I$，以及对每个 $i\in I$ 给出的一个靶为 $U$ 的 $\mathcal{C}$ 中态射
$U_i \to U$ 构成。我们用记号 $\{U_i \to U\}_{i\in I}$ 表示它。
\end{definition}

\noindent
集合 $I$ 可能为空！这一记号意在提示拓扑学中的开覆盖。

\begin{definition}
\label{sites-definition-site}
一个{\it 位点}\footnote{如引言所述，这一记号与 \cite{SGA4} 中的不同。}
由一个范畴 $\mathcal{C}$ 和一族具有固定靶的态射族
$\{U_i \to U\}_{i \in I}$ 所成的集合 $\text{Cov}(\mathcal{C})$ 给出；这些
态射族称为 $\mathcal{C}$ 的{\it 覆盖}，并满足下列公理：
\begin{enumerate}
\item 若 $V \to U$ 是同构，则
$\{V \to U\} \in \text{Cov}(\mathcal{C})$。
\item 若 $\{U_i \to U\}_{i\in I} \in \text{Cov}(\mathcal{C})$，且对每个
$i$ 都有 $\{V_{ij} \to U_i\}_{j\in J_i} \in \text{Cov}(\mathcal{C})$，则
$\{V_{ij} \to U\}_{i \in I, j\in J_i} \in \text{Cov}(\mathcal{C})$。
\item 若 $\{U_i \to U\}_{i\in I}\in \text{Cov}(\mathcal{C})$，且
$V \to U$ 是 $\mathcal{C}$ 的态射，则对所有 $i$，纤维积
$U_i \times_U V$ 都存在，并且
$\{U_i \times_U V \to V \}_{i\in I} \in \text{Cov}(\mathcal{C})$。
\end{enumerate}
\end{definition}

\noindent
作几点说明。在公理 (1) 中，我们要求存在 $\mathcal{C}$ 的一个覆盖
$\{U_i \to U\}_{i \in I}$，使得 $I = \{i\}$ 是单元素集，且态射
$U_i \to U$ 等于 (1) 中给定的态射 $V \to U$。下文中，我们经常用
$\{V \to U\}$ 表示指标集为单元素集的具有固定靶的态射族。在公理 (3) 中，
我们要求：对 $i \in I$ 的纤维积 $U_i \times_U V$ 的某种选择，所述覆盖存在。

\begin{remark}
\label{sites-remark-no-big-sites}
（关于集合论问题——初读时可跳过。）引入位点的主要原因，是研究位点上的层
范畴；它推广了在代数几何中极其重要的拓扑空间上的层范畴。为避免考虑“类的类”
之类的对象，与我们对范畴和 $2$-范畴的做法不同，我们不允许位点是“大全畴”。

\medskip\noindent
假设 $\mathcal{C}$ 是范畴，且 $\text{Cov}(\mathcal{C})$ 是满足上述
(1)、(2)、(3) 的覆盖真类。我们也不允许它作为位点，主要因为接下来要在覆盖
上取极限。不过，有几种自然方法可以把 $\text{Cov}(\mathcal{C})$ 替换成覆盖
的一个集合或稍有不同的结构，而所得层范畴不变。例如：
\begin{enumerate}
\item 在 Sets，第 \ref{sets-section-coverings-site} 节中，我们说明如何选取一个
适当的覆盖集，使所得层范畴相同。
\item 也可以取相伴拓扑（见定义 \ref{sites-definition-topology-associated-site}）。
$\mathcal{C}$ 上所得拓扑具有相同的层范畴。两个拓扑的层范畴相同，当且仅当
它们相等；见定理 \ref{sites-theorem-topology-and-topos}。范畴上的拓扑由对象上的筛
的选择给出。$\mathcal{C}$ 上所有可能的筛，乃至所有可能的拓扑，所成的集合
都是集合。
\item 还可以略微修改位点的概念，见下面的备注
\ref{sites-remark-shrink-coverings}，从而得到一个典范的覆盖集。
\end{enumerate}
每种方案都有一些小缺点。对第一种方案，需要验证后续构造不依赖于覆盖集的选择。
对第二种方案，需要学习拓扑，并为位点重新完成许多论证。对第三种方案，见备注
\ref{sites-remark-shrink-coverings} 的最后一句。

\medskip\noindent
我们将采用上述定义 \ref{sites-definition-site} 中的位点。给定一个范畴
$\mathcal{C}$ 以及如上的覆盖真类，我们用 Sets，引理
\ref{sets-lemma-coverings-site} 将其替换为一个覆盖集，从而产生一个位点。
下面的引理 \ref{sites-lemma-choice-set-coverings-immaterial} 将证明，所得层范畴
（托扑斯）与这一选择无关。若读者愿意，可自行采用另外两种策略之一处理这些
问题。
\end{remark}

\begin{example}
\label{sites-example-site-topological}
设 $X$ 为拓扑空间。令 $X_{Zar}$ 为这样的范畴：其对象是 $X$ 中的所有开集
$U$，态射仅为包含映射。也就是说，$X_{Zar}$ 的任意两个对象之间至多有一个
态射。现在定义
$\{U_i \to U\}_{i \in I}\in \text{Cov}(X_{Zar})$，当且仅当
$\bigcup U_i = U$。上述条件 (1) 和 (2) 显然成立；一旦注意到在
$X_{Zar}$ 中 $U \times V = U \cap V$，条件 (3) 也显然成立。特别要注意，
空集必须是 $X_{Zar}$ 的元素，否则一般而言这一构造无法成立。此外，正如稍后
将看到的，同样重要的是允许{\it 空集的空覆盖作为一个覆盖}！我们按照 Sets，
引理 \ref{sets-lemma-coverings-site} 选取适当的覆盖集
$\text{Cov}(X_{Zar})_{\kappa, \alpha}$，从而把 $X_{Zar}$ 变成位点。
% P01-ERRATA: P01-E0058 -- the frozen source has the ungrammatical plural phrase "a presheaves".
位点 $X_{Zar}$ 上的预层和层（定义见下文），恰好与 Sheaves，第
\ref{sheaves-section-introduction} 节所定义的拓扑空间上的通常预层和层概念
一致。
\end{example}

\begin{example}
\label{sites-example-site-on-group}
设 $G$ 为群。考虑范畴 $G\textit{-Sets}$，其对象是带左 $G$-作用的集合
$X$，态射是 $G$-等变映射。一个重要例子是 ${}_GG$：它是一个 $G$-集合，
其底层集合是 $G$，作用由左乘给出。这个范畴有纤维积，见 Categories，第
\ref{categories-section-example-fibre-products} 节。若
$\bigcup_{i\in I} \varphi_i(U_i) = U$，则称
$\{\varphi_i : U_i \to U\}_{i\in I}$ 为覆盖。这在
$G\textit{-Sets}$ 上给出一个覆盖类，容易看出它满足定义
\ref{sites-definition-site} 的条件 (1)、(2)、(3)。所得结构还不是位点，因为底层
范畴的对象汇集与覆盖汇集都是真类。
% P01-ERRATA: P01-E0059 -- the frozen source says "replace by ... by".
首先按照 Sets，引理 \ref{sets-lemma-sets-with-group-action}，把
$G\textit{-Sets}$ 替换为一个全子范畴 $G\textit{-Sets}_\alpha$。然后如
Sets，引理 \ref{sets-lemma-coverings-site} 那样适当限制覆盖汇集，所得位点
$(G\textit{-Sets}_\alpha,
\text{Cov}_{\kappa, \alpha'}(G\textit{-Sets}_\alpha))$
记作 $\mathcal{T}_G$。

\medskip\noindent
特别地，若群 $G$ 可数，则可令 $\mathcal{T}_G$ 为可数 $G$-集合的范畴，并
把所有指标集 $I$ 可数的联合满态射族
$\{\varphi_i : U_i \to U\}_{i \in I}$ 取作覆盖。
\end{example}

\begin{example}
\label{sites-example-indiscrete}
设 $\mathcal{C}$ 为范畴。有一种典范方式把它变成位点，其中 $U$ 的覆盖正是
$\{f : V \to U \mid f\text{ 为同构}\}$。这个位点上的层就是
$\mathcal{C}$ 上的预层。
% P01-ERRATA: P01-E0060 -- the frozen source says "This corresponding topology".
相应的拓扑称为{\it 混沌拓扑}或{\it 不可分拓扑}。
\end{example}
















\section{层}
\label{sites-section-sheaves}

\noindent
设 $\mathcal{C}$ 为位点。在引入取值于一个范畴的层这一概念之前，我们先说明
一个集合预层成为层是什么意思。设 $\mathcal{F}$ 是 $\mathcal{C}$ 上的集合
预层，并设
$\{U_i \to U\}_{i\in I}$ 是 $\text{Cov}(\mathcal{C})$ 的元素。根据假设，
所有纤维积 $U_i \times_U U_j$ 都在 $\mathcal{C}$ 中存在。有两个自然映射
$$
\xymatrix{
\prod\nolimits_{i\in I}
\mathcal{F}(U_i)
\ar@<1ex>[r]^-{\text{pr}_0^*} \ar@<-1ex>[r]_-{\text{pr}_1^*}
&
\prod\nolimits_{(i_0, i_1) \in I \times I}
\mathcal{F}(U_{i_0} \times_U U_{i_1})
}
$$
如所示，记为 $\text{pr}^*_i$，$i = 0, 1$。事实上，左端的一个元素对应于
一族 $(s_i)_{i\in I}$，其中每个 $s_i$ 都是 $\mathcal{F}$ 在 $U_i$ 上的
截面。对每一对 $(i_0, i_1) \in I \times I$，都有投影态射
$$
\text{pr}^{(i_0, i_1)}_{i_0} :
U_{i_0} \times_U U_{i_1}
\longrightarrow
U_{i_0}
\text{ 且 }
\text{pr}^{(i_0, i_1)}_{i_1} :
U_{i_0} \times_U U_{i_1}
\longrightarrow
U_{i_1}.
$$
因此，可以通过第一个映射拉回截面 $s_{i_0}$，也可以通过第二个映射拉回截面
$s_{i_1}$。明确地说，上述映射为
$$
\text{pr}_0^* :
(s_i)_{i\in I}
\longmapsto
\Big(
\text{pr}^{(i_0, i_1), *}_{i_0}(s_{i_0})
\Big)_{(i_0, i_1) \in I \times I}
$$
和
$$
\text{pr}_1^* :
(s_i)_{i\in I}
\longmapsto
\Big(
\text{pr}^{(i_0, i_1), *}_{i_1}(s_{i_1})
\Big)_{(i_0, i_1) \in I \times I}.
$$
最后考虑自然映射
$$
\mathcal{F}(U)
\longrightarrow
\prod\nolimits_{i\in I}
\mathcal{F}(U_i), \quad
s
\longmapsto
(s|_{U_i})_{i \in I}
$$
这里用 $s|_{U_i}$ 表示 $s$ 沿映射 $U_i \to U$ 的拉回。由
$\mathcal{F}$ 的函子性及纤维积图表的交换性，显然有
$\text{pr}_0^*( (s|_{U_i})_{i \in I} ) =
\text{pr}_1^*( (s|_{U_i})_{i \in I} )$。

\begin{definition}
\label{sites-definition-sheaf-sets}
设 $\mathcal{C}$ 为位点，$\mathcal{F}$ 为 $\mathcal{C}$ 上的集合预层。若对
每个覆盖 $\{U_i \to U\}_{i \in I} \in \text{Cov}(\mathcal{C})$，图表
\begin{equation}
\label{sites-equation-sheaf-condition}
\xymatrix{
\mathcal{F}(U) \ar[r]
&
\prod\nolimits_{i\in I}
\mathcal{F}(U_i)
\ar@<1ex>[r]^-{\text{pr}_0^*} \ar@<-1ex>[r]_-{\text{pr}_1^*}
&
\prod\nolimits_{(i_0, i_1) \in I \times I}
\mathcal{F}(U_{i_0} \times_U U_{i_1})
}
\end{equation}
把第一个箭头表示为 $\text{pr}_0^*$ 与 $\text{pr}_1^*$ 的等化子，则称
$\mathcal{F}$ 是一个{\it 层}。
\end{definition}

\noindent
粗略地说，这意味着：若给定截面 $s_i \in \mathcal{F}(U_i)$，使得对所有
$(i, j) \in I \times I$，在 $\mathcal{F}(U_i \times_U U_j)$ 中都有
$$
s_i|_{U_i \times_U U_j} = s_j|_{U_i \times_U U_j}
$$
则存在唯一的 $s \in \mathcal{F}(U)$，使得 $s_i = s|_{U_i}$。

\begin{remark}
\label{sites-remark-sheaf-condition-empty-covering}
若覆盖 $\{U_i \to U\}_{i \in I}$ 是空族（即 $I = \emptyset$），则层条件
表示 $\mathcal{F}(U) = \{*\}$ 是单元素集。这是因为在
(\ref{sites-equation-sheaf-condition}) 中，第二、第三个集合都是集合范畴中的空积；
它们是集合范畴的终对象，因而是单元素集。
\end{remark}

\begin{example}
\label{sites-example-sheaves-topological}
设 $X$ 为拓扑空间，令 $X_{Zar}$ 为例
\ref{sites-example-site-topological} 中构造的位点。$X_{Zar}$ 上的层概念与
Sheaves，定义 \ref{sheaves-definition-sheaf} 中引入的 $X$ 上层的概念一致。
\end{example}

\begin{example}
\label{sites-example-topological-wrong}
设 $X$ 为拓扑空间。考虑位点 $X'_{Zar}$：除了不允许空集的空覆盖之外，它与例
\ref{sites-example-site-topological} 中的位点 $X_{Zar}$ 相同。换言之，我们允许覆盖
$\{\emptyset \to \emptyset\}$，但不允许指标集为空的覆盖。容易证明，这仍然
定义一个位点。不过，我们断言 $X'_{Zar}$ 上的层与 $X_{Zar}$ 上的层不同。
例如，作为一个极端情形，考虑 $X = \{p\}$ 为单元素集的情况。此时
$X'_{Zar}$ 的对象是 $\emptyset, X$；$\emptyset$ 的每个覆盖都可由
$\{\emptyset \to \emptyset\}$ 加细，而 $X$ 的每个覆盖都可由
$\{X \to X\}$ 加细。显然，这里的一个层可由集合
$\mathcal{F}(\emptyset)$ 的任意选择、集合 $\mathcal{F}(X)$ 的任意选择，
以及限制映射 $\mathcal{F}(X) \to \mathcal{F}(\emptyset)$ 的任意选择给出。
因此，$X'_{Zar}$ 上的层等同于二点空间 $\{\eta, p\}$ 上的通常层，其中开集为
$\{\emptyset, \{\eta\}, \{p, \eta\}\}$。一般而言，$X'_{Zar}$ 上的层
等同于空间 $X \amalg \{\eta\}$ 上的层；该空间的开集是空集以及形如
$U \cup \{\eta\}$ 的任意集合，其中 $U \subset X$ 为开集。
\end{example}


\begin{definition}
\label{sites-definition-category-sheaves-sets}
集合层范畴 {\it $\Sh(\mathcal{C})$} 是
$\textit{PSh}(\mathcal{C})$ 的全子范畴，其对象为集合层。
\end{definition}

\noindent
设 $\mathcal{A}$ 为范畴。若对每个作为覆盖族指标集出现的 $I$，
$\mathcal{A}$ 中都存在以 $I$ 和 $I \times I$ 为指标的积，则上述定义
\ref{sites-definition-sheaf-sets} 有意义，并定义了 $\mathcal{C}$ 上取值于
$\mathcal{A}$ 的层这一概念。注意，$\mathcal{A}$ 中的图表
$$
\xymatrix{
\mathcal{F}(U) \ar[r]
&
\prod\nolimits_{i\in I}
\mathcal{F}(U_i)
\ar@<1ex>[r]^-{\text{pr}_0^*} \ar@<-1ex>[r]_-{\text{pr}_1^*}
&
\prod\nolimits_{(i_0, i_1) \in I \times I}
\mathcal{F}(U_{i_0} \times_U U_{i_1})
}
$$
是等化子图表，当且仅当对 $\mathcal{A}$ 的每个对象 $X$，集合图表
$$
\xymatrix{
\Mor_\mathcal{A}(X, \mathcal{F}(U)) \ar[r]
&
\prod
\Mor_\mathcal{A}(X, \mathcal{F}(U_i))
\ar@<1ex>[r]^-{\text{pr}_0^*} \ar@<-1ex>[r]_-{\text{pr}_1^*}
&
\prod
\Mor_\mathcal{A}(X, \mathcal{F}(U_{i_0} \times_U U_{i_1}))
}
$$
是等化子图表。

\medskip\noindent
现在假设 $\mathcal{A}$ 任意。设 $\mathcal{F}$ 为取值于 $\mathcal{A}$ 的
预层。任取对象 $X\in \Ob(\mathcal{A})$。由规则
$$
\mathcal{F}_X(U) = \Mor_\mathcal{A}(X, \mathcal{F}(U)).
$$
定义一个集合预层 $\mathcal{F}_X$。由上可知，要求所有预层
$\mathcal{F}_X$ 都是集合层，便得到一个良好定义。

\begin{definition}
\label{sites-definition-sheaf}
设 $\mathcal{C}$ 为位点，$\mathcal{A}$ 为范畴，$\mathcal{F}$ 是
$\mathcal{C}$ 上取值于 $\mathcal{A}$ 的预层。若对 $\mathcal{A}$ 的所有对象
$X$，上面定义的集合预层 $\mathcal{F}_X$ 都是层，则称 $\mathcal{F}$ 是一个
{\it 层}。
\end{definition}












\section{具有固定靶的态射族}
\label{sites-section-refinements}

\noindent
本节引入若干与具有相同靶的态射族有关的概念。

\begin{definition}
\label{sites-definition-morphism-coverings}
设 $\mathcal{C}$ 为范畴。设
$\mathcal{U} = \{U_i \to U\}_{i\in I}$ 是 $\mathcal{C}$ 中一个具有固定靶的
态射族，$\mathcal{V} = \{V_j \to V\}_{j\in J}$ 是另一个这样的态射族。
\begin{enumerate}
\item 从 $\mathcal{U}$ 到 $\mathcal{V}$ 的一个
{\it $\mathcal{C}$ 中具有固定靶的映射族的态射}，简称
{\it 从 $\mathcal{U}$ 到 $\mathcal{V}$ 的态射}，由一个态射 $U \to V$、
一个集合映射 $\alpha : I \to J$，以及对每个 $i\in I$ 给出的一个态射
$U_i \to V_{\alpha(i)}$ 构成，使得图表
$$
\xymatrix{
U_i \ar[r] \ar[d]
&
V_{\alpha(i)} \ar[d]
\\
U \ar[r]
&
V
}
$$
交换。
\item 在 $U = V$ 且 $U \to V$ 是恒等态射的特殊情形下，我们称
$\mathcal{U}$ 是态射族 $\mathcal{V}$ 的一个{\it 加细}。
\end{enumerate}
\end{definition}

\noindent
一个平凡但重要的观察是：若
$\mathcal{V} = \{V_j \to V\}_{j \in J}$ 是{\it 空映射族}，即
$J = \emptyset$，则任何满足 $I \not = \emptyset$ 的族
$\mathcal{U} = \{U_i \to V\}_{i \in I}$ 都不能加细 $\mathcal{V}$！

\begin{definition}
\label{sites-definition-combinatorial-tautological}
设 $\mathcal{C}$ 为范畴。设
$\mathcal{U} = \{\varphi_i : U_i \to U\}_{i\in I}$ 和
$\mathcal{V} = \{\psi_j : V_j \to U\}_{j\in J}$ 是两个具有固定靶的态射族。
\begin{enumerate}
\item 若存在映射 $\alpha : I \to J$ 和 $\beta : J\to I$，使得
$\varphi_i = \psi_{\alpha(i)}$ 且 $\psi_j = \varphi_{\beta(j)}$，则称
$\mathcal{U}$ 与 $\mathcal{V}$ {\it 组合等价}。
\item 若存在映射 $\alpha : I \to J$ 和 $\beta : J\to I$，并且对所有
$i\in I$ 和 $j \in J$ 都有交换图表
$$
\xymatrix{
U_i \ar[rd] \ar[rr] & &
V_{\alpha(i)} \ar[ld] & &
V_j \ar[rd] \ar[rr] & &
U_{\beta(j)} \ar[ld] \\
&
U & & & &
U &
}
$$
其中水平箭头都是同构，则称 $\mathcal{U}$ 与 $\mathcal{V}$
{\it 自明等价}。
\end{enumerate}
\end{definition}

\begin{lemma}
\label{sites-lemma-tautological-combinatorial}
设 $\mathcal{C}$ 为范畴。设
$\mathcal{U} = \{\varphi_i : U_i \to U\}_{i\in I}$ 和
$\mathcal{V} = \{\psi_j : V_j \to U\}_{j\in J}$ 是两个具有相同固定靶的
态射族。
\begin{enumerate}
\item 若 $\mathcal{U}$ 与 $\mathcal{V}$ 组合等价，则它们自明等价。
\item 若 $\mathcal{U}$ 与 $\mathcal{V}$ 自明等价，则 $\mathcal{U}$ 是
$\mathcal{V}$ 的加细，且 $\mathcal{V}$ 是 $\mathcal{U}$ 的加细。
\item “组合等价”关系是所有具有固定靶的态射族上的等价关系。
\item “自明等价”关系是所有具有固定靶的态射族上的等价关系。
\item “$\mathcal{U}$ 加细 $\mathcal{V}$ 且 $\mathcal{V}$ 加细
$\mathcal{U}$”关系是所有具有固定靶的态射族上的等价关系。
\end{enumerate}
\end{lemma}

\begin{proof}
从略。
\end{proof}

\noindent
在下面的引理中，给定一个范畴 $\mathcal{C}$、$\mathcal{C}$ 上的预层
$\mathcal{F}$，以及满足所有纤维积 $U_i \times_U U_{i'}$ 都存在的态射族
$\mathcal{U} = \{U_i \to U\}_{i\in I}$，若图表
(\ref{sites-equation-sheaf-condition}) 是等化子图表，我们就说
{\it $\mathcal{F}$ 关于 $\mathcal{U}$ 的层条件}成立。

\begin{lemma}
\label{sites-lemma-tautological-same-sheaf}
设 $\mathcal{C}$ 为范畴。设
$\mathcal{U} = \{\varphi_i : U_i \to U\}_{i\in I}$ 和
$\mathcal{V} = \{\psi_j : V_j \to U\}_{j\in J}$ 是两个具有相同固定靶的
态射族。假设纤维积 $U_i \times_U U_{i'}$ 和
$V_j \times_U V_{j'}$ 都存在。若 $\mathcal{U}$ 与 $\mathcal{V}$ 自明等价，
则对 $\mathcal{C}$ 上任意预层 $\mathcal{F}$，$\mathcal{F}$ 关于
$\mathcal{U}$ 的层条件等价于 $\mathcal{F}$ 关于 $\mathcal{V}$ 的层条件。
\end{lemma}

\begin{proof}
首先注意，若 $\varphi : A \to B$ 是 $\mathcal{C}$ 中的同构，则
$\varphi^* : \mathcal{F}(B) \to \mathcal{F}(A)$ 是同构。设
$\beta : J \to I$ 为映射，并设
$\chi_j : V_j \to U_{\beta(j)}$ 是假设所保证的 $U$ 上同构。我们证明，
$\mathcal{V}$ 的层条件蕴含 $\mathcal{U}$ 的层条件。假设给定截面
$s_i \in \mathcal{F}(U_i)$，使得对所有
$(i, i') \in I \times I$，在 $\mathcal{F}(U_i \times_U U_{i'})$ 中有
$$
s_i|_{U_i \times_U U_{i'}} = s_{i'}|_{U_i \times_U U_{i'}}
$$
定义 $s_j = \chi_j^*s_{\beta(j)}$。对任意一对
$(j, j') \in J \times J$，态射
$\chi_j \times_{\text{id}_U} \chi_{j'} : V_j \times_U V_{j'} \to
U_{\beta(j)} \times_U U_{\beta(j')}$ 也是同构。因此，由结构传递可知
$$
s_j|_{V_j \times_U V_{j'}} = s_{j'}|_{V_j \times_U V_{j'}}
$$
$\mathcal{V}$ 的层条件蕴含存在唯一的 $s$，使得对所有 $j \in J$，都有
$s|_{V_j} = s_j$。由证明开头的观察，这也蕴含对所有
$i \in \Im(\beta)$，都有 $s|_{U_i} = s_i$。现设 $i \in I$ 且
$i \not \in \Im(\beta)$。对此 $i$，在 $U$ 上有同构
$U_i \to V_{\alpha(i)} \to U_{\beta(\alpha(i))}$。它给出态射
$U_i \to U_i \times_U U_{\beta(\alpha(i))}$，且此态射是投影的截面。
由于 $s_i$ 和 $s_{\beta(\alpha(i))}$ 在纤维积上限制为同一元素，可知
$s_{\beta(\alpha(i))}$ 沿
$U_i \to U_{\beta(\alpha(i))}$ 拉回为 $s_i$。因此也有
$s_i = s|_{U_i}$，正合所需。
\end{proof}

\begin{lemma}
\label{sites-lemma-compare-separated-presheaf-condition}
设 $\mathcal{C}$ 为范畴。设
$\mathcal{V} = \{V_j \to U\}_{j \in J} \to
\mathcal{U} = \{U_i \to U\}_{i \in I}$ 是 $\mathcal{C}$ 中具有固定靶的映射
族的一个态射，由 $\text{id} : U \to U$、$\alpha : J \to I$ 和
$f_j : V_j \to U_{\alpha(j)}$ 给出。设 $\mathcal{F}$ 是
$\mathcal{C}$ 上的预层。若
$\mathcal{F}(U) \to \prod_{j \in J} \mathcal{F}(V_j)$ 是单射，则
$\mathcal{F}(U) \to \prod_{i \in I} \mathcal{F}(U_i)$ 也是单射。
\end{lemma}

\begin{proof}
从略。
\end{proof}

\begin{lemma}
\label{sites-lemma-compare-sheaf-condition}
设 $\mathcal{C}$ 为范畴。设
$\mathcal{V} = \{V_j \to U\}_{j \in J} \to
\mathcal{U} = \{U_i \to U\}_{i \in I}$ 是 $\mathcal{C}$ 中具有固定靶的映射
族的一个态射，由 $\text{id} : U \to U$、$\alpha : J \to I$ 和
$f_j : V_j \to U_{\alpha(j)}$ 给出。设 $\mathcal{F}$ 是
$\mathcal{C}$ 上的预层。假设
\begin{enumerate}
\item 纤维积 $U_i \times_U U_{i'}$、$U_i \times_U V_j$、
$V_j \times_U V_{j'}$ 均存在；
\item $\mathcal{F}$ 满足关于 $\mathcal{V}$ 的层条件；
\item 对每个 $i \in I$，映射
$\mathcal{F}(U_i) \to \prod_{j \in J} \mathcal{F}(V_j \times_U U_i)$
都是单射。
\end{enumerate}
那么 $\mathcal{F}$ 满足关于 $\mathcal{U}$ 的层条件。
\end{lemma}

\begin{proof}
由引理 \ref{sites-lemma-compare-separated-presheaf-condition}，映射
$\mathcal{F}(U) \to \prod \mathcal{F}(U_i)$ 是单射。假设给定
$s_i \in \mathcal{F}(U_i)$，使得对所有 $i, i' \in I$，都有
$s_i|_{U_i \times_U U_{i'}} = s_{i'}|_{U_i \times_U U_{i'}}$。令
$s_j = f_j^*(s_{\alpha(j)}) \in \mathcal{F}(V_j)$。由于态射 $f_j$ 都是
$U$ 上的态射，我们得到与投影映射下的 $f_j, f_{j'}$ 相容的诱导态射
% P01-ERRATA: P01-E0061 -- the frozen target uses undefined i,i' instead of j,j'.
$f_{jj'} : V_j \times_U V_{j'} \to
U_{\alpha(i)} \times_U U_{\alpha(i')}$。由此可得
$$
s_j|_{V_j \times_U V_{j'}}
= f_{jj'}^*(s_{\alpha(j)}|_{U_{\alpha(j)} \times_U U_{\alpha(j')}})
= f_{jj'}^*(s_{\alpha(j')}|_{U_{\alpha(j)} \times_U U_{\alpha(j')}})
= s_{j'}|_{V_j \times_U V_{j'}}
$$
对所有 $j, j' \in J$ 成立。因此，由 $\mathcal{F}$ 关于
$\mathcal{V}$ 的层条件，得到截面 $s \in \mathcal{F}(U)$，其在每个 $V_j$
上的限制为 $s_j$。若能证明对任意 $i \in I$，$s$ 在 $U_i$ 上的限制为
$s_i$，则证明完成。由于 $\mathcal{F}$ 满足 (3)，只须证明对所有
$j \in J$，$s$ 和 $s_i$ 在 $U_i \times_U V_j$ 上限制为同一元素。为此使用
$$
s|_{U_i \times_U V_j} = s_j|_{U_i \times_U V_j} =
(\text{id} \times f_j)^*s_{\alpha(j)}|_{U_i \times_U U_{\alpha(j)}} =
(\text{id} \times f_j)^*s_i|_{U_i \times_U U_{\alpha(j)}} =
s_i|_{U_i \times_U V_j}
$$
正合所需。
\end{proof}

\begin{lemma}
\label{sites-lemma-refine-same-topology}
设 $\mathcal{C}$ 为范畴。设 $\text{Cov}_i$，$i = 1, 2$，是两组具有固定靶的
态射族集合，并且各自在 $\mathcal{C}$ 上定义一个位点结构。
\begin{enumerate}
\item 若每个 $\mathcal{U} \in \text{Cov}_1$ 都与某个
$\mathcal{V} \in \text{Cov}_2$ 自明等价，则
$\Sh(\mathcal{C}, \text{Cov}_2) \subset
\Sh(\mathcal{C}, \text{Cov}_1)$。
若还对每个 $\mathcal{U} \in \text{Cov}_2$，都存在某个与之自明等价的
$\mathcal{V} \in \text{Cov}_1$，则两个层范畴相等。
% P01-ERRATA: P01-E0062 -- the frozen source has subject-verb disagreement in "category ... are".
\item 假设对每个 $\mathcal{U} \in \text{Cov}_1$，都存在
$\mathcal{V} \in \text{Cov}_2$，使得 $\mathcal{V}$ 加细 $\mathcal{U}$。
在这种情况下
$\Sh(\mathcal{C}, \text{Cov}_2) \subset
\Sh(\mathcal{C}, \text{Cov}_1)$。
若还对每个 $\mathcal{U} \in \text{Cov}_2$，都存在
$\mathcal{V} \in \text{Cov}_1$，使得 $\mathcal{V}$ 加细 $\mathcal{U}$，
则两个层范畴相等。
\end{enumerate}
\end{lemma}

\begin{proof}
(1) 直接由引理 \ref{sites-lemma-tautological-same-sheaf} 和定义得到。

\medskip\noindent
证明 (2)。设 $\mathcal{F}$ 是位点
$(\mathcal{C}, \text{Cov}_2)$ 上的集合层。取
$\mathcal{U} \in \text{Cov}_1$，写成
$\mathcal{U} = \{U_i \to U\}_{i \in I}$。由假设，可以选择
$\mathcal{U}$ 的一个加细 $\mathcal{V} \in \text{Cov}_2$；写成
$\mathcal{V} = \{V_j \to U\}_{j \in J}$，其加细数据由
$\alpha : J \to I$ 和 $f_j : V_j \to U_{\alpha(j)}$ 给出。注意，
$\mathcal{F}$ 满足 $\mathcal{V}$ 的层条件，也满足覆盖
$\{V_j \times_U U_i \to U_i\}_{j \in J}$ 的层条件，因为这些覆盖属于
$\text{Cov}_2$。故由引理 \ref{sites-lemma-compare-sheaf-condition}，
$\mathcal{F}$ 满足 $\mathcal{U}$ 的层条件。
\end{proof}

\begin{lemma}
\label{sites-lemma-choice-set-coverings-immaterial}
设 $\mathcal{C}$ 为范畴。设 $\text{Cov}(\mathcal{C})$ 是满足定义
\ref{sites-definition-site} 中条件 (1)、(2)、(3) 的覆盖真类。设
$\text{Cov}_1, \text{Cov}_2 \subset \text{Cov}(\mathcal{C})$ 是
$\text{Cov}(\mathcal{C})$ 的两个子集，它们都赋予 $\mathcal{C}$ 位点结构。
若每个覆盖 $\mathcal{U} \in \text{Cov}(\mathcal{C})$ 都与
$\text{Cov}_1$ 中的某个覆盖组合等价，并且与 $\text{Cov}_2$ 中的某个覆盖
组合等价，则
$\Sh(\mathcal{C}, \text{Cov}_1) =
\Sh(\mathcal{C}, \text{Cov}_2)$。
\end{lemma}

\begin{proof}
由上述引理 \ref{sites-lemma-refine-same-topology} 和
\ref{sites-lemma-tautological-combinatorial} 即得，因为假设蕴含每个覆盖
$\mathcal{U} \in \text{Cov}_1 \subset \text{Cov}(\mathcal{C})$
都与 $\text{Cov}_2$ 的一个元素组合等价；交换 $\text{Cov}_1$ 与
$\text{Cov}_2$ 的角色亦然。
\end{proof}

















\section{G-集合的例子}
\label{sites-section-example-sheaf-G-sets}

\noindent
作为例子，考虑例 \ref{sites-example-site-on-group} 中的位点 $\mathcal{T}_G$。
我们将描述 $\mathcal{T}_G$ 上的层范畴。结果将与定义 $\mathcal{T}_G$ 时所作
的选择无关。事实上，在证明中只需要位点 $\mathcal{T}_G$ 的下列性质：
\begin{enumerate}
\item[(a)] $\mathcal{T}_G$ 是 $G\textit{-Sets}$ 的全子范畴；
\item[(b)] $\mathcal{T}_G$ 包含 $G$-集合 ${}_GG$；
\item[(c)] $\mathcal{T}_G$ 有纤维积，且它们与 $G\textit{-Sets}$ 中的相同；
\item[(d)] 给定 $U \in \Ob(\mathcal{T}_G)$ 和 $G$-不变子集
$O \subset U$，存在 $\mathcal{T}_G$ 的一个与 $O$ 同构的对象；
\item[(e)] 任何满射映射族 $\{U_i \to U\}_{i \in I}$，其中
$U, U_i \in \Ob(\mathcal{T}_G)$，都与 $\mathcal{T}_G$ 的某个覆盖组合等价。
\end{enumerate}
这些性质由 Sets，引理 \ref{sets-lemma-what-is-in-it-G-sets} 和
\ref{sets-lemma-coverings-site} 成立。

\medskip\noindent
注意，映射
$$
\Hom_G({}_GG, {}_GG)
\longrightarrow
G^{opp},
\varphi
\longmapsto
\varphi(1)
$$
是群同构。其逆映射把 $g \in G$ 送到映射
$R_g : s \mapsto sg$（即右乘）。注意
$R_{g_1g_2} = R_{g_2} \circ R_{g_1}$，故必须取对群。

\medskip\noindent
这意味着，对 $\mathcal{T}_G$ 上的每个预层 $\mathcal{F}$，集合
$\mathcal{F}({}_GG)$ 按如下方式继承一个 $G$-集合结构：对
$g \in G$ 和 $s \in \mathcal{F}({}_GG)$，把 $g \cdot s$ 定义为
$\mathcal{F}(R_g)(s)$。这是左作用，因为
$$
(g_1g_2) \cdot s  = \mathcal{F}(R_{g_1g_2})(s) =
\mathcal{F}(R_{g_2}\circ R_{g_1})(s) =
\mathcal{F}(R_{g_1})( \mathcal{F}(R_{g_2})(s)) =
g_1 \cdot (g_2 \cdot s).
$$
这里使用了 $\mathcal{F}$ 的反变性。注意，若
$\mathcal{F} \to \mathcal{G}$ 是 $\mathcal{T}_G$ 上集合预层的态射，则得到
映射 $\mathcal{F}({}_GG) \to \mathcal{G}({}_GG)$，它与刚定义的
$G$-作用相容。综上，我们构造了函子
$$
\textit{PSh}(\mathcal{T}_G)
\longrightarrow
G\textit{-Sets}, \quad
\mathcal{F}
\longmapsto
\mathcal{F}({}_GG).
$$
留给读者验证，这一构造具有如下良好性质：可表预层 $h_U$ 被映到一个与 $U$
典范同构的对象。用公式表示，即
$h_U({}_GG) = \Hom_G({}_GG, U) \cong U$。

\medskip\noindent
假设 $S$ 是一个 $G$-集合。用公式\footnote{它看起来可能是由 $S$ 定义的可表
预层，但未必如此，因为 $S$ 可能不是 $\mathcal{T}_G$ 的对象；该位点被选成
一个足够大的 $G$-集合之集。}
$$
\mathcal{F}_S(U)
=
\Mor_{G\textit{-Sets}}(U, S).
$$
定义预层 $\mathcal{F}_S$。这显然是一个预层。另一方面，假设
$\{U_i \to U\}_{i\in I}$ 是 $\mathcal{T}_G$ 中的覆盖。这蕴含
$\coprod_i U_i \to U$ 是满射。因此映射
$$
\mathcal{F}_S(U)
=
\Mor_{G\textit{-Sets}}(U, S)
\longrightarrow
\prod \mathcal{F}_S(U_i)
=
\prod \Mor_{G\textit{-Sets}}(U_i, S)
$$
显然是单射。再给定一族 $G$-等变映射 $s_i : U_i \to S$，使得所有图表
$$
\xymatrix{
U_i \times_U U_j \ar[d] \ar[r]
&
U_j \ar[d]^{s_j}
\\
U_i \ar[r]^{s_i}
&
S
}
$$
都交换，则存在唯一的 $G$-等变映射 $s : U \to S$，使得 $s_i$ 是复合
$U_i \to U \to S$。事实上，定义 $s(u) = s_i(u_i)$，其中 $i\in I$ 是任一
满足存在某个 $u_i \in U_i$ 经映射 $U_i \to U$ 映到 $u$ 的指标。上述图表
的交换性恰好蕴含这一构造是良定义的。综上，我们构造了函子
$$
G\textit{-Sets}
\longrightarrow
\Sh(\mathcal{T}_G), \quad
S
\longmapsto
\mathcal{F}_S
.
$$

\medskip\noindent
现在得到如下范畴与函子的图表：
$$
\xymatrix{
\textit{PSh}(\mathcal{T}_G) \ar[rr]^{\mathcal{F} \mapsto \mathcal{F}({}_GG)}
&
&
G\textit{-Sets} \ar[ld]_{S \mapsto \mathcal{F}_S}
\\
&
\Sh(\mathcal{T}_G) \ar[lu]
&
}
$$
由定义立即可知
$\mathcal{F}_S({}_GG) = \Mor_G({}_GG, S) = S$，最后一个等号来自在 $1$
处取值。这几乎证明了下面的命题。

\begin{proposition}
\label{sites-proposition-sheaves-on-group}
函子 $\mathcal{F} \mapsto \mathcal{F}({}_GG)$ 和
$S \mapsto \mathcal{F}_S$ 给出 $\Sh(\mathcal{T}_G)$ 与
$G\textit{-Sets}$ 之间的一对拟逆等价。
\end{proposition}

\begin{proof}
我们已经看到，按一个方向复合这两个函子得到与恒等函子同构的函子。在另一个
方向上，对任意层 $\mathcal{H}$，存在自然的层映射
$$
can :
\mathcal{H}
\longrightarrow
\mathcal{F}_{\mathcal{H}({}_GG)}.
$$
事实上，对 $\mathcal{T}_G$ 的任意对象 $U$，令 $can_U$ 为映射
$$
\begin{matrix}
\mathcal{H}(U)
&
\longrightarrow
&
\mathcal{F}_{\mathcal{H}({}_GG)}(U)
=
\Mor_G(U, \mathcal{H}({}_GG))
\\
s
&
\longmapsto
&
(u \mapsto \alpha_u^*s).
\end{matrix}
$$
这里 $\alpha_u : {}_GG \to U$ 是映射 $\alpha_u(g) = gu$，而
$\alpha_u^* : \mathcal{H}(U) \to \mathcal{H}({}_GG)$ 是拉回映射。
一个平凡但容易混淆的验证表明，这确实是预层映射。我们需要证明 $can$ 是同构。
为此，证明对所有 $U \in \Ob(\mathcal{T}_G)$，$can_U$ 都是同构。
把 $U = {}_GG$ 这一重要但容易的情形留给读者。$\mathcal{T}_G$ 的一般对象
$U$ 是若干 $G$-轨道的不交并：$U = \coprod_{i\in I} O_i$。映射族
$\{O_i \to U\}_{i \in I}$ 与 $\mathcal{T}_G$ 中的某个覆盖自明等价（由本节
开头列出的 $\mathcal{T}_G$ 的性质）。故由引理
\ref{sites-lemma-tautological-same-sheaf}，层 $\mathcal{H}$ 满足关于
$\{O_i \to U\}_{i \in I}$ 的层性质。这个覆盖的层性质蕴含
$\mathcal{H}(U) = \prod_i \mathcal{H}(O_i)$。因此，只须证明当 $U$ 由单个
$G$-轨道构成时，$can_U$ 是同构。取 $u \in U$，并令 $H \subset G$ 为其
稳定子。显然，$\Mor_G(U, \mathcal{H}({}_GG)) = \mathcal{H}({}_GG)^H$
等于 $H$-不变元素组成的子集。另一方面，考虑由 $g \mapsto gu$ 给出的覆盖
$\{{}_GG \to U\}$（它同样只是与 $\mathcal{T}_G$ 中的某个覆盖组合等价，
而这同样无关紧要）。注意，纤维积 $({}_GG)\times_U ({}_GG)$ 等于
$\{(g, gh), g\in G, h\in H\} \cong \coprod_{h\in H} {}_GG$。因此，这个覆盖
的层性质表述为
$$
\xymatrix{
\mathcal{H}(U) \ar[r]
&
\mathcal{H}({}_GG)
\ar@<1ex>[r]^-{\text{pr}_0^*} \ar@<-1ex>[r]_-{\text{pr}_1^*}
&
\prod_{h \in H}
\mathcal{H}({}_GG).
}
$$
到因子 $\mathcal{H}({}_GG)$ 的两个映射 $\text{pr}_i^*$ 相差乘以 $h$。
现在，结论由此以及 $can$ 在 $U = {}_GG$ 时是同构这一事实得到。
\end{proof}






















\section{层化}
\label{sites-section-sheafification}

\noindent
为了定义层化，我们研究一个覆盖的第零 {\v C}ech 上同调群及其函子性。
% P01-ERRATA: P01-E0063 -- set-valued H^0 need not be a group.

\medskip\noindent
设 $\mathcal{F}$ 是 $\mathcal{C}$ 上的集合预层，并设
$\mathcal{U} = \{U_i \to U\}_{i \in I}$ 是 $\mathcal{C}$ 的一个覆盖。
用记号 $\mathcal{F}(\mathcal{U})$ 表示等化子
$$
H^0(\mathcal{U}, \mathcal{F})
=
\{
(s_i)_{i\in I} \in \prod\nolimits_i \mathcal{F}(U_i)
\mid
s_i|_{U_i \times_U U_j} = s_j|_{U_i \times_U U_j}
\ \forall i, j \in I
\}.
$$
稍后将看到，这是 $\mathcal{F}$ 在 $U$ 上关于覆盖 $\mathcal{U}$ 的第零
{\v C}ech 上同调。顺便指出，只要所有态射 $U_i \to U$ 都可表，便可以定义
$H^0(\mathcal{U}, \mathcal{F})$；也就是说，$\mathcal{U}$ 不必是该位点的
覆盖。存在典范映射
$\mathcal{F}(U) \to H^0(\mathcal{U}, \mathcal{F})$。显然，覆盖的一个态射
$\mathcal{U} \to \mathcal{V}$ 诱导交换图表
$$
\xymatrix{
& U_i \ar[rr] & & V_{\alpha(i)} \\
U_i \times_U U_j \ar[rr] \ar[ur] \ar[dr] & &
V_{\alpha(i)} \times_V V_{\alpha(j)} \ar[ur] \ar[dr] & \\
& U_j \ar[rr] & & V_{\alpha(j)}
}.
$$
它进而产生映射 $H^0(\mathcal{V}, \mathcal{F}) \to
H^0(\mathcal{U}, \mathcal{F})$，并与映射
$\mathcal{F}(V) \to \mathcal{F}(U)$ 相容。

\medskip\noindent
由构造，预层 $\mathcal{F}$ 是层，当且仅当对 $\mathcal{C}$ 的每个覆盖
$\mathcal{U}$，自然映射
$\mathcal{F}(U) \to H^0(\mathcal{U}, \mathcal{F})$ 都是双射。我们将用这一
概念证明下面关于层的极限的简单引理。

\begin{lemma}
\label{sites-lemma-limit-sheaf}
\begin{slogan}
层图表的极限存在，并且与作为预层的极限一致
\end{slogan}
设 $\mathcal{F} : \mathcal{I} \to \Sh(\mathcal{C})$ 是一个图表。那么
$\lim_\mathcal{I} \mathcal{F}$ 存在，并且等于预层范畴中的极限。
\end{lemma}

\begin{proof}
设 $\lim_i \mathcal{F}_i$ 为作为预层的极限。我们证明它是层；由此立即可知，
它也是层范畴中的极限。为证明层性质，设
$\mathcal{V} = \{V_j \to V\}_{j\in J}$ 为覆盖。取
$(s_j)_{j\in J}$ 为 $H^0(\mathcal{V}, \lim_i \mathcal{F}_i)$ 的元素。
利用投影映射，得到 $H^0(\mathcal{V}, \mathcal{F}_i)$ 中的元素
$(s_{j, i})_{j\in J}$。由 $\mathcal{F}_i$ 的层性质，存在唯一的
$s_i \in \mathcal{F}_i(V)$，使得 $s_{j, i} = s_i|_{V_j}$。设
$\phi : i \to i'$ 是指标范畴的态射。我们要证明
$\mathcal{F}(\phi) : \mathcal{F}_i \to \mathcal{F}_{i'}$ 把 $s_i$ 映到
$s_{i'}$。我们知道，对所有 $j$，这对截面 $s_{i, j}$ 和 $s_{i', j}$ 成立；
故由 $\mathcal{F}_{i'}$ 的层性质，所需结论成立。此时得到
$(\lim_i \mathcal{F}_i)(V)$ 的一个元素
$s = (s_i)_{i \in \Ob(\mathcal{I})}$。留给读者验证，该元素具有所需性质
$s_j = s|_{V_j}$。
\end{proof}

\begin{example}
\label{sites-example-singleton-sheaf}
一个特殊例子是空图表上的极限。它给出（预）层范畴中的终对象。该预层把一个
单元素集赋给 $\mathcal{C}$ 的每个对象 $U$，其限制映射均为唯一映射，并且
这个预层还是层。我们常用 $*$ 表示这个层。
\end{example}

\noindent
令 $\mathcal{J}_U$ 为 $U$ 的所有覆盖组成的范畴。换言之，
$\mathcal{J}_U$ 的对象是 $\mathcal{C}$ 中 $U$ 的覆盖，态射是加细。根据我们
对位点的约定，这确实是一个范畴，也就是说，其对象和态射的汇集都是集合。
注意，$\Ob(\mathcal{J}_U)$ 非空，因为 $\{\text{id}_U\}$ 是其中一个对象。
由上述讨论，构造
$\mathcal{U} \mapsto H^0(\mathcal{U}, \mathcal{F})$ 是
$\mathcal{J}_U$ 上的反变函子。定义
$$
\mathcal{F}^{+}(U)
=
\colim_{\mathcal{J}_U^{opp}}
H^0(\mathcal{U}, \mathcal{F})
$$
关于极限和余极限的讨论，见 Categories，第
\ref{categories-section-limits} 节。我们指出，稍后将看到
$\mathcal{F}^{+}(U)$ 是 $\mathcal{F}$ 在 $U$ 上的第零 {\v C}ech 上同调。

\medskip\noindent
在进一步讨论这个余极限的结构之前，先把集合族
$\mathcal{F}^{+}(U)$，$U \in \Ob(\mathcal{C})$，变成一个预层。具体地，
设 $V \to U$ 是 $\mathcal{C}$ 的态射。由位点公理，存在函子\footnote{这一
构造实际上涉及对纤维积 $U_i \times_U V$ 的选择，因而涉及选择公理。所得映射
与所作选择无关；见下文。}
$$
\mathcal{J}_U
\longrightarrow
\mathcal{J}_V, \quad
\{U_i \to U\}
\longmapsto
\{U_i \times_U V \to V\}.
$$
注意，投影映射给出函子性的覆盖态射
$\{U_i \times_U V \to V\} \to \{U_i \to U\}$，从而由上述构造给出
函子性的集合映射
$H^0(\{U_i \to U\}, \mathcal{F}) \to
H^0(\{U_i \times_U V \to V\}, \mathcal{F})$。换言之，存在从函子
$H^0(-, \mathcal{F}) : \mathcal{J}_U^{opp} \to \textit{Sets}$ 到复合
$\mathcal{J}_U^{opp} \to \mathcal{J}_V^{opp}
\xrightarrow{H^0(-, \mathcal{F})} \textit{Sets}$ 的函子变换。因此，由余极限
的一般性质，得到典范映射 $\mathcal{F}^+(U) \to \mathcal{F}^+(V)$。
用上面对集合 $\mathcal{F}^+(U)$ 的描述来说，它把由
$s = (s_i) \in H^0(\{U_i \to U\}, \mathcal{F})$ 所代表的元素送到由
$(s_i|_{V \times_U U_i})
\in H^0(\{U_i \times_U V \to V\}, \mathcal{F})$ 所代表的元素。

\begin{lemma}
\label{sites-lemma-plus-presheaf}
上述构造定义一个预层 $\mathcal{F}^+$，并带有一个典范的预层映射
$\mathcal{F} \to \mathcal{F}^+$。
\end{lemma}

\begin{proof}
只须证明：给定态射 $W \to V \to U$，复合
$\mathcal{F}^+(U) \to \mathcal{F}^+(V) \to \mathcal{F}^+(W)$ 等于映射
$\mathcal{F}^+(U) \to \mathcal{F}^+(W)$。可以直接证明这一点：验证对于给定
的覆盖 $\{U_i \to U\}$ 和
$s = (s_i) \in H^0(\{U_i \to U\}, \mathcal{F})$，存在典范同构
$W \times_U U_i \cong W \times_V (V \times_U U_i)$，并且经此同构，
$s_i|_{W \times_U U_i}$ 对应于
$(s_i|_{V \times_U U_i})|_{W \times_V (V \times_U U_i)}$。
\end{proof}

\noindent
更间接地，引理 \ref{sites-lemma-independent-refinement} 的结果表明，可以沿从
$W$ 的任一覆盖到 $\{U_i \to U\}$ 的任意态射拉回上述元素 $s$，而所得总是
$\mathcal{F}^+(W)$ 中的同一元素。

\begin{lemma}
\label{sites-lemma-plus-functorial}
对应 $\mathcal{F} \mapsto (\mathcal{F} \to \mathcal{F}^+)$ 是一个函子。
\end{lemma}

\begin{proof}
我们不作证明，只准确陈述需要证明的内容。设
$\mathcal{F} \to \mathcal{G}$ 是预层映射。证明下图交换：
$$
\xymatrix{
\mathcal{F} \ar[r] \ar[d]
&
\mathcal{F}^{+} \ar[d]
\\
\mathcal{G} \ar[r]
&
\mathcal{G}^{+}
}
$$
\end{proof}

\noindent
下面两个引理表明，上述余极限是定向集上的余极限。

\begin{lemma}
\label{sites-lemma-common-refinement}
给定位点 $\mathcal{C}$ 的某个对象 $U$ 的一对覆盖
$\{U_i \to U\}$ 和 $\{V_j \to U\}$，存在一个共同加细的覆盖。
\end{lemma}

\begin{proof}
由于 $\mathcal{C}$ 是位点，对每个 $i$，族
$\{V_j \times_U U_i \to U_i\}_j$ 都是覆盖。于是，另一个公理蕴含
$\{V_j \times_U U_i \to U\}_{i, j}$ 是 $U$ 的覆盖。显然，这个覆盖加细
两个给定覆盖。
\end{proof}

\begin{lemma}
\label{sites-lemma-independent-refinement}
覆盖的任意两个态射 $f, g: \mathcal{U} \to \mathcal{V}$，若诱导同一个态射
$U \to V$，则诱导同一个映射
$H^0(\mathcal{V}, \mathcal{F}) \to  H^0(\mathcal{U}, \mathcal{F})$。
\end{lemma}

\begin{proof}
设 $\mathcal{U} = \{U_i \to U\}_{i\in I}$，
$\mathcal{V} = \{V_j \to V\}_{j\in J}$。态射 $f$ 由映射
$U\to V$、映射 $\alpha : I \to J$ 以及映射
$f_i : U_i \to V_{\alpha(i)}$ 构成。类似地，$g$ 决定映射
$\beta : I \to J$ 和映射 $g_i : U_i \to V_{\beta(i)}$。由于 $f$ 与 $g$
诱导同一个映射 $U\to V$，图表
$$
\xymatrix{
&
V_{\alpha(i)} \ar[dr]
\\
U_i \ar[ur]^{f_i} \ar[dr]_{g_i}
&
&
V
\\
&
V_{\beta(i)} \ar[ur]
}
$$
对每个 $i\in I$ 都交换。因此，$f$ 与 $g$ 通过纤维积分解：
$$
\xymatrix{
&
V_{\alpha(i)}
\\
U_i \ar[r]^-\varphi \ar[ur]^{f_i} \ar[dr]_{g_i}
&
V_{\alpha(i)} \times_V V_{\beta(i)} \ar[u]_{\text{pr}_1} \ar[d]^{\text{pr}_2}
\\
&
V_{\beta(i)}.
}
$$
现在取 $s = (s_j)_j \in H^0(\mathcal{V}, \mathcal{F})$。则对所有
$i\in I$：
$$
(f^*s)_i
=
f_i^*(s_{\alpha(i)})
=
\varphi^*\text{pr}_1^*(s_{\alpha(i)})
=
\varphi^*\text{pr}_2^*(s_{\beta(i)})
=
g_i^*(s_{\beta(i)})
=
(g^*s)_i,
$$
其中中间的等号来自 $H^0(\mathcal{V}, \mathcal{F})$ 的定义。这表明由
$f$ 与 $g$ 诱导的映射
$H^0(\mathcal{V}, \mathcal{F}) \to H^0(\mathcal{U}, \mathcal{F})$ 相等。
\end{proof}

\begin{remark}
\label{sites-remark-both-refine-same-H0}
特别地，这个引理表明，若 $\mathcal{U}$ 是 $\mathcal{V}$ 的加细，并且
$\mathcal{V}$ 是 $\mathcal{U}$ 的加细，则存在典范等同
$H^0(\mathcal{U}, \mathcal{F}) =
H^0(\mathcal{V}, \mathcal{F})$。
\end{remark}

\noindent
由这两个引理以及 $\mathcal{J}_U$ 非空这一事实，可知图表
$H^0(-, \mathcal{F}) : \mathcal{J}_U^{opp}
\to \textit{Sets}$ 是滤过的；见 Categories，定义
\ref{categories-definition-directed}。因此，由 Categories，第
\ref{categories-section-directed-colimits} 节，余极限
$\mathcal{F}^{+}(U)$ 可以按如下直接方式描述。集合 $\mathcal{F}^{+}(U)$
中的每个元素，都来自某个 $U$ 的覆盖 $\mathcal{U}$ 的一个元素
$s \in H^0(\mathcal{U}, \mathcal{F})$。再给定一个元素
$s' \in H^0(\mathcal{U}', \mathcal{F})$，则 $s$ 与 $s'$ 决定余极限中的
同一元素，当且仅当存在 $U$ 的一个覆盖 $\mathcal{V}$，以及加细
$f : \mathcal{V} \to \mathcal{U}$ 和
$f' : \mathcal{V} \to \mathcal{U}'$，使得在
$H^0(\mathcal{V}, \mathcal{F})$ 中有 $f^*s = (f')^*s'$。由于平凡覆盖
$\{\text{id}_U\}$ 是 $\mathcal{J}_U$ 的对象，得到典范映射
$\mathcal{F}(U) \to \mathcal{F}^+(U)$。

\begin{lemma}
\label{sites-lemma-plus-surjective}
映射 $\theta : \mathcal{F} \to \mathcal{F}^+$ 具有如下性质：对
$\mathcal{C}$ 的每个对象 $U$ 及每个截面 $s \in \mathcal{F}^+(U)$，存在
覆盖 $\{U_i \to U\}$，使得 $s|_{U_i}$ 属于
$\theta : \mathcal{F}(U_i) \to \mathcal{F}^{+}(U_i)$ 的像。
\end{lemma}

\begin{proof}
事实上，取一个覆盖 $\{U_i \to U\}$，使 $s$ 来自元素
$(s_i) \in H^0(\{U_i \to U\}, \mathcal{F})$。由引理
\ref{sites-lemma-independent-refinement}，为计算 $s$ 拉回到
$\mathcal{F}^+(U_i)$ 的元素，可以使用覆盖 $\{U_i \to U_i\}$ 及显然的覆盖
态射 $\{U_i \to U_i\} \to \{U_i \to U\}$。而使用这个覆盖时，$s$ 在
$U_i$ 上的限制恰好是 $\theta(s_i)$。
\end{proof}

\begin{definition}
\label{sites-definition-separated}
若对所有覆盖 $\{U_i \rightarrow U\}$，映射
$\mathcal{F}(U) \to \prod \mathcal{F}(U_i)$ 都是单射，则称位点
$\mathcal{C}$ 上的集合预层 $\mathcal{F}$ 是{\it 分离的}。
% P01-ERRATA: P01-E0064 -- the frozen source says "for all coverings of {family}".
\end{definition}

\begin{theorem}
\label{sites-theorem-plus}
对上述 $\mathcal{F}$，有：
\begin{enumerate}
\item
\label{sites-item-sep}
预层 $\mathcal{F}^+$ 是分离的。
\item
\label{sites-item-sheaf}
若 $\mathcal{F}$ 是分离的，则 $\mathcal{F}^+$ 是层，并且预层映射
$\mathcal{F} \to \mathcal{F}^+$ 是单射。
\item
\label{sites-item-plus-iso}
若 $\mathcal{F}$ 是层，则 $\mathcal{F} \to \mathcal{F}^+$ 是同构。
\item
\label{sites-item-plusplus}
预层 $\mathcal{F}^{++}$ 总是层。
\end{enumerate}
\end{theorem}

\begin{proof}
证明 (\ref{sites-item-sep})。假设 $s, s' \in \mathcal{F}^+(U)$，并假设存在某个
覆盖 $\{U_i \to U\}$，使得对所有 $i$，都有
$s|_{U_i} = s'|_{U_i}$。现在有 $U$ 的三个覆盖：上述覆盖
$\{U_i \to U\}$，如引理 \ref{sites-lemma-plus-surjective} 中用于 $s$ 的覆盖
$\mathcal{U}$，以及用于 $s'$ 的类似覆盖 $\mathcal{U}'$。由引理
\ref{sites-lemma-common-refinement}，可以找到一个共同加细，记作
$\{W_j \to U\}$。这意味着存在
$s_j, s'_j \in \mathcal{F}(W_j)$，使得
$s|_{W_j} = \theta(s_j)$，对 $s'|_{W_j}$ 亦然，并且
$\theta(s_j) = \theta(s'_j)$。最后一个等式意味着存在某个覆盖
$\{W_{jk} \to W_j\}$，使得
$s_j|_{W_{jk}} = s'_j|_{W_{jk}}$。由于 $\{W_{jk} \to U\}$ 是覆盖，
可知 $s, s'$ 映到 $H^0(\{W_{jk} \to U\}, \mathcal{F})$ 的同一元素，
正合所需。

\medskip\noindent
证明 (\ref{sites-item-sheaf})。$\mathcal{F} \to \mathcal{F}^+$ 显然是单射，因为
所有映射 $\mathcal{F}(U) \to H^0(\mathcal{U}, \mathcal{F})$ 都是单射。
同样显然，若 $\mathcal{U} \to \mathcal{U}'$ 是加细，则
$H^0(\mathcal{U}', \mathcal{F})
\to H^0(\mathcal{U}, \mathcal{F})$ 是单射。现在假设
$\{U_i \to U\}$ 是覆盖，并设 $(s_i)$ 是一族
$\mathcal{F}^+(U_i)$ 的元素，满足对所有 $i, i' \in I$ 的层条件
$s_i|_{U_i \times_U U_{i'}} = s_{i'}|_{U_i \times_U U_{i'}}$。
选择如引理 \ref{sites-lemma-plus-surjective} 中的覆盖
$\{U_{ij} \to U_i\}$，使得 $s_i|_{U_{ij}}$ 是唯一元素
$s_{ij} \in \mathcal{F}(U_{ij})$ 的像。层条件蕴含 $s_{ij}$ 与
$s_{i'j'}$ 在 $U_{ij} \times_U U_{i'j'}$ 上相等，因为该对象映到
$U_i \times_U U_{i'}$，而我们在那里已有等式。因此
$(s_{ij}) \in H^0(\{U_{ij} \to U\}, \mathcal{F})$ 产生一个元素
$s \in \mathcal{F}^+(U)$。留给读者验证 $s|_{U_i} = s_i$。

\medskip\noindent
证明 (\ref{sites-item-plus-iso})。这由定义立即得到，因为层性质恰好说明，对
$U$ 的每个覆盖 $\mathcal{U}$，每个映射
$\mathcal{F} \to H^0(\mathcal{U}, \mathcal{F})$ 都是双射。
% P01-ERRATA: P01-E0065 -- the frozen source omits (U) from the domain of this map.

\medskip\noindent
陈述 (\ref{sites-item-plusplus}) 现在是显然的。
\end{proof}

\begin{definition}
\label{sites-definition-associated-sheaf}
设 $\mathcal{C}$ 为位点，$\mathcal{F}$ 是 $\mathcal{C}$ 上的集合预层。
层 $\mathcal{F}^\# := \mathcal{F}^{++}$ 连同典范映射
$\mathcal{F} \to \mathcal{F}^\#$，称为{\it 与 $\mathcal{F}$ 相伴的层}。
\end{definition}

\begin{proposition}
\label{sites-proposition-sheafification-adjoint}
典范映射 $\mathcal{F} \to \mathcal{F}^\#$ 具有如下泛性质：对任意映射
$\mathcal{F} \to \mathcal{G}$，其中 $\mathcal{G}$ 是集合层，存在唯一映射
$\mathcal{F}^\# \to \mathcal{G}$，使得复合
$\mathcal{F} \to \mathcal{F}^\# \to \mathcal{G}$ 等于给定映射。
\end{proposition}

\begin{proof}
由引理 \ref{sites-lemma-plus-functorial}，得到交换图表
$$
\xymatrix{
\mathcal{F} \ar[r] \ar[d]
&
\mathcal{F}^{+} \ar[r] \ar[d]
&
\mathcal{F}^{++} \ar[d]
\\
\mathcal{G} \ar[r]
&
\mathcal{G}^{+} \ar[r]
&
\mathcal{G}^{++}
}
$$
而由定理 \ref{sites-theorem-plus}，下方的水平映射都是同构。唯一性来自引理
\ref{sites-lemma-plus-surjective}；该引理说明，$\mathcal{F}^\#$ 的每个截面在局部
都来自 $\mathcal{F}$ 的截面。
\end{proof}

\noindent
由此结果容易看出，函子 $\mathcal{F}
\mapsto (\mathcal{F} \to \mathcal{F}^\#)$ 在函子的唯一同构意义下唯一。
事实上，暂用 $i : \Sh(\mathcal{C}) \to \textit{PSh}(\mathcal{C})$ 表示
包含函子。上述结果实际上说明
$$
\Mor_{\textit{PSh}(\mathcal{C})}(\mathcal{F}, i(\mathcal{G}))
=
\Mor_{\Sh(\mathcal{C})}(\mathcal{F}^\#, \mathcal{G}).
$$
换言之，层化函子是包含函子 $i$ 的左伴随。

\begin{example}
\label{sites-example-constant-sheaf}
设 $\mathcal{C}$ 为位点，$E$ 为集合。{\it 取常值 $E$ 的常值层
$\underline{E}$}，是取常值 $E$ 的常值函子所给预层的层化。它由下列规则刻画：
对 $\mathcal{C}$ 上的每个层 $\mathcal{F}$，
$\Mor_{\Sh(\mathcal{C})}(\underline{E}, \mathcal{F}) =
\Mor_{\textit{Sets}}(E, \Gamma(\mathcal{C}, \mathcal{F}))$；这里的记号来自定义
\ref{sites-definition-global-sections}。
\end{example}

\noindent
以几个引理结束本节。

\begin{lemma}
\label{sites-lemma-colimit-sheaves}
\begin{slogan}
层范畴中的余极限等于预层范畴中余极限的层化。
\end{slogan}
设 $\mathcal{F} : \mathcal{I} \to \Sh(\mathcal{C})$ 是一个图表。那么
$\colim_\mathcal{I} \mathcal{F}$ 存在，并且是预层范畴中余极限的层化。
\end{lemma}

\begin{proof}
由于层化函子是左伴随，它与所有余极限交换；见 Categories，引理
\ref{categories-lemma-adjoint-exact}。所以，由于
$\textit{PSh}(\mathcal{C})$ 有余极限，可知 $\Sh(\mathcal{C})$ 也有余极限
（它们是预层中余极限的层化）。
\end{proof}

\begin{lemma}
\label{sites-lemma-sheafification-exact}
函子 $\textit{PSh}(\mathcal{C}) \to \Sh(\mathcal{C})$，
$\mathcal{F} \mapsto \mathcal{F}^\#$，是正合的。
\end{lemma}

\begin{proof}
由于它是左伴随，故右正合；见 Categories，引理
\ref{categories-lemma-exact-adjoint}。另一方面，由引理
\ref{sites-lemma-common-refinement} 和引理 \ref{sites-lemma-independent-refinement}，
$\mathcal{F}^+$ 的构造中的余极限实际上取在定向集
$\Ob(\mathcal{J}_U)$ 上，其中
$\mathcal{U} \geq \mathcal{U}'$ 当且仅当 $\mathcal{U}$ 是
$\mathcal{U}'$ 的加细。因此，由 Categories，引理
\ref{categories-lemma-directed-commutes}，可知
$\mathcal{F} \to \mathcal{F}^+$ 与有限极限交换（视为从预层到预层的函子）。
再用引理 \ref{sites-lemma-limit-sheaf} 即得结论。
\end{proof}

\begin{lemma}
\label{sites-lemma-sections-sheafification}
设 $\mathcal{C}$ 为位点，$\mathcal{F}$ 为 $\mathcal{C}$ 上的集合预层。记
$\theta^2 : \mathcal{F} \to \mathcal{F}^\#$ 为从 $\mathcal{F}$ 到其层化的
典范映射。设 $U$ 为 $\mathcal{C}$ 的对象，并设
$s \in \mathcal{F}^\#(U)$。存在覆盖 $\{U_i \to U\}$ 和截面
$s_i \in \mathcal{F}(U_i)$，使得
\begin{enumerate}
\item $s|_{U_i} = \theta^2(s_i)$；
\item 对每个 $i, j$，存在 $\mathcal{C}$ 的一个覆盖
$\{U_{ijk} \to U_i \times_U U_j\}$，使得 $s_i$ 与 $s_j$ 拉回到每个
$U_{ijk}$ 后相等。
\end{enumerate}
反过来，给定任意覆盖 $\{U_i \to U\}$ 和满足 (2) 的元素
$s_i \in \mathcal{F}(U_i)$，则存在唯一的截面
$s \in \mathcal{F}^\#(U)$，使得 (1) 成立。
\end{lemma}

\begin{proof}
从略。
\end{proof}

\begin{lemma}
\label{sites-lemma-isomorphism-sheafifications}
设 $\mathcal{C}$ 为位点。设 $\mathcal{F} \to \mathcal{G}$ 是
$\mathcal{C}$ 上集合预层的映射。记 $\mathcal{B}$ 为所有满足
$\mathcal{F}(U) \to \mathcal{G}(U)$ 是双射的
$U \in \Ob(\mathcal{C})$ 所成的集合。若 $\mathcal{C}$ 的每个对象都有一个
由 $\mathcal{B}$ 的元素组成的覆盖，则
$\mathcal{F}^\# \to \mathcal{G}^\#$ 是同构。
\end{lemma}

\begin{proof}
取 $U \in \Ob(\mathcal{C})$。我们证明
$\mathcal{F}^\#(U) \to \mathcal{G}^\#(U)$ 是满射。为此，我们将使用引理
\ref{sites-lemma-sections-sheafification}，不再另行说明。对任意
$s \in \mathcal{G}^\#(U)$，存在覆盖 $\{U_i \to U\}$ 和截面
$s_i \in \mathcal{G}(U_i)$，使得
\begin{enumerate}
\item $s|_{U_i}$ 是 $s_i$ 经 $\mathcal{G} \to \mathcal{G}^\#$ 的像；
\item 对每个 $i, j$，存在 $\mathcal{D}$ 的一个覆盖
$\{U_{ijk} \to U_i \times_U U_j\}$，使得 $s_i$ 与 $s_j$ 拉回到每个
$U_{ijk}$ 后相等。
% P01-ERRATA: P01-E0066 -- the frozen source names undefined site D instead of C.
\end{enumerate}
由假设，对每个 $i$，可以选择一个覆盖 $\{U_{i, a} \to U_i\}$，其中
$U_{i, a} \in \mathcal{B}$。于是 $\{U_{i, a} \to U\}$ 是覆盖。用
$s_{i, a}$ 表示 $s_i$ 在 $\mathcal{G}(U_{i, a})$ 中的像，可知
$s_{i, a}$ 与 $s_{j, b}$ 拉回到覆盖
$$
\{(U_{i, a} \times_U U_{j, b}) \times_{U_i \times_U U_j} U_{ijk}
\to U_{i, a} \times_U U_{j, b}\}
$$
的成员上后相等。因此，可以假设 $U_i \in \mathcal{B}$。重复此论证，还可
以假设对所有 $i, j, k$ 都有 $U_{ijk} \in \mathcal{B}$（细节从略）。于是，
由于引理陈述中对 $\mathcal{B}$ 的定义保证
$\mathcal{F}(U_i) \to \mathcal{G}(U_i)$ 是双射，得到唯一的
$t_i \in \mathcal{F}(U_i)$ 映到 $s_i$。$t_i$ 与 $t_j$ 拉回到每个
$U_{ijk}$ 后在 $\mathcal{F}(U_{ijk})$ 中相等，因为
$U_{ijk} \in \mathcal{B}$，并且 $t_i$ 与 $t_j$ 已有相等性。
% P01-ERRATA: P01-E0067 -- the frozen rationale is circular; the known agreement is for s_i,s_j.
于是该引理表明，存在唯一的 $t \in \mathcal{F}^\#(U)$，使得
$t|_{U_i}$ 是 $t_i$ 的像。由此 $t$ 映到 $s$。我们略去
$\mathcal{F}^\#(U) \to \mathcal{G}^\#(U)$ 为单射的证明。
\end{proof}









\section{层的单射与满射}
\label{sites-section-sheaves-injective}

\begin{definition}
\label{sites-definition-sheaves-injective-surjective}
设 $\mathcal{C}$ 为位点，并设 $\varphi : \mathcal{F}
\to \mathcal{G}$ 是集合层的映射。
\begin{enumerate}
\item 若对 $\mathcal{C}$ 的每个对象 $U$，映射
$\varphi : \mathcal{F}(U) \to \mathcal{G}(U)$ 都是单射，则称
$\varphi$ 是{\it 单射的}。
\item 若对 $\mathcal{C}$ 的每个对象 $U$ 和每个截面
$s\in \mathcal{G}(U)$，都存在一个覆盖 $\{U_i \to U\}$，使得对所有 $i$，
限制 $s|_{U_i}$ 都属于
$\varphi : \mathcal{F}(U_i) \to \mathcal{G}(U_i)$ 的像，则称
$\varphi$ 是{\it 满射的}。
\end{enumerate}
\end{definition}

\begin{lemma}
\label{sites-lemma-mono-epi-sheaves}
上面定义的单射（相应地，满射）恰好是范畴 $\Sh(\mathcal{C})$ 中的单态射
（相应地，满态射）。层的映射是同构，当且仅当它既是单射又是满射。
\end{lemma}

\begin{proof}
从略。
\end{proof}

\begin{lemma}
\label{sites-lemma-coequalizer-surjection}
设 $\mathcal{C}$ 为位点。设 $\mathcal{F} \to \mathcal{G}$ 是集合层的满射。
则图表
$$
\xymatrix{
\mathcal{F} \times_\mathcal{G} \mathcal{F}
\ar@<1ex>[r] \ar@<-1ex>[r]
&
\mathcal{F} \ar[r]
&
\mathcal{G}}
$$
把 $\mathcal{G}$ 表示为余等化子。
\end{lemma}

\begin{proof}
设 $\mathcal{H}$ 为集合层，并设 $\varphi : \mathcal{F} \to \mathcal{H}$
是等化两个映射
$\mathcal{F} \times_\mathcal{G} \mathcal{F} \to \mathcal{F}$ 的层映射。
设 $\mathcal{G}' \subset \mathcal{G}$ 为映射
$\mathcal{F} \to \mathcal{G}$ 的预层像。由于积
$\mathcal{F} \times_\mathcal{G} \mathcal{F}$ 可以在预层范畴中计算，可知它
等于预层积 $\mathcal{F} \times_{\mathcal{G}'} \mathcal{F}$。因此
$\varphi$ 诱导唯一的预层映射
$\psi' : \mathcal{G}' \to \mathcal{H}$。由引理
\ref{sites-lemma-mono-epi-sheaves}，$\mathcal{G}$ 是 $\mathcal{G}'$ 的层化，故
$\psi'$ 唯一延拓为层映射 $\psi : \mathcal{G} \to \mathcal{H}$。我们略去
验证 $\varphi$ 等于 $\psi$ 与给定映射的复合。
\end{proof}














\section{可表层}
\label{sites-section-representable-sheaves}

\noindent
设 $\mathcal{C}$ 为范畴。典范拓扑是使所有可表预层都成为层的最细拓扑
（它在定义 \ref{sites-definition-canonical-topology} 中有形式定义，但这里不需要）。
该拓扑未必总是 $\mathcal{C}$ 上某个位点结构的相伴拓扑。若
$\mathcal{C}$ 有纤维积，我们将给出生成此拓扑的一类覆盖。先给出下面的一般
定义。

\begin{definition}
\label{sites-definition-universal-effective-epimorphisms}
设 $\mathcal{C}$ 为范畴。若族 $\{U_i \to U\}_{i \in I}$ 中所有态射
$U_i \to U$ 都可表（见 Categories，定义
\ref{categories-definition-representable-morphism}），并且对任意
$X\in \Ob(\mathcal{C})$，序列
$$
\xymatrix{
\Mor_\mathcal{C}(U, X) \ar[r]
&
\prod\nolimits_{i \in I} \Mor_\mathcal{C}(U_i, X)
\ar@<1ex>[r] \ar@<-1ex>[r]
&
\prod\nolimits_{(i, j) \in I^2} \Mor_\mathcal{C}(U_i \times_U U_j, X)
}
$$
是等化子图表，则称该族是一个{\it 有效满态射}。若对任意态射
$V \to U$，基变换 $\{U_i \times_U V \to V\}$ 都是有效满态射，则称族
$\{U_i \to U\}$ 是一个{\it 泛有效满态射}。
\end{definition}

\noindent
泛有效满态射族所成的类满足定义 \ref{sites-definition-site} 的公理。若
$\mathcal{C}$ 有纤维积，则相伴拓扑就是典范拓扑。（在这种情况下，要得到一个
位点，可按 Sets，引理 \ref{sets-lemma-coverings-site} 论证。）

\medskip\noindent
反过来，假设 $\mathcal{C}$ 是一个位点，且所有可表预层都是层。于是显然，
所有覆盖都是泛有效满态射。因此，在位点语境中，下面的定义才是“正确”的。

\begin{definition}
\label{sites-definition-weaker-than-canonical}
若 $\mathcal{C}$ 的所有覆盖都是泛有效满态射，则称位点 $\mathcal{C}$ 上的
拓扑{\it 弱于典范拓扑}，或称该拓扑是{\it 次典范的}。
\end{definition}

\noindent
可表层是同时为层的可表预层。由于当拓扑不是次典范时最好避免这一术语，我们
只在次典范情形下给出其形式定义。

\begin{definition}
\label{sites-definition-representable-sheaf}
设 $\mathcal{C}$ 为拓扑次典范的位点。Yoneda 嵌入 $h$（见 Categories，
第 \ref{categories-section-opposite} 节）把 $\mathcal{C}$ 表示为
$\mathcal{C}$ 的层范畴的全子范畴。在这种情况下，我们称形如 $h_U$ 的层为
$\mathcal{C}$ 上的{\it 可表层}，其中 $U \in \Ob(\mathcal{C})$。
记号：有时，与 $U$ 相伴的可表层 $h_U$ 记作 {\it $\underline{U}$}。
\end{definition}

\noindent
注意，在该定义的情形下，对每个层 $\mathcal{F}$，都有
$$
\Mor_{\Sh(\mathcal{C})}(h_U, \mathcal{F}) = \mathcal{F}(U)
$$
因为该等式对预层成立；见
(\ref{sites-equation-map-representable-into-presheaf})。一般而言，预层 $h_U$ 不是
层；要得到层，必须对它们作层化。在这种情况下，仍有
\begin{equation}
\label{sites-equation-map-representable-into-sheaf}
\Mor_{\Sh(\mathcal{C})}(h_U^\#, \mathcal{F}) =
\Mor_{\textit{PSh}(\mathcal{C})}(h_U, \mathcal{F}) =
\mathcal{F}(U)
\end{equation}
对每个层 $\mathcal{F}$ 成立。事实上，第一个等式来自 $\#$ 的伴随性质，
第二个等式就是
(\ref{sites-equation-map-representable-into-presheaf})。

\begin{lemma}
\label{sites-lemma-covering-surjective-after-sheafification}
\begin{slogan}
覆盖经层化后成为满射。
\end{slogan}
设 $\mathcal{C}$ 为位点。若 $\{U_i \to U\}_{i \in I}$ 是位点
$\mathcal{C}$ 的一个覆盖，则集合预层的态射
$$
\coprod\nolimits_{i \in I} h_{U_i} \to h_U
$$
经层化后成为满射。
\end{lemma}

\begin{proof}
由上述引理 \ref{sites-lemma-mono-epi-sheaves}，需要证明
$\coprod\nolimits_{i \in I} h_{U_i}^\# \to h_U^\#$ 是满态射。设
$\mathcal{F}$ 为集合层。态射 $h_U^\# \to \mathcal{F}$ 对应于截面
$s \in \mathcal{F}(U)$。因此，映射
$\Mor(h_U^\#, \mathcal{F})
\to \prod_i \Mor(h_{U_i}^\#, \mathcal{F})$ 的单射性直接来自
$\mathcal{F}$ 的层性质。
\end{proof}

\noindent
下一个引理说明：在拓扑弱于典范拓扑的情况下，每个层在某种意义下都由可表层
组成。

\begin{lemma}
\label{sites-lemma-sheaf-coequalizer-representable}
设 $\mathcal{C}$ 为位点。设 $E \subset \Ob(\mathcal{C})$ 是一个子集，使得
$\mathcal{C}$ 的每个对象都有由 $E$ 的元素组成的覆盖。设 $\mathcal{F}$ 为
集合层。存在集合层的图表
$$
\xymatrix{
\mathcal{F}_1 \ar@<1ex>[r] \ar@<-1ex>[r] &
\mathcal{F}_0 \ar[r] &
\mathcal{F}
}
$$
把 $\mathcal{F}$ 表示为余等化子，并且对 $i = 0, 1$，$\mathcal{F}_i$ 都是
形如 $h_U^\#$（其中 $U \in E$）的层的余积。
\end{lemma}

\begin{proof}
首先证明存在所需类型的满态射 $\mathcal{F}_0 \to \mathcal{F}$。事实上，只须取
$$
\mathcal{F}_0 =
\coprod\nolimits_{U \in E, s \in \mathcal{F}(U)}
(h_U)^\# \longrightarrow \mathcal{F}
$$
这里，箭头限制在与 $(U, s)$ 对应的分量上时，把元素
$\text{id}_U \in h_U^\#(U)$ 映到截面 $s \in \mathcal{F}(U)$。由引理
\ref{sites-lemma-mono-epi-sheaves} 以及对 $E$ 的条件，这是一条满态射。为构造
$\mathcal{F}_1$，先令
$\mathcal{G} = \mathcal{F}_0 \times_\mathcal{F} \mathcal{F}_0$，再像上面那样
构造满态射 $\mathcal{F}_1 \to \mathcal{G}$。见引理
\ref{sites-lemma-coequalizer-surjection}。
\end{proof}

\begin{lemma}
\label{sites-lemma-colimit-representable}
设 $\mathcal{C}$ 为位点，$\mathcal{F}$ 为 $\mathcal{C}$ 上的集合层。则存在
图表 $\mathcal{I} \to \mathcal{C}$，$i \mapsto U_i$，使得
$$
\mathcal{F} = \colim_{i \in \mathcal{I}} h_{U_i}^\#
$$
此外，若 $E \subset \Ob(\mathcal{C})$ 是一个子集，使 $\mathcal{C}$ 的每个
对象都有由 $E$ 的元素组成的覆盖，则可假设对所有
$i \in \Ob(\mathcal{I})$，$U_i$ 都是 $E$ 的元素。
\end{lemma}

\begin{proof}
令 $\mathcal{I}$ 为如下范畴：其对象是二元组 $(U, s)$，其中
$U \in \Ob(\mathcal{C})$ 且 $s \in \mathcal{F}(U)$；态射
$(U, s) \to (U', s')$ 是 $\mathcal{C}$ 中满足 $f^*s' = s$ 的
态射 $f : U \to U'$。对 $\mathcal{I}$ 的每个对象 $(U, s)$，由 Yoneda 引理，
元素 $s$ 定义预层映射 $s : h_U \to \mathcal{F}$，从而由层化的泛性质得到
映射 $h_U^\# \to \mathcal{F}$。立即可见，这些映射与范畴
$\mathcal{I}$ 中的态射相容。因此，得到预层映射
$\colim_{(U, s)} h_U \to \mathcal{F}$（其中余极限在预层范畴中取），以及层
映射 $\colim_{(U, s)} (h_U)^\# \to \mathcal{F}$（其中余极限在层范畴中取）。
由于层化是包含函子
$\Sh(\mathcal{C}) \to \textit{PSh}(\mathcal{C})$ 的左伴随（命题
\ref{sites-proposition-sheafification-adjoint}），由 Categories，引理
\ref{categories-lemma-adjoint-exact}，有
$\colim (h_U)^\# = (\colim h_U)^\#$。因此只须证明
$\colim_{(U, s)} h_U \to \mathcal{F}$ 是预层同构。为此，证明对
$\mathcal{C}$ 的每个对象 $X$，映射
$\colim_{(U, s)} h_U(X) \to \mathcal{F}(X)$ 是双射。事实上，它有如下逆映射：
把元素 $t \in \mathcal{F}(X)$ 送到集合
$\colim_{(U, s)} h_U(X)$ 中由 $(X, t)$ 和
$\text{id}_X \in h_X(X)$ 对应的元素。

\medskip\noindent
最后一个陈述的证明从略。
\end{proof}





\section{连续函子}
\label{sites-section-continuous-functors}

\begin{definition}
\label{sites-definition-continuous}
设 $\mathcal{C}$ 和 $\mathcal{D}$ 为位点。若对每个
$\{V_i \to V\}_{i\in I} \in \text{Cov}(\mathcal{C})$，函子
$u : \mathcal{C} \to \mathcal{D}$ 都满足下列条件，则称它是
{\it 连续的}：
\begin{enumerate}
\item $\{u(V_i) \to u(V)\}_{i\in I}$ 属于 $\text{Cov}(\mathcal{D})$；
\item 对 $\mathcal{C}$ 中任意态射 $T \to V$，态射
$u(T \times_V V_i) \to u(T) \times_{u(V)} u(V_i)$ 是同构。
\end{enumerate}
\end{definition}

\noindent
回忆一下，给定上述函子 $u$ 和 $\mathcal{D}$ 上的集合预层
$\mathcal{F}$，我们已把 $u^p\mathcal{F}$ 定义为预层
$\mathcal{F} \circ u$，换言之，对 $\mathcal{C}$ 的每个对象 $V$，有
$$
u^p\mathcal{F} (V) = \mathcal{F}(u(V))
$$

\begin{lemma}
\label{sites-lemma-pushforward-sheaf}
设 $\mathcal{C}$ 和 $\mathcal{D}$ 为位点。设
$u : \mathcal{C} \to \mathcal{D}$ 为连续函子。若 $\mathcal{F}$ 是
$\mathcal{D}$ 上的层，则 $u^p\mathcal{F}$ 也是层。
\end{lemma}

\begin{proof}
设 $\{V_i \to V\}$ 为覆盖。由假设，$\{u(V_i) \to u(V)\}$ 是
$\mathcal{D}$ 中的覆盖，并且
$u(V_i \times_V V_j) = u(V_i)\times_{u(V)}u(V_j)$。所以，
$u^p\mathcal{F}$ 关于覆盖 $\{V_i \to V\}$ 的层条件，恰好就是
$\mathcal{F}$ 关于覆盖 $\{u(V_i) \to u(V)\}$ 的层条件。
\end{proof}

\noindent
为避免混淆，有时把 $u^p$ 限制在集合层子范畴上所得的函子记作
$$
u^s :
\Sh(\mathcal{D})
\longrightarrow
\Sh(\mathcal{C})
$$
回忆一下，$u^p$ 有左伴随
$u_p : \textit{PSh}(\mathcal{C}) \to \textit{PSh}(\mathcal{D})$；见第
\ref{sites-section-functoriality-PSh} 节。

\begin{lemma}
\label{sites-lemma-adjoint-sheaves}
在引理 \ref{sites-lemma-pushforward-sheaf} 的情形下，函子
$u_s : \mathcal{G} \mapsto (u_p \mathcal{G})^\#$ 是 $u^s$ 的左伴随。
\end{lemma}

\begin{proof}
直接由引理 \ref{sites-lemma-adjoints-u} 和命题
\ref{sites-proposition-sheafification-adjoint} 得到。
\end{proof}

\noindent
下面是一个技术性引理。

\begin{lemma}
\label{sites-lemma-technical-up}
在引理 \ref{sites-lemma-pushforward-sheaf} 的情形下，对 $\mathcal{C}$ 上的任意预层
$\mathcal{G}$，都有
$(u_p\mathcal{G})^\# = (u_p(\mathcal{G}^\#))^\#$。
\end{lemma}

\begin{proof}
对 $\mathcal{D}$ 上的任意层 $\mathcal{F}$，有
\begin{eqnarray*}
\Mor_{\Sh(\mathcal{D})}(u_s(\mathcal{G}^\#), \mathcal{F})
& = &
\Mor_{\Sh(\mathcal{C})}(\mathcal{G}^\#, u^s\mathcal{F}) \\
& = &
\Mor_{\textit{PSh}(\mathcal{C})}(\mathcal{G}^\#, u^p\mathcal{F}) \\
& = &
\Mor_{\textit{PSh}(\mathcal{C})}(\mathcal{G}, u^p\mathcal{F}) \\
& = &
\Mor_{\textit{PSh}(\mathcal{D})}(u_p\mathcal{G}, \mathcal{F}) \\
& = &
\Mor_{\Sh(\mathcal{D})}((u_p\mathcal{G})^\#, \mathcal{F})
\end{eqnarray*}
结论由 Yoneda 引理得到。
\end{proof}

\begin{lemma}
\label{sites-lemma-pullback-representable-sheaf}
设 $u : \mathcal{C} \to \mathcal{D}$ 为位点之间的连续函子。对
$\mathcal{C}$ 的任意对象 $U$，都有 $u_sh_U^\# = h_{u(U)}^\#$。
\end{lemma}

\begin{proof}
由引理 \ref{sites-lemma-pullback-representable-presheaf} 和
\ref{sites-lemma-technical-up} 得到。
\end{proof}


\begin{remark}
\label{sites-remark-quasi-continuous}
（初读时可跳过。）设 $\mathcal{C}$ 和 $\mathcal{D}$ 为位点。我们利用自明
等价映射族的定义（见定义 \ref{sites-definition-combinatorial-tautological}），
（略微）削弱连续性的定义条件。设
$u : \mathcal{C} \to \mathcal{D}$ 为函子。若对每个
$\mathcal{V} = \{V_i \to V\}_{i\in I} \in \text{Cov}(\mathcal{C})$，
都有下列性质，则称 $u$ 是{\it 拟连续的}：
\begin{enumerate}
\item[(1')] 映射族 $\{u(V_i) \to u(V)\}_{i\in I}$ 与
$\text{Cov}(\mathcal{D})$ 的一个元素自明等价；
\item[(2)] 对 $\mathcal{C}$ 中任意态射 $T \to V$，态射
$u(T \times_V V_i) \to u(T) \times_{u(V)} u(V_i)$ 是同构。
\end{enumerate}
我们将看到，当 $u$ 拟连续时，引理 \ref{sites-lemma-pushforward-sheaf} 和
\ref{sites-lemma-adjoint-sheaves} 同样成立。

\medskip\noindent
首先注意，态射 $u(V_i) \to u(V)$ 是可表的，因为由第一个条件，它们与可表
态射同构。特别地，族
$u(\mathcal{V}) = \{u(V_i) \to u(V)\}_{i\in I}$ 对 $\mathcal{D}$ 上的任意
预层 $\mathcal{F}$ 都给出一个第零 {\v C}ech 上同调群
$H^0(u(\mathcal{V}), \mathcal{F})$。
% P01-ERRATA: P01-E0063 -- this is the same set-valued H^0/group mismatch.
设 $\mathcal{U} = \{U_j \to u(V)\}_{j \in J}$ 是
$\text{Cov}(\mathcal{D})$ 的一个元素，并且与
$\{u(V_i) \to u(V)\}_{i \in I}$ 自明等价。注意，$u(\mathcal{V})$ 是
$\mathcal{U}$ 的加细，反之亦然。因此，由备注
\ref{sites-remark-both-refine-same-H0}，有
$H^0(u(\mathcal{V}), \mathcal{F}) = H^0(\mathcal{U}, \mathcal{F})$。
特别地，若 $\mathcal{F}$ 是层，则由用第零 {\v C}ech 上同调群表述的层性质，
有 $\mathcal{F}(u(V)) = H^0(u(\mathcal{V}), \mathcal{F})$。由此可知，若
$\mathcal{F}$ 是层，则 $u^p\mathcal{F}$ 是层，因为
$H^0(\mathcal{V}, u^p\mathcal{F}) =
H^0(u(\mathcal{V}), \mathcal{F})$，而刚才已经看到，它等于
$\mathcal{F}(u(V)) = u^p\mathcal{F}(V)$。这样便证明了引理
\ref{sites-lemma-pushforward-sheaf}。引理 \ref{sites-lemma-adjoint-sheaves} 随即成立。
\end{remark}







\section{位点态射}
\label{sites-section-morphism-sites}

\begin{definition}
\label{sites-definition-morphism-sites}
设 $\mathcal{C}$ 和 $\mathcal{D}$ 为位点。一个{\it 位点态射}
$f : \mathcal{D} \to \mathcal{C}$，由连续函子
$u : \mathcal{C} \to \mathcal{D}$ 给出，并要求函子 $u_s$ 正合。
\end{definition}

\noindent
注意，函子 $u$ 的方向与态射 $f$ 的方向{\it 相反}。若
$f \leftrightarrow u$ 是位点态射，则采用记号
$f^{-1} = u_s$ 和 $f_* = u^s$。函子 $f^{-1}$ 称为{\it 逆像函子}，函子
$f_*$ 称为{\it 直像函子}。与拓扑学中一样，有下列伴随性质：
$$
\Mor_{\Sh(\mathcal{D})}(f^{-1}\mathcal{G}, \mathcal{F})
=
\Mor_{\Sh(\mathcal{C})}(\mathcal{G}, f_*\mathcal{F})
$$
这一定义的动机来自下面的例子。

\begin{example}
\label{sites-example-continuous-map}
设 $f : X  \to Y$ 是拓扑空间的连续映射。回忆一下，我们有位点
$X_{Zar}$ 和 $Y_{Zar}$；见例 \ref{sites-example-site-topological}。考虑函子
$u : Y_{Zar} \to X_{Zar}$，$V \mapsto f^{-1}(V)$。这个函子显然连续，因为
开覆盖的逆像仍是开覆盖。（实际上，这取决于如何为 $X_{Zar}$ 与 $Y_{Zar}$
选择覆盖集。但无论如何，该函子是拟连续的；见备注
\ref{sites-remark-quasi-continuous}。）容易验证，函子 $u^s$ 等于拓扑学中通常的
直像函子 $f_*$。因此，由于 $u_s$ 是伴随函子，而通常的拓扑逆像函子
$f^{-1}$ 也是伴随函子，得到典范同构 $f^{-1} \cong u_s$。由于 $f^{-1}$
正合，可知 $u_s$ 正合。故 $u$ 定义位点态射
$f : X_{Zar} \to Y_{Zar}$；我们也可以把它记作 $f$，因为已经看到函子
$u_s, u^s$ 与它们通常的概念一致。
\end{example}

\begin{example}
\label{sites-example-finer-topology}
设 $\mathcal{C}$ 为范畴。设
$$
\text{Cov}(\mathcal{C}) \supset \text{Cov}'(\mathcal{C})
$$
是两组具有固定靶的态射族集合，并且都把 $\mathcal{C}$ 变成位点。记
$\mathcal{C}_\tau$ 为对应于 $\text{Cov}(\mathcal{C})$ 的位点，记
$\mathcal{C}_{\tau'}$ 为对应于 $\text{Cov}'(\mathcal{C})$ 的位点。我们断言，
$\mathcal{C}$ 上的恒等函子定义位点态射
$$
\epsilon : \mathcal{C}_\tau \longrightarrow \mathcal{C}_{\tau'}
$$
事实上，注意 $\text{id} : \mathcal{C}_{\tau'} \to \mathcal{C}_\tau$ 连续，
因为每个 $\tau'$-覆盖都是 $\tau$-覆盖。因此，函子
$\epsilon_* = \text{id}^s$ 是底层预层上的恒等函子。故 $\epsilon_*$ 的左伴随
$\epsilon^{-1}$ 把 $\tau'$-层 $\mathcal{F}$ 送到 $\mathcal{F}$ 的
$\tau$-层化（由层化的泛性质）。$\tau'$-层的有限极限与预层的有限极限一致
（引理 \ref{sites-lemma-limit-sheaf}），而 $\tau$-层化与有限极限交换（引理
\ref{sites-lemma-sheafification-exact}）。因此，$\epsilon^{-1}$ 左正合。又因为
$\epsilon^{-1}$ 是左伴随，它也右正合（Categories，引理
\ref{categories-lemma-exact-adjoint}）。所以 $\epsilon^{-1}$ 正合，我们已经
验证了定义 \ref{sites-definition-morphism-sites} 的全部条件。
\end{example}

\begin{lemma}
\label{sites-lemma-composition-morphisms-sites}
\begin{slogan}
位点函子的复合保持位点函子的连续性。
\end{slogan}
设 $\mathcal{C}_i$，$i = 1, 2, 3$，为位点。设
$u : \mathcal{C}_2 \to \mathcal{C}_1$ 和
$v : \mathcal{C}_3 \to \mathcal{C}_2$ 是诱导位点态射的连续函子。那么函子
$u \circ v : \mathcal{C}_3 \to \mathcal{C}_1$ 连续，并定义位点态射
$\mathcal{C}_1 \to \mathcal{C}_3$。
\end{lemma}

\begin{proof}
由定义立即可知 $u \circ v$ 是连续函子。此外，显然有
$(u \circ v)^p = v^p \circ u^p$，从而
$(u \circ v)^s = v^s \circ u^s$。因此，函子 $(u \circ v)_s$ 与
$u_s \circ v_s$ 都是 $(u \circ v)^s$ 的左伴随。故
$(u \circ v)_s \cong u_s \circ v_s$，于是 $(u \circ v)_s$ 作为正合函子的
复合是正合的。
\end{proof}

\begin{definition}
\label{sites-definition-composition-morphisms-sites}
设 $\mathcal{C}_i$，$i = 1, 2, 3$，为位点。设
$f : \mathcal{C}_1 \to \mathcal{C}_2$ 和
$g : \mathcal{C}_2 \to \mathcal{C}_3$ 是位点态射，分别由连续函子
$u : \mathcal{C}_2 \to \mathcal{C}_1$ 和
$v : \mathcal{C}_3 \to \mathcal{C}_2$ 给出。{\it 复合} $g \circ f$ 是对应于
函子 $u \circ v$ 的位点态射。
\end{definition}

\noindent
在这种情况下，有 $(g \circ f)_* = g_* \circ f_*$ 及
$(g \circ f)^{-1} = f^{-1} \circ g^{-1}$（见引理
\ref{sites-lemma-composition-morphisms-sites} 的证明）。

\begin{lemma}
\label{sites-lemma-directed-morphism}
设 $\mathcal{C}$ 和 $\mathcal{D}$ 为位点，且
$u : \mathcal{C} \to \mathcal{D}$ 连续。假设第
\ref{sites-section-functoriality-PSh} 节中的所有范畴
$(\mathcal{I}_V^u)^{opp}$ 都是滤过的。那么 $u$ 定义位点态射
$\mathcal{D} \to \mathcal{C}$；换言之，$u_s$ 正合。
\end{lemma}

\begin{proof}
由于 $u_s$ 是 $u^s$ 的左伴随，$u_s$ 右正合；见 Categories，引理
\ref{categories-lemma-exact-adjoint}。因此只须证明 $u_s$ 左正合，换言之，
只须证明 $u_s$ 与有限极限交换。由于范畴
$\mathcal{I}_Y^{opp}$ 是滤过的，由 Categories，引理
\ref{categories-lemma-directed-commutes}，可知 $u_p$ 与有限极限交换（这里也
使用了 $\textit{PSh}$ 中极限的描述，见第
\ref{sites-section-limits-colimits-PSh} 节）。又因层化也与有限极限交换（引理
\ref{sites-lemma-sheafification-exact}），结合 $u_s = \# \circ u_p$ 即得结论。
\end{proof}

\begin{proposition}
\label{sites-proposition-get-morphism}
设 $\mathcal{C}$ 和 $\mathcal{D}$ 为位点。设
$u : \mathcal{C} \to \mathcal{D}$ 连续，并进一步假设：
\begin{enumerate}
\item 范畴 $\mathcal{C}$ 有终对象 $X$，且 $u(X)$ 是 $\mathcal{D}$ 的终对象；
\item 范畴 $\mathcal{C}$ 有纤维积，且 $u$ 与纤维积交换。
\end{enumerate}
那么 $u$ 定义位点态射 $\mathcal{D} \to \mathcal{C}$；换言之，$u_s$ 正合。
\end{proposition}

\begin{proof}
这由引理 \ref{sites-lemma-directed} 和 \ref{sites-lemma-directed-morphism} 得到。
\end{proof}

\begin{remark}
\label{sites-remark-explain-left-exact}
上述命题 \ref{sites-proposition-get-morphism} 的条件等价于说 $u$ 左正合，即与有限
极限交换。见 Categories，引理
\ref{categories-lemma-finite-limits-exist} 和
\ref{categories-lemma-characterize-left-exact}。用终对象和纤维积来表述似乎更
自然，因为在例子中，这种表述似乎有更多几何意义。
\end{remark}

\noindent
引理 \ref{sites-lemma-continuous-with-continuous-left-adjoint} 将提供另一种方法，
证明连续函子产生位点态射。

\begin{remark}
\label{sites-remark-quasi-continuous-morphism-sites}
（初读时可跳过。）设 $\mathcal{C}$ 和 $\mathcal{D}$ 为位点。仿照定义
\ref{sites-definition-morphism-sites}，我们说一个{\it 拟位点态射
$f : \mathcal{D} \to \mathcal{C}$}，由一个拟连续函子
$u : \mathcal{C} \to \mathcal{D}$（见备注
\ref{sites-remark-quasi-continuous}）给出，并要求 $u_s$ 正合。命题
\ref{sites-proposition-get-morphism} 在此情形下的类似命题，只须把“连续”替换为
“拟连续”，并把“态射”替换为“拟态射”。证明完全相同。
\end{remark}

\noindent
定义 \ref{sites-definition-morphism-sites} 中 $u_s$ 正合这一条件不能省略。例如，
若只假设 $u$ 连续，下面引理的结论未必成立。

\begin{lemma}
\label{sites-lemma-morphism-of-sites-covering}
设 $f : \mathcal{D} \to \mathcal{C}$ 是由函子
$u : \mathcal{C} \to \mathcal{D}$ 给出的位点态射。给定
$\mathcal{D}$ 的任意对象 $V$，存在覆盖 $\{V_j \to V\}$，使得对每个 $j$，
都存在态射 $V_j \to u(U_j)$，其中 $U_j$ 是 $\mathcal{C}$ 的某个对象。
\end{lemma}

\begin{proof}
由于 $f^{-1} = u_s$ 正合，有 $f^{-1}* = *$，其中 $*$ 表示层范畴的终对象
（例 \ref{sites-example-singleton-sheaf}）。由于 $f^{-1}* = u_s*$ 是 $u_p*$ 的
层化，可知存在覆盖 $\{V_j \to V\}$，使得 $(u_p*)(V_j)$ 非空。又因为
$(u_p*)(V_j)$ 是范畴 $\mathcal{I}^u_{V_j}$ 上的余极限，而该范畴的对象是态射
$V_j \to u(U)$，引理得证。
\end{proof}


























\section{托扑斯}
\label{sites-section-topoi}

\noindent
下面给出适合本文用途的托扑斯定义。托扑斯就是一个位点上的层范畴；要指定一个
托扑斯，只须指定该位点。托扑斯与位点之间真正的区别在于态射的定义。事实上，
有许多托扑斯态射并非来自底层位点的态射。

\begin{definition}[托扑斯]
\label{sites-definition-topos}
一个{\it 托扑斯}是某个位点 $\mathcal{C}$ 上的层范畴
$\Sh(\mathcal{C})$。
\begin{enumerate}
\item 设 $\mathcal{C}$、$\mathcal{D}$ 为位点。从
$\Sh(\mathcal{D})$ 到 $\Sh(\mathcal{C})$ 的一个{\it 托扑斯态射} $f$，
由一对函子
$f_* : \Sh(\mathcal{D}) \to \Sh(\mathcal{C})$ 和
$f^{-1} : \Sh(\mathcal{C}) \to \Sh(\mathcal{D})$ 给出，并满足：
\begin{enumerate}
\item 双函子自然地有
$$
\Mor_{\Sh(\mathcal{D})}(f^{-1}\mathcal{G}, \mathcal{F})
=
\Mor_{\Sh(\mathcal{C})}(\mathcal{G}, f_*\mathcal{F})
$$
\item 函子 $f^{-1}$ 与有限极限交换，即它左正合。
\end{enumerate}
\item 设 $\mathcal{C}$、$\mathcal{D}$、$\mathcal{E}$ 为位点。给定托扑斯态射
$f :\Sh(\mathcal{D}) \to \Sh(\mathcal{C})$ 和
$g :\Sh(\mathcal{E}) \to \Sh(\mathcal{D})$，{\it 复合
$f\circ g$} 是由函子
$(f \circ g)_* = f_* \circ g_*$ 和
$(f \circ g)^{-1} = g^{-1} \circ f^{-1}$ 定义的托扑斯态射。
\end{enumerate}
\end{definition}

\noindent
假设 $\alpha : \mathcal{S}_1 \to \mathcal{S}_2$ 是（可能为“大的”）范畴的
等价。若 $\mathcal{S}_1$、$\mathcal{S}_2$ 是托扑斯，则令
$f_* = \alpha$，并令 $f^{-1}$ 等于 $\alpha$ 的一个拟逆，便得到托扑斯态射
$f : \mathcal{S}_1 \to \mathcal{S}_2$。此外，这个态射是托扑斯
$2$-范畴中的等价（见第 \ref{sites-section-2-category} 节）。所以，若
$\mathcal{S}$ 等价于某个位点上的层范畴（而不一定等于某个位点上的层范畴），
说“$\mathcal{S}$ 是一个托扑斯”是有意义的。我们偶尔采用这种记号滥用。

\medskip\noindent
{\it 空托扑斯}是位点 $\mathcal{C}$ 上的层所成的托扑斯，其中
$\mathcal{C}$ 是空范畴。我们有时用 $\emptyset$ 表示这个位点。这个位点只有
唯一一个层（因为 $\emptyset$
没有对象）。所以 $\Sh(\emptyset)$ 等价于只有一个对象和一个态射的范畴。

\medskip\noindent
{\it 点托扑斯}是如下位点 $\mathcal{C}$ 上的层所成的托扑斯：它只有一个对象
$pt$ 和一个态射 $\text{id}_{pt}$，唯一的覆盖是
$\{\text{id}_{pt}\}$。我们简记这个位点为 $pt$。显然，层范畴
$ = $ 预层范畴 $ = $ 集合范畴。用公式表示：
$\Sh(pt) = \textit{Sets}$。

\medskip\noindent
设 $\mathcal{C}$ 和 $\mathcal{D}$ 为位点。设
$f : \Sh(\mathcal{D}) \to \Sh(\mathcal{C})$ 为托扑斯态射。注意，$f_*$
与所有极限交换，而 $f^{-1}$ 与所有余极限交换；见 Categories，引理
\ref{categories-lemma-adjoint-exact}。特别地，上述定义中对 $f^{-1}$ 的条件
保证 $f^{-1}$ 正合。托扑斯态射通常使用引理
\ref{sites-lemma-cocontinuous-morphism-topoi} 或下一个引理来构造。

\begin{lemma}
\label{sites-lemma-morphism-sites-topoi}
给定与函子 $u : \mathcal{C} \to \mathcal{D}$ 对应的位点态射
$f : \mathcal{D} \to \mathcal{C}$，函子对
$(f^{-1} = u_s, f_* = u^s)$ 是托扑斯态射
$\Sh(\mathcal{D}) \to \Sh(\mathcal{C})$。
\end{lemma}

\begin{proof}
这由定义 \ref{sites-definition-morphism-sites} 显然成立。
\end{proof}

\begin{remark}
\label{sites-remark-pt-topos}
有许多位点产生托扑斯 $\Sh(pt)$。下面是一个有用的例子。假设 $S$ 是一个
（由集合组成的）集合，并且至少含有一个非空元素。令 $\mathcal{S}$ 为这样的
范畴：对象是 $S$ 的元素，态射是任意集合映射。假设 $\mathcal{S}$ 有纤维积。
例如，若 $S = \mathcal{P}(\text{无限集})$ 是任意无限集合的幂集，则
此条件成立（集合论练习）。把满射映射族声明为覆盖，从而使
$\mathcal{S}$ 成为位点（并如 Sets，第 \ref{sets-section-coverings-site} 节
那样选择一个适当且充分大的覆盖族集合）。我们断言
$\Sh(\mathcal{S})$ 等价于集合范畴。

\medskip\noindent
先证明 $S$ 含有一个单元素集 $e \in S$ 的情形。在此情形下，存在托扑斯等价
$i : \Sh(pt) \to \Sh(\mathcal{S})$，由函子
\begin{equation}
\label{sites-equation-sheaves-pt-sets}
i^{-1}\mathcal{F} = \mathcal{F}(e), \quad
i_*E = (U \mapsto \Mor_{\textit{Sets}}(U, E))
\end{equation}
给出。事实上，假设 $\mathcal{F}$ 是 $\mathcal{S}$ 上的层。对任意
$U \in \Ob(\mathcal{S}) = S$，都可以找到覆盖
$\{\varphi_u : e \to U\}_{u \in U}$，其中 $\varphi_u$ 把 $e$ 的唯一元素
映到 $u \in U$。此时层条件蕴含
$\mathcal{F}(U) = \prod_{u \in U} \mathcal{F}(e)$。换言之，
$\mathcal{F}(U) = \Mor_{\textit{Sets}}(U, \mathcal{F}(e))$。此外，这条规则
与限制映射相容。因此函子
$$
i_* :
\textit{Sets} = \Sh(pt)
\longrightarrow
\Sh(\mathcal{S}), \quad
E \longmapsto (U \mapsto \Mor_{\textit{Sets}}(U, E))
$$
是范畴等价，其逆就是上面给出的函子 $i^{-1}$。

\medskip\noindent
若 $\mathcal{S}$ 不含单元素集，则上述定义的函子 $i_*$ 仍有意义。为证明它在
此情形下仍为等价，任取一个非空的 $\tilde e \in S$，并取一个像为单元素集的
映射 $\varphi : \tilde e \to \tilde e$。对任意层 $\mathcal{F}$，令
$$
\mathcal{F}(e) :=
\Im(
\mathcal{F}(\varphi) :
\mathcal{F}(\tilde e)
\longrightarrow
\mathcal{F}(\tilde e)
)
$$
并证明这给出 $i_*$ 的拟逆。细节从略。
\end{remark}

\begin{remark}
\label{sites-remark-morphism-topoi-big}
（与托扑斯态射有关的集合论问题。初读时可跳过。）上述定义的一个托扑斯态射
不是集合，而是类。换言之，它由数学公式而非数学对象给出。虽然可以考虑两个
给定托扑斯之间所有态射的汇集，但把它作为数学对象引入并非好主意。另一方面，
假设 $\mathcal{C}$ 和 $\mathcal{D}$ 是给定的位点。考虑函子
$\Phi : \mathcal{C} \to \Sh(\mathcal{D})$。这样的函子是集合，换言之，它
是数学对象。可以依次对 $\Phi$ 提出下列问题：
\begin{enumerate}
\item 给定 $\mathcal{D}$ 上的层 $\mathcal{F}$，规则
$U \mapsto \Mor_{\Sh(\mathcal{D})}(\Phi(U), \mathcal{F})$ 是否定义
$\mathcal{C}$ 上的层？若是，它定义函子
$\Phi_* : \Sh(\mathcal{D}) \to \Sh(\mathcal{C})$。
\item $\Phi_*$ 是否有左伴随？若有，把这个左伴随记作 $\Phi^{-1}$。
\item $\Phi^{-1}$ 是否正合？
\end{enumerate}
若最后一个问题的答案仍为“是”，便得到托扑斯态射
$(\Phi_*, \Phi^{-1})$。此外，给定任意托扑斯态射
$(f_*, f^{-1})$，可令 $\Phi(U) = f^{-1}(h_U^\#)$，并得到上述函子
$\Phi$，使得 $f_* \cong \Phi_*$ 且 $f^{-1} \cong \Phi^{-1}$（与伴随性质
相容）。结论是：通过使用 $\Phi$ 的汇集而非托扑斯态射，我们 (a) 把托扑斯
态射概念替换为数学对象，并且 (b) $\Phi$ 的汇集构成一个类（而不是类的
汇集）。当然，还可以说得更多，例如可以更准确地阐明上述条件 (2) 和 (3) 的
意义；我们在第 \ref{sites-section-points} 节对托扑斯的点讨论这一点。
\end{remark}

\begin{remark}
\label{sites-remark-quasi-continuous-morphism-topoi}
（初读时可跳过。）设 $\mathcal{C}$ 和 $\mathcal{D}$ 为位点。拟位点态射
$f : \mathcal{D} \to \mathcal{C}$（见备注
\ref{sites-remark-quasi-continuous-morphism-sites}）完全如引理
\ref{sites-lemma-morphism-sites-topoi} 那样，产生从
$\Sh(\mathcal{D})$ 到 $\Sh(\mathcal{C})$ 的托扑斯态射 $f$。
\end{remark}














\section{G-集合与态射}
\label{sites-section-G-sets-morphisms}

\noindent
设 $\varphi : G \to H$ 为群同态。按照例 \ref{sites-example-site-on-group} 和第
\ref{sites-section-example-sheaf-G-sets} 节，选择（适当的）位点
$\mathcal{T}_G$ 和 $\mathcal{T}_H$。令
$u : \mathcal{T}_H \to \mathcal{T}_G$ 为这样的函子：它把 $H$-集合 $U$
送到具有相同底层集合的 $G$-集合 $U_\varphi$，其 $G$-作用由
$g \cdot \xi = \varphi(g)\xi$ 定义，其中 $g \in G$ 且 $\xi \in U$。
显然，$u$ 与有限极限交换并且连续\footnote{集合论说明：先选择
$\mathcal{T}_H$，再选择 $\mathcal{T}_G$，使其包含 $u(\mathcal{T}_H)$，
并使 $\mathcal{T}_H$ 中的每个覆盖都对应于 $\mathcal{T}_G$ 中的覆盖。这可由
Sets，引理 \ref{sets-lemma-sets-with-group-action}、
\ref{sets-lemma-what-is-in-it-G-sets} 和
\ref{sets-lemma-coverings-site} 实现。}。应用命题
\ref{sites-proposition-get-morphism} 和引理 \ref{sites-lemma-morphism-sites-topoi}，得到
与 $\varphi$ 相伴的托扑斯态射
$$
f : \Sh(\mathcal{T}_G) \longrightarrow \Sh(\mathcal{T}_H)
$$
由命题 \ref{sites-proposition-sheaves-on-group}，得到一对伴随函子
$$
f_* : G\textit{-Sets} \longrightarrow H\textit{-Sets}, \quad
f^{-1} : H\textit{-Sets} \longrightarrow G\textit{-Sets}.
$$
下面具体求出这两个函子。

\medskip\noindent
先求 $f_*$ 的公式。回忆一下，给定 $G$-集合 $S$，$\mathcal{T}_G$ 上的相应
层 $\mathcal{F}_S$ 由规则
$\mathcal{F}_S(U) = \Mor_G(U, S)$ 给出。另一方面，给定
$\mathcal{T}_H$ 上的层 $\mathcal{G}$，相应的 $H$-集合由规则
$\mathcal{G}({}_HH)$ 给出。因此
$$
f_*S = \Mor_{G\textit{-Sets}}(({}_HH)_\varphi, S).
$$
进一步展开，得到
$$
f_*S = \{ a : H \to S \mid a(gh) = ga(h) \}
$$
其左 $H$-作用由
$(h \cdot a)(h') = a(h'h)$ 给出，其中 $a \in f_*S$ 任意。
% P01-ERRATA: P01-E0068 -- g is a G-element but a's argument lies in H; the formula omits \varphi(g).

\medskip\noindent
接着显式计算 $f^{-1}$。注意，由于 $\mathcal{T}_G$ 和
$\mathcal{T}_H$ 上的拓扑都是次典范的，所有可表预层都是层。此外，给定
$\mathcal{T}_H$ 的对象 $V$，可知 $f^{-1}h_V$ 等于 $h_{u(V)}$（见引理
\ref{sites-lemma-pullback-representable-sheaf}）。所以，对可表层有
$f^{-1}S = S_\varphi$。由于 $\mathcal{T}_H$ 上的每个层都是可表层的余积，
可知一般情形下也成立。因此，对任意 $H$-集合 $T$，有
$$
f^{-1}T = T_\varphi.
$$
$f^{-1}$ 与 $f_*$ 之间的伴随性由公式
$$
\Mor_{G\textit{-Sets}}(T_\varphi, S) =
\Mor_{H\textit{-Sets}}(T, f_*S)
$$
体现，其中 $f_*S$ 如上。这也可以直接证明。此外，由此显然可知
$(f^{-1}, f_*)$ 构成伴随对，且 $f^{-1}$ 正合。因此，也可以不用上述构造，
而直接以这种方式定义托扑斯态射
$f : G\textit{-Sets} \to H\textit{-Sets}$。














\section{拟紧对象与余极限}
\label{sites-section-quasi-compact}

\noindent
为能使用与拓扑空间情形相同的语言，我们引入如下术语。

\begin{definition}
\label{sites-definition-quasi-compact}
设 $\mathcal{C}$ 为位点。若对 $\mathcal{C}$ 中任一覆盖
$\mathcal{U} = \{U_i \to U\}_{i \in I}$，都存在另一个覆盖
$\mathcal{V} = \{V_j \to U\}_{j \in J}$，以及一个由
$\text{id} : U \to U$、$\alpha : J \to I$ 和
$V_j \to U_{\alpha(j)}$ 给出的具有固定靶的态射族之态射
$\mathcal{V} \to \mathcal{U}$（见定义
\ref{sites-definition-morphism-coverings}），使得 $\alpha$ 的像是 $I$ 的有限
子集，则称 $\mathcal{C}$ 的对象 $U$ 是{\it 拟紧的}。
\end{definition}

\noindent
当然，通常的概念足以推出 $U$ 拟紧。

\begin{lemma}
\label{sites-lemma-conclude-quasi-compact}
设 $\mathcal{C}$ 为位点，$U$ 为 $\mathcal{C}$ 的对象。考虑下列条件：
\begin{enumerate}
\item $U$ 拟紧；
\item 对 $\mathcal{C}$ 中每个覆盖 $\{U_i \to U\}_{i \in I}$，都存在
$\mathcal{C}$ 中加细 $\mathcal{U}$ 的有限覆盖
$\{V_j \to U\}_{j = 1, \ldots, m}$；
\item 对 $\mathcal{C}$ 中每个覆盖 $\{U_i \to U\}_{i \in I}$，都存在
有限子集 $I' \subset I$，使得
$\{U_i \to U\}_{i \in I'}$ 是 $\mathcal{C}$ 中的覆盖。
\end{enumerate}
则总有 (3) $\Rightarrow$ (2) $\Rightarrow$ (1)，但反向蕴含一般不成立。
\end{lemma}

\begin{proof}
这些蕴含由定义立即可得。令 $X = [0, 1] \subset \mathbf{R}$ 为拓扑空间
（带通常的 $\epsilon$-$\delta$ 拓扑）。令 $\mathcal{C}$ 为 $X$ 的开子空间
所成的范畴，以包含映射为态射，并取通常的开覆盖（与例
\ref{sites-example-site-topological} 比较）。不过，我们随后改变 $\mathcal{C}$ 中
覆盖的概念：除形如 $\{U \to U\}$ 的覆盖外，排除所有有限覆盖。容易看出，
这会成为定义 \ref{sites-definition-site} 所述的位点。在这个位点中，对象
$X = U = [0, 1]$ 在定义 \ref{sites-definition-quasi-compact} 的意义下拟紧，但
$U$ 不满足 (2)。留给读者构造满足 (2) 而不满足 (3) 的例子。
\end{proof}

\noindent
下面给出拟紧对象的托扑斯论含义。

\begin{lemma}
\label{sites-lemma-quasi-compact}
设 $\mathcal{C}$ 为位点，$U$ 为 $\mathcal{C}$ 的对象。下列条件等价：
\begin{enumerate}
\item $U$ 拟紧；
\item 对每个层的满射
$\coprod_{i \in I} \mathcal{F}_i \to h_U^\#$，都存在有限子集
$J \subset I$，使得
$\coprod_{i \in J} \mathcal{F}_i \to h_U^\#$ 是满射。
\end{enumerate}
\end{lemma}

\begin{proof}
假设 (1) 成立，并设 $\coprod_{i \in I} \mathcal{F}_i \to h_U^\#$ 为满射。
则 $\text{id}_U$ 是 $h_U^\#$ 在 $U$ 上的截面。因此存在覆盖
$\{U_a \to U\}_{a \in A}$，并且对每个 $a \in A$，存在
$\coprod_{i \in I} \mathcal{F}_i$ 在 $U_a$ 上的截面 $s_a$，它映到
$\text{id}_U$。根据把预层余积层化以构造余积的方法（引理
\ref{sites-lemma-colimit-sheaves}），对每个 $a$ 都存在覆盖
$\{U_{ab} \to U_a\}_{b \in B_a}$；对所有 $b \in B_a$，都存在
$\iota(b) \in I$ 以及 $\mathcal{F}_{\iota(b)}$ 在 $U_{ab}$ 上的截面
$s_{b}$，它映到 $\text{id}_U|_{U_{ab}}$。因此，用
$\{U_{ab} \to U\}_{a \in A, b \in B_a}$ 代替覆盖
$\{U_a \to U\}_{a \in A}$ 后，可以假设有映射 $\iota : A \to I$，
且对每个 $a \in A$，有 $\mathcal{F}_{\iota(a)}$ 在 $U_a$ 上的截面
$s_a$，它映到 $\text{id}_U$。
由于 $U$ 拟紧，存在覆盖 $\{V_c \to U\}_{c \in C}$、像为有限集的映射
$\alpha : C \to A$，以及 $U$ 上的态射
$V_c \to U_{\alpha(c)}$。于是 $J = \Im(\iota \circ \alpha) \subset I$
满足要求，因为
$\coprod_{c \in C} h_{V_c}^\# \to h_U^\#$ 是满射（引理
\ref{sites-lemma-covering-surjective-after-sheafification}），并且经由
$\coprod_{i \in J} \mathcal{F}_i \to h_U^\#$ 分解。（这里用到复合
$h_{V_c}^\# \to h_{U_{\alpha(c)}}
\xrightarrow{s_{\alpha(c)}} \mathcal{F}_{\iota(\alpha(c))} \to h_U^\#$
就是由态射 $V_c \to U$ 诱导的映射 $h_{V_c}^\# \to h_U^\#$，因为
$s_{\alpha(c)}$ 映到 $\text{id}_U|_{U_{\alpha(c)}}$。）
% P01-ERRATA: P01-E0069 -- the intermediate representable is missing sheafification marker \#.

\medskip\noindent
假设 (2) 成立。令 $\{U_i \to U\}_{i \in I}$ 为覆盖。由引理
\ref{sites-lemma-covering-surjective-after-sheafification} 可知
$\coprod_{i \in I} h_{U_i}^\# \to h_U^\#$ 是满射。因此可找到有限子集
$J \subset I$，使得
$\coprod_{j \in J} h_{U_j}^\# \to h_U^\#$ 是满射。再按上述论证，得到
$\mathcal{C}$ 中 $U$ 的覆盖 $\{V_c \to U\}_{c \in C}$ 和映射
$\iota : C \to J$，使得 $\text{id}_U$ 可提升为
$h_{U_{\iota(c)}}^\#$ 在 $V_c$ 上的截面 $s_c$。
% P01-ERRATA: P01-E0070 -- source says "a section of s_c of"; translated by its only type-correct reading.
进一步加细该覆盖后，可以假设
$s_c \in h_{U_{\iota(c)}}(V_c)$ 映到 $\text{id}_U$。于是
$s_c : V_c \to U_{\iota(c)}$ 是 $U$ 上的态射，结论成立。
\end{proof}

\noindent
上述引理启发如下定义。

\begin{definition}
\label{sites-definition-quasi-compact-topos}
若对 $\Sh(\mathcal{C})$ 中任一满射
$\coprod_{i \in I} \mathcal{F}_i \to \mathcal{F}$，都存在有限子集
$J \subset I$，使得
$\coprod_{i \in J} \mathcal{F}_i \to \mathcal{F}$ 是满射，则称对象
$\mathcal{F}$ 是{\it 拟紧的}。若托扑斯
$\Sh(\mathcal{C})$ 的终对象 $*$ 是拟紧对象，则称
$\Sh(\mathcal{C})$ 是{\it 拟紧的}。
\end{definition}

\noindent
由引理 \ref{sites-lemma-quasi-compact}，若位点 $\mathcal{C}$ 有终对象 $X$，则
$\Sh(\mathcal{C})$ 拟紧当且仅当 $X$ 拟紧。

\begin{lemma}
\label{sites-lemma-image-of-quasi-compact}
\begin{slogan}
层的满射传递拟紧性。
\end{slogan}
设 $\mathcal{C}$ 为位点。
\begin{enumerate}
\item 若 $U \to V$ 是 $\mathcal{C}$ 的态射，且
$h_U^\# \to h_V^\#$ 是满射，而 $U$ 拟紧，则 $V$ 拟紧。
\item 若 $\mathcal{F} \to \mathcal{G}$ 是集合层的满射，且
$\mathcal{F}$ 拟紧，则 $\mathcal{G}$ 拟紧。
\end{enumerate}
\end{lemma}

\begin{proof}
略。
\end{proof}

\begin{lemma}
\label{sites-lemma-coproduct-quasi-compact}
\begin{slogan}
层的有限余积保持拟紧性。
\end{slogan}
% P01-ERRATA: P01-E0071 -- source slogan misspells "coproducts" as "copoducts".
设 $\mathcal{C}$ 为位点。若 $n \geq 1$，且
$\mathcal{F}_1, \ldots, \mathcal{F}_n$ 是 $\mathcal{C}$ 上的拟紧层，则
$\coprod_{i = 1, \ldots, n} \mathcal{F}_i$ 拟紧。
\end{lemma}

\begin{proof}
略。
\end{proof}

\noindent
下列两个引理是 Sheaves，引理
\ref{sheaves-lemma-directed-colimits-sections} 对位点的类似结果。

\begin{lemma}
\label{sites-lemma-directed-colimits-sections}
设 $\mathcal{C}$ 为位点，且
$\mathcal{I} \to \Sh(\mathcal{C})$，$i \mapsto \mathcal{F}_i$
是集合层的滤过图表。设 $U \in \Ob(\mathcal{C})$。考虑典范映射
$$
\Psi :
\colim_i \mathcal{F}_i(U)
\longrightarrow
\left(\colim_i \mathcal{F}_i\right)(U)
$$
采用上面引入的术语：
\begin{enumerate}
\item 若所有转移映射都是单射，则对任意 $U$，$\Psi$ 都是单射。
\item 若 $U$ 拟紧，则 $\Psi$ 是单射。
\item 若 $U$ 拟紧，且所有转移映射都是单射，则 $\Psi$ 是同构。
\item 若 $U$ 有一个由覆盖 $\{U_j \to U\}_{j \in J}$ 组成的共尾系，
其中 $J$ 有限，且对所有 $j, j' \in J$，
$U_j \times_U U_{j'}$ 都拟紧，则 $\Psi$ 是双射。
\end{enumerate}
\end{lemma}

\begin{proof}
假设所有转移映射都是单射。此时预层
$\mathcal{F}' : V \mapsto \colim_i \mathcal{F}_i(V)$ 是分离的（见定义
\ref{sites-definition-separated}）。由引理 \ref{sites-lemma-colimit-sheaves} 有
$(\mathcal{F}')^\# = \colim_i \mathcal{F}_i$。由定理
\ref{sites-theorem-plus} 可知
$\mathcal{F}' \to (\mathcal{F}')^\#$ 是单射。这证明了 (1)。

\medskip\noindent
假设 $U$ 拟紧。设 $s \in \mathcal{F}_i(U)$ 和
$s' \in \mathcal{F}_{i'}(U)$ 在左边给出的元素经 $\Psi$ 后具有相同的像。
这意味着可以选择覆盖 $\{U_a \to U\}_{a \in A}$，并且对每个
$a \in A$ 选择指标 $i_a \in I$，其中 $i_a \geq i$、$i_a \geq i'$，
使得 $\varphi_{ii_a}(s) = \varphi_{i'i_a}(s')$。因为 $U$ 拟紧，可以选择
覆盖 $\{V_b \to U\}_{b \in B}$、像为有限集的映射
$\alpha : B \to A$，以及 $U$ 上的态射
$V_b \to U_{\alpha(b)}$。选择 $i''\in I$，使它不小于（即 $\geq$）所有
$i_{\alpha(b)}$；这是可能的，因为 $\alpha$ 的像有限。于是对所有
$b \in B$，$\varphi_{ii''}(s)$ 和 $\varphi_{i'i''}(s)$ 在 $V_b$ 上相等，从而
$\varphi_{ii''}(s) = \varphi_{i'i''}(s)$。这证明了 (2)。
% P01-ERRATA: P01-E0072 -- both second terms use s but the setup requires s'; formula preserved.

\medskip\noindent
假设 $U$ 拟紧，且所有转移映射都是单射。令 $s$ 为 $\Psi$ 的靶中的元素。
存在覆盖 $\{U_a \to U\}_{a \in A}$；对每个 $a \in A$，存在指标
$i_a \in I$ 和截面 $s_a \in \mathcal{F}_{i_a}(U_a)$，使得对所有
$a \in A$，$s|_{U_a}$ 都来自 $s_a$。因为 $U$ 拟紧，可以选择覆盖
$\{V_b \to U\}_{b \in B}$、像为有限集的映射
$\alpha : B \to A$，以及 $U$ 上的态射
$V_b \to U_{\alpha(b)}$。选择 $i \in I$，使它不小于（即 $\geq$）所有
$i_{\alpha(b)}$；这是可能的，因为 $\alpha$ 的像有限。由 (1)，截面
$s_b = \varphi_{i_{\alpha(b)} i}(s_{\alpha(b)})|_{V_b}$
在 $V_b \times_U V_{b'}$ 上相符。因此它们粘合为截面
$s' \in \mathcal{F}_i(U)$，该截面经 $\Psi$ 映到 $s$。这证明了 (3)。

\medskip\noindent
假设 (4) 的假设成立。由引理 \ref{sites-lemma-conclude-quasi-compact}，对象
$U$ 拟紧，故由 (2)，$\Psi$ 是单射。为证明满射性，令 $s$ 为 $\Psi$ 的靶中
的元素。根据假设，存在有限覆盖
$\{U_j \to U\}_{j = 1, \ldots, m}$，使得对所有
$1 \leq j, j' \leq m$，$U_j \times_U U_{j'}$ 都拟紧；并且对每个 $j$
存在指标 $i_j \in I$ 和 $s_j \in \mathcal{F}_{i_j}(U_j)$，使得
$s|_{U_j}$ 是 $s_j$ 的像，对所有 $j$ 都成立。由于
$U_j \times_U U_{j'}$ 拟紧，可以应用
(2)，从而存在 $i_{jj'} \in I$，满足 $i_{jj'} \geq i_j$、
$i_{jj'} \geq i_{j'}$，并使
$\varphi_{i_ji_{jj'}}(s_j)$ 与
$\varphi_{i_{j'}i_{jj'}}(s_{j'})$ 在
$U_j \times_U U_{j'}$ 上相等。选择指标 $i \in I$，使其不小于所有
$i_{jj'}$。于是 $\mathcal{F}_i$ 的截面
$\varphi_{i_ji}(s_j)$ 粘合成 $\mathcal{F}_i$ 在 $U$ 上的截面。
该截面按要求映到元素 $s$。
\end{proof}

\begin{lemma}
\label{sites-lemma-directed-colimits-global-sections}
设 $\mathcal{C}$ 为位点，且
$\mathcal{I} \to \Sh(\mathcal{C})$，$i \mapsto \mathcal{F}_i$
是集合层的滤过图表。考虑典范映射
$$
\Psi :
\colim_i \Gamma(\mathcal{C}, \mathcal{F}_i)
\longrightarrow
\Gamma(\mathcal{C}, \colim_i \mathcal{F}_i)
$$
有下列结论：
\begin{enumerate}
\item 若所有转移映射都是单射，则 $\Psi$ 是单射。
\item 若 $\Sh(\mathcal{C})$ 拟紧，则 $\Psi$ 是单射。
\item 若 $\Sh(\mathcal{C})$ 拟紧，且所有转移映射都是单射，则
$\Psi$ 是同构。
\item 假设存在集合 $S \subset \Ob(\Sh(\mathcal{C}))$，具有如下性质：
\begin{enumerate}
\item 对每个满射 $\mathcal{F} \to *$，都存在
$\mathcal{K} \in S$ 和映射 $\mathcal{K} \to \mathcal{F}$，使得
$\mathcal{K} \to *$ 是满射；
\item 对 $\mathcal{K} \in S$，积
$\mathcal{K} \times \mathcal{K}$ 拟紧。
\end{enumerate}
则 $\Psi$ 是双射。
\end{enumerate}
\end{lemma}

\begin{proof}
(1) 的证明。假设所有转移映射都是单射。此时预层
$\mathcal{F}' : V \mapsto \colim_i \mathcal{F}_i(V)$ 是分离的（见定义
\ref{sites-definition-separated}）。由引理 \ref{sites-lemma-colimit-sheaves} 有
$(\mathcal{F}')^\# = \colim_i \mathcal{F}_i$。由定理
\ref{sites-theorem-plus} 可知
$\mathcal{F}' \to (\mathcal{F}')^\#$ 是单射。这证明了 (1)。

\medskip\noindent
(2) 的证明。假设 $\Sh(\mathcal{C})$ 拟紧。回忆一下，对
$\Sh(\mathcal{C})$ 中的所有 $\mathcal{F}$，都有
$\Gamma(\mathcal{C}, \mathcal{F}) = \Mor(*, \mathcal{F})$。
令 $a_i, b_i : * \to \mathcal{F}_i$；对 $i' \geq i$，以
$a_{i'}, b_{i'} : * \to \mathcal{F}_{i'}$ 表示它们与系统转移映射的复合。
令 $a = \colim_{i' \geq i} a_{i'}$，并对 $b$ 作类似定义。对
$i' \geq i$，记
$$
E_{i'} = \text{Equalizer}(a_{i'}, b_{i'}) \subset *
\quad\text{且}\quad
E = \text{Equalizer}(a, b) \subset *
$$
由 Categories，引理 \ref{categories-lemma-directed-commutes}，有
$E = \colim_{i' \geq i} E_{i'}$。从而
$\coprod_{i' \geq i} E_{i'} \to E$ 是层的满射。因此，若 $E = *$，
即 $a = b$，则因为 $*$ 拟紧，可知对某个 $i' \geq i$ 有
$E_{i'} = *$，进而对某个 $i' \geq i$ 有 $a_{i'} = b_{i'}$。
这证明了 (2)。

\medskip\noindent
(3) 的证明。假设 $\Sh(\mathcal{C})$ 拟紧，且所有转移映射都是单射。
令 $a : * \to \colim \mathcal{F}_i$ 为映射。则
$E_i = a^{-1}(\mathcal{F}_i) \subset *$ 是子层，并且
$\colim E_i = *$（由上面的引用）。故对某个 $i$ 有 $E_i = *$，从而
$a$ 的像按要求包含在 $\mathcal{F}_i$ 中。

\medskip\noindent
(4) 的证明。设 $S \subset \Ob(\Sh(\mathcal{C}))$ 满足 (4)(a)、(b)。
把 (4)(a) 应用于 $\text{id} : * \to *$，可找到
$\mathcal{K} \in S$，使得 $\mathcal{K} \to *$ 是满射。映射
$\mathcal{K} \times \mathcal{K} \to \mathcal{K} \to *$ 都是满射。
由 (4)(b) 和引理 \ref{sites-lemma-image-of-quasi-compact}，可知
$\mathcal{K}$ 与 $\Sh(\mathcal{C})$ 都拟紧。因此由 (2)，$\Psi$ 是单射。
令 $\mathcal{F} = \colim \mathcal{F}_i$，并令
$s : * \to \mathcal{F}$ 为该余极限的整体截面。由于
$\coprod \mathcal{F}_i \to \mathcal{F}$ 是满射，投影
$$
\coprod\nolimits_{i \in I} * \times_{s, \mathcal{F}} \mathcal{F}_i \to *
$$
是满射。由 (4)(a)，得到 $\mathcal{K} \in S$ 和映射
$\mathcal{K} \to \coprod_{i \in I} * \times_{s, \mathcal{F}} \mathcal{F}_i$，
其中 $\mathcal{K} \to *$ 是满射。因为 $\mathcal{K}$ 拟紧，对某个有限子集
$I' \subset I$，可得到分解
$\mathcal{K} \to \coprod_{i' \in I'} * \times_{s, \mathcal{F}}
\mathcal{F}_{i'}$。令 $i \in I$ 为有限子集 $I'$ 的一个上界。转移映射定义
映射 $\coprod_{i' \in I'} \mathcal{F}_{i'} \to \mathcal{F}_i$，后者又产生
映射 $\mathcal{K} \to * \times_{s, \mathcal{F}} \mathcal{F}_i$。换言之，
得到 $\mathcal{K} \in S$，其中 $\mathcal{K} \to *$ 是满射，并有交换图
$$
\xymatrix{
\mathcal{K} \times \mathcal{K} \ar@<1ex>[r] \ar@<-1ex>[r] &
\mathcal{K} \ar[d] \ar[r] & {*} \ar[d]^s \\
& \mathcal{F}_i \ar[r] & \mathcal{F} \ar@{=}[r] & \colim \mathcal{F}_i
}
$$
注意该图的顶行是余等化子。因此，只需证明增大 $i$ 后，两个诱导映射
$a_i, b_i : \mathcal{K} \times \mathcal{K} \to \mathcal{F}_i$
相等即可。下一段使用与 (2) 的证明完全相同的论证来说明这一点；建议读者
跳过证明的剩余部分。
% P01-ERRATA: P01-E0073 -- source has the duplicated phrase "This is done shown".

\medskip\noindent
对 $i' \geq i$，以
$a_{i'}, b_{i'} : \mathcal{K} \times \mathcal{K} \to \mathcal{F}_{i'}$
表示 $a_i, b_i$ 与系统转移映射的复合。令
$a = \colim_{i' \geq i} a_{i'} :
\mathcal{K} \times \mathcal{K} \to \mathcal{F}$，并对 $b$ 作类似定义。
由上图的交换性，有 $a = b$。对 $i' \geq i$，记
$$
E_{i'} = \text{Equalizer}(a_{i'}, b_{i'}) \subset
\mathcal{K} \times \mathcal{K}
\quad\text{且}\quad
E = \text{Equalizer}(a, b) \subset
\mathcal{K} \times \mathcal{K}
$$
由 Categories，引理 \ref{categories-lemma-directed-commutes}，有
$E = \colim_{i' \geq i} E_{i'}$。从而
$\coprod_{i' \geq i} E_{i'} \to E$ 是层的满射。因为 $a = b$，有
$E = \mathcal{K} \times \mathcal{K}$。又因为由 (4)(b)，
$\mathcal{K} \times \mathcal{K}$ 拟紧，所以对某个 $i' \geq i$ 有
$E_{i'} = \mathcal{K} \times \mathcal{K}$，进而对某个 $i' \geq i$ 有
$a_{i'} = b_{i'}$。
\end{proof}

\begin{remark}
\label{sites-remark-stronger-conditions}
设 $\mathcal{C}$ 为位点。可以用多种方式保证引理
\ref{sites-lemma-directed-colimits-global-sections} 第 (4) 部分的假设成立。下面
列出几种。
\begin{enumerate}
\item 假设存在集合 $\mathcal{B} \subset \Ob(\mathcal{C})$，具有如下性质：
\begin{enumerate}
\item 对每个满射 $\mathcal{F} \to *$，都存在 $m \geq 0$ 和
$U_1, \ldots, U_m \in \mathcal{B}$，使得 $\mathcal{F}(U_j)$ 非空，且
$\coprod h_{U_j}^\# \to *$ 是满射；
\item 对 $U, U' \in \mathcal{B}$，层
$h_U^\# \times h_{U'}^\#$ 拟紧。
\end{enumerate}
\item 假设存在集合 $\mathcal{B} \subset \Ob(\mathcal{C})$，具有如下性质：
\begin{enumerate}
\item 存在 $m \geq 0$ 和 $U_1, \ldots, U_m \in \mathcal{B}$，使得
$\coprod h_{U_j}^\# \to *$ 是满射；
\item 对 $U \in \mathcal{B}$，$U$ 的任意覆盖都可由有限覆盖
$\{U_j \to U\}_{j = 1, \ldots, m}$ 加细，其中
$U_j \in \mathcal{B}$；
\item 对 $U, U' \in \mathcal{B}$，存在 $m \geq 0$、
$U_1, \ldots, U_m \in \mathcal{B}$，以及态射
$U_j \to U$ 和 $U_j \to U'$，使得
$\coprod h_{U_j}^\# \to h_U^\# \times h_{U'}^\#$ 是满射。
\end{enumerate}
\item 假设：
\begin{enumerate}
\item $\Sh(\mathcal{C})$ 拟紧；
\item $\mathcal{C}$ 的每个对象都有一个覆盖，其成员均为拟紧对象；
\item 若 $U$ 和 $U'$ 拟紧，则层
$h_U^\# \times h_{U'}^\#$ 拟紧。
\end{enumerate}
\end{enumerate}
在情形 (1) 和 (2) 中，令 $S \subset \Ob(\Sh(\mathcal{C}))$ 为由
$U \in \mathcal{B}$ 的层 $h_U^\#$ 的有限余积组成的集合。在情形 (3) 中，
令 $S \subset \Ob(\Sh(\mathcal{C}))$ 为由拟紧对象
$U \in \Ob(\mathcal{C})$ 的层 $h_U^\#$ 的有限余积组成的集合。
\end{remark}

\noindent
以后将需要如下关于余极限可能情形的界。

\begin{lemma}
\label{sites-lemma-colimit-over-ordinal-sections}
设 $\mathcal{C}$ 为位点，$\beta$ 为序数。设
$\beta \to \Sh(\mathcal{C})$，$\alpha \mapsto \mathcal{F}_\alpha$
为 $\beta$ 上的层系。对 $U \in \Ob(\mathcal{C})$，考虑典范映射
$$
\colim_{\alpha < \beta} \mathcal{F}_\alpha(U)
\longrightarrow
\left(\colim_{\alpha < \beta} \mathcal{F}_\alpha \right)(U)
$$
若 $\beta$ 的共尾性足够大，则该映射对所有 $U$ 都是双射。
\end{lemma}

\begin{proof}
左边是预层范畴中所取余极限 $\mathcal{F}_{\colim}$ 在 $U$ 上的值，见第
\ref{sites-section-limits-colimits-PSh} 节。回忆一下，
$\colim_{\alpha < \beta} \mathcal{F}_\alpha$ 是
$\mathcal{F}_{\colim}$ 的层化 $\mathcal{F}_{\colim}^\#$，见引理
\ref{sites-lemma-colimit-sheaves}。令
$\mathcal{U} = \{U_i \to U\}_{i \in I}$ 为 $\mathcal{C}$ 的覆盖所成集合
$\text{Cov}(\mathcal{C})$ 的元素。若 $\beta$ 的共尾性大于 $I$ 的基数，
则断言
$$
H^0(\mathcal{U}, \mathcal{F}_{\colim}) =
\colim H^0(\mathcal{U}, \mathcal{F}_\alpha) =
\colim \mathcal{F}_\alpha(U) = \mathcal{F}_{\colim}(U)
$$
第二、第三个等号显然。对第一个等号，设
$s = (s_i) \in H^0(\mathcal{U}, \mathcal{F}_{\colim})$。则对每个 $i$，
元素 $s_i$ 来自某个 $\alpha_i < \beta$ 的元素
$s_{i, \alpha_i} \in \mathcal{F}_{\alpha_i}(U_i)$。由共尾性假设，可选择
与 $i$ 无关的 $\alpha_i = \alpha$。
% P01-ERRATA: P01-E0074 -- the source elides replacing all alpha_i by a common upper bound alpha.
于是对某个 $\alpha_{i, j} < \beta$，$s_i$ 和 $s_j$ 映到
$\mathcal{F}_{\alpha_{i, j}}(U_i \times_U U_j)$ 中的同一元素。由于
$I \times I$ 的基数也小于 $\beta$ 的共尾性，增大 $\alpha$ 后，可以假设
对所有 $i, j$ 都有 $\alpha_{i, j} = \alpha$。这证明自然映射
$\colim H^0(\mathcal{U}, \mathcal{F}_\alpha)
\to H^0(\mathcal{U}, \mathcal{F}_{\colim})$ 是满射。非常类似的论证表明
它是单射。特别地，$\mathcal{F}_{\colim}$ 对 $\mathcal{U}$ 满足层条件。
因此，若 $\beta$ 的共尾性大于所有覆盖的指标集 $I$ 的基数所成集合的上确界，
便得到结论。
% P01-ERRATA: P01-E0075 -- source has "cardinality if I x I" instead of "of".
\end{proof}







\section{位点的余极限}
\label{sites-section-colimit-sites}

\noindent
当位点是一个位点逆系统的极限时，我们需要引理
\ref{sites-lemma-directed-colimits-sections} 的类似结果。为简单起见，只说明覆盖的
指标集有限时的构造。

\begin{situation}
\label{sites-situation-inverse-limit-sites}
这里给定：
\begin{enumerate}
\item 余滤过指标范畴 $\mathcal{I}$；
\item 对 $i \in \Ob(\mathcal{I})$，给定位点 $\mathcal{C}_i$，使得
$\mathcal{C}_i$ 中每个覆盖都有有限指标集；
\item 对 $\mathcal{I}$ 中的态射 $a : i \to j$，给定由连续函子
$u_a : \mathcal{C}_j \to \mathcal{C}_i$ 所确定的位点态射
$f_a : \mathcal{C}_i \to \mathcal{C}_j$。
\end{enumerate}
并且只要 $\mathcal{I}$ 中 $c = a \circ b$，就有
$f_a \circ f_b = f_c$。
\end{situation}

\begin{lemma}
\label{sites-lemma-colimit-sites}
在情形 \ref{sites-situation-inverse-limit-sites} 中，可以如下构造位点
$(\mathcal{C}, \text{Cov}(\mathcal{C}))$：
\begin{enumerate}
\item 作为范畴，$\mathcal{C} = \colim \mathcal{C}_i$；
\item $\text{Cov}(\mathcal{C})$ 是各个
$\text{Cov}(\mathcal{C}_i)$ 在
$u_i : \mathcal{C}_i \to \mathcal{C}$ 下之像的并集。
\end{enumerate}
\end{lemma}

\begin{proof}
我们对位点态射复合的定义蕴含：只要 $\mathcal{I}$ 中
$c = a \circ b$，就有 $u_b \circ u_a = u_c$。公式
$\mathcal{C} = \colim \mathcal{C}_i$ 意味着
$\Ob(\mathcal{C}) = \colim \Ob(\mathcal{C}_i)$ 且
$\text{Arrows}(\mathcal{C}) = \colim \text{Arrows}(\mathcal{C}_i)$。
于是源、靶和复合都从各个 $\text{Arrows}(\mathcal{C}_i)$ 上的源、靶和
复合继承而来。这样便得到一个范畴。以
$u_i : \mathcal{C}_i \to \mathcal{C}$ 表示显然的函子。注意，给定
$\mathcal{C}$ 中任意有限图表，都存在某个 $i$，使该图表是
$\mathcal{C}_i$ 中一个图表的像。

\medskip\noindent
令 $\{U^t \to U\}$ 为 $\mathcal{C}$ 的覆盖。先证明：若
$V \to U$ 是 $\mathcal{C}$ 的态射，则 $U^t \times_U V$ 存在。由上述说明
和覆盖的定义，可以找到 $i$、$\mathcal{C}_i$ 的覆盖
$\{U_i^t \to U_i\}$ 以及态射 $V_i \to U_i$，它们在 $u_i$ 下的像就是
给定数据。断言 $U^t \times_U V$ 是
$U^t_i \times_{U_i} V_i$ 在 $u_i$ 下的像。事实上，对
$\mathcal{I}$ 中每个 $a : j \to i$，函子 $u_a$ 连续，所以
$u_a(U^t_i \times_{U_i} V_i) = u_a(U^t_i) \times_{u_a(U_i)} u_a(V_i)$。
特别地，如有需要，可以用 $j$ 代替 $i$。因此，若 $W$ 是
$\mathcal{C}$ 的另一个对象，则可以假设 $W = u_i(W_i)$，并得到
\begin{align*}
& \Mor_\mathcal{C}(W, u_i(U^t_i \times_{U_i} V_i)) \\
& =
\colim_{a : j \to i}
\Mor_{\mathcal{C}_j}(u_a(W_i), u_a(U^t_i \times_{U_i} V_i)) \\
& =
\colim_{a : j \to i}
\Mor_{\mathcal{C}_j}(u_a(W_i), u_a(U^t_i))
\times_{\Mor_{\mathcal{C}_j}(u_a(W_i), u_a(U_i))}
\Mor_{\mathcal{C}_j}(u_a(W_i), u_a(V_i)) \\
& =
\Mor_\mathcal{C}(W, U^t)
\times_{\Mor_\mathcal{C}(W, U)}
\Mor_\mathcal{C}(W, V)
\end{align*}
因为滤过余极限与有限极限交换（Categories，引理
\ref{categories-lemma-directed-commutes}）。同时可知
$\{U^t \times_U V \to V\}$ 是 $\mathcal{C}$ 中的覆盖。这样就验证了
定义 \ref{sites-definition-site} 的公理 (3)。

\medskip\noindent
为验证定义 \ref{sites-definition-site} 的公理 (2)，令
$\{U^t \to U\}_{t \in T}$ 为 $\mathcal{C}$ 的覆盖，并且对每个 $t$，令
$\{U^{ts} \to U^t\}$ 为 $\mathcal{C}$ 的覆盖。可以找到 $i$ 和
$\mathcal{C}_i$ 的覆盖 $\{U^t_i \to U_i\}_{t \in T}$，其在 $u_i$ 下的像
是 $\{U^t \to U\}$。因为 $T$ 是{\bf 有限的}，可以选择
$\mathcal{I}$ 中的 $a : j \to i$ 和 $\mathcal{C}_j$ 中的覆盖
$\{U^{ts}_j \to u_a(U^t_i)\}$，使它们在 $u_j$ 下的像给出
$\{U^{ts} \to U^t\}$。于是，把公理 (2) 应用于位点 $\mathcal{C}_j$，
可知 $\{U^{ts} \to U\}$ 是 $\mathcal{C}$ 的覆盖。

\medskip\noindent
定义 \ref{sites-definition-site} 的公理 (1) 的证明从略。
\end{proof}

\begin{lemma}
\label{sites-lemma-compute-pullback-to-limit}
在情形 \ref{sites-situation-inverse-limit-sites} 中，令
$u_i : \mathcal{C}_i \to \mathcal{C}$ 如引理
\ref{sites-lemma-colimit-sites} 中所构造。则 $u_i$ 定义位点态射
$f_i : \mathcal{C} \to \mathcal{C}_i$。对
$U_i \in \Ob(\mathcal{C}_i)$ 以及 $\mathcal{C}_i$ 上的层
$\mathcal{F}$，有
\begin{equation}
\label{sites-equation-compute-pullback-to-limit}
f_i^{-1}\mathcal{F}(u_i(U_i)) =
\colim_{a : j \to i} f_a^{-1}\mathcal{F}(u_a(U_i))
\end{equation}
\end{lemma}

\begin{proof}
由引理 \ref{sites-lemma-colimit-sites} 证明中的论证，立即可知各函子 $u_i$ 连续。
为完成证明，必须说明 $f_i^{-1} := u_{i, s}$ 是正合函子
$\Sh(\mathcal{C}_i) \to \Sh(\mathcal{C})$。事实上，只需证明
$f_i^{-1}$ 左正合，因为它作为左伴随是右正合的（Categories，引理
\ref{categories-lemma-exact-adjoint}）。先证明
(\ref{sites-equation-compute-pullback-to-limit})，然后由此推出正合性。

\medskip\noindent
对 $\mathcal{C}$ 的任意对象 $V$，可以选取 $a : j \to i$ 和对象
$V_j \in \Ob(\mathcal{C})$，使得 $V = u_j(V_j)$。然后可令
% P01-ERRATA: P01-E0076 -- source says "pick a a" before the arrow a : j -> i.
% P01-ERRATA: P01-E0077 -- V_j is typed as an object of C, but must lie in C_j for u_j(V_j).
$$
\mathcal{G}(V) = \colim_{b : k \to j} f_{a \circ b}^{-1}\mathcal{F}(u_b(V_j))
$$
该余极限的值 $\mathcal{G}(V)$ 与 $b : j \to i$ 的选择以及满足
$u_j(V_j) = V$ 的对象 $V_j$ 的选择无关；验证从略。
% P01-ERRATA: P01-E0078 -- the choice was denoted a : j -> i; source switches to b and reuses b incompatibly.
此外，若 $\alpha : V \to V'$ 是 $\mathcal{C}$ 的态射，则可以选择
$b : j \to i$ 和态射 $\alpha_j : V_j \to V'_j$，使得
$u_j(\alpha_j) = \alpha$。利用沿态射
$u_b(\alpha_j) : u_b(V_j) \to u_b(V'_j)$ 的限制映射，可诱导映射
$\mathcal{G}(V') \to \mathcal{G}(V)$。验证表明 $\mathcal{G}$ 是预层
（从略）。事实上，$\mathcal{G}$ 满足层条件。任一
$\mathcal{C}$ 中的覆盖 $\mathcal{U} = \{U^t \to U\}$ 都来自某个有限层级。
设 $\mathcal{U}_j = \{U^t_j \to U_j\}$ 在 $\mathcal{I}$ 中某个
$a : j \to i$ 下由 $u_j$ 映到 $\mathcal{U}$。则
$$
H^0(\mathcal{U}, \mathcal{G}) =
\colim_{b : k \to j} H^0(u_b(\mathcal{U}_j), f_{b \circ a}^{-1}\mathcal{F}) =
\colim_{b : k \to j} f_{b \circ a}^{-1}\mathcal{F}(u_b(U_j)) =
\mathcal{G}(U)
$$
如所需。第一个等号成立是因为滤过余极限与有限极限交换（Categories，引理
\ref{categories-lemma-directed-commutes}）。
% P01-ERRATA: P01-E0079 -- b \circ a is ill-typed here; the setup and earlier formula require a \circ b.
根据构造，$\mathcal{G}(U)$ 由
(\ref{sites-equation-compute-pullback-to-limit}) 的右边给出。因此，如果能证明
$\mathcal{G}$ 等于 $f_i^{-1}\mathcal{F}$，则
(\ref{sites-equation-compute-pullback-to-limit}) 成立。

\medskip\noindent
本段验证 $\mathcal{G}$ 与 $f_i^{-1}\mathcal{F}$ 典范同构。强烈建议读者
跳过本段。为验证这一点，必须证明有双射
$\Mor_{\Sh(\mathcal{C})}(\mathcal{G}, \mathcal{H}) =
\Mor_{\Sh(\mathcal{C}_i)}(\mathcal{F}, f_{i, *}\mathcal{H})$，并且它对
$\mathcal{C}$ 上的层 $\mathcal{H}$ 是函子性的，其中
$f_{i, *} = u_i^p$。映射 $\mathcal{G} \to \mathcal{H}$ 等价于相容的映射系
$$
\varphi_{a, b, V_j} :
f_{a \circ b}^{-1}\mathcal{F}(u_b(V_j))
\longrightarrow
\mathcal{H}(u_j(V_j))
$$
这里 $a : j \to i$、$b : k \to j$ 且
$V_j \in \Ob(\mathcal{C}_j)$ 任意。相容性迫使映射
$\varphi_{a, b, V_j}$ 等于
$\varphi_{a \circ b, \text{id}, u_b(V_j)}$。给定
$a : j \to i$，映射族 $\varphi_{a, \text{id}, V_j}$ 对应于层的映射
$\varphi_a : f_a^{-1}\mathcal{F} \to f_{j, *}\mathcal{H}$。
$\varphi_{a, \text{id}, u_a(V_i)}$ 与
$\varphi_{\text{id}, \text{id}, V_i}$ 之间的相容性蕴含：
$\varphi_a$ 是映射 $\varphi_{id}$ 关于下列等式给出的伴随：
$$
\Mor_{\Sh(\mathcal{C}_j)}(f_a^{-1}\mathcal{F}, f_{j, *}\mathcal{H}) =
\Mor_{\Sh(\mathcal{C}_i)}(\mathcal{F}, f_{a, *}f_{j, *}\mathcal{H}) =
\Mor_{\Sh(\mathcal{C}_i)}(\mathcal{F}, f_{i, *}\mathcal{H})
$$
% P01-ERRATA: P01-E0080 -- source has subject-verb disagreement: "compatibilities ... implies".
最终可知，整个映射系 $\varphi_{a, b, V_j}$ 由映射
$\varphi_{id} : \mathcal{F} \to f_{i, *}\mathcal{H}$ 决定。反过来，给定
这样的映射 $\psi : \mathcal{F} \to f_{i, *}\mathcal{H}$，反向阅读刚才的
论证即可构造映射族 $\varphi_{a, b, V_j}$。这便完成了
$\mathcal{G} = f_i^{-1}\mathcal{F}$ 的证明。

\medskip\noindent
假设 (\ref{sites-equation-compute-pullback-to-limit}) 成立。则函子
$\mathcal{F} \mapsto f_i^{-1}\mathcal{F}(U)$ 与有限极限交换，因为层的
有限极限在预层范畴中计算（引理 \ref{sites-lemma-limit-sheaf}），各函子
$f_a^{-1}$ 与有限极限交换，并且滤过余极限与有限极限交换。为说明对
$\mathcal{C}$ 的一般对象 $V$，函子
$\mathcal{F} \mapsto f_i^{-1}\mathcal{F}(V)$ 与有限极限交换，可以对上面
给出的 $f_i^{-1}\mathcal{F}(V) = \mathcal{G}(V)$ 公式使用同一论证。因此
$f_i^{-1}$ 左正合，引理得证。
\end{proof}

\begin{lemma}
\label{sites-lemma-colimit}
在情形 \ref{sites-situation-inverse-limit-sites} 中，假设给定：
\begin{enumerate}
\item 对所有 $i \in \Ob(\mathcal{I})$，给定 $\mathcal{C}_i$ 上的层
$\mathcal{F}_i$；
\item 对 $a : j \to i$，给定 $\mathcal{C}_j$ 上的层映射
$\varphi_a : f_a^{-1}\mathcal{F}_i \to \mathcal{F}_j$。
\end{enumerate}
并且只要 $c = a \circ b$，就有
$\varphi_c = \varphi_b \circ f_b^{-1}\varphi_a$。在引理
\ref{sites-lemma-colimit-sites} 的位点 $\mathcal{C}$ 上，令
$\mathcal{F} = \colim f_i^{-1}\mathcal{F}_i$。设
$i \in \Ob(\mathcal{I})$ 且 $X_i \in \text{Ob}(\mathcal{C}_i)$。则
$$
\colim_{a : j \to i} \mathcal{F}_j(u_a(X_i)) = \mathcal{F}(u_i(X_i))
$$
\end{lemma}

\begin{proof}
形式论证表明
$$
\colim_{a : j \to i} \mathcal{F}_i(u_a(X_i)) =
\colim_{a : j \to i} \colim_{b : k \to j}
f_b^{-1}\mathcal{F}_j(u_{a \circ b}(X_i))
$$
% P01-ERRATA: P01-E0081 -- the left side uses F_i on an object of C_j; statement and typing require F_j.
由 (\ref{sites-equation-compute-pullback-to-limit})，内层余极限等于
$f_j^{-1}\mathcal{F}_j(u_i(X_i))$，故由引理
\ref{sites-lemma-directed-colimits-sections} 得到结论。
\end{proof}

\begin{lemma}
\label{sites-lemma-colimit-push-pull}
在情形 \ref{sites-situation-inverse-limit-sites} 中，假设有
$\mathcal{C}$ 上的层 $\mathcal{F}$。则
$$
\mathcal{F} = \colim f_i^{-1}f_{i, *}\mathcal{F}
$$
其中对 $a : j \to i$，转移映射是 $f_j^{-1}\varphi_a$；这里
$\varphi_a : f_a^{-1}f_{i, *}\mathcal{F} \to f_{j, *}\mathcal{F}$
是典范映射，满足与引理 \ref{sites-lemma-colimit} 中相同的余圈条件。
\end{lemma}

\begin{proof}
对态射
$$
\varphi_a : f_a^{-1}f_{i, *}\mathcal{F} \to f_{j, *}\mathcal{F}
$$
我们取恒等映射
$$
f_{i, *}\mathcal{F} \to f_{a, *}f_{j, *}\mathcal{F}
$$
的伴随。因此 $\varphi_a$ 是伴随对
$(f_a^{-1}, f_{a, *})$ 的余单位。必须证明，对所有
$a : j \to i$ 和 $b : k \to i$，若复合 $c = a \circ b$，则
$\varphi_c = \varphi_b \circ f_b^{-1}\varphi_a$。
% P01-ERRATA: P01-E0082 -- b is typed k -> i, so c = a \circ b is undefined; b must target j.
这由 Categories，引理 \ref{categories-lemma-compose-counits} 得出。最后，必须
证明由伴随给出的映射
$$
\colim_{i \in I} f_i^{-1}f_{i, *}\mathcal{F}
\longrightarrow
\mathcal{F}
$$
是同构。对 $\mathcal{C}$ 的对象 $U$，需要证明映射
$$
(\colim_{i \in I} f_i^{-1}\mathcal{F}_i)(U) \to \mathcal{F}(U)
$$
是双射。选择 $i$ 和 $\mathcal{C}_i$ 的对象 $U_i$，使得
$u_i(U_i) = U$。则由引理 \ref{sites-lemma-colimit}，左边等于
% P01-ERRATA: P01-E0083 -- F_i is undefined in this lemma; both displayed occurrences should be f_{i,*}F.
$$
(\colim_{i \in I} f_i^{-1}\mathcal{F}_i)(U) =
\colim_{a : j \to i} f_{j, *}\mathcal{F}(u_a(U_i))
$$
因为 $u_j(u_a(U_i)) = U$，根据定义，对所有
$a : j \to i$ 都有
$f_{j, *}\mathcal{F}(u_a(U_i)) = \mathcal{F}(U)$。因此该余极限的值是
$\mathcal{F}(U)$，证明完成。
\end{proof}












\section{预层的进一步函子性}
\label{sites-section-more-functoriality-PSh}

\noindent
本节重新讨论第 \ref{sites-section-functoriality-PSh} 节的内容。设
$u : \mathcal{C} \to \mathcal{D}$ 为范畴之间的函子。回忆
$$
u^p :
\textit{PSh}(\mathcal{D})
\longrightarrow
\textit{PSh}(\mathcal{C})
$$
是这样的函子：它把 $\mathcal{D}$ 上的 $\mathcal{G}$ 送到预层
$u^p\mathcal{G} = \mathcal{G} \circ u$。事实表明，该函子不仅有左伴随
（即 $u_p$），而且还有右伴随。

\medskip\noindent
具体地，对任意 $V \in \Ob(\mathcal{D})$，定义范畴
${}_V\mathcal{I} = {}_V^u\mathcal{I}$。其对象是序对
$(U, \psi : u(U) \to V)$。注意，这里的箭头方向与第
\ref{sites-section-functoriality-PSh} 节定义范畴 $\mathcal{I}_V^u$ 时所用的
箭头方向相反。态射 $(U, \psi) \to (U', \psi')$ 由满足
$\psi = \psi' \circ u(\alpha)$ 的态射 $\alpha : U \to U'$ 给出。
此外，给定 $\mathcal{C}$ 上任意集合预层 $\mathcal{F}$，引入函子
${}_V\mathcal{F} : {}_V\mathcal{I}^{opp} \to \textit{Sets}$，由规则
${}_V\mathcal{F}(U, \psi) = \mathcal{F}(U)$ 定义。令
$$
{}_pu(\mathcal{F})(V) := \lim_{{}_V\mathcal{I}^{opp}} {}_V\mathcal{F}
$$
作为极限，对 ${}_V\mathcal{I}$ 的每个对象 $(U, \psi)$，都有投影映射
$c(\psi) : {}_pu(\mathcal{F})(V) \to \mathcal{F}(U)$。事实上，
$$
{}_pu(\mathcal{F})(V)
=
\left\{
\begin{matrix}
\text{族 }
s_{(U, \psi)} \in \mathcal{F}(U) \\
\forall \beta : (U_1, \psi_1) \to (U_2, \psi_2)
\text{ 于 }{}_V\mathcal{I} \\
\text{ 有 } \beta^*s_{(U_2, \psi_2)} = s_{(U_1, \psi_1)}
\end{matrix}
\right\}
$$
其中对应由 $s \mapsto s_{(U, \psi)} = c(\psi)(s)$ 给出。留给读者定义与
$\mathcal{D}$ 的任意态射 $V' \to V$ 相伴的限制映射
${}_pu(\mathcal{F})(V) \to {}_pu(\mathcal{F})(V')$。所得预层记为
${}_pu\mathcal{F}$。

\begin{lemma}
\label{sites-lemma-recover-pu}
存在与限制映射相容的典范映射
${}_pu\mathcal{F}(u(U)) \to \mathcal{F}(U)$。
\end{lemma}

\begin{proof}
这就是上面的投影映射 $c(\text{id}_{u(U)})$。
\end{proof}

\noindent
注意，预层的任意映射 $\mathcal{F} \to \mathcal{F}'$ 都产生函子间相容的
映射系 ${}_V\mathcal{F} \to {}_V\mathcal{F}'$，从而产生预层的映射
${}_pu\mathcal{F} \to {}_pu\mathcal{F}'$。换言之，我们定义了函子
$$
{}_pu :
\textit{PSh}(\mathcal{C})
\longrightarrow
\textit{PSh}(\mathcal{D})
$$

\begin{lemma}
\label{sites-lemma-adjoints-pu}
函子 ${}_pu$ 是函子 $u^p$ 的右伴随。换言之，公式
$$
\Mor_{\textit{PSh}(\mathcal{C})}(u^p\mathcal{G}, \mathcal{F})
=
\Mor_{\textit{PSh}(\mathcal{D})}(\mathcal{G}, {}_pu\mathcal{F})
$$
关于 $\mathcal{F}$ 和 $\mathcal{G}$ 双函子地成立。
\end{lemma}

\begin{proof}
证明与引理 \ref{sites-lemma-adjoints-u} 的证明完全相同。注意，引理
\ref{sites-lemma-recover-pu} 中的映射
$u^p{}_pu \mathcal{F} \to \mathcal{F}$ 正是用于从右边得到左边的映射。

\medskip\noindent
另一种证法是：把 $\mathcal{C}$ 上的集合预层 $\mathcal{F}$ 看作
$\mathcal{C}^{opp}$ 上取值于 $\textit{Sets}^{opp}$ 的预层
$\mathcal{F}'$，并对 $\mathcal{D}$ 作类似处理。验证
$({}_pu \mathcal{F})' = u_p(\mathcal{F}')$ 且
$(u^p\mathcal{G})' = u^p(\mathcal{G}')$。由备注
\ref{sites-remark-functoriality-presheaves-values}，对取值于
$\textit{Sets}^{opp}$ 的预层，$u_p$ 与 $u^p$ 互为伴随。结论由此形式地
得到。
\end{proof}

\noindent
因此，给定范畴的函子 $u : \mathcal{C} \to \mathcal{D}$，可在预层范畴
之间得到函子序列
$$
u_p, u^p, {}_pu
$$
其中每一对相邻函子中，前者是后者的左伴随。

\begin{lemma}
\label{sites-lemma-adjoint-functors}
设 $u : \mathcal{C} \to \mathcal{D}$ 和
$v : \mathcal{D} \to \mathcal{C}$ 为范畴的函子。假设 $v$ 是 $u$ 的
右伴随。则：
\begin{enumerate}
\item 对 $\mathcal{D}$ 中任意 $V$，有 $u^ph_V = h_{v(V)}$；
\item 范畴 $\mathcal{I}^v_U$ 有始对象；
\item 范畴 ${}_V^u\mathcal{I}$ 有终对象；
\item ${}_pu = v^p$；
\item $u^p = v_p$。
\end{enumerate}
\end{lemma}

\begin{proof}
(1) 的证明。设 $V$ 为 $\mathcal{D}$ 的对象。由假设，
$u^ph_V = h_{v(V)}$，因为
$u^ph_V(U) = \Mor_\mathcal{D}(u(U), V) = \Mor_\mathcal{C}(U, v(V))$。

\medskip\noindent
(2) 的证明。设 $U$ 为 $\mathcal{C}$ 的对象。令
$\eta : U \to v(u(U))$ 为映射 $\text{id} : u(U) \to u(U)$ 的伴随映射。
断言 $(u(U), \eta)$ 是 $\mathcal{I}_U^v$ 的始对象。事实上，给定
$\mathcal{I}_U^v$ 的对象 $(V, \phi : U \to v(V))$，态射 $\phi$ 的伴随是
映射 $\psi : u(U) \to V$，后者定义态射
$(u(U), \eta) \to (V, \phi)$。

\medskip\noindent
(3) 的证明。设 $V$ 为 $\mathcal{D}$ 的对象。令
$\xi : u(v(V)) \to V$ 为映射 $\text{id} : v(V) \to v(V)$ 的伴随映射。
断言 $(v(V), \xi)$ 是 ${}_V^u\mathcal{I}$ 的终对象。事实上，给定
${}_V^u\mathcal{I}$ 的对象 $(U, \psi : u(U) \to V)$，态射 $\psi$ 的
伴随是映射 $\phi : U \to v(V)$，后者定义态射
$(U, \psi) \to (v(V), \xi)$。

\medskip\noindent
因此，对 $\mathcal{C}$ 上任意预层 $\mathcal{F}$，有
\begin{eqnarray*}
v^p\mathcal{F}(V)
& = &
\mathcal{F}(v(V)) \\
& = &
\Mor_{\textit{PSh}(\mathcal{C})}(h_{v(V)}, \mathcal{F}) \\
& = &
\Mor_{\textit{PSh}(\mathcal{C})}(u^ph_V, \mathcal{F}) \\
& = &
\Mor_{\textit{PSh}(\mathcal{D})}(h_V, {}_pu\mathcal{F}) \\
& = &
{}_pu\mathcal{F}(V)
\end{eqnarray*}
这证明了 (4)。由伴随函子的唯一性，(5) 随之成立。
\end{proof}





\section{余连续函子}
\label{sites-section-cocontinuous-functors}

\noindent
还有另一种构造托扑斯态射的方法。它使用位点之间的余连续函子，定义如下。

\begin{definition}
\label{sites-definition-cocontinuous}
设 $\mathcal{C}$ 和 $\mathcal{D}$ 为位点，并设
$u : \mathcal{C} \to \mathcal{D}$ 为函子。若对每个
$U \in \Ob(\mathcal{C})$ 以及 $\mathcal{D}$ 的每个覆盖
$\{V_j \to u(U)\}_{j \in J}$，都存在 $\mathcal{C}$ 的覆盖
$\{U_i \to U\}_{i\in I}$，使得态射族
$\{u(U_i) \to u(U)\}_{i \in I}$ 加细覆盖
$\{V_j \to u(U)\}_{j \in J}$，则称函子 $u$ 是{\it 余连续的}。
\end{definition}

\noindent
注意，$\{u(U_i) \to u(U)\}_{i \in I}$ 一般{\it 不是}位点
$\mathcal{D}$ 的覆盖。

\begin{lemma}
\label{sites-lemma-pu-sheaf}
设 $\mathcal{C}$ 和 $\mathcal{D}$ 为位点，
$u : \mathcal{C} \to \mathcal{D}$ 余连续，并设 $\mathcal{F}$ 为
$\mathcal{C}$ 上的层。则 ${}_pu\mathcal{F}$ 是 $\mathcal{D}$ 上的层，
将其记为 ${}_su\mathcal{F}$。
\end{lemma}

\begin{proof}
令 $\{V_j \to V\}_{j \in J}$ 为位点 $\mathcal{D}$ 的覆盖。必须证明
$$
\xymatrix{
&\ \phantom{}_pu\mathcal{F}(V) \ar[r] &
\prod {}_pu\mathcal{F}(V_j) \ar@<1ex>[r] \ar@<-1ex>[r] &
\prod {}_pu\mathcal{F}(V_j \times_V V_{j'}) &
}
$$
是等化子图。因为 ${}_pu$ 是 $u^p$ 的右伴随，有
$$
{}_pu\mathcal{F}(V) =
\Mor_{\textit{PSh}(\mathcal{D})}(h_V, {}_pu\mathcal{F}) =
\Mor_{\textit{PSh}(\mathcal{C})}(u^ph_V, \mathcal{F}) =
\Mor_{\Sh(\mathcal{C})}((u^ph_V)^\#, \mathcal{F})
$$
因此，只需证明
\begin{equation}
\label{sites-equation-coequalizer}
\xymatrix{
\coprod u^p h_{V_j \times_V V_{j'}} \ar@<1ex>[r] \ar@<-1ex>[r] &
\coprod u^p h_{V_j} \ar[r] &
u^p h_V
}
\end{equation}
在层化后成为余等化子图。（回忆一下，层范畴中的余积是预层范畴中余积的
层化，见引理 \ref{sites-lemma-colimit-sheaves}。）

\medskip\noindent
先证明 (\ref{sites-equation-coequalizer}) 的第二个箭头在层化后成为满射。为此使用
引理 \ref{sites-lemma-mono-epi-sheaves}。因此，只需证明：$u^ph_V$ 在 $U$ 上的
截面 $s$ 可在 $U$ 的某个覆盖的各成员上提升为
$\coprod u^p h_{V_j}$ 的截面。注意，$s$ 是态射
$s : u(U) \to V$。于是
$\{V_j \times_{V, s} u(U) \to u(U)\}$ 是 $\mathcal{D}$ 的覆盖。因为 $u$
余连续，存在覆盖 $\{U_i \to U\}$，使得
$\{u(U_i) \to u(U)\}$ 加细
$\{V_j \times_{V, s} u(U) \to u(U)\}$。这意味着每个限制
$s|_{U_i} : u(U_i) \to V$ 都经由某个 $j$ 的态射
$s_i : u(U_i) \to V_j$ 分解；换言之，$s|_{U_i}$ 按要求在
$u^ph_{V_j}(U_i) \to u^ph_V(U_i)$ 的像中。

\medskip\noindent
设 $s, s' \in (\coprod u^ph_{V_j})^\#(U)$ 映到
$(u^ph_V)^\#(U)$ 中同一元素。为完成引理的证明，说明把 $U$ 替换为某个
覆盖的各成员后，$s, s'$ 是 $\coprod u^p h_{V_j \times_V V_{j'}}$ 的同一
截面在 (\ref{sites-equation-coequalizer}) 的两个映射下的像。先把 $U$ 替换为
某个覆盖的各成员，从而可以假设 $s \in u^ph_{V_j}(U)$ 且
$s' \in u^ph_{V_{j'}}(U)$。再次作这样的替换，可以保证 $s$ 和 $s'$ 在
$u^ph_V(U)$ 中，而不只是在层化中，具有相同的像。因此
$s : u(U) \to V_j$ 和 $s' : u(U) \to V_{j'}$ 是 $\mathcal{D}$ 的态射，
使得
$$
\xymatrix{
u(U) \ar[r]_{s'} \ar[d]_s & V_{j'} \ar[d] \\
V_j \ar[r] & V
}
$$
交换。于是得到 $t = (s, s') : u(U) \to V_j \times_V V_{j'}$，即截面
$t \in u^ph_{V_j \times_V V_{j'}}(U)$，它按要求映到 $s, s'$。
\end{proof}

\begin{lemma}
\label{sites-lemma-exact-cocontinuous}
设 $\mathcal{C}$ 和 $\mathcal{D}$ 为位点，
$u : \mathcal{C} \to \mathcal{D}$ 余连续。函子
$\Sh(\mathcal{D}) \to \Sh(\mathcal{C})$，
$\mathcal{G} \mapsto (u^p\mathcal{G})^\#$ 是上面引理
\ref{sites-lemma-pu-sheaf} 中引入的函子 ${}_su$ 的左伴随。此外，它是正合的。
\end{lemma}

\begin{proof}
如下证明伴随性质：
\begin{eqnarray*}
\Mor_{\Sh(\mathcal{C})}
((u^p\mathcal{G})^\#, \mathcal{F})
& = &
\Mor_{\textit{PSh}(\mathcal{C})}
(u^p\mathcal{G}, \mathcal{F}) \\
& = &
\Mor_{\textit{PSh}(\mathcal{D})}
(\mathcal{G}, {}_pu\mathcal{F}) \\
& = &
\Mor_{\Sh(\mathcal{D})}
(\mathcal{G}, {}_su\mathcal{F}).
\end{eqnarray*}
因此它是左伴随，从而右正合，见 Categories，引理
\ref{categories-lemma-exact-adjoint}。已经看到层化是左正合的，见引理
\ref{sites-lemma-sheafification-exact}。此外，由引理
\ref{sites-lemma-limit-sheaf}，包含函子
$i : \Sh(\mathcal{D}) \to \textit{PSh}(\mathcal{D})$ 左正合。最后，函子
$u^p$ 左正合，因为它是右伴随（即 $u_p$ 的右伴随）。因此，该函子是左正合
函子之复合 ${}^\# \circ u^p \circ i$，所以左正合。
\end{proof}

\noindent
以一个技术性引理结束本节。

\begin{lemma}
\label{sites-lemma-technical-pu}
在引理 \ref{sites-lemma-exact-cocontinuous} 的情形中。对 $\mathcal{D}$ 上任意
预层 $\mathcal{G}$，有
$(u^p\mathcal{G})^\# = (u^p(\mathcal{G}^\#))^\#$。
\end{lemma}
% P01-ERRATA: P01-E0084 -- the first sentence of the lemma is a fragment ending with a period.

\begin{proof}
对 $\mathcal{C}$ 上任意层 $\mathcal{F}$，有
\begin{eqnarray*}
\Mor_{\Sh(\mathcal{C})}((u^p(\mathcal{G}^\#))^\#, \mathcal{F})
& = &
\Mor_{\Sh(\mathcal{D})}(\mathcal{G}^\#, {}_su\mathcal{F}) \\
& = &
\Mor_{\Sh(\mathcal{D})}(\mathcal{G}^\#, {}_pu\mathcal{F}) \\
& = &
\Mor_{\textit{PSh}(\mathcal{D})}(\mathcal{G}, {}_pu\mathcal{F}) \\
& = &
\Mor_{\textit{PSh}(\mathcal{C})}(u^p\mathcal{G}, \mathcal{F}) \\
& = &
\Mor_{\Sh(\mathcal{C})}((u^p\mathcal{G})^\#, \mathcal{F})
\end{eqnarray*}
结论由米田引理得出。
\end{proof}

\begin{remark}
\label{sites-remark-cartesian-cocontinuous}
设 $u : \mathcal{C} \to \mathcal{D}$ 为范畴之间的函子。给定
$\mathcal{D}$ 中的态射 $g : u(U) \to V$ 和 $f : W \to V$，可以考虑函子
$$
\mathcal{C}^{opp} \longrightarrow \textit{Sets},\quad
T \longmapsto
\Mor_\mathcal{C}(T, U)
\times_{\Mor_\mathcal{D}(u(T), V)}
\Mor_\mathcal{D}(u(T), W)
$$
若该函子可表，以 $U \times_{g, V, f} W$ 表示 $\mathcal{C}$ 中相应的对象。
假设 $\mathcal{C}$ 和 $\mathcal{D}$ 是位点。考虑性质 $P$：对
$\mathcal{D}$ 的每个覆盖 $\{f_j : V_j \to V\}$ 和任意态射
$g : u(U) \to V$，有：
\begin{enumerate}
\item 对所有 $i$，$U \times_{g, V, f_i} V_i$ 存在；
\item $\{U \times_{g, V, f_i} V_i \to U\}$ 是 $\mathcal{C}$ 的覆盖。
\end{enumerate}
% P01-ERRATA: P01-E0085 -- the covering is indexed by j, but both property clauses switch to undefined i.
请注意这与连续函子定义的相似之处。若 $u$ 具有 $P$，则 $u$ 余连续
（细节从略）。我们将遇到的许多余连续函子都满足 $P$。
\end{remark}



\section{余连续函子与托扑斯态射}
\label{sites-section-cocontinuous-morphism-topoi}

\noindent
由上文清楚可知，余连续函子 $u$ 给出一个与 $u$ 同方向的托扑斯态射。
因此，这与定义 \ref{sites-definition-morphism-sites}、命题
\ref{sites-proposition-get-morphism} 和引理
\ref{sites-lemma-morphism-sites-topoi} 中在某些条件下与连续 $u$ 相伴的托扑斯态射
方向相反。

\begin{lemma}
\label{sites-lemma-cocontinuous-morphism-topoi}
设 $\mathcal{C}$ 和 $\mathcal{D}$ 为位点，并设
$u : \mathcal{C} \to \mathcal{D}$ 余连续。函子
$g_* = {}_su$ 和 $g^{-1} = (u^p\ )^\#$ 定义从
$\Sh(\mathcal{C})$ 到 $\Sh(\mathcal{D})$ 的托扑斯态射 $g$。
\end{lemma}

\begin{proof}
这正是引理 \ref{sites-lemma-exact-cocontinuous} 的内容。
\end{proof}

\begin{lemma}
\label{sites-lemma-composition-cocontinuous}
\begin{slogan}
位点函子的复合保持位点函子的余连续性。
\end{slogan}
设 $u : \mathcal{C} \to \mathcal{D}$ 和
$v : \mathcal{D} \to \mathcal{E}$ 为余连续函子。则
$v \circ u$ 余连续，并且 $h = g \circ f$，其中
$f : \Sh(\mathcal{C}) \to \Sh(\mathcal{D})$、相应地
$g : \Sh(\mathcal{D}) \to \Sh(\mathcal{E})$、相应地
$h : \Sh(\mathcal{C}) \to \Sh(\mathcal{E})$，分别是与 $u$、相应地
$v$、相应地 $v \circ u$ 相伴的托扑斯态射。
\end{lemma}

\begin{proof}
令 $U \in \Ob(\mathcal{C})$。设 $\{E_i \to v(u(U))\}$ 是
$\mathcal{E}$ 中 $U$ 的一个覆盖。根据假设，存在 $\mathcal{D}$ 中的覆盖
$\{D_j \to u(U)\}$，使得 $\{v(D_j) \to v(u(U))\}$ 加细
$\{E_i \to v(u(U))\}$。同样根据假设，存在 $\mathcal{C}$ 中的覆盖
$\{C_l \to U\}$，使得 $\{u(C_l) \to u(U)\}$ 加细
$\{D_j \to u(U)\}$。于是
$\{v(u(C_l)) \to v(u(U))\}$ 确实加细覆盖
$\{E_i \to v(u(U))\}$。这证明 $v \circ u$ 余连续。
% P01-ERRATA: P01-E0086 -- the source calls {E_i -> v(u(U))} a covering of U in E; U is not an E-object.
为证明最后的断言，只需说明
${}_sv \circ {}_su = {}_s(v \circ u)$。由引理 \ref{sites-lemma-pu-sheaf}，只需
证明 ${}_pv \circ {}_pu = {}_p(v \circ u)$。因为 ${}_pu$、相应地
${}_pv$、相应地 ${}_p(v \circ u)$，分别是 $u^p$、相应地 $v^p$、
相应地 $(v \circ u)^p$ 的右伴随，所以只需证明
$u^p \circ v^p = (v \circ u)^p$。这直接由定义得出。
\end{proof}

\begin{example}
\label{sites-example-open-immersion-cocontinuous}
设 $X$ 为拓扑空间，$j : U  \to X$ 为开子空间的包含映射。回忆一下，
我们有位点 $X_{Zar}$ 和 $U_{Zar}$，见例
\ref{sites-example-site-topological}。还记得，与 $j$ 相伴的函子
$u : X_{Zar} \to U_{Zar}$ 连续，并产生位点态射
$U_{Zar} \to X_{Zar}$，见例 \ref{sites-example-continuous-map}。它也给出托扑斯
态射 $(j_*, j^{-1})$。接着考虑函子
$v : U_{Zar} \to X_{Zar}$，$V \mapsto v(V) = V$（仍是同一个开集，但现在
把它看作 $X_{Zar}$ 的对象）。该函子余连续。事实上，若
$v(V) = \bigcup_{j \in J} W_j$ 是 $X$ 中的开覆盖，则每个 $W_j$ 必为
$U$ 的子集，因而形如 $v(V_j)$，而且显然
$V = \bigcup_{j \in J} V_j$ 是 $U$ 中的开覆盖。由上面的引理
\ref{sites-lemma-cocontinuous-morphism-topoi}，可知有与 $v$ 相伴的托扑斯态射
$$
\Sh(U) \longrightarrow \Sh(X)
$$
由 ${}_sv$ 和 $(v^p\ )^\#$ 给出。断言事实上
$(v^p\ )^\# = j^{-1}$ 且 ${}_sv = j_*$；换言之，这与上面给出的托扑斯
态射相同。看出这一点的最简便方法或许是注意：对 $X$ 上任意层
$\mathcal{G}$，有 $v^p\mathcal{G}(V) = \mathcal{G}(V)$；根据 Sheaves，引理
\ref{sheaves-lemma-j-pullback}，这描述了 $j^{-1}\mathcal{G}$（因此在此情形
无需层化）。${}_sv$ 与 $j_*$ 的相等由伴随函子的唯一性得出（也可直接
计算）。
\end{example}

\begin{example}
\label{sites-example-open-map-cocontinuous}
这是例 \ref{sites-example-open-immersion-cocontinuous} 的一个小推广。设
$f : X \to Y$ 为拓扑空间的连续映射，并假设 $f$ 是开映射。回忆一下，
我们有位点 $X_{Zar}$ 和 $Y_{Zar}$，见例
\ref{sites-example-site-topological}。还记得，与 $f$ 相伴的函子
$u : Y_{Zar} \to X_{Zar}$ 连续，并产生位点态射
$X_{Zar} \to Y_{Zar}$，见例 \ref{sites-example-continuous-map}。它也给出托扑斯
态射 $(f_*, f^{-1})$。接着考虑函子
$v : X_{Zar} \to Y_{Zar}$，$U \mapsto v(U) = f(U)$。该函子余连续。
事实上，若 $f(U) = \bigcup_{j \in J} V_j$ 是 $Y$ 中的开覆盖，则令
$U_j = f^{-1}(V_j) \cap U$，可得开覆盖 $U = \bigcup U_j$，使得
$f(U) = \bigcup f(U_j)$ 加细 $f(U) = \bigcup V_j$。由上面的引理
\ref{sites-lemma-cocontinuous-morphism-topoi}，可知有与 $v$ 相伴的托扑斯态射
$$
\Sh(X) \longrightarrow \Sh(Y)
$$
由 ${}_sv$ 和 $(v^p\ )^\#$ 给出。断言事实上
$(v^p\ )^\# = f^{-1}$ 且 ${}_sv = f_*$；换言之，这与上面给出的托扑斯
态射相同。对 $Y$ 上任意层 $\mathcal{G}$，有
$v^p\mathcal{G}(U) = \mathcal{G}(f(U))$。另一方面，可以计算
$u_p\mathcal{G}(U) = \colim_{f(U) \subset V} \mathcal{G}(V)
= \mathcal{G}(f(U))$，因为显然
$(f(U), U \subset f^{-1}(f(U)))$ 是第
\ref{sites-section-functoriality-PSh} 节中范畴 $\mathcal{I}_U^u$ 的始对象。
因此 $u_p = v^p$，从而 $f^{-1} = u_s = (v^p\ )^\#$。${}_sv$ 与 $f_*$
的相等由伴随函子的唯一性得出（也可直接计算）。
\end{example}

\noindent
在第一个例 \ref{sites-example-open-immersion-cocontinuous} 中，函子 $v$ 也连续。
但在第二个例 \ref{sites-example-open-map-cocontinuous} 中，它一般不连续，因为
定义 \ref{sites-definition-continuous} 的条件 (2) 可能不成立。因此下述引理适用于
第一个例，但不适用于第二个例。

\begin{lemma}
\label{sites-lemma-when-shriek}
\begin{slogan}
对同时连续且余连续的位点函子，相伴托扑斯态射中的层化是多余的。
\end{slogan}
设 $\mathcal{C}$ 和 $\mathcal{D}$ 为位点，并设
$u : \mathcal{C} \to \mathcal{D}$ 为函子。假设：
\begin{enumerate}
\item[(a)] $u$ 余连续；
\item[(b)] $u$ 连续。
\end{enumerate}
令 $g : \Sh(\mathcal{C}) \to \Sh(\mathcal{D})$ 为相伴的托扑斯态射。则：
\begin{enumerate}
\item 公式 $g^{-1} = (u^p\ )^\#$ 中无需层化，换言之，
$g^{-1}(\mathcal{G})(U) = \mathcal{G}(u(U))$；
\item $g^{-1}$ 有左伴随 $g_{!} = (u_p\ )^\#$；
\item $g^{-1}$ 与任意极限和余极限交换。
\end{enumerate}
\end{lemma}

\begin{proof}
由引理 \ref{sites-lemma-pushforward-sheaf}，对 $\mathcal{D}$ 上任意层
$\mathcal{G}$，预层 $u^p\mathcal{G}$ 是 $\mathcal{C}$ 上的层（记为
$u^s\mathcal{G}$）。由引理 \ref{sites-lemma-adjoint-sheaves}，函子
$\mathcal{F} \mapsto (u_p\mathcal{F})^\#$ 是 $u^s$ 的左伴随。关于极限和
余极限的断言由 Categories，第 \ref{categories-section-adjoint} 节中的讨论
得出。
\end{proof}

\noindent
在上述引理 \ref{sites-lemma-when-shriek} 的情形中，可知有伴随函子序列
$$
g_{!}, \ g^{-1}, \ g_*.
$$
函子 $g_!$ 一般{\it 不}正合，因为一般而言，它不把
$\Sh(\mathcal{C})$ 的终对象变为 $\Sh(\mathcal{D})$ 的终对象。见
Sheaves，备注 \ref{sheaves-remark-j-shriek-not-exact}。另一方面，在例
\ref{sites-example-open-immersion-cocontinuous} 的拓扑情形中，函子 $j_!$ 在
阿贝尔层上正合，见 Modules，引理 \ref{modules-lemma-j-shriek-exact}。
下述引理把它推广到位点情形。

\begin{lemma}
\label{sites-lemma-preserve-equalizers}
设 $\mathcal{C}$ 和 $\mathcal{D}$ 为位点，并设
$u : \mathcal{C} \to \mathcal{D}$ 为函子。假设：
\begin{enumerate}
\item[(a)] $u$ 余连续；
\item[(b)] $u$ 连续；
\item[(c)] $\mathcal{C}$ 中存在纤维积和等化子，并且 $u$ 与它们交换。
\end{enumerate}
此时上述函子 $g_!$ 与纤维积和等化子交换（更一般地，与有限连通极限
交换）。
\end{lemma}

\begin{proof}
假设 (a)、(b)、(c)。有 $g_! = (u_p\ )^\#$。回忆一下（引理
\ref{sites-lemma-limit-sheaf}），层的极限等于相应的预层极限，而层化与有限极限
交换（引理 \ref{sites-lemma-sheafification-exact}）。因此，只需说明 $u_p$ 与
纤维积和等化子交换。为此，只需第 \ref{sites-section-functoriality-PSh} 节中的
范畴 $(\mathcal{I}_V^u)^{opp}$ 上的余极限与纤维积和等化子交换。这由引理
\ref{sites-lemma-almost-directed} 和 Categories，引理
\ref{categories-lemma-almost-directed-commutes-equalizers} 得出。
\end{proof}

\noindent
下述引理处理一种甚至更像与拓扑空间开浸入相伴之态射的情形。

\begin{lemma}
\label{sites-lemma-back-and-forth}
设 $\mathcal{C}$ 和 $\mathcal{D}$ 为位点，并设
$u : \mathcal{C} \to \mathcal{D}$ 为函子。假设：
\begin{enumerate}
\item[(a)] $u$ 余连续；
\item[(b)] $u$ 连续；
\item[(c)] $u$ 全忠实。
\end{enumerate}
对上述 $g_!, g^{-1}, g_*$，任意 $\mathcal{C}$ 上的层 $\mathcal{F}$ 的
典范映射 $\mathcal{F} \to g^{-1}g_!\mathcal{F}$ 和
$g^{-1}g_*\mathcal{F} \to \mathcal{F}$ 都是同构。
\end{lemma}

\begin{proof}
设 $X$ 为 $\mathcal{C}$ 的对象。引理 \ref{sites-lemma-pu-sheaf} 和
\ref{sites-lemma-when-shriek} 已经表明，对函子
$g^{-1} = (u^p\ )^\#$ 和 $g_{*} = ({}_pu\ )^\#$，层化不是必需的。
可以计算
$(g^{-1}g_{*}\mathcal{F})(X) = g_{*}\mathcal{F}(u(X))
= \lim \mathcal{F}(Y)$。这里的极限遍历序对
$(Y, u(Y) \to u(X))$ 所成的范畴，其中不要求态射
$u(Y) \to u(X)$ 形如 $u(\alpha)$，$\alpha$ 为 $\mathcal{C}$ 的态射。
由假设 (c)，它们自动来自 $\mathcal{C}$ 的态射，因而该极限是在
$(X, u(\text{id}_X))$ 上的值，即 $\mathcal{F}(X)$。这证明
$g^{-1}g_{*}\mathcal{F} = \mathcal{F}$。

\medskip\noindent
另一方面，$(g^{-1}g_{!}\mathcal{F})(X) =
g_{!}\mathcal{F}(u(X)) = (u_p\mathcal{F})^\#(u(X))$，并且
$u_p\mathcal{F}(u(X)) = \colim \mathcal{F}(Y)$。这里的余极限遍历序对
$(Y, u(X) \to u(Y))$ 所成的范畴，其中不要求态射
$u(X) \to u(Y)$ 形如 $u(\alpha)$，$\alpha$ 为 $\mathcal{C}$ 的态射。
由假设 (c)，它们自动来自 $\mathcal{C}$ 的态射，因而该余极限是在
$(X, u(\text{id}_X))$ 上的值，即 $\mathcal{F}(X)$。因此，对每个
$X \in \Ob(\mathcal{C})$，有
$u_p\mathcal{F}(u(X)) = \mathcal{F}(X)$。由于 $u$ 余连续且连续，
$\mathcal{D}$ 中 $u(X)$ 的任意覆盖都可由 $\mathcal{D}$ 的一个覆盖
(!) $\{u(X_i) \to u(X)\}$ 加细，其中
$\{X_i \to X\}$ 是 $\mathcal{C}$ 中的覆盖。这也蕴含
$(u_p\mathcal{F})^+(u(X)) = \mathcal{F}(X)$，因为在定义
$(u_p\mathcal{F})^+$ 在 $u(X)$ 上之值的余极限中，可以限制到上述覆盖
$\{u(X_i) \to u(X)\}$ 所成的共尾系。因此，对 $\mathcal{C}$ 的所有对象
$X$，同样有 $(u_p\mathcal{F})^+(u(X)) = \mathcal{F}(X)$。再重复一次该
论证，便对 $\mathcal{C}$ 的所有对象 $X$ 得到
$(u_p\mathcal{F})^\#(u(X)) = \mathcal{F}(X)$。这产生所需等式
$g^{-1}g_!\mathcal{F} = \mathcal{F}$。
\end{proof}

\noindent
最后给出一种没有相应拓扑例子的情形。我们将用该引理考察在保持相同拓扑的
同时扩大概形的“部分宇宙”时会发生什么。在引理的情形中，托扑斯态射
$g : \Sh(\mathcal{C}) \to \Sh(\mathcal{D})$ 把
$\Sh(\mathcal{C})$ 识别为 $\Sh(\mathcal{D})$ 的子托扑斯（第
\ref{sites-section-subtopoi} 节），而且给定嵌入还有一个收缩。

\begin{lemma}
\label{sites-lemma-bigger-site}
设 $\mathcal{C}$ 和 $\mathcal{D}$ 为位点，并设
$u : \mathcal{C} \to \mathcal{D}$ 为函子。假设：
\begin{enumerate}
\item[(a)] $u$ 余连续；
\item[(b)] $u$ 连续；
\item[(c)] $u$ 全忠实；
\item[(d)] $\mathcal{C}$ 中存在纤维积，且 $u$ 与纤维积交换；
\item[(e)] 存在终对象
$e_\mathcal{C} \in \Ob(\mathcal{C})$、
$e_\mathcal{D} \in \Ob(\mathcal{D})$，使得
$u(e_\mathcal{C}) = e_\mathcal{D}$。
\end{enumerate}
令 $g_!, g^{-1}, g_*$ 如上。则 $u$ 定义位点态射
$f : \mathcal{D} \to \mathcal{C}$，并且 $f_* = g^{-1}$、
$f^{-1} = g_!$。复合
$$
\xymatrix{
\Sh(\mathcal{C}) \ar[r]^g &
\Sh(\mathcal{D}) \ar[r]^f &
\Sh(\mathcal{C})
}
$$
同构于托扑斯 $\Sh(\mathcal{C})$ 的恒等态射。此外，函子 $f^{-1}$
全忠实。
\end{lemma}

\begin{proof}
由假设，函子 $u$ 满足命题 \ref{sites-proposition-get-morphism} 的假设。因此，
$u$ 定义位点态射，从而定义引理 \ref{sites-lemma-morphism-sites-topoi} 中的
托扑斯态射 $f$。公式 $f_* = g^{-1}$ 和 $f^{-1} = g_!$ 由所引引理及引理
\ref{sites-lemma-when-shriek} 清楚可得。由引理
\ref{sites-lemma-back-and-forth}，有
$f_* \circ g_* = g^{-1} \circ g_* \cong \text{id}$，以及
$g^{-1} \circ f^{-1} = g^{-1} \circ g_! \cong \text{id}$。

\medskip\noindent
还须证明 $f^{-1}$ 全忠实。令
$\mathcal{F}, \mathcal{G} \in \Ob(\Sh(\mathcal{C}))$。必须证明映射
$$
\Mor_{\Sh(\mathcal{C})}(\mathcal{F}, \mathcal{G})
\longrightarrow
\Mor_{\Sh(\mathcal{D})}(f^{-1}\mathcal{F}, f^{-1}\mathcal{G})
$$
是双射。但右边等于
\begin{align*}
\Mor_{\Sh(\mathcal{D})}(f^{-1}\mathcal{F}, f^{-1}\mathcal{G})
& =
\Mor_{\Sh(\mathcal{C})}(\mathcal{F}, f_*f^{-1}\mathcal{G}) \\
& =
\Mor_{\Sh(\mathcal{C})}(\mathcal{F}, g^{-1}f^{-1}\mathcal{G}) \\
& =
\Mor_{\Sh(\mathcal{C})}(\mathcal{F}, \mathcal{G})
\end{align*}
（第一个等号由伴随得到），这证明了所需结论。
\end{proof}

\begin{example}
\label{sites-example-closed-map-cocontinuous-false}
设 $X$ 为拓扑空间，$i : Z  \to X$ 为一个子集（赋以诱导拓扑）的包含映射。
考虑函子 $u : X_{Zar} \to Z_{Zar}$，
$U \mapsto u(U) = Z \cap U$。乍看之下，该函子似乎也余连续。毕竟，既然
$Z$ 带有诱导拓扑，$U\cap Z$ 的任意覆盖难道不应来自 $X$ 中 $U$ 的某个
覆盖吗？并非如此！例如，若 $U \cap Z = \emptyset$ 会怎样？此时空覆盖是
$U \cap Z$ 的覆盖，而空覆盖只能由空覆盖加细。因此得到
$u$ 余连续 $\Rightarrow$ $X$ 的每个非空开集 $U$ 都与 $Z$ 有非空交。
% P01-ERRATA: P01-E0087 -- source has the malformed phrase "shouldn't any covering ... it come".
% P01-ERRATA: P01-E0088 -- source omits the verb in "$u$ cocontinuous".
但这并不充分。例如，若 $X = \mathbf{R}$ 是带通常拓扑的实数直线，且
$Z = \mathbf{R} \setminus \{0\}$，则 $Z$ 有开覆盖
$Z = \{x < 0\} \cup \bigcup_n \{1/n < x\}$，它不可能由 $X$ 的任何开覆盖
之限制来加细。
\end{example}







\section{具有右伴随的余连续函子}
\label{sites-section-cocontinuous-adjoint}

\noindent
设 $\mathcal{C}$ 和 $\mathcal{D}$ 为位点。设
$u : \mathcal{C} \to \mathcal{D}$ 和
$v : \mathcal{D} \to \mathcal{C}$ 为底层范畴的函子，并且 $v$ 是 $u$ 的
右伴随。在此情形中，若 $v$ 连续，则 $u$ 余连续（引理
\ref{sites-lemma-continuous-with-continuous-left-adjoint}）。若 $u$ 余连续，则常常
（但并非总是如此，见例 \ref{sites-example-not-equivalent}）$v$ 连续；若确实如此，
则 $v$ 定义一个位点态射，其相伴托扑斯态射与 $u$ 所定义的相同。

\begin{lemma}
\label{sites-lemma-have-functor-other-way}
设 $\mathcal{C}$ 和 $\mathcal{D}$ 为位点，设
$u : \mathcal{C} \to \mathcal{D}$ 和
$v : \mathcal{D} \to \mathcal{C}$ 为函子。假设 $u$ 余连续，且 $v$ 是
$u$ 的右伴随。令 $g : \Sh(\mathcal{C}) \to \Sh(\mathcal{D})$ 为与 $u$
相伴的托扑斯态射，见引理 \ref{sites-lemma-cocontinuous-morphism-topoi}。则：
\begin{enumerate}
\item 对 $\mathcal{C}$ 上的层 $\mathcal{F}$，层 $g_*\mathcal{F}$ 等于
预层 $v^p\mathcal{F}$；换言之，
$(g_*\mathcal{F})(V) = \mathcal{F}(v(V))$；
\item 对 $\mathcal{D}$ 上的层 $\mathcal{G}$，有
$g^{-1}\mathcal{G} = (v_p\mathcal{G})^\#$。
\end{enumerate}
\end{lemma}

\begin{proof}
对 (1) 中的 $\mathcal{F}$，有
$$
g_*\mathcal{F} =
{}_su\mathcal{F} =
{}_pu\mathcal{F} =
v^p\mathcal{F} =
\mathcal{F} \circ v
$$
第一个等号是引理 \ref{sites-lemma-cocontinuous-morphism-topoi}。第二个等号是
引理 \ref{sites-lemma-pu-sheaf}。第三个等号是引理
\ref{sites-lemma-adjoint-functors}。最后一个等号是第
\ref{sites-section-functoriality-PSh} 节中 $v^p$ 的定义。这证明 (1)。
对 (2) 中的 $\mathcal{G}$，有
$$
g^{-1}\mathcal{G} = (u^p\mathcal{G})^\# = (v_p\mathcal{G})^\#
$$
第一个等号是引理 \ref{sites-lemma-cocontinuous-morphism-topoi}。第二个等号是
引理 \ref{sites-lemma-adjoint-functors}。
\end{proof}

\begin{lemma}
\label{sites-lemma-have-functor-other-way-morphism}
记号与假设同引理 \ref{sites-lemma-have-functor-other-way}。若另外 $v$ 连续，则
$v$ 定义位点态射 $f : \mathcal{C} \to \mathcal{D}$，其相伴托扑斯态射
等于 $g$。
\end{lemma}

\begin{proof}
使用引理 \ref{sites-lemma-have-functor-other-way} 的结果而不再另行说明。为证明
$v$ 定义引理所述的位点态射 $f$，必须说明 $v_s$ 是正合函子（见定义
\ref{sites-definition-morphism-sites}）。由于
$v_s\mathcal{G} = (v_p\mathcal{G})^\# = g^{-1}\mathcal{G}$，这由 $g$ 是
托扑斯态射这一事实得出。于是 $f^{-1} = v_s = g^{-1}$，从而发现 $f = g$
作为托扑斯态射成立。
\end{proof}

\begin{example}
\label{sites-example-finer-topology-bis}
本例延续例 \ref{sites-example-finer-topology} 的讨论，并沿用其中的记号
$\mathcal{C}, \tau, \tau', \epsilon$。注意，恒等函子
$v : \mathcal{C}_{\tau'} \to \mathcal{C}_\tau$ 连续，而恒等函子
$u : \mathcal{C}_\tau \to \mathcal{C}_{\tau'}$ 余连续。此外，$u$ 是 $v$ 的
左伴随。因此，引理 \ref{sites-lemma-have-functor-other-way} 和
\ref{sites-lemma-have-functor-other-way-morphism} 的结果适用，从而 $v$ 定义位点
态射，即
$$
\epsilon : \mathcal{C}_\tau \longrightarrow \mathcal{C}_{\tau'}
$$
其相应托扑斯态射与余连续函子 $u$ 所伴随的托扑斯态射相同。
\end{example}

\begin{lemma}
\label{sites-lemma-continuous-with-continuous-left-adjoint}
设 $\mathcal{C}$ 和 $\mathcal{D}$ 为位点，并设
$v : \mathcal{D} \to \mathcal{C}$ 为连续函子。假设 $v$ 有左伴随
$u : \mathcal{C} \to \mathcal{D}$。则：
\begin{enumerate}
\item $u$ 余连续；
\item 引理 \ref{sites-lemma-have-functor-other-way} 和
\ref{sites-lemma-have-functor-other-way-morphism} 的结论成立。
\end{enumerate}
特别地，$v$ 定义位点态射 $f : \mathcal{C} \to \mathcal{D}$。
\end{lemma}

\begin{proof}
令 $U$ 为 $\mathcal{C}$ 的对象，并令
$\{V_j \to u(U)\}_{j \in J}$ 为 $\mathcal{D}$ 中的覆盖。则
$\{v(V_j) \to v(u(U))\}_{i \in I}$ 是 $\mathcal{C}$ 中的覆盖。通过伴随
映射 $U \to v(u(U))$，可以基变换为覆盖
$\{W_j \to U\}_{j \in J}$，其中
$W_j = v(V_j) \times_{v(u(U))} U$。以 $p_j : W_j \to v(V_j)$ 表示第一
投影，得到映射
% P01-ERRATA: P01-E0089 -- the pulled-forward covering switches from j in J to undefined i in I.
$$
u(W_j) \xrightarrow{u(p_j)} u(v(V_j)) \longrightarrow V_j
$$
其中第二个箭头是伴随映射。这确定具有固定靶的态射族之态射
$\{u(W_j) \to u(U)\} \to \{V_j \to u(U)\}$，表明 $u$ 的确余连续。
其余断言由引理 \ref{sites-lemma-have-functor-other-way} 和
\ref{sites-lemma-have-functor-other-way-morphism} 立即得出。
\end{proof}

\begin{example}
\label{sites-example-not-equivalent}
设 $\mathcal{C}$ 和 $\mathcal{D}$ 为位点。设
$u : \mathcal{C} \to \mathcal{D}$ 和
$v : \mathcal{D} \to \mathcal{C}$ 为底层范畴的函子，并且 $v$ 是 $u$ 的
右伴随。引理 \ref{sites-lemma-continuous-with-continuous-left-adjoint} 表明，若
$v$ 连续，则 $u$ 余连续。反过来，若 $u$ 余连续，不能由此断定 $v$ 连续。
% P01-ERRATA: P01-E0090 -- source misspells "cocontinuous" as "cocontinous".
将用概形的大 \'etale 位点和平滑位点给出这一现象的例子，不过大概也有
初等例子。具体而言，考虑概形 $S$ 以及位点
$(\Sch/S)_\etale$ 和 $(\Sch/S)_{smooth}$。可以假设这些位点具有同一底层
范畴，见 Topologies，备注 \ref{topologies-remark-choice-sites}。令
$u = v = \text{id}$。则作为从 $(\Sch/S)_\etale$ 到
$(\Sch/S)_{smooth}$ 的函子，$u$ 余连续，因为概形的每个平滑覆盖都可由
\'etale 覆盖加细，见 More on Morphisms，引理
\ref{more-morphisms-lemma-etale-dominates-smooth}。反之，从
$(\Sch/S)_{smooth}$ 到 $(\Sch/S)_\etale$ 的函子 $v$ 不连续，因为平滑
覆盖一般不是 \'etale 覆盖。
\end{example}





\section{具有左伴随的余连续函子}
\label{sites-section-cocontinuous-left-adjoint}

\noindent
余连续函子 $u$ 可能有左伴随 $w$。

\begin{lemma}
\label{sites-lemma-have-left-adjoint}
设 $\mathcal{C}$ 和 $\mathcal{D}$ 为位点。令
$g : \Sh(\mathcal{C}) \to \Sh(\mathcal{D})$ 为与连续且余连续的函子
$u : \mathcal{C} \to \mathcal{D}$ 相伴的托扑斯态射，见引理
\ref{sites-lemma-cocontinuous-morphism-topoi} 和 \ref{sites-lemma-when-shriek}。
\begin{enumerate}
\item 若 $w : \mathcal{D} \to \mathcal{C}$ 是 $u$ 的左伴随，则：
\begin{enumerate}
\item $g_!\mathcal{F}$ 是与预层 $w^p\mathcal{F}$ 相伴的层；
\item $g_!$ 正合。
\end{enumerate}
\item 若 $w$ 是连续的左伴随，则 $g_!$ 有左伴随。
\item 若 $w$ 是余连续的左伴随，则 $g_! = h^{-1}$ 且
$g^{-1} = h_*$，其中
$h : \Sh(\mathcal{D}) \to \Sh(\mathcal{C})$ 是与 $w$ 相伴的托扑斯态射。
\end{enumerate}
\end{lemma}

\begin{proof}
回忆一下，$g_!\mathcal{F}$ 是 $u_p\mathcal{F}$ 的层化。因此，(1)(a) 由
引理 \ref{sites-lemma-adjoint-functors} 中 $u_p = w^p$ 这一事实得出。

\medskip\noindent
为说明 (1)(b)，注意 $g_!$ 作为左伴随与所有余极限交换（Categories，引理
\ref{categories-lemma-adjoint-exact}）。令
$i \mapsto \mathcal{F}_i$ 为 $\Sh(\mathcal{C})$ 中的有限图表。则
$\lim \mathcal{F}_i$ 在预层范畴中计算（引理
\ref{sites-lemma-limit-sheaf}）。因为 $w^p$ 是右伴随（引理
\ref{sites-lemma-adjoints-u}），有
$w^p \lim \mathcal{F}_i = \lim w^p\mathcal{F}_i$。层化正合（引理
\ref{sites-lemma-sheafification-exact}），故结论由 (1)(a) 得出。

\medskip\noindent
假设 $w$ 连续。则 $g_! = (w^p\ )^\# = w^s$，但层化不是必需的，并且
$g_!$ 有左伴随 $w_s$，见引理 \ref{sites-lemma-pushforward-sheaf} 和
\ref{sites-lemma-adjoint-sheaves}。

\medskip\noindent
假设 $w$ 余连续。等式 $g_! = h^{-1}$ 由 (1)(a) 和定义得出。等式
$g^{-1} = h_*$ 由 $g_! = h^{-1}$ 和伴随函子的唯一性得出。或者，也可由
引理 \ref{sites-lemma-have-functor-other-way} 推出。
\end{proof}









\section{下惊叹号的存在性}
\label{sites-section-lower-shriek}

\noindent
本节讨论托扑斯态射 $f$ 的一些情形；在这些情形中，$f^{-1}$ 有左伴随
$f_!$。

\begin{lemma}
\label{sites-lemma-existence-lower-shriek}
设 $\mathcal{C}$、$\mathcal{D}$ 为两个位点。设
$f : \Sh(\mathcal{D}) \to \Sh(\mathcal{C})$ 为托扑斯态射。设
$E \subset \Ob(\mathcal{D})$ 为子集，满足：
\begin{enumerate}
\item 对 $V \in E$，存在 $\mathcal{C}$ 上的层 $\mathcal{G}$，使得
$f^{-1}\mathcal{F}(V) =
\Mor_{\Sh(\mathcal{C})}(\mathcal{G}, \mathcal{F})$ 对
$\Sh(\mathcal{C})$ 中的 $\mathcal{F}$ 函子地成立；
\item $\mathcal{D}$ 的每个对象都有一个由 $E$ 中对象组成的覆盖。
\end{enumerate}
则 $f^{-1}$ 有左伴随 $f_!$。
\end{lemma}

\begin{proof}
由米田引理（Categories，引理 \ref{categories-lemma-yoneda}），与
$V \in E$ 对应的层 $\mathcal{G}_V$ 由公式
$f^{-1}\mathcal{F}(V) =
\Mor_{\Sh(\mathcal{C})}(\mathcal{G}_V, \mathcal{F})$ 在唯一同构意义下
确定。回忆一下，
$f^{-1}\mathcal{F}(V) =
\Mor_{\Sh(\mathcal{D})}(h_V^\#, f^{-1}\mathcal{F})$。以
$i_V : h_V^\# \to f^{-1}\mathcal{G}_V$ 表示与
$\Mor(\mathcal{G}_V, \mathcal{G}_V)$ 中的 $\text{id}$ 对应的映射。(1) 中的
函子性蕴含该双射由下式给出：
$$
\Mor_{\Sh(\mathcal{C})}(\mathcal{G}_V, \mathcal{F}) \to
\Mor_{\Sh(\mathcal{D})}(h_V^\#, f^{-1}\mathcal{F}),\quad
\varphi \mapsto f^{-1}\varphi \circ i_V
$$
对任意 $V_1, V_2 \in E$，存在典范映射
$$
\Mor_{\Sh(\mathcal{D})}(h^\#_{V_2}, h^\#_{V_1})
\to
\Hom_{\Sh(\mathcal{C})}(\mathcal{G}_{V_2}, \mathcal{G}_{V_1}),\quad
\varphi \mapsto f_!(\varphi)
$$
其特征为
$f^{-1}(f_!(\varphi)) \circ i_{V_2} = i_{V_1} \circ \varphi$。注意，
$\varphi \mapsto f_!(\varphi)$ 与复合相容；这可直接由上述刻画看出。因此，
$h_V^\# \mapsto \mathcal{G}_V$ 和
$\varphi \mapsto f_!\varphi$ 构成一个以 $\Sh(\mathcal{D})$ 的全子范畴
为定义域的函子，其中该全子范畴的对象是各个 $h_V^\#$。

\medskip\noindent
设 $J$ 为集合，并设 $J \to E$，$j \mapsto V_j$ 为映射。利用上述双射的
乘积，得到函子性双射
$$
\Mor_{\Sh(\mathcal{C})}(\coprod \mathcal{G}_{V_j}, \mathcal{F})
\longrightarrow
\Mor_{\Sh(\mathcal{D})}(\coprod h_{V_j}^\#, f^{-1}\mathcal{F})
$$
因此，可以把函子 $f_!$ 扩张到 $\Sh(\mathcal{D})$ 的全子范畴，该全子
范畴的对象是 $V \in E$ 的 $h_V^\#$ 的余积。

\medskip\noindent
给定 $\mathcal{D}$ 上任意层 $\mathcal{H}$，选择余等化子图
$$
\xymatrix{
\mathcal{H}_1 \ar@<1ex>[r] \ar@<-1ex>[r] &
\mathcal{H}_0 \ar[r] &
\mathcal{H}
}
$$
其中 $\mathcal{H}_i = \coprod h_{V_{i, j}}^\#$ 是一个余积，且
$V_{i, j} \in E$。由假设 (2)，这是可能的，见引理
\ref{sites-lemma-sheaf-coequalizer-representable}（对于担心集合论问题的读者，
注意引理 \ref{sites-lemma-sheaf-coequalizer-representable} 中给出的构造是典范的）。
把 $f_!(\mathcal{H})$ 定义为 $\mathcal{C}$ 上使
$$
\xymatrix{
f_!\mathcal{H}_1 \ar@<1ex>[r] \ar@<-1ex>[r] &
f_!\mathcal{H}_0 \ar[r] &
f_!\mathcal{H}
}
$$
成为余等化子图的层。于是
\begin{align*}
\Mor(f_!\mathcal{H}, \mathcal{F})
& =
\text{Equalizer}(
\xymatrix{
\Mor(f_!\mathcal{H}_0, \mathcal{F}) \ar@<1ex>[r] \ar@<-1ex>[r] &
\Mor(f_!\mathcal{H}_1, \mathcal{F})
}
) \\
& =
\text{Equalizer}(
\xymatrix{
\Mor(\mathcal{H}_0, f^{-1}\mathcal{F}) \ar@<1ex>[r] \ar@<-1ex>[r] &
\Mor(\mathcal{H}_1, f^{-1}\mathcal{F})
}
) \\
& =
\Hom(\mathcal{H}, f^{-1}\mathcal{F})
\end{align*}
因此可把 $f_!$ 扩张到 $\mathcal{D}$ 上整个层范畴。
\end{proof}








\section{局部化}
\label{sites-section-localize}

\noindent
设 $\mathcal{C}$ 为位点，并设 $U \in \Ob(\mathcal{C})$。关于 $U$ 上对象
所成范畴 $\mathcal{C}/U$ 的定义，见 Categories，例
\ref{categories-example-category-over-X}。如下把 $\mathcal{C}/U$ 变为位点：
当且仅当 $U$ 上对象的态射族 $\{V_j \to V\}$ 是 $\mathcal{C}$ 中的覆盖时，
才宣称它是 $\mathcal{C}/U$ 的覆盖。考虑遗忘函子
$$
j_U : \mathcal{C}/U \longrightarrow \mathcal{C}.
$$
它显然余连续且连续。因此，由引理 \ref{sites-lemma-when-shriek}，得到托扑斯态射
$$
j_U : \Sh(\mathcal{C}/U) \longrightarrow \Sh(\mathcal{C})
$$
由 $j_U^{-1}$ 和 $j_{U*}$ 给出，并有 $j_U^{-1}$ 的左伴随函子 $j_{U!}$。

\begin{definition}
\label{sites-definition-localize}
设 $\mathcal{C}$ 为位点，并设 $U \in \Ob(\mathcal{C})$。
\begin{enumerate}
\item 位点 $\mathcal{C}/U$ 称为{\it 位点 $\mathcal{C}$ 在对象 $U$ 处的
局部化}。
\item 托扑斯态射
$j_U : \Sh(\mathcal{C}/U) \to \Sh(\mathcal{C})$ 称为{\it 局部化态射}。
\item 函子 $j_{U*}$ 称为{\it 直像函子}。
\item 对 $\mathcal{C}$ 上的层 $\mathcal{F}$，层
$j_U^{-1}\mathcal{F}$ 称为 $\mathcal{F}$ {\it 到 $\mathcal{C}/U$ 的限制}。
\item 对 $\mathcal{C}/U$ 上的层 $\mathcal{G}$，层
$j_{U!}\mathcal{G}$ 称为 $\mathcal{G}$ {\it 以空集的延拓}。
\end{enumerate}
\end{definition}

\noindent
如所预期，限制 $j_U^{-1}\mathcal{F}$ 是由规则
$j_U^{-1}\mathcal{F}(X/U) = \mathcal{F}(X)$ 定义的层。在此情形中，以空集
的延拓也有非常简单的描述，如下。

\begin{lemma}
\label{sites-lemma-describe-j-shriek}
设 $\mathcal{C}$ 为位点，$U \in \Ob(\mathcal{C})$，并设
$\mathcal{G}$ 为 $\mathcal{C}/U$ 上的预层。则
$j_{U!}(\mathcal{G}^\#)$ 是与下述预层相伴的层：
$$
V
\longmapsto
\coprod\nolimits_{\varphi \in \Mor_\mathcal{C}(V, U)}
\mathcal{G}(V \xrightarrow{\varphi} U)
$$
其限制映射显然。
\end{lemma}

\begin{proof}
由引理 \ref{sites-lemma-when-shriek}，有
$j_{U!}(\mathcal{G}^\#) = ((j_U)_p\mathcal{G}^\#)^\#$。由引理
\ref{sites-lemma-technical-up}，它等于 $((j_U)_p\mathcal{G})^\#$。因此，只需证明
对 $\mathcal{C}/U$ 上任意预层 $\mathcal{G}$，$(j_U)_p$ 由上述公式给出。
现在，根据第 \ref{sites-section-functoriality-PSh} 节中的定义，有
% P01-ERRATA: P01-E0091 -- source begins the sentence with the editing remnant "OK, and".
$$
(j_U)_p\mathcal{G}(V)
=
\colim_{(W/U, V \to W)} \mathcal{G}(W)
$$
显然，对每个 $\varphi : V \to U$，序对 $(W/U, V \to W)$ 所成范畴都有
对象 $O_\varphi = (\varphi : V \to U, \text{id} : V \to V)$；而且从某个
$O_\varphi$ 到任意对象恰有一个态射。结论随之成立。
\end{proof}

\begin{lemma}
\label{sites-lemma-describe-j-shriek-representable}
设 $\mathcal{C}$ 为位点，$U \in \Ob(\mathcal{C})$，并设 $X/U$ 为
$\mathcal{C}/U$ 的对象。则 $j_{U!}(h_{X/U}^\#) = h_X^\#$。
\end{lemma}

\begin{proof}
以 $p : X \to U$ 表示 $X$ 的结构态射。由引理
\ref{sites-lemma-describe-j-shriek}，$j_{U!}(h_{X/U}^\#)$ 是与下述预层相伴的层：
$$
V
\longmapsto
\coprod\nolimits_{\varphi \in \Mor_\mathcal{C}(V, U)}
\{\psi : V \to X \mid p \circ \psi = \varphi\}
$$
这显然就是 $\Mor_\mathcal{C}(V, X)$，故引理成立。
\end{proof}

\noindent
由上述任一引理，有 $j_{U!}(*) = h_U^\#$。因此，对 $\mathcal{C}/U$ 上的
每个层 $\mathcal{G}$，都有层的典范映射
$j_{U!}\mathcal{G} \to h_U^\#$。这刻画了 $j_{U!}$ 本质像中的层。

\begin{lemma}
\label{sites-lemma-essential-image-j-shriek}
设 $\mathcal{C}$ 为位点，$U \in \Ob(\mathcal{C})$。函子 $j_{U!}$ 给出
范畴等价
$$
\Sh(\mathcal{C}/U)
\longrightarrow
\Sh(\mathcal{C})/h_U^\#
$$
\end{lemma}

\begin{proof}
把 $\mathcal{C}/U$ 的对象记作序对 $(X, a)$，其中 $X$ 是
$\mathcal{C}$ 的对象，$a : X \to U$ 是 $\mathcal{C}$ 的态射。类似地，
$\Sh(\mathcal{C})/h_U^\#$ 的对象是序对 $(\mathcal{F}, \varphi)$。函子
$\Sh(\mathcal{C}/U) \to \Sh(\mathcal{C})/h_U^\#$ 把 $\mathcal{G}$ 送到
序对 $(j_{U!}\mathcal{G}, \gamma)$，其中 $\gamma$ 是
$j_{U!}\mathcal{G} \to j_{U!}*$ 与等同
$j_{U!}* = h_U^\#$ 的复合。

\medskip\noindent
构造从 $\Sh(\mathcal{C})/h_U^\#$ 到 $\Sh(\mathcal{C}/U)$ 的函子。假设
给定 $(\mathcal{F}, \varphi)$。对 $\mathcal{C}/U$ 的对象 $(X, a)$，考虑
集合 $\mathcal{F}_\varphi(X, a)$；它由 $s \in \mathcal{F}(X)$ 中经
$\varphi$ 映到 $a \in \Mor_\mathcal{C}(X, U) = h_U(X)$ 在
$h_U^\#(X)$ 中之像的元素组成。容易看出，
$(X, a) \mapsto \mathcal{F}_\varphi(X, a)$ 是 $\mathcal{C}/U$ 上的层。
显然，规则 $(\mathcal{F}, \varphi) \mapsto \mathcal{F}_\varphi$ 定义函子
$\Sh(\mathcal{C})/h_U^\# \to \Sh(\mathcal{C}/U)$。

\medskip\noindent
还考虑函子
$\textit{PSh}(\mathcal{C})/h_U \to \textit{PSh}(\mathcal{C}/U)$，
$(\mathcal{F}, \varphi) \mapsto \mathcal{F}_\varphi$，其中
$\mathcal{F}_\varphi(X, a)$ 定义为 $\mathcal{F}(X)$ 中映到
$a \in h_U(X)$ 的元素所成的集合。断言图
$$
\xymatrix{
\textit{PSh}(\mathcal{C})/h_U \ar[r] \ar[d] &
\textit{PSh}(\mathcal{C}/U) \ar[d] \\
\Sh(\mathcal{C})/h_U^\# \ar[r] &
\Sh(\mathcal{C}/U)
}
$$
交换，其中竖直箭头由层化给出。为看出这一点\footnote{另一种做法是，
对预层用笛卡尔图
$$
\vcenter{
\xymatrix{
\mathcal{F}_\varphi \ar[r] \ar[d] & {*} \ar[d] \\
\mathcal{F}|_{\mathcal{C}/U} \ar[r] & h_U|_{\mathcal{C}/U}
}
}
\quad\text{对预层以及}\quad
\vcenter{
\xymatrix{
\mathcal{F}_\varphi \ar[r] \ar[d] & {*} \ar[d] \\
\mathcal{F}|_{\mathcal{C}/U} \ar[r] & h_U^\#|_{\mathcal{C}/U}
}
}
$$
描述 $\mathcal{F}_\varphi$，并对层使用后一图，再利用到
$\mathcal{C}/U$ 的限制与层化交换这一事实。}，只需证明该构造与引理
\ref{sites-lemma-plus-presheaf}、\ref{sites-lemma-plus-functorial} 及定理
\ref{sites-theorem-plus} 中的函子
$\mathcal{F} \mapsto \mathcal{F}^+$ 交换。与
$\mathcal{F} \mapsto \mathcal{F}^+$ 的交换性来自如下事实：给定
$(X, a)$，$\mathcal{C}/U$ 中 $(X, a)$ 的覆盖范畴与 $\mathcal{C}$ 中 $X$
的覆盖范畴典范等同。

\medskip\noindent
接着，令 $\textit{PSh}(\mathcal{C}/U) \to
\textit{PSh}(\mathcal{C})/h_U$ 把 $\mathcal{G}$ 送到序对
$(j_{U!}^{PSh}\mathcal{G}, \gamma)$，其中
$j_{U!}^{PSh}\mathcal{G}$ 是由引理 \ref{sites-lemma-describe-j-shriek} 中公式
定义的预层，而 $\gamma$ 是
$j_{U!}^{PSh}\mathcal{G} \to j_{U!}*$ 与等同
$j_{U!}^{PSh}* = h_U$ 的复合（由公式显然可知）。
% P01-ERRATA: P01-E0092 -- source omits "is" in "where j_{U!}^{PSh}G the presheaf".
% P01-ERRATA: P01-E0093 -- the target of the first map drops the PSh superscript, inconsistent with the stated identification.
于是立即清楚可知图
$$
\xymatrix{
\textit{PSh}(\mathcal{C}/U) \ar[r] \ar[d] &
\textit{PSh}(\mathcal{C})/h_U \ar[d] \\
\Sh(\mathcal{C}/U) \ar[r] &
\Sh(\mathcal{C})/h_U^\#
}
$$
交换，其中竖直箭头是层化。综合上述各点，只需证明：对
$\textit{PSh}(\mathcal{C}/U)$ 中的 $\mathcal{G}$，有函子性同构
$(j_{U!}^{PSh}\mathcal{G})_\gamma = \mathcal{G}$；并且对
$\textit{PSh}(\mathcal{C})/h_U$ 中的 $(\mathcal{F}, \varphi)$，有
$j_{U!}^{PSh}\mathcal{F}_\varphi = \mathcal{F}$。预层
$(j_{U!}^{PSh}\mathcal{G})_\gamma$ 在 $(X, a)$ 上的值是映射
$$
\coprod\nolimits_{a' : X \to U} \mathcal{G}(X, a') \to \Mor_\mathcal{C}(X, U)
$$
在 $a$ 上的纤维，即 $\mathcal{G}(X, a)$。这证明第一个等式。预层
$j_{U!}^{PSh}\mathcal{F}_\varphi$ 在 $X$ 上的值是
% P01-ERRATA: P01-E0094 -- source has the duplicated phrase "is on X is".
$$
\coprod\nolimits_{a : X \to U} \mathcal{F}_\varphi(X, a) =
\mathcal{F}(X)
$$
因为给定集合映射 $S \to S'$，集合 $S$ 是其各纤维的无交并。
\end{proof}

\noindent
引理 \ref{sites-lemma-essential-image-j-shriek} 说明，函子 $j_{U!}$ 是复合
$$
\Sh(\mathcal{C}/U) \rightarrow
\Sh(\mathcal{C})/h_U^\# \rightarrow
\Sh(\mathcal{C})
$$
其中第一个箭头是等价。

\begin{lemma}
\label{sites-lemma-j-shriek-commutes-equalizers-fibre-products}
设 $\mathcal{C}$ 为位点，$U \in \Ob(\mathcal{C})$。函子 $j_{U!}$ 与纤维积和
等化子交换（更一般地，与有限连通极限交换）。特别地，若
$\mathcal{F} \subset \mathcal{F}'$ 于 $\Sh(\mathcal{C}/U)$ 中，则
$j_{U!}\mathcal{F} \subset j_{U!}\mathcal{F}'$。
\end{lemma}

\begin{proof}
利用引理 \ref{sites-lemma-essential-image-j-shriek} 以及范畴等价与所有极限交换这一
事实，问题归结为函子
$\Sh(\mathcal{C})/h_U^\# \rightarrow \Sh(\mathcal{C})$ 与纤维积和等化子
交换。也可以使用引理 \ref{sites-lemma-describe-j-shriek} 中对 $j_{U!}$ 的描述，
并利用层化是正合的，直接证明这一点。（此外，当 $\mathcal{C}$ 有纤维积和
等化子时，结论也由引理 \ref{sites-lemma-preserve-equalizers} 得出。）
\end{proof}

\begin{lemma}
\label{sites-lemma-j-shriek-reflects-injections-surjections}
设 $\mathcal{C}$ 为位点，$U \in \Ob(\mathcal{C})$。函子 $j_{U!}$ 反映单射与
满射。
\end{lemma}

\begin{proof}
我们必须证明 $j_{U!}$ 反映单态射与满态射，见引理
\ref{sites-lemma-mono-epi-sheaves}。利用引理
\ref{sites-lemma-essential-image-j-shriek}，这归结为函子
$\Sh(\mathcal{C})/h_U^\# \to \Sh(\mathcal{C})$ 反映单态射与满态射这一
事实。
\end{proof}

\begin{lemma}
\label{sites-lemma-compute-j-shriek-restrict}
设 $\mathcal{C}$ 为位点，$U \in \Ob(\mathcal{C})$。对 $\mathcal{C}$ 上的任意
层 $\mathcal{F}$，有
$j_{U!}j_U^{-1}\mathcal{F} = \mathcal{F} \times h_U^\#$。
\end{lemma}

\begin{proof}
这由引理 \ref{sites-lemma-describe-j-shriek} 中对 $j_{U!}$ 的描述显然可知。
\end{proof}

\begin{lemma}
\label{sites-lemma-relocalize}
设 $\mathcal{C}$ 为位点。设 $f : V \to U$ 是 $\mathcal{C}$ 的态射。则存在
连续且余连续函子构成的交换图
$$
\xymatrix{
\mathcal{C}/V \ar[rd]_{j_V} \ar[rr]_j & &
\mathcal{C}/U \ar[ld]^{j_U} \\
& \mathcal{C} &
}
$$
函子 $j : \mathcal{C}/V \to \mathcal{C}/U$，
$(a : W \to V) \mapsto (f \circ a : W \to U)$，通过等同
$(\mathcal{C}/U)/(V/U) = \mathcal{C}/V$ 与函子
$j_{V/U} : (\mathcal{C}/U)/(V/U) \to \mathcal{C}/U$ 等同。此外，有
$j_{V!} = j_{U!} \circ j_!$、
$j_V^{-1} = j^{-1} \circ j_U^{-1}$ 以及
$j_{V*} = j_{U*} \circ j_*$。
\end{lemma}

\begin{proof}
图的交换性是立即的。$j$ 与 $j_{V/U}$ 相符由定义得出。由引理
\ref{sites-lemma-composition-cocontinuous}，可见下列拓扑斯态射图
\begin{equation}
\label{sites-equation-relocalize}
\vcenter{
\xymatrix{
\Sh(\mathcal{C}/V) \ar[rd]_{j_V} \ar[rr]_j & &
\Sh(\mathcal{C}/U) \ar[ld]^{j_U} \\
& \Sh(\mathcal{C}) &
}
}
\end{equation}
交换。这证明
$j_V^{-1} = j^{-1} \circ j_U^{-1}$ 和 $j_{V*} = j_{U*} \circ j_*$。
等式 $j_{V!} = j_{U!} \circ j_!$ 由伴随性质形式地得出。
\end{proof}

\begin{lemma}
\label{sites-lemma-relocalize-explicit}
记号 $\mathcal{C}$、$f : V \to U$、$j_U$、$j_V$ 和 $j$ 如引理
\ref{sites-lemma-relocalize}。通过引理
\ref{sites-lemma-essential-image-j-shriek} 的等同
$\Sh(\mathcal{C}/V) = \Sh(\mathcal{C})/h_V^\#$ 与
$\Sh(\mathcal{C}/U) = \Sh(\mathcal{C})/h_U^\#$，有
\begin{enumerate}
\item 函子 $j^{-1}$ 有如下描述：
$$
j^{-1}(\mathcal{H} \xrightarrow{\varphi} h_U^\#)
=
(\mathcal{H} \times_{\varphi, h_U^\#, f} h_V^\# \to h_V^\#).
$$
\item 函子 $j_!$ 有如下描述：
$$
j_!(\mathcal{H} \xrightarrow{\varphi} h_V^\#) =
(\mathcal{H} \xrightarrow{h_f \circ \varphi} h_U^\#)
$$
\end{enumerate}
\end{lemma}

\begin{proof}
先证明 (2)。回忆等同
$\Sh(\mathcal{C}/V) \to \Sh(\mathcal{C})/h_V^\#$ 把 $\mathcal{G}$ 送到
$j_{V!}\mathcal{G} \to j_{V!}(*) = h_V^\#$，而
$\Sh(\mathcal{C}/U) \to \Sh(\mathcal{C})/h_U^\#$ 的情形也类似。因此，
$j_!\mathcal{G}$ 被送到
$j_{U!}(j_!\mathcal{G}) \to j_{U!}(*) = h_U^\#$；又由引理
\ref{sites-lemma-relocalize} 有 $j_{U!}j_! = j_{V!}$，故 (2) 成立。

\medskip\noindent
读者现在可以利用 $j^{-1}$ 是 $j_!$ 的右伴随，并利用 (1) 中的规则是 (2)
中规则的右伴随，来证明 (1)。下面给出直接证明。假设
$\varphi : \mathcal{H} \to h_U^\#$ 是 $\Sh(\mathcal{C})/h_U^\#$ 的对象。
由引理 \ref{sites-lemma-essential-image-j-shriek} 的证明，它对应于
$\mathcal{C}/U$ 上按规则
$$
(a : W \to U)
\longmapsto
\{ s \in \mathcal{H}(W) \mid \varphi(s) = a\}
$$
定义的层 $\mathcal{H}_\varphi$。这个 $\mathcal{C}/U$ 上的层到
$\mathcal{C}/V$ 的拉回
$j^{-1}\mathcal{H}_\varphi$ 由规则
$$
(a : W \to V)
\longmapsto
\{ s \in \mathcal{H}(W) \mid \varphi(s) = f \circ a\}
$$
给出，这是因为把 $j^{-1} = j_{U/V}^{-1}$ 描述为
$\mathcal{H}_\varphi$ 到 $\mathcal{C}/V$ 的限制。
% P01-ERRATA: P01-E0095 -- source writes j_{U/V}^{-1}, while the defined localization functor is j=j_{V/U}.
另一方面，把该规则应用于 $\Sh(\mathcal{C})/h_V^\#$ 的对象
$$
\xymatrix{
\mathcal{H}' = \mathcal{H} \times_{\varphi, h_U^\#, f} h_V^\#
\ar[rr]^-{\varphi'} & & h_V^\#
}
$$
便得到 $\mathcal{H}'_{\varphi'}$，它由
\begin{align*}
(a : W \to V)
\longmapsto
&  \{ s' \in \mathcal{H}'(W) \mid \varphi'(s') = a\} \\
= &
\{ (s, a') \in \mathcal{H}(W) \times h_V^\#(W) \mid
a' = a \text{ 且 } \varphi(s) = f \circ a'\}
\end{align*}
给出。这恰与上面对 $j^{-1}\mathcal{H}_\varphi$ 的描述是同一规则。
\end{proof}

\begin{remark}
\label{sites-remark-localize-presheaves}
局部化与预层。设 $\mathcal{C}$ 为范畴，$U$ 是 $\mathcal{C}$ 的对象。严格地说，
函子 $j_U^{-1}$、$j_{U*}$ 与 $j_{U!}$ 尚未对预层定义。不过，我们当然可以把
预层视为 $\mathcal{C}$ 上混沌拓扑的层（见例 \ref{sites-example-indiscrete}）。因此也
得到函子
$$
j_U^{-1} :
\textit{PSh}(\mathcal{C})
\longrightarrow
\textit{PSh}(\mathcal{C}/U)
$$
以及函子
$$
j_{U*}, j_{U!} :
\textit{PSh}(\mathcal{C}/U)
\longrightarrow
\textit{PSh}(\mathcal{C})
$$
它们分别是 $j_U^{-1}$ 的右、左伴随。由引理
\ref{sites-lemma-describe-j-shriek} 可见，$j_{U!}\mathcal{G}$ 是预层
$$
V \longmapsto
\coprod\nolimits_{\varphi \in \Mor_\mathcal{C}(V, U)}
\mathcal{G}(V \xrightarrow{\varphi} U)
$$
此外，函子 $j_{U!}$ 与纤维积和等化子交换。
\end{remark}

\begin{remark}
\label{sites-remark-localization-cartesian-cocontinuous}
设 $\mathcal{C}$ 为位点，$U \to V$ 是 $\mathcal{C}$ 的态射。余连续函子
$\mathcal{C}/U \to \mathcal{C}$ 与
$j : \mathcal{C}/U \to \mathcal{C}/V$（引理 \ref{sites-lemma-relocalize}）满足
备注 \ref{sites-remark-cartesian-cocontinuous} 的性质 $P$。例如，若有对象
$(X/U)$、$(W/V)$、态射 $g : j(X/U) \to (W/V)$ 以及覆盖
$\{f_i  : (W_i/V) \to (W/V)\}$，则
$(X \times_W W_i/U)$ 是 $(X/U) \times_{g, (W/V), f_i} (W_i/V)$ 的一个化身，
而族 $\{(X \times_W W_i/U) \to (X/U)\}$ 是 $\mathcal{C}/U$ 的覆盖。
\end{remark}






\section{层的粘合}
\label{sites-section-glueing-sheaves}

\noindent
本节类似于《层》，第 \ref{sheaves-section-glueing-sheaves} 节。

\begin{lemma}
\label{sites-lemma-glue-maps}
\begin{slogan}
层的映射可以粘合。
\end{slogan}
设 $\mathcal{C}$ 为位点，$\{U_i \to U\}$ 是 $\mathcal{C}$ 的覆盖，
$\mathcal{F}$、$\mathcal{G}$ 是 $\mathcal{C}$ 上的层。给定一族层映射
$$
\varphi_i :
\mathcal{F}|_{\mathcal{C}/U_i}
\longrightarrow
\mathcal{G}|_{\mathcal{C}/U_i}
$$
使得对所有 $i, j \in I$，映射 $\varphi_i, \varphi_j$ 在相应交叠上的
限制为同一映射
$\varphi_{ij} : \mathcal{F}|_{\mathcal{C}/U_i \times_U U_j} \to
\mathcal{G}|_{\mathcal{C}/U_i \times_U U_j}$，则存在唯一的层映射
$$
\varphi :
\mathcal{F}|_{\mathcal{C}/U}
\longrightarrow
\mathcal{G}|_{\mathcal{C}/U}
$$
其到每个 $\mathcal{C}/U_i$ 的限制都与 $\varphi_i$ 相符。
\end{lemma}

\begin{proof}
引理中所用的限制是引理 \ref{sites-lemma-relocalize} 中的限制。设 $V/U$ 是
$\mathcal{C}/U$ 的对象。令 $V_i = U_i \times_U V$，并记
$\mathcal{V} = \{V_i \to V\}$。注意
$(U_i \times_U U_j) \times_U V = V_i \times_V V_j$。于是有
$\mathcal{F}|_{\mathcal{C}/U_i}(V_i/U_i) = \mathcal{F}(V_i)$，以及
$\mathcal{F}|_{\mathcal{C}/U_i \times_U U_j}(V_i \times_V V_j/U_i \times_U U_j)
= \mathcal{F}(V_i \times_V V_j)$，对 $\mathcal{G}$ 也类似。因此，可以把
$V$ 上截面间的 $\varphi$ 定义为图中的虚线箭头：
$$
\xymatrix{
\mathcal{F}(V) \ar@{=}[r] &
H^0(\mathcal{V}, \mathcal{F}) \ar@{..>}[d] \ar[r] &
\prod \mathcal{F}(V_i)
\ar[d]_{\prod \varphi_i}
\ar@<1ex>[r] \ar@<-1ex>[r] &
\prod \mathcal{F}(V_i \times_V V_j) \ar[d]_{\prod \varphi_{ij}} \\
\mathcal{G}(V) \ar@{=}[r] &
H^0(\mathcal{V}, \mathcal{G}) \ar[r] &
\prod \mathcal{G}(V_i)
\ar@<1ex>[r] \ar@<-1ex>[r] &
\prod \mathcal{G}(V_i \times_V V_j)
}
$$
等号来自层条件；记号 $H^0(\mathcal{V}, -)$ 见第
\ref{sites-section-sheafification} 节。我们略去这些映射与限制映射相容的验证。
\end{proof}

\noindent
前一引理蕴含：给定位点 $\mathcal{C}$ 上的两个层 $\mathcal{F}$、
$\mathcal{G}$，规则
$$
U \longmapsto
\Mor_{\Sh(\mathcal{C}/U)}(
\mathcal{F}|_{\mathcal{C}/U}, \mathcal{G}|_{\mathcal{C}/U})
$$
定义一个层 $\SheafHom(\mathcal{F},\mathcal{G})$。这是一种
{\it 内 Hom 层}。在集合层的语境中很少使用它，而在模层的语境中更常用；
见《位点上的模》，第 \ref{sites-modules-section-internal-hom} 节。

\begin{lemma}
\label{sites-lemma-internal-hom-sheaf}
\begin{slogan}
位点上的层范畴是笛卡尔闭的
\end{slogan}
设 $\mathcal{C}$ 为位点，$\mathcal{F}$、$\mathcal{G}$ 与 $\mathcal{H}$ 是
$\mathcal{C}$ 上的层。存在典范双射
$$
\Mor_{\Sh(\mathcal{C})}(\mathcal{F}\times\mathcal{G},\mathcal{H}) =
\Mor_{\Sh(\mathcal{C})}(\mathcal{F},\SheafHom(\mathcal{G},\mathcal{H}))
$$
并且它关于三个变元都是函子性的。
\end{lemma}

\begin{proof}
引理断言函子 $-\times\mathcal{G}$ 与 $\SheafHom(\mathcal{G},-)$ 互为伴随。
为证明这一点，我们使用单位和余单位的概念，见《范畴》，第
\ref{categories-section-adjoint} 节。单位
$$
\eta_\mathcal{F} :
\mathcal{F}
\longrightarrow
\SheafHom(\mathcal{G},\mathcal{F}\times\mathcal{G})
$$
把 $s \in \mathcal{F}(U)$ 送到映射
$\mathcal{G}|_{\mathcal{C}/U} \to
\mathcal{F}|_{\mathcal{C}/U}\times\mathcal{G}|_{\mathcal{C}/U}$；该映射在
$V/U$ 上由
$$
\mathcal{G}(V) \longrightarrow \mathcal{F}(V)\times \mathcal{G}(V), \quad
t \longmapsto (s|_{V},t).
$$
给出。余单位
$$
\epsilon_{\mathcal{H}} :
\SheafHom(\mathcal{G}, \mathcal{H}) \times \mathcal{G}
\longrightarrow
\mathcal{H}
$$
是求值映射。它由规则
$$
\Mor_{\Sh(\mathcal{C}/U)}(
\mathcal{G}|_{\mathcal{C}/U}, \mathcal{H}|_{\mathcal{C}/U})
\times \mathcal{G}(U)
\longrightarrow
\mathcal{H}(U),\quad
(\varphi, s) \longmapsto \varphi(s).
$$
给出。于是，对每个 $\varphi : \mathcal{F} \times \mathcal{G} \to \mathcal{H}$，
相应的态射 $\mathcal{F} \to \SheafHom(\mathcal{G},\mathcal{H})$ 把每个截面
$s \in \mathcal{F}(U)$ 送到 $\mathcal{C}/U$ 上的层态射；该态射在 $V/U$
上的截面间由
$$
\mathcal{G}(V) \longrightarrow \mathcal{H}(V),\quad
t \longmapsto \varphi(s|_V, t).
$$
给出。反过来，对每个
$\psi : \mathcal{F} \to \SheafHom(\mathcal{G}, \mathcal{H})$，相应的态射
$\mathcal{F} \times \mathcal{G} \to \mathcal{H}$ 在对象 $U$ 上的截面间由
$$
\mathcal{F}(U) \times \mathcal{G}(U) \longrightarrow \mathcal{H}(U),\quad
(s, t) \longmapsto \psi(s)(t)
$$
给出。我们略去证明这些构造互逆的细节。
\end{proof}

\begin{lemma}
\label{sites-lemma-hom-sheaf-hU}
设 $\mathcal{C}$ 为位点，$U \in \Ob(\mathcal{C})$。则对
$\Sh(\mathcal{C})$ 中的 $\mathcal{F}$，有
$\SheafHom(h_U^\#, \mathcal{F}) = j_*(\mathcal{F}|_{\mathcal{C}/U})$。
\end{lemma}

\begin{proof}
可以直接构造层的同构来证明。这里改用如下论证。设 $\mathcal{G}$ 是
$\mathcal{C}$ 上的层。则
\begin{align*}
\Mor(\mathcal{G}, j_*(\mathcal{F}|_{\mathcal{C}/U}))
& =
\Mor(\mathcal{G}|_{\mathcal{C}/U}, \mathcal{F}|_{\mathcal{C}/U}) \\
& =
\Mor(j_!(\mathcal{G}|_{\mathcal{C}/U}), \mathcal{F}) \\
& =
\Mor(\mathcal{G} \times h_U^\#, \mathcal{F}) \\
& =
\Mor(\mathcal{G}, \SheafHom(h_U^\#, \mathcal{F}))
\end{align*}
由米田引理即可得出结论。这里使用了引理
\ref{sites-lemma-internal-hom-sheaf} 和
\ref{sites-lemma-compute-j-shriek-restrict}。
\end{proof}

\noindent
设 $\mathcal{C}$ 为位点，$\{U_i \to U\}_{i \in I}$ 是
$\mathcal{C}$ 的覆盖。对每个 $i \in I$，设 $\mathcal{F}_i$ 是
$\mathcal{C}/U_i$ 上的集合层。对每一对 $i, j \in I$，设
$$
\varphi_{ij} :
\mathcal{F}_i|_{\mathcal{C}/U_i \times_U U_j}
\longrightarrow
\mathcal{F}_j|_{\mathcal{C}/U_i \times_U U_j}
$$
是集合层的同构。再假设，对指标的每个三元组 $i, j, k \in I$，下图交换：
$$
\xymatrix{
\mathcal{F}_i|_{\mathcal{C}/U_i \times_U U_j \times_U U_k}
\ar[rr]_{\varphi_{ik}}
\ar[rd]_{\varphi_{ij}} & &
\mathcal{F}_k|_{\mathcal{C}/U_i \times_U U_j \times_U U_k} \\
&
\mathcal{F}_j|_{\mathcal{C}/U_i \times_U U_j \times_U U_k}
\ar[ru]_{\varphi_{jk}}
}
$$
我们把这样一族数据 $(\mathcal{F}_i, \varphi_{ij})$ 称为
{\it 关于覆盖 $\{U_i \to U\}_{i \in I}$ 的集合层粘合数据}。

\begin{lemma}
\label{sites-lemma-glue-sheaves}
设 $\mathcal{C}$ 为位点，$\{U_i \to U\}_{i \in I}$ 是
$\mathcal{C}$ 的覆盖。给定关于覆盖 $\{U_i \to U\}_{i \in I}$ 的任意
集合层粘合数据 $(\mathcal{F}_i, \varphi_{ij})$，存在
$\mathcal{C}/U$ 上的集合层 $\mathcal{F}$ 以及同构
$$
\varphi_i : \mathcal{F}|_{\mathcal{C}/U_i} \to \mathcal{F}_i
$$
使得下列各图交换：
$$
\xymatrix{
\mathcal{F}|_{\mathcal{C}/U_i \times_U U_j}
\ar[d]_{\text{id}} \ar[r]_{\varphi_i} &
\mathcal{F}_i|_{\mathcal{C}/U_i \times_U U_j} \ar[d]^{\varphi_{ij}} \\
\mathcal{F}|_{\mathcal{C}/U_i \times_U U_j} \ar[r]^{\varphi_j} &
\mathcal{F}_j|_{\mathcal{C}/U_i \times_U U_j}
}
$$
\end{lemma}

\begin{proof}
说明如何构造 $\mathcal{C}/U$ 上的层 $\mathcal{F}$。设
$a : V \to U$ 是 $\mathcal{C}/U$ 的对象。则
$$
\mathcal{F}(V/U) = \{
(s_i)_{i \in I} \in \prod_{i \in I} \mathcal{F}_i(U_i \times_U V/U_i)
\mid
\varphi_{ij}(s_i|_{U_i \times_U U_j \times_U V})
=
s_j|_{U_i \times_U U_j \times_U V}
\}
$$
我们略去限制映射的构造，略去这是一个层的验证，也略去同构
$\varphi_i$ 的构造以及引理中各图交换性的证明。
\end{proof}

\noindent
设 $\mathcal{C}$ 为位点，$\{U_i \to U\}_{i \in I}$ 是
$\mathcal{C}$ 的覆盖，$\mathcal{F}$ 是 $\mathcal{C}/U$ 上的层。与
$\mathcal{F}$ 相伴的{\it 典范粘合数据}由限制
$\mathcal{F}|_{\mathcal{C}/U_i}$ 以及典范同构
$$
\left(\mathcal{F}|_{\mathcal{C}/U_i}\right)|_{\mathcal{C}/U_i \times_U U_j}
=
\left(\mathcal{F}|_{\mathcal{C}/U_j}\right)|_{\mathcal{C}/U_i \times_U U_j}
$$
给出；这些同构来自如下事实：复合函子
$\mathcal{C}/U_i \times_U U_j \to \mathcal{C}/U_i \to \mathcal{C}/U$
与
$\mathcal{C}/U_i \times_U U_j \to \mathcal{C}/U_j \to \mathcal{C}/U$
相等。

\begin{lemma}
\label{sites-lemma-mapping-property-glue}
设 $\mathcal{C}$ 为位点，$\{U_i \to U\}_{i \in I}$ 是
$\mathcal{C}$ 的覆盖。把 $\mathcal{C}/U$ 上的 $\mathcal{F}$ 送到其典范
粘合数据的函子，给出范畴 $\Sh(\mathcal{C}/U)$ 与粘合数据范畴之间的等价。
\end{lemma}

\begin{proof}
在引理 \ref{sites-lemma-glue-maps} 中，我们看到该函子是全忠实的；在引理
\ref{sites-lemma-glue-sheaves} 中，我们证明它是本质满的（方法是显式构造一个
拟逆函子）。
\end{proof}

\noindent
设 $\mathcal{C}$ 为位点。下面讨论在整个位点 $\mathcal{C}$ 上粘合层的一个
版本。对 $\mathcal{C}$ 的每个对象 $U$，设 $\mathcal{F}_U$ 是
$\mathcal{C}/U$ 上的层。回忆：对 $\mathcal{C}$ 中的每个态射
$f : V \rightarrow U$，都有相伴函子
$j_f : \mathcal{C}/V \rightarrow \mathcal{C}/U$，它由
$(a : W \to V) \mapsto (f \circ a : W \to U)$ 给出。对每个这样的 $f$，
设
$$
c_f : j_f^{-1} \mathcal{F}_U \to \mathcal{F}_V
$$
是层的同构。再假设，对 $\mathcal{C}$ 中任意两个箭头
$f : V \to U$ 与 $g : W \to V$，复合
$c_g \circ j_{g}^{-1} c_f$ 等于 $c_{f \circ g}$。我们把这样一族数据
$(\mathcal{F}_{U}, c_f)$ 称为{\it $\mathcal{C}$ 上集合层的绝对粘合数据}。
绝对粘合数据间的态射
$(\mathcal{F}_{U}, c_f) \to (\mathcal{G}_{U}, c'_f)$ 由一族层态射
$(\varphi_U)$ 给出，其中 $\varphi_U : \mathcal{F}_U \to \mathcal{G}_U$，
并且
$$
\xymatrix{
j_f^{-1} \mathcal{F}_U \ar[r]_{c_f} \ar[d]_{j_f^{-1} \varphi_U} &
\mathcal{F}_V \ar[d]^{\varphi_V} \\
j_f^{-1} \mathcal{G}_U \ar[r]^{c'_f} &
\mathcal{G}_V
}
$$
对 $\mathcal{C}$ 中每个态射 $f : V \to U$ 都交换。

\medskip\noindent
与 $\mathcal{C}$ 上任意层 $\mathcal{F}$ 相伴的，是其{\it 典范绝对粘合数据}
$(\mathcal{F}|_{\mathcal{C}/U},c_f)$；其中，对 $f : V \to U$，典范同构
$c_f : j_f^{-1} \mathcal{F}|_{\mathcal{C}/U} \to
\mathcal{F}|_{\mathcal{C}/V}$ 来自引理 \ref{sites-lemma-relocalize} 中的关系
$j_V = j_U \circ j_f$。$\mathcal{C}$ 的层之间的任意态射
$\varphi : \mathcal{F} \to \mathcal{G}$ 都诱导典范绝对粘合数据间的态射
$(\varphi|_{\mathcal{C}/U})$。

\begin{lemma}
\label{sites-lemma-glue-sheaves-absolute}
设 $\mathcal{C}$ 为位点。把 $\mathcal{C}$ 上的 $\mathcal{F}$ 送到其典范
绝对粘合数据的函子，给出范畴 $\Sh(\mathcal{C})$ 与绝对粘合数据范畴之间的
等价。
\end{lemma}

\begin{proof}
给定绝对粘合数据 $(\mathcal{F}_U, c_f)$，令
$\mathcal{F}(U) = \mathcal{F}_U(U)$，从而构造 $\mathcal{C}$ 上的层
$\mathcal{F}$；沿 $f : V \to U$ 的限制由交换图给出：
% P01-ERRATA: P01-E0096 -- source sentence omits "is" before "given by the commutative diagram".
$$
\xymatrix{
\mathcal{F}_U(U) \ar[r] \ar@{=}[d] &
\mathcal{F}_U(V) \ar[r]^{c_f} &
\mathcal{F}_V(V) \ar@{=}[d] \\
\mathcal{F}(U) \ar[rr] &  &
\mathcal{F}(V)
}
$$
相容条件 $c_g \circ j_{g}^{-1} c_f = c_{f \circ g}$ 保证
$\mathcal{F}$ 是预层，也保证映射
$c_f : \mathcal{F}_U(V) \to \mathcal{F}(V)$ 对每个 $U$ 都定义同构
$\mathcal{F}_U \to \mathcal{F}|_{\mathcal{C}/U}$。由于每个
$\mathcal{F}_U$ 都是层，这也蕴含 $\mathcal{F}$ 是层。刚才构造的函子
$(\mathcal{F}_U, c_f) \mapsto \mathcal{F}$ 是把 $\mathcal{C}$ 上的层送到
其典范粘合数据的函子的拟逆。其余细节略去。
\end{proof}

\begin{remark}
\label{sites-remark-variant-glue-sheaves-absolute}
引理 \ref{sites-lemma-glue-sheaves-absolute} 有一个会在代数几何中出现的变体。
具体地，假设 $\mathcal{C}$ 是具有所有纤维积的位点，并且对每个
$U \in \Ob(\mathcal{C})$，都给定满足下列性质的满子范畴
$U_\tau \subset \mathcal{C}/U$：
\begin{enumerate}
\item $U/U$ 属于 $U_\tau$；
\item 对 $U_\tau$ 中的 $V/U$ 以及 $\mathcal{C}$ 的覆盖
$\{V_j \to V\}$，有 $V_j/U$ 属于 $U_\tau$；且
\item 对 $\mathcal{C}$ 的态射 $U' \to U$ 以及 $U_\tau$ 中的 $V/U$，
基变换 $V' = V \times_U U'$ 属于 $U'_\tau$。
\end{enumerate}
在此情形，对 $\mathcal{C}$ 中所有 $U$，$U_\tau$ 都是位点；并且对
$\mathcal{C}$ 的每个态射 $f : U' \to U$，基变换函子
$U_\tau \to U'_\tau$ 定义位点态射 $f_\tau : U_\tau \to U'_\tau$。
此时所得粘合陈述如下：$\mathcal{C}$ 上的层 $\mathcal{F}$ 由下列数据给出：
\begin{enumerate}
\item 对每个 $U \in \Ob(\mathcal{C})$，给定 $U_\tau$ 上的层
$\mathcal{F}_U$；
\item 对 $\mathcal{C}$ 中每个 $f : U' \to U$，给定映射
$c_f : f_\tau^{-1}\mathcal{F}_U \to \mathcal{F}_{U'}$。
\end{enumerate}
这些数据须满足下列条件：
\begin{enumerate}
\item[(a)] 给定 $\mathcal{C}$ 中的 $f : U' \to U$ 和
$g : U'' \to U'$，复合 $c_g \circ g_\tau^{-1}c_f$ 等于
$c_{f \circ g}$；且
\item[(b)] 若 $f : U' \to U$ 属于 $U_\tau$，则 $c_f$ 是同构。
\end{enumerate}
如果以后需要这一结果，我们会在此精确陈述并证明。（注意，这一结果与上述
陈述略有不同，因为这里并不要求所有映射 $c_f$ 都是同构！）
\end{remark}



\section{进一步的局部化}
\label{sites-section-more-localize}

\noindent
本节证明几个关于局部化的引理，其中会对位点或作局部化的对象施加一些附加
假设。
% P01-ERRATA: P01-E0097 -- source has "hypotheses on the site on or the object" instead of "on the site or on the object".

\begin{lemma}
\label{sites-lemma-describe-j-shriek-good-site}
设 $\mathcal{C}$ 为位点，$U \in \Ob(\mathcal{C})$。若
$\mathcal{C}$ 上的拓扑是次典范的，见定义
\ref{sites-definition-weaker-than-canonical}，并且 $\mathcal{G}$ 是
$\mathcal{C}/U$ 上的层，则
$$
j_{U!}(\mathcal{G})(V)
=
\coprod\nolimits_{\varphi \in \Mor_\mathcal{C}(V, U)}
\mathcal{G}(V \xrightarrow{\varphi} U),
$$
换言之，引理 \ref{sites-lemma-describe-j-shriek} 中无需层化。
\end{lemma}

\begin{proof}
设 $\mathcal{V} = \{V_i \to V\}_{i \in I}$ 是位点 $\mathcal{C}$ 中
$V$ 的覆盖。我们直接验证引理 \ref{sites-lemma-describe-j-shriek} 中预层
$\mathcal{H}$ 的层条件。设
$(s_i, \varphi_i)_{i \in I} \in \prod_i \mathcal{H}(V_i)$。这意味着
$\varphi_i : V_i \to U$ 是 $\mathcal{C}$ 中的态射，并且
$s_i \in \mathcal{G}(V_i \xrightarrow{\varphi_i} U)$。
% P01-ERRATA: P01-E0098 -- source joins these sentences with a comma before capitalized "This".
序对 $(s_i, \varphi_i)$ 到 $V_i \times_V V_j$ 的限制是
$(s_i|_{V_i \times_V V_j/U}, \text{pr}_1 \circ \varphi_i)$；类似地，序对
$(s_j, \varphi_j)$ 到 $V_i \times_V V_j$ 的限制是
$(s_j|_{V_i \times_V V_j/U}, \text{pr}_2 \circ \varphi_j)$。因此，若族
$(s_i, \varphi_i)$ 属于 $\check{H}^0(\mathcal{V}, \mathcal{H})$，则可见
$\text{pr}_1 \circ \varphi_i = \text{pr}_2 \circ \varphi_j$。
$\mathcal{C}$ 上的拓扑弱于典范拓扑这一条件于是蕴含：存在唯一态射
$\varphi : V \to U$，使 $\varphi_i$ 是 $V_i \to V$ 与 $\varphi$ 的复合。
此时，$\mathcal{G}$ 的层条件保证各截面 $s_i$ 粘合为唯一截面
$s \in \mathcal{G}(V \xrightarrow{\varphi} U)$。故如所需，
$(s, \varphi) \in \mathcal{H}(V)$。
\end{proof}

\begin{lemma}
\label{sites-lemma-localize-given-products}
设 $\mathcal{C}$ 为位点，$U \in \Ob(\mathcal{C})$。假设
$\mathcal{C}$ 有任意两个对象的积。则
\begin{enumerate}
\item 函子 $j_U$ 有连续的右伴随，即函子 $v(X) = X \times U / U$；
\item 函子 $v$ 定义位点态射 $\mathcal{C}/U \to \mathcal{C}$，其相伴的
拓扑斯态射等于 $j_U : \Sh(\mathcal{C}/U) \to \Sh(\mathcal{C})$；且
\item 有 $j_{U*}\mathcal{F}(X) = \mathcal{F}(X \times U/U)$。
\end{enumerate}
\end{lemma}

\begin{proof}
函子 $v$ 是 $j_U$ 的右伴随，意即给定 $Y/U$ 和 $X$ 时有
$$
\Mor_\mathcal{C}(Y, X)
=
\Mor_{\mathcal{C}/U}(Y/U, X \times U/U)
$$
这是显然的。为验证 $v$ 连续，设 $\{X_i \to X\}$ 是
$\mathcal{C}$ 的覆盖。由位点的第三条公理（定义 \ref{sites-definition-site}），
可见
$$
\{X_i \times_X (X \times U) \to X \times_X (X \times U)\}
=
\{X_i \times U \to X \times U\}
$$
也是 $\mathcal{C}$ 的覆盖。因此 $v$ 连续。引理的其余陈述由引理
\ref{sites-lemma-have-functor-other-way} 和
\ref{sites-lemma-have-functor-other-way-morphism} 得出。
\end{proof}

\begin{lemma}
\label{sites-lemma-relocalize-given-fibre-products}
设 $\mathcal{C}$ 为位点，$U \to V$ 是 $\mathcal{C}$ 的态射。假设
$\mathcal{C}$ 有纤维积。设 $j$ 如引理 \ref{sites-lemma-relocalize}。则
\begin{enumerate}
\item 函子 $j : \mathcal{C}/U \to \mathcal{C}/V$ 有连续的右伴随，即函子
$v : (X/V) \mapsto (X \times_V U/U)$；
\item 函子 $v$ 定义位点态射 $\mathcal{C}/U \to \mathcal{C}/V$，其相伴的
拓扑斯态射等于 $j$；且
\item 有 $j_*\mathcal{F}(X/V) = \mathcal{F}(X \times_V U/U)$。
\end{enumerate}
\end{lemma}

\begin{proof}
由引理 \ref{sites-lemma-localize-given-products} 得出，因为根据引理
\ref{sites-lemma-relocalize}，$j$ 可以视为局部化函子。
\end{proof}

\noindent
开浸入的一条基本性质是：先限制直像，或者先作空集延拓再限制，都会还原原来
的层。对上面与 $j_U$ 相伴的函子，这并不总成立。当 $U$ 是“终对象的子对象”
时则成立。
% P01-ERRATA: P01-E0099 -- source compound subject "the restriction ... and the restriction ..." takes singular "produces".

\begin{lemma}
\label{sites-lemma-restrict-back}
设 $\mathcal{C}$ 为位点，$U \in \Ob(\mathcal{C})$。假设
$\mathcal{C}$ 中每个 $X$ 至多有一个到 $U$ 的态射。设 $\mathcal{F}$ 是
$\mathcal{C}/U$ 上的层。典范映射
$\mathcal{F} \to j_U^{-1}j_{U!}\mathcal{F}$ 与
$j_U^{-1}j_{U*}\mathcal{F} \to \mathcal{F}$ 都是同构。
\end{lemma}

\begin{proof}
这是引理 \ref{sites-lemma-back-and-forth} 的特殊情形，因为对 $U$ 的假设等价于
局部化函子 $\mathcal{C}/U \to \mathcal{C}$ 的全忠实性。
% P01-ERRATA: P01-E0100 -- source uses the malformed noun phrase "the fully faithfulness".
\end{proof}

\begin{lemma}
\label{sites-lemma-localize-cartesian-square}
设 $\mathcal{C}$ 为位点。设
$$
\xymatrix{
U' \ar[d] \ar[r] & U \ar[d] \\
V' \ar[r] & V
}
$$
是 $\mathcal{C}$ 的交换图。引理 \ref{sites-lemma-relocalize} 的态射给出连续且
余连续函子的交换图以及拓扑斯的交换图
$$
\vcenter{
\xymatrix{
\mathcal{C}/U' \ar[d]_{j_{U'/V'}} \ar[r]_{j_{U'/U}} &
\mathcal{C}/U \ar[d]^{j_{U/V}} \\
\mathcal{C}/V' \ar[r]^{j_{V'/V}} & \mathcal{C}/V
}
}
\quad\text{且}\quad
\vcenter{
\xymatrix{
\Sh(\mathcal{C}/U') \ar[d]_{j_{U'/V'}} \ar[r]_{j_{U'/U}} &
\Sh(\mathcal{C}/U) \ar[d]^{j_{U/V}} \\
\Sh(\mathcal{C}/V') \ar[r]^{j_{V'/V}} &
\Sh(\mathcal{C}/V)
}
}
$$
此外，若 $\mathcal{C}$ 的初始图是笛卡尔的，则
$j_{V'/V}^{-1} \circ j_{U/V, *} = j_{U'/V', *} \circ j_{U'/U}^{-1}$。
\end{lemma}

\begin{proof}
引理第一条陈述中左方形的交换性由定义立即得出。由引理
\ref{sites-lemma-composition-cocontinuous}，它蕴含拓扑斯图的交换性。假设该图
是笛卡尔的。由伴随函子的唯一性，要证明
$j_{V'/V}^{-1} \circ j_{U/V, *} = j_{U'/V', *} \circ j_{U'/U}^{-1}$，
等价于证明
$j_{U/V}^{-1} \circ j_{V'/V!} = j_{U'/U!} \circ j_{U'/V'}^{-1}$。
借助引理 \ref{sites-lemma-essential-image-j-shriek} 的等同，可以把拓扑斯图看作
$$
\xymatrix{
\Sh(\mathcal{C})/h_{U'}^\# \ar[d] \ar[r] &
\Sh(\mathcal{C})/h_U^\# \ar[d] \\
\Sh(\mathcal{C})/h_{V'}^\# \ar[r] &
\Sh(\mathcal{C})/h_V^\#
}
$$
而由引理 \ref{sites-lemma-relocalize-explicit}，我们知道如何解释函子
$j^{-1}$ 与 $j_!$。因此必须证明：给定 $\mathcal{F} \to h_{V'}^\#$，有
$$
\mathcal{F} \times_{h_{V'}^\#} h_{U'}^\# =
\mathcal{F} \times_{h_V^\#} h_U^\#
$$
作为带有到 $h_U^\#$ 的映射的层。这是真的，因为
$h_{U'} = h_{V'} \times_{h_V} h_U$；而且由于层化是正合的，也有
$$
h_{U'}^\# = h_{V'}^\# \times_{h_V^\#} h_U^\#
$$
\end{proof}












\section{局部化与态射}
\label{sites-section-localize-morphisms}

\noindent
为理解局部化与位点及拓扑斯的态射之间的关系，下列引理很重要。
% P01-ERRATA: P01-E0101 -- source omits "the" in "understand relation between".

\begin{lemma}
\label{sites-lemma-localize-morphism}
设 $f : \mathcal{C} \to \mathcal{D}$ 是对应于连续函子
$u : \mathcal{D} \to \mathcal{C}$ 的位点态射。设
$V \in \Ob(\mathcal{D})$，并令 $U = u(V)$。则函子
$u' : \mathcal{D}/V \to \mathcal{C}/U$，
$V'/V \mapsto u(V')/U$，确定一个位点态射
$f' : \mathcal{C}/U \to \mathcal{D}/V$。态射 $f'$ 嵌入拓扑斯的交换图
$$
\xymatrix{
\Sh(\mathcal{C}/U) \ar[r]_{j_U} \ar[d]_{f'} &
\Sh(\mathcal{C}) \ar[d]^f \\
\Sh(\mathcal{D}/V) \ar[r]^{j_V} &
\Sh(\mathcal{D}).
}
$$
利用引理 \ref{sites-lemma-essential-image-j-shriek} 的等同
$\Sh(\mathcal{C}/U) = \Sh(\mathcal{C})/h_U^\#$ 与
$\Sh(\mathcal{D}/V) = \Sh(\mathcal{D})/h_V^\#$，函子
$(f')^{-1}$ 可由规则
$$
(f')^{-1}(\mathcal{H} \xrightarrow{\varphi} h_V^\#)
=
(f^{-1}\mathcal{H} \xrightarrow{f^{-1}\varphi} h_U^\#).
$$
描述。最后，有 $f'_*j_U^{-1} = j_V^{-1}f_*$。
\end{lemma}

\begin{proof}
显然 $u'$ 连续，因而得到函子
$f'_* = (u')^s = (u')^p$（见第 \ref{sites-section-functoriality-PSh} 节与第
\ref{sites-section-continuous-functors} 节）以及伴随
$(f')^{-1} = (u')_s = ((u')_p\ )^\#$。断言
$f'_*j_U^{-1} = j_V^{-1}f_*$ 来自
$$
(j_V^{-1}f_*\mathcal{F})(V'/V)
= f_*\mathcal{F}(V') = \mathcal{F}(u(V'))
= (j_U^{-1}\mathcal{F})(u(V')/U)
= (f'_*j_U^{-1}\mathcal{F})(V'/V)
$$
这甚至对预层也成立。先验并不清楚的是：$(f')^{-1}$ 是正合的、该图交换，
并且对 $(f')^{-1}$ 的描述成立。

\medskip\noindent
设 $\mathcal{H}$ 是 $\mathcal{D}/V$ 上的层。计算
$j_{U!}(f')^{-1}\mathcal{H}$。有
\begin{eqnarray*}
j_{U!}(f')^{-1}\mathcal{H}
& =
((j_U)_p(u'_p\mathcal{H})^\#)^\# \\
& =
((j_U)_pu'_p\mathcal{H})^\# \\
& =
(u_p(j_V)_p\mathcal{H})^\# \\
& =
f^{-1}j_{V!}\mathcal{H}
\end{eqnarray*}
第一个等式由展开定义得出。第二个等式由引理
\ref{sites-lemma-technical-up} 得出。第三个等式来自
$u \circ j_V = j_U \circ u'$。第四个等式再次由引理
\ref{sites-lemma-technical-up} 得出。
% P01-ERRATA: P01-E0102 -- source presents all four explanations as sentence fragments "The ... equality by/because ...".
上述等式全都是函子同构，因此可以把它解释为下列范畴与函子图交换：
$$
\xymatrix{
\Sh(\mathcal{C}/U) \ar[r]_{j_{U!}} &
\Sh(\mathcal{C})/h_U^\# \ar[r] &
\Sh(\mathcal{C}) \\
\Sh(\mathcal{D}/V) \ar[r]^{j_{V!}} \ar[u]^{(f')^{-1}} &
\Sh(\mathcal{D})/h_V^\# \ar[r] \ar[u]^{f^{-1}} &
\Sh(\mathcal{D}) \ar[u]^{f^{-1}}
}
$$
中间箭头有意义，因为
$f^{-1}h_V^\# = (h_{u(V)})^\# = h_U^\#$；见引理
\ref{sites-lemma-pullback-representable-sheaf}。特别地，这证明了引理陈述中对
$(f')^{-1}$ 的描述。根据引理 \ref{sites-lemma-essential-image-j-shriek}，左侧
水平箭头是等价；又因为根据假设 $f^{-1}$ 是正合的，故得出
$(f')^{-1} = u'_s$ 是正合的。具体地，因为它是左伴随，所以已经右正合
（《范畴》，引理 \ref{categories-lemma-adjoint-exact}）。因此只需证明它把
终对象变为终对象，并且与纤维积交换（《范畴》，引理
\ref{categories-lemma-characterize-left-exact}）。对诱导函子
$f^{-1} : \Sh(\mathcal{D})/h_V^\# \to \Sh(\mathcal{C})/h_U^\#$，两点都显然。
这证明 $f'$ 是位点态射。

\medskip\noindent
还须验证 $(f')^{-1}j_V^{-1} = j_U^{-1}f^{-1}$。为此，使用上面的公式以及
引理 \ref{sites-lemma-compute-j-shriek-restrict} 中的描述。具体地，二者结合给出：
对 $\mathcal{D}$ 上的任意层 $\mathcal{G}$，有
$$
j_{U!}(f')^{-1}j_V^{-1}\mathcal{G}
=
f^{-1}j_{V!}j_V^{-1}\mathcal{G}
=
f^{-1}(\mathcal{G} \times h_V^\#)
=
f^{-1}\mathcal{G} \times h_U^\#
=
j_{U!}j_U^{-1}f^{-1}\mathcal{G}.
$$
由于函子 $j_{U!}$ 诱导等价
$\Sh(\mathcal{C}/U) \to \Sh(\mathcal{C})/h_U^\#$，结论成立。
\end{proof}

\noindent
下列引理是上面更一般的引理 \ref{sites-lemma-localize-morphism} 的特殊情形。

\begin{lemma}
\label{sites-lemma-localize-morphism-strong}
设 $\mathcal{C}$、$\mathcal{D}$ 为位点，
$u : \mathcal{D} \to \mathcal{C}$ 为函子。设
$V \in \Ob(\mathcal{D})$，令 $U = u(V)$。假设
\begin{enumerate}
\item $\mathcal{C}$ 和 $\mathcal{D}$ 有所有有限极限；
\item $u$ 连续；且
\item $u$ 与有限极限交换。
\end{enumerate}
存在位点态射的交换图
$$
\xymatrix{
\mathcal{C}/U \ar[r]_{j_U} \ar[d]_{f'} & \mathcal{C} \ar[d]^f \\
\mathcal{D}/V \ar[r]^{j_V} & \mathcal{D}
}
$$
其中右竖直箭头对应于 $u$，左竖直箭头对应于函子
$u' : \mathcal{D}/V \to \mathcal{C}/U$，
$V'/V \mapsto u(V')/u(V)$；水平箭头则对应于引理
\ref{sites-lemma-localize-given-products} 中的函子
$\mathcal{C} \to \mathcal{C}/U$，$X \mapsto X \times U$，以及
$\mathcal{D} \to \mathcal{D}/V$，$Y \mapsto Y \times V$。
% P01-ERRATA: P01-E0103 -- source omits the slice structure notation /U and /V from these two functor values, unlike the cited lemma.
此外，相伴的拓扑斯态射图等于引理 \ref{sites-lemma-localize-morphism} 的图。
特别地，有 $f'_*j_U^{-1} = j_V^{-1}f_*$。
\end{lemma}

\begin{proof}
注意，$u$ 满足命题 \ref{sites-proposition-get-morphism} 的假设，因此由该命题诱导
位点态射 $f : \mathcal{C} \to \mathcal{D}$。显然，$u$ 诱导所述函子
$u'$。显然，这一函子也满足命题 \ref{sites-proposition-get-morphism} 的假设。
因此得到位点态射 $f' : \mathcal{C}/U \to \mathcal{D}/V$。根据我们对位点
态射复合的定义（见定义 \ref{sites-definition-composition-morphisms-sites}），并且
由于
$$
u(Y \times V) = u(Y) \times u(V) = u(Y) \times U
$$
表明与引理之图相反的范畴与函子图交换，故原图交换。
\end{proof}

\noindent
至此，我们可以局部化一个位点，知道如何再次局部化，也可以在下方位点的一个
对象处局部化位点态射。把这些结合起来，就得到下列类型的图。

\begin{lemma}
\label{sites-lemma-relocalize-morphism}
设 $f : \mathcal{C} \to \mathcal{D}$ 是对应于连续函子
$u : \mathcal{D} \to \mathcal{C}$ 的位点态射。设
$V \in \Ob(\mathcal{D})$、$U \in \Ob(\mathcal{C})$，且
$c : U \to u(V)$ 是 $\mathcal{C}$ 的态射。存在拓扑斯的交换图
$$
\xymatrix{
\Sh(\mathcal{C}/U) \ar[r]_{j_U} \ar[d]_{f_c} &
\Sh(\mathcal{C}) \ar[d]^f \\
\Sh(\mathcal{D}/V) \ar[r]^{j_V} &
\Sh(\mathcal{D}).
}
$$
有 $f_c = f' \circ j_{U/u(V)}$，其中
$f' : \Sh(\mathcal{C}/u(V)) \to \Sh(\mathcal{D}/V)$ 如引理
\ref{sites-lemma-localize-morphism}，而
$j_{U/u(V)} : \Sh(\mathcal{C}/U) \to \Sh(\mathcal{C}/u(V))$ 如引理
\ref{sites-lemma-relocalize}。利用引理
\ref{sites-lemma-essential-image-j-shriek} 的等同
$\Sh(\mathcal{C}/U) = \Sh(\mathcal{C})/h_U^\#$ 与
$\Sh(\mathcal{D}/V) = \Sh(\mathcal{D})/h_V^\#$，函子
$(f_c)^{-1}$ 可由规则
$$
(f_c)^{-1}(\mathcal{H} \xrightarrow{\varphi} h_V^\#)
=
(f^{-1}\mathcal{H} \times_{f^{-1}\varphi, h_{u(V)}^\#, c} h_U^\#
\rightarrow h_U^\#).
$$
描述。最后，给定任意态射 $b : V' \to V$、$a : U' \to U$ 以及
$c' : U' \to u(V')$，使得
$$
\xymatrix{
U' \ar[r]_-{c'} \ar[d]_a & u(V') \ar[d]^{u(b)} \\
U \ar[r]^-c & u(V)
}
$$
交换，则图
$$
\xymatrix{
\Sh(\mathcal{C}/U') \ar[r]_{j_{U'/U}} \ar[d]_{f_{c'}} &
\Sh(\mathcal{C}/U) \ar[d]^{f_c} \\
\Sh(\mathcal{D}/V') \ar[r]^{j_{V'/V}} &
\Sh(\mathcal{D}/V).
}
$$
交换。
\end{lemma}

\begin{proof}
这一引理可谓自证，它更多是把理论发展到此处已经知道的事项汇集起来。例如，
第一个方形的交换性来自图 (\ref{sites-equation-relocalize}) 的交换性以及引理
\ref{sites-lemma-localize-morphism} 中图的交换性。对 $f_c^{-1}$ 的描述来自合并
引理 \ref{sites-lemma-relocalize-explicit} 与引理
\ref{sites-lemma-localize-morphism}。最后一个方形的交换性随即由等式
$$
f^{-1}\mathcal{H} \times_{h_{u(V)}^\#, c} h_U^\# \times_{h_U^\#} h_{U'}^\#
=
f^{-1}(\mathcal{H} \times_{h_V^\#} h_{V'}^\#)
\times_{h_{u(V'), c'}^\#} h_{U'}^\#
$$
得出；这在利用
$f^{-1}h_V^\# = h_{u(V)}^\#$ 和
$f^{-1}h_{V'}^\# = h_{u(V')}^\#$ 后是形式的，见引理
\ref{sites-lemma-pullback-representable-sheaf}。
% P01-ERRATA: P01-E0104 -- source puts c' inside h_{u(V'),c'}^#; expected fibre-product indexing separates the base h_{u(V')}^# and map c'.
\end{proof}

\noindent
下列引理给出另一种局部化的函子性，适用于拓扑斯态射来自余连续函子的情形。
这里的图与引理 \ref{sites-lemma-localize-morphism} 中的图不同；特别地，引理
\ref{sites-lemma-localize-morphism} 中的等式
$f'_*j_U^{-1} = j_V^{-1}f_*$ 在下列引理的情形中一般不成立。

\begin{lemma}
\label{sites-lemma-localize-cocontinuous}
设 $\mathcal{C}$、$\mathcal{D}$ 为位点，
$u : \mathcal{C} \to \mathcal{D}$ 是余连续函子。设 $U$ 是
$\mathcal{C}$ 的对象，并令 $V = u(U)$。有交换图
$$
\xymatrix{
\mathcal{C}/U \ar[r]_{j_U} \ar[d]_{u'} & \mathcal{C} \ar[d]^u \\
\mathcal{D}/V \ar[r]^-{j_V} & \mathcal{D}
}
$$
其中左竖直箭头为
$u' : \mathcal{C}/U \to \mathcal{D}/V$，$U'/U \mapsto V'/V$。
% P01-ERRATA: P01-E0105 -- V' is undefined; the functor value used in the proof is u(U')/V.
则 $u'$ 也余连续，并且得到拓扑斯的交换图
$$
\xymatrix{
\Sh(\mathcal{C}/U) \ar[r]_{j_U} \ar[d]_{f'} &
\Sh(\mathcal{C}) \ar[d]^f \\
\Sh(\mathcal{D}/V) \ar[r]^-{j_V} &
\Sh(\mathcal{D})
}
$$
其中 $f$（相应地 $f'$）对应于 $u$（相应地 $u'$）。
\end{lemma}

\begin{proof}
第一个图的交换性是显然的。只要证明 $u'$ 余连续，它就蕴含第二个图交换。

\medskip\noindent
设 $U'/U$ 是 $\mathcal{C}/U$ 的对象。设
$\{V_j/V \to u(U')/V\}_{j \in J}$ 是 $\mathcal{D}/V$ 中
$u(U')/V$ 的覆盖。由于 $u$ 余连续，存在覆盖
$\{U_i' \to U'\}_{i \in I}$，使族
$\{u(U_i') \to u(U')\}$ 加细 $\mathcal{D}$ 中的覆盖
$\{V_j \to u(U')\}$。换言之，存在指标集映射
$\alpha : I \to J$ 以及在 $U'$ 上的态射
$\phi_i : u(U_i') \to V_{\alpha(i)}$。
% P01-ERRATA: P01-E0106 -- these are morphisms in D and must be over u(U'), not over U' (an object of C).
通过复合 $U'_i \to U' \to U$，把 $U_i'$ 看作 $U$ 上的对象。则
$\{U'_i/U \to U'/U\}$ 是 $\mathcal{C}/U$ 的覆盖，使
$\{u(U_i')/V \to u(U')/V\}$ 加细
$\{V_j/V \to u(U')/V\}$（使用同一个 $\alpha$ 和同一些映射
$\phi_i$）。因此 $u' : \mathcal{C}/U \to \mathcal{D}/V$ 余连续。
\end{proof}

\begin{lemma}
\label{sites-lemma-localize-cocontinuous-downstairs}
设 $\mathcal{C}$、$\mathcal{D}$ 为位点，
$u : \mathcal{C} \to \mathcal{D}$ 是余连续函子。设 $V$ 是
$\mathcal{D}$ 的对象。令 ${}^u_V\mathcal{I}$ 为第
\ref{sites-section-more-functoriality-PSh} 节引入的范畴。有交换图
$$
\vcenter{
\xymatrix{
\,_V^u\mathcal{I} \ar[r]_j \ar[d]_{u'} &
\mathcal{C} \ar[d]^u \\
\mathcal{D}/V \ar[r]^-{j_V} &
\mathcal{D}
}
}
\quad\text{其中}\quad
\begin{matrix}
j : (U, \psi) \mapsto U \\
u' : (U, \psi) \mapsto (\psi : u(U) \to V)
\end{matrix}
$$
规定 ${}^u_V\mathcal{I}$ 的一族态射
$\{(U_i, \psi_i) \to (U, \psi)\}$ 是覆盖，当且仅当
$\{U_i \to U\}$ 是 $\mathcal{C}$ 中的覆盖。则
\begin{enumerate}
\item ${}^u_V\mathcal{I}$ 是位点；
\item $j$ 连续且余连续；
\item $u'$ 余连续；
\item 得到拓扑斯的交换图
$$
\xymatrix{
\Sh({}^u_V\mathcal{I}) \ar[r]_j \ar[d]_{f'} &
\Sh(\mathcal{C}) \ar[d]^f \\
\Sh(\mathcal{D}/V) \ar[r]^-{j_V} &
\Sh(\mathcal{D})
}
$$
其中 $f$（相应地 $f'$）对应于 $u$（相应地 $u'$）；且
\item 有 $f'_*j^{-1} = j_V^{-1}f_*$。
\end{enumerate}
\end{lemma}

\begin{proof}
(1)、(2)、(3) 和 (4) 是定义以及函子 $j$ 与纤维积交换这一事实的直接推论。
细节略去。为证明 (5)，回忆 $f_*$ 由 ${}_su = {}_pu$ 给出。因此，
$j_V^{-1}f_*\mathcal{F}$ 在 $V'/V$ 上的值，是
${}_pu\mathcal{F}$ 在 $V'$ 上的值，也就是 $\mathcal{F}$ 在范畴
${}^u_{V'}\mathcal{I}$ 上各值的极限。显然有范畴等价
$$
{}^u_{V'}\mathcal{I} \to {}^{u'}_{V'/V}\mathcal{I}
$$
另一方面，$f'_*j^{-1}\mathcal{F}$ 在 $V'/V$ 上的值由
$j^{-1}\mathcal{F}$ 在范畴 ${}^{u'}_{V'/V}\mathcal{I}$ 上各值的极限给出；
而且 $j^{-1}\mathcal{F}$ 在 ${}^u_V\mathcal{I}$ 的对象上的值正是
$\mathcal{F}$ 的值（由引理 \ref{sites-lemma-when-shriek}，因为 $j$ 连续且
余连续），故 (5) 成立。
\end{proof}

\noindent
下列两个结果的性质略有不同。

\begin{lemma}
\label{sites-lemma-special-square-cocontinuous}
假设给定位点 $\mathcal{C}', \mathcal{C}, \mathcal{D}', \mathcal{D}$ 以及函子
$$
\xymatrix{
\mathcal{C}' \ar[r]_{v'} \ar[d]_{u'} &
\mathcal{C} \ar[d]^u \\
\mathcal{D}' \ar[r]^v &
\mathcal{D}
}
$$
假设
\begin{enumerate}
\item $u$、$u'$、$v$ 与 $v'$ 余连续，并由引理
\ref{sites-lemma-cocontinuous-morphism-topoi} 分别产生拓扑斯态射
$f$、$f'$、$g$ 与 $g'$；
\item $v \circ u' = u \circ v'$；
\item $v$ 与 $v'$ 既连续又余连续；且
\item 对 $\mathcal{D}'$ 的任意对象 $V'$，由 $v$ 给出的函子
${}^{u'}_{V'}\mathcal{I} \to {}^{\ \ \ u}_{v(V')}\mathcal{I}$ 是余终的。
\end{enumerate}
则 $f'_* \circ (g')^{-1} = g^{-1} \circ f_*$，并且
$g'_! \circ (f')^{-1} = f^{-1} \circ g_!$。
\end{lemma}

\begin{proof}
范畴 ${}^{u'}_{V'}\mathcal{I}$ 与
${}^{\ \ \ u}_{v(V')}\mathcal{I}$ 定义于第
\ref{sites-section-more-functoriality-PSh} 节。条件 (4) 中的函子把
${}^{u'}_{V'}\mathcal{I}$ 的对象 $\psi : u'(U') \to V'$ 送到
${}^{\ \ \ u}_{v(V')}\mathcal{I}$ 的对象
$v(\psi) : uv'(U') = vu'(U') \to v(V')$。回忆，$g^{-1}$ 由
$v^p$ 给出（引理 \ref{sites-lemma-when-shriek}），而 $f_*$ 由
${}_su = {}_pu$ 给出。因此，$g^{-1}f_*\mathcal{F}$ 在 $V'$ 上的值，是
${}_pu\mathcal{F}$ 在 $v(V')$ 上的值，即极限
$$
\lim_{u(U) \to v(V') \in \Ob({}^{\ \ \ u}_{v(V')}\mathcal{I}^{opp})}
\mathcal{F}(U)
$$
同理，$f'_*(g')^{-1}\mathcal{F}$ 在 $V'$ 上的值由极限
$$
\lim_{u'(U') \to V' \in \Ob({}^{u'}_{V'}\mathcal{I}^{opp})} \mathcal{F}(v'(U'))
$$
给出。因此，假设 (4) 及《范畴》，引理 \ref{categories-lemma-initial}
表明二者相同，从而证明引理的第一个等式。第二个等式由第一个等式及伴随函子
的唯一性得出。
\end{proof}

\begin{lemma}
\label{sites-lemma-special-square-continuous}
假设给定位点 $\mathcal{C}', \mathcal{C}, \mathcal{D}', \mathcal{D}$ 以及函子
$$
\xymatrix{
\mathcal{C}' \ar[r]_{v'} &
\mathcal{C} \\
\mathcal{D}' \ar[r]^v \ar[u]^{u'} &
\mathcal{D} \ar[u]_u
}
$$
沿用第 \ref{sites-section-morphism-sites} 节与第
\ref{sites-section-cocontinuous-morphism-topoi} 节的记号，假设
\begin{enumerate}
\item $u$ 与 $u'$ 连续，并产生位点态射 $f$ 与 $f'$；
\item $v$ 与 $v'$ 余连续，并产生拓扑斯态射 $g$ 与 $g'$；
\item $u \circ v = v' \circ u'$；且
\item $v$ 与 $v'$ 既连续又余连续。
\end{enumerate}
则\footnote{在这一一般性下，我们不知道 $f \circ g'$ 与
$g \circ f'$ 作为拓扑斯态射是否相等（从前者到后者有典范 $2$-箭头，但它未必
是同构）。}
$f'_* \circ (g')^{-1} = g^{-1} \circ f_*$，并且
$g'_! \circ (f')^{-1} = f^{-1} \circ g_!$。
\end{lemma}

\begin{proof}
具体地，有
$$
f'_*(g')^{-1}\mathcal{F} = (u')^p((v')^p\mathcal{F})^\# =
(u')^p(v')^p\mathcal{F}
$$
第一个等式由定义得出，第二个等式由引理 \ref{sites-lemma-when-shriek} 得出。又有
$$
g^{-1}f_*\mathcal{F} = (v^pu^p\mathcal{F})^\# =
((u')^p(v')^p\mathcal{F})^\# = (u')^p(v')^p\mathcal{F}
$$
第一个等式由定义得出，第二个等式来自
$u \circ v = v' \circ u'$，第三个等式则因为已经看到
$(u')^p(v')^p\mathcal{F}$ 是层。这证明
$f'_* \circ (g')^{-1} = g^{-1} \circ f_*$；而等式
$g'_! \circ (f')^{-1} = f^{-1} \circ g_!$ 由左伴随的唯一性得出。
% P01-ERRATA: P01-E0107 -- source proof again uses fragmentary "The first equality by ..." clauses for both displays.
\end{proof}
















\section{拓扑斯的态射}
\label{sites-section-morphisms-topoi}

\noindent
本节证明，每个拓扑斯态射都等价于一个来自位点态射的拓扑斯态射。请与
\cite[Expos\'e III, Proposition 4.9.4]{SGA4} 比较。

\begin{lemma}
\label{sites-lemma-equivalence}
设 $\mathcal{C}$、$\mathcal{D}$ 为位点，
$u : \mathcal{C} \to \mathcal{D}$ 为函子。假设
\begin{enumerate}
\item $u$ 余连续；
\item $u$ 连续；
\item 给定 $\mathcal{C}$ 中的 $a, b : U' \to U$，若
$u(a) = u(b)$，则 $\mathcal{C}$ 中存在覆盖
$\{f_i : U'_i \to U'\}$，使 $a \circ f_i = b \circ f_i$；
\item 给定 $U', U \in \Ob(\mathcal{C})$ 以及 $\mathcal{D}$ 中的态射
$c : u(U') \to u(U)$，存在 $\mathcal{C}$ 中的覆盖
$\{f_i : U_i' \to U'\}$ 以及态射 $c_i : U_i' \to U$，使
$u(c_i) = c \circ u(f_i)$；且
\item 给定 $V \in \Ob(\mathcal{D})$，$\mathcal{D}$ 中存在一个形如
$\{u(U_i) \to V\}_{i \in I}$ 的 $V$ 的覆盖。
\end{enumerate}
则由引理 \ref{sites-lemma-cocontinuous-morphism-topoi} 与余连续函子 $u$ 相伴的
拓扑斯态射
$$
g : \Sh(\mathcal{C}) \longrightarrow \Sh(\mathcal{D})
$$
是等价。
\end{lemma}

\begin{proof}
假设 $u$ 满足性质 (1) -- (5)。将证明伴随映射
$$
\mathcal{G} \longrightarrow g_*g^{-1}\mathcal{G}
\quad\text{且}\quad
g^{-1}g_*\mathcal{F} \longrightarrow \mathcal{F}
$$
是同构。

\medskip\noindent
注意，引理 \ref{sites-lemma-when-shriek} 适用，并且对 $\mathcal{D}$ 上任意层
$\mathcal{G}$，有 $g^{-1}\mathcal{G}(U) = \mathcal{G}(u(U))$。接着，设
$\mathcal{F}$ 是 $\mathcal{C}$ 上的层，$V$ 是 $\mathcal{D}$ 的对象。由定义，
有 $g_*\mathcal{F}(V) = \lim_{u(U) \to V} \mathcal{F}(U)$。因此
$$
g^{-1}g_*\mathcal{F}(U) = \lim_{U', u(U') \to u(U)} \mathcal{F}(U')
$$
其中态射 $\psi : u(U') \to u(U)$ 不一定具有形式 $u(\alpha)$。这样一些序对
$(U', \psi)$ 所成的范畴有终对象，即 $(U, \text{id})$，它产生从该极限到
$\mathcal{F}(U)$ 的映射。
% P01-ERRATA: P01-E0109 -- (U,id) need not be final: a morphism from (U',psi) to it would require psi=u(alpha), which the preceding sentence explicitly says need not hold.
设 $(s_{(U', \psi)})$ 是该极限的一个元素。要证明
$s_{(U', \psi)}$ 由值
$s_{(U, \text{id})} \in \mathcal{F}(U)$ 唯一确定。由性质 (4)，给定任意
$(U', \psi)$，都存在覆盖 $\{U'_i \to U'\}$，使复合
$u(U'_i) \to u(U') \to u(U)$ 具有形式 $u(c_i)$，其中
$c_i : U'_i \to U$ 是 $\mathcal{C}$ 中的态射。因此
$$
s_{(U', \psi)}|_{U'_i} = c_i^*(s_{(U, \text{id})}).
$$
因为 $\mathcal{F}$ 是层，确实可得 $s_{(U', \psi)}$ 由
$s_{(U, \text{id})}$ 确定。这证明唯一性。

为证明存在性，任取
$s \in \mathcal{F}(U)$、$\psi : u(U') \to u(U)$、
$\{f_i : U_i' \to U'\}$ 以及如上的 $c_i : U_i' \to U$，使
$\psi \circ u(f_i) = u(c_i)$。我们断言存在（唯一的）元素
$s_{(U', \psi)} \in \mathcal{F}(U')$，使
$$
s_{(U', \psi)}|_{U'_i} = c_i^*(s).
$$
事实上，先验并不清楚元素
$c_i^*(s)|_{U_i' \times_{U'} U_j'}$ 与
$c_j^*(s)|_{U_i' \times_{U'} U_j'}$ 是否相等，因为图
$$
\xymatrix{
U_i' \times_{U'} U_j' \ar[r]_-{\text{pr}_2} \ar[d]_{\text{pr}_1} &
U_j' \ar[d]^{c_j} \\
U_i' \ar[r]^{c_i} & U}
$$
不一定交换。但引理的条件 (3) 保证存在覆盖
$\{f_{ijk} : U'_{ijk} \to U_i' \times_{U'} U_j'\}_{k \in K_{ij}}$，使
$c_i \circ \text{pr}_1 \circ f_{ijk} = c_j \circ \text{pr}_2 \circ f_{ijk}$。
因此
$$
f_{ijk}^* \left(c_i^*s|_{U_i' \times_{U'} U_j'}\right)
=
f_{ijk}^* \left(c_j^*s|_{U_i' \times_{U'} U_j'}\right)
$$
所以，由 $\mathcal{F}$ 的层条件，
$c_i^*(s)|_{U_i' \times_{U'} U_j'} = c_j^*(s)|_{U_i' \times_{U'} U_j'}$；
再由 $\mathcal{F}$ 的层条件，$s_{(U', \psi)}$ 存在。唯一性保证如此得到的族
$(s_{(U', \psi)})$ 是该极限的元素，并且 $s_{(U, \psi)} = s$。这证明
$g^{-1}g_*\mathcal{F} \to \mathcal{F}$ 是同构。
% P01-ERRATA: P01-E0108 -- source concludes s_(U,psi)=s although the distinguished component throughout is s_(U,id).

\medskip\noindent
设 $\mathcal{G}$ 是 $\mathcal{D}$ 上的层，$V$ 是 $\mathcal{D}$ 的对象。于是
$$
g_*g^{-1}\mathcal{G}(V) =
\lim_{U, \psi : u(U) \to V} \mathcal{G}(u(U))
$$
由前一段可见，层 $g_*g^{-1}\mathcal{G}$ 在形如 $V = u(U)$ 的对象
$V$ 上的值等于 $\mathcal{G}(u(U))$。（形式地，这是因为
$g^{-1}g_*g^{-1} \cong g^{-1}$ 以及证明开头给出的 $g^{-1}$ 的描述；非形式地，
只需比较这里和上面的极限。）因此，伴随映射
$\mathcal{G} \to g_*g^{-1}\mathcal{G}$ 在每个形如 $u(U)$ 的对象上的截面间
都是双射。由公理 (5)，存在由形如 $u(U)$ 的对象构成的 $V$ 的覆盖，因而容易
看出该伴随映射是同构。
\end{proof}

\noindent
给引理 \ref{sites-lemma-equivalence} 中那样的余连续函子一个名称会很方便。

\begin{definition}
\label{sites-definition-special-cocontinuous-functor}
设 $\mathcal{C}$、$\mathcal{D}$ 为位点。{\it 从 $\mathcal{C}$ 到
$\mathcal{D}$ 的特殊余连续函子 $u$}，是满足引理
\ref{sites-lemma-equivalence} 的假设与结论的余连续函子
$u : \mathcal{C} \to \mathcal{D}$。
\end{definition}

\begin{lemma}
\label{sites-lemma-localize-special-cocontinuous}
设 $\mathcal{C}$、$\mathcal{D}$ 为位点，
$u : \mathcal{C} \to \mathcal{D}$ 是特殊余连续函子。对
$\mathcal{C}$ 的每个对象 $U$，都有引理
\ref{sites-lemma-localize-cocontinuous} 中那样的交换图
$$
\xymatrix{
\mathcal{C}/U \ar[r]_{j_U} \ar[d] & \mathcal{C} \ar[d]^u \\
\mathcal{D}/u(U) \ar[r]^-{j_{u(U)}} & \mathcal{D}
}
$$
左竖直箭头是特殊余连续函子。因此，在拓扑斯的交换图
$$
\xymatrix{
\Sh(\mathcal{C}/U) \ar[r]_{j_U} \ar[d] &
\Sh(\mathcal{C}) \ar[d]^u \\
\Sh(\mathcal{D}/u(U)) \ar[r]^-{j_{u(U)}} &
\Sh(\mathcal{D})
}
$$
中，竖直箭头都是等价。
\end{lemma}

\begin{proof}
引理 \ref{sites-lemma-localize-cocontinuous} 已经给出这些图的存在性与交换性。
我们必须对诱导函子 $u : \mathcal{C}/U \to \mathcal{D}/u(U)$ 验证引理
\ref{sites-lemma-equivalence} 的假设 (1) -- (5)。这完全是机械的。

\medskip\noindent
性质 (1)。这是引理 \ref{sites-lemma-localize-cocontinuous}。

\medskip\noindent
性质 (2)。设 $\{U_i'/U \to U'/U\}_{i \in I}$ 是
$\mathcal{C}/U$ 中 $U'/U$ 的覆盖。由于 $u$ 连续，
$\{u(U_i')/u(U) \to u(U')/u(U)\}_{i \in I}$ 是
$\mathcal{D}/u(U)$ 中 $u(U')/u(U)$ 的覆盖。因此，(2) 对
$u : \mathcal{C}/U \to \mathcal{D}/u(U)$ 成立。

\medskip\noindent
性质 (3)。设 $a, b : U''/U \to U'/U$ 是 $\mathcal{C}/U$ 中满足
$u(a) = u(b)$ 于 $\mathcal{D}/u(U)$ 中的态射。由于 $u$ 满足 (3)，存在
$\mathcal{C}$ 中的覆盖 $\{f_i : U''_i \to U''\}$，使
$a \circ f_i = b \circ f_i$。这给出 $\mathcal{C}/U$ 中的覆盖
$\{f_i : U''_i/U \to U''/U\}$，使 $a \circ f_i = b \circ f_i$。
因此，(3) 对 $u : \mathcal{C}/U \to \mathcal{D}/u(U)$ 成立。

\medskip\noindent
性质 (4)。给定 $U''/U, U'/U \in \Ob(\mathcal{C}/U)$ 以及
$\mathcal{D}/u(U)$ 中的态射
$c : u(U'')/u(U) \to u(U')/u(U)$。由于 $u$ 满足性质 (4)，存在
$\mathcal{C}$ 中的覆盖 $\{f_i : U_i'' \to U''\}$ 以及态射
$c_i : U_i'' \to U'$，使 $u(c_i) = c \circ u(f_i)$。通过复合
$U_i'' \to U'' \to U$，把 $U_i''$ 看作 $U$ 上的对象。$c_i$ 未必是
$U$ 上的态射！但由于 $u(c_i)$ 是 $u(U)$ 上的态射，可以应用 $u$ 的性质
(3)，找到覆盖 $\{f_{ik} : U''_{ik} \to U''_i\}$，使
$c_{ik} = c_i \circ f_{ik} : U''_{ik} \to U'$ 是 $U$ 上的态射。因此，
$\{f_i \circ f_{ik} : U''_{ik}/U \to U''/U\}$ 是
$\mathcal{C}/U$ 中的覆盖，使 $u(c_{ik}) = c \circ u(f_{ik})$。
% P01-ERRATA: P01-E0110 -- source conclusion omits f_i: it should use u(f_i circ f_ik), matching the displayed covering and preceding equality.
因此，(4) 对 $u : \mathcal{C}/U \to \mathcal{D}/u(U)$ 成立。

\medskip\noindent
性质 (5)。设 $h : V \to u(U)$ 是 $\mathcal{D}/u(U)$ 的对象。由于
$u$ 满足性质 (5)，$\mathcal{D}$ 中存在覆盖
$\{c_i : u(U_i) \to V\}$。由性质 (4)，可以找到覆盖
$\{f_{ij} : U_{ij} \to U_i\}$ 和态射 $c_{ij} : U_{ij} \to U$，使
$u(c_{ij}) = h \circ c_i \circ u(f_{ij})$。因此，
$\{u(U_{ij})/u(U) \to V/u(U)\}$ 是 $\mathcal{D}/u(U)$ 中具有所需形式的
覆盖，从而得出 (5) 对 $u : \mathcal{C}/U \to \mathcal{D}/u(U)$ 成立。
\end{proof}

\begin{lemma}
\label{sites-lemma-special-equivalence}
设 $\mathcal{C}$ 为位点。设
$\mathcal{C}' \subset \Sh(\mathcal{C})$ 是满足下列条件的满子范畴（其对象构成
一个集合）：
\begin{enumerate}
\item 对所有 $U \in \Ob(\mathcal{C})$，有
$h_U^\# \in \Ob(\mathcal{C}')$；且
\item $\mathcal{C}'$ 在 $\Sh(\mathcal{C})$ 中的纤维积下保持不变。
\end{enumerate}
规定 $\mathcal{C}'$ 的覆盖是任意一族映射
$\{\mathcal{F}_i \to \mathcal{F}\}_{i \in I}$，使
$\coprod_{i \in I} \mathcal{F}_i \to \mathcal{F}$ 是层的满射。则
\begin{enumerate}
\item $\mathcal{C}'$ 是位点（选择一个覆盖集之后，见《集合》，引理
\ref{sets-lemma-coverings-site}）；
\item $\mathcal{C}'$ 上的可表预层是层（即 $\mathcal{C}'$ 上的拓扑是
次典范的，见定义 \ref{sites-definition-weaker-than-canonical}）；
\item 函子 $v : \mathcal{C} \to \mathcal{C}'$，$U \mapsto h_U^\#$ 是特殊
余连续函子，因而诱导等价
$g : \Sh(\mathcal{C}) \to \Sh(\mathcal{C}')$；
\item 对任意 $\mathcal{F} \in \Ob(\mathcal{C}')$，有
$g^{-1}h_\mathcal{F} = \mathcal{F}$；且
\item 对任意 $U \in \Ob(\mathcal{C})$，有
$g_*h_U^\# = h_{v(U)} = h_{h_U^\#}$。
\end{enumerate}
\end{lemma}

\begin{proof}
提醒：上面的某些陈述乍看可能有些令人困惑；这是因为 $\mathcal{C}'$ 的对象也
可以看作 $\mathcal{C}$ 上的层！
% P01-ERRATA: P01-E0111 -- source has the ungrammatical phrase "may look be a bit confusing".
我们略去证明引理所述的 $\mathcal{C}'$ 的覆盖满足定义
\ref{sites-definition-site} 的条件。

\medskip\noindent
假设 $\{\mathcal{F}_i \to \mathcal{F}\}$ 是层态射的满族。设
$\mathcal{G}$ 是另一个层。引理的 (2) 断言下图
$$
\xymatrix{
\Mor_{\Sh(\mathcal{C})}(
\coprod_{i \in I} \mathcal{F}_i, \mathcal{G})
\ar@<1ex>[r] \ar@<-1ex>[r]
&
\Mor_{\Sh(\mathcal{C})}(
\coprod_{(i_0, i_1) \in I \times I}
\mathcal{F}_{i_0} \times_\mathcal{F} \mathcal{F}_{i_1}, \mathcal{G})
}
$$
的等化子是 $\Mor_{\Sh(\mathcal{C})}(\mathcal{F}, \mathcal{G}).$
这是显然的（例如使用引理 \ref{sites-lemma-coequalizer-surjection}）。

\medskip\noindent
为证明 (3)，必须验证引理 \ref{sites-lemma-equivalence} 的条件 (1) -- (5)。
$v$ 余连续这一事实等价于引理 \ref{sites-lemma-mono-epi-sheaves} 中对层满射的
描述。函子 $v$ 连续，因为 $U \mapsto h_U^\#$ 与纤维积交换，并且把覆盖
变为覆盖（见引理 \ref{sites-lemma-sheafification-exact} 和引理
\ref{sites-lemma-covering-surjective-after-sheafification}）。引理
\ref{sites-lemma-equivalence} 的性质 (3)、(4) 是关于态射
$f : h_{U'}^\# \to h_U^\#$ 的陈述。这样的态射等同于
$h_U^\#(U')$ 的元素。因此，(3) 和 (4) 由层化的构造立即得出。引理
\ref{sites-lemma-equivalence} 的性质 (5) 是引理
\ref{sites-lemma-sheaf-coequalizer-representable}。以
$g : \Sh(\mathcal{C}) \to \Sh(\mathcal{C}')$ 表示由引理
\ref{sites-lemma-equivalence} 与 $v$ 相伴的拓扑斯等价。
% P01-ERRATA: P01-E0112 -- source "Denote g ... the equivalence" omits "by" or otherwise needs recasting.

\medskip\noindent
设 $\mathcal{F}$ 如引理的 (4)。对任意 $U \in \Ob(\mathcal{C})$，有
$$
g^{-1}h_\mathcal{F}(U) = h_\mathcal{F}(v(U))
= \Mor_{\Sh(\mathcal{C})}(h_U^\#, \mathcal{F})
= \mathcal{F}(U)
$$
第一个等式由引理 \ref{sites-lemma-when-shriek} 得出。因此 (4) 成立。
% P01-ERRATA: P01-E0113 -- source explanation is a sentence fragment: "The first equality by ...".

\medskip\noindent
设 $\mathcal{F} \in \Ob(\mathcal{C}')$，并设
$U \in \Ob(\mathcal{C})$。则
\begin{align*}
g_*h_U^\#(\mathcal{F})
& =
\Mor_{\Sh(\mathcal{C}')}(h_\mathcal{F}, g_*h_U^\#) \\
& =
\Mor_{\Sh(\mathcal{C})}(g^{-1}h_\mathcal{F}, h_U^\#) \\
& =
\Mor_{\Sh(\mathcal{C})}(\mathcal{F}, h_U^\#) \\
& =
\Mor_{\mathcal{C}'}(\mathcal{F}, h_U^\#)
\end{align*}
如所需（其中第三个等式已在上面证明）。
\end{proof}

\noindent
利用这一点，可以把任意拓扑斯整理为定义在具有所有有限极限的位点上。

\begin{lemma}
\label{sites-lemma-topos-good-site}
设 $\Sh(\mathcal{C})$ 为拓扑斯，$\{\mathcal{F}_i\}_{i \in I}$ 是
$\mathcal{C}$ 上的一个层集。存在拓扑斯等价
$g : \Sh(\mathcal{C}) \to \Sh(\mathcal{C}')$；它由位点间的特殊余连续函子
$u : \mathcal{C} \to \mathcal{C}'$ 诱导，并且
\begin{enumerate}
\item $\mathcal{C}'$ 具有次典范拓扑；
\item $\mathcal{C}'$ 的一族态射 $\{V_j \to V\}$ 是（组合等价于）
$\mathcal{C}'$ 的覆盖，当且仅当 $\coprod h_{V_j} \to h_V$ 是满射；
\item $\mathcal{C}'$ 有纤维积和终对象（即 $\mathcal{C}'$ 有所有有限极限）；
\item $\mathcal{C}'$ 上可表层的每个子层都是可表的；且
\item 每个 $g_*\mathcal{F}_i$ 都是可表层。
\end{enumerate}
\end{lemma}

\begin{proof}
考虑满子范畴 $\mathcal{C}_1 \subset \Sh(\mathcal{C})$，其对象包括所有
$U \in \Ob(\mathcal{C})$ 对应的 $h_U^\#$、给定的层
$\mathcal{F}_i$，以及终层 $*$（见例 \ref{sites-example-singleton-sheaf}）。
我们将归纳定义满子范畴
$$
\mathcal{C}_1 \subset \mathcal{C}_2 \subset \mathcal{C}_2 \subset \ldots
\subset \Sh(\mathcal{C})
$$
% P01-ERRATA: P01-E0114 -- source chain repeats C_2; the third term should be C_3.
具体地，给定 $\mathcal{C}_n$，令 $\mathcal{C}_{n + 1}$ 为由
$\mathcal{C}_n$ 的对象的一切纤维积和子层组成的满子范畴。（注意，
$\mathcal{C}_{n + 1}$ 有一个对象集。）令
$\mathcal{C}' = \bigcup_{n \geq 1} \mathcal{C}_n$。
$\mathcal{C}'$ 中的覆盖是 $\mathcal{C}'$ 的对象间的任意一族态射
$\{\mathcal{G}_j \to \mathcal{G}\}_{j \in J}$，使
$\coprod \mathcal{G}_j \to \mathcal{G}$ 作为 $\mathcal{C}$ 上的层映射是
满射。函子 $v : \mathcal{C} \to \mathcal{C'}$ 由
$U \mapsto h_U^\#$ 给出。应用引理 \ref{sites-lemma-special-equivalence}。
\end{proof}

\noindent
下面是本节的目标。

\begin{lemma}
\label{sites-lemma-morphism-topoi-comes-from-morphism-sites}
\begin{reference}
这一陈述与 \cite[Proposition 4.9.4. Expos\'e IV]{SGA4} 密切相关。为得到完整
结果，还应使用 \cite[Remarque 4.7.4, Expos\'e IV]{SGA4}。
\end{reference}
设 $\mathcal{C}$、$\mathcal{D}$ 为位点。设
$f : \Sh(\mathcal{C}) \to \Sh(\mathcal{D})$ 是拓扑斯态射。则存在位点
$\mathcal{C}'$ 以及函子图
$$
\xymatrix{
\mathcal{C} \ar[r]_v & \mathcal{C}' & \mathcal{D} \ar[l]^u
}
$$
使得
\begin{enumerate}
\item 函子 $v$ 是特殊余连续函子；
\item 函子 $u$ 与纤维积交换、连续，并定义位点态射
$\mathcal{C}' \to \mathcal{D}$；且
\item 拓扑斯态射 $f$ 与拓扑斯态射的复合
$$
\Sh(\mathcal{C}) \longrightarrow
\Sh(\mathcal{C}') \longrightarrow
\Sh(\mathcal{D})
$$
相符，其中第一个箭头由引理 \ref{sites-lemma-equivalence} 从 $v$ 得到，第二个
箭头由引理 \ref{sites-lemma-morphism-sites-topoi} 从 $u$ 得到。
\end{enumerate}
\end{lemma}

\begin{proof}
考虑满子范畴 $\mathcal{C}_1 \subset \Sh(\mathcal{C})$，其对象包括对所有
$U \in \Ob(\mathcal{C})$ 和所有 $V \in \Ob(\mathcal{D})$ 的
$h_U^\#$ 与 $f^{-1}h_V^\#$。令 $\mathcal{C}_{n + 1}$ 为由
$\mathcal{C}_n$ 的对象的一切纤维积组成的满子范畴。令
$\mathcal{C}' = \bigcup_{n \geq 1} \mathcal{C}_n$。
$\mathcal{C}'$ 中的覆盖是任意一族
$\{\mathcal{F}_i \to \mathcal{F}\}_{i \in I}$，使
$\coprod_{i \in I} \mathcal{F}_i \to \mathcal{F}$ 作为
$\mathcal{C}$ 上的层映射是满射。函子
$v : \mathcal{C} \to \mathcal{C'}$ 由 $U \mapsto h_U^\#$ 给出。函子
$u : \mathcal{D} \to \mathcal{C'}$ 由 $V \mapsto f^{-1}h_V^\#$ 给出。

\medskip\noindent
(1) 由引理 \ref{sites-lemma-special-equivalence} 得出。

\medskip\noindent
证明引理的 (2) 与 (3)。函子 $u$ 与纤维积交换，因为
$V \mapsto h_V^\#$ 和 $f^{-1}$ 都如此。此外，由于 $f^{-1}$ 正合并且与任意
余极限交换，可见它把覆盖变为层态射的满族。因此 $u$ 连续。为证明它定义
位点态射，还须证明 $u_s$ 正合。为此，将证明
$g^{-1} \circ u_s = f^{-1}$。事实上，届时因为 $g^{-1}$ 是等价并且
$f^{-1}$ 正合，即可得出结论。因为 $g^{-1}$ 伴随于 $g_*$，$u_s$ 伴随于
$u^s$，而 $f^{-1}$ 伴随于 $f_*$，所以证明
$u^s \circ g_* = f_*$ 也足够。设 $\mathcal{F}$ 是 $\mathcal{C}$ 上的层，
$V$ 是 $\mathcal{D}$ 的对象。则
\begin{align*}
(u^s g_{\ast}\mathcal{F})(V)
& =
(g_{\ast}\mathcal{F})(f^{-1}h_V^\#) \\
& =
\Mor_{Sh(\mathcal{C}')}(h_{f^{-1}h_V^\#},g_{\ast}\mathcal{F}) \\
& =
\Mor_{Sh(\mathcal{C})}(g^{-1}h_{f^{-1}h_V^\#},\mathcal{F}) \\
& =
\Mor_{Sh(\mathcal{C})}(f^{-1}h_V^\#,\mathcal{F}) \\
& =
\Mor_{Sh(\mathcal{D})}(h_V^\#,f_{\ast}\mathcal{F}) \\
& =
f_{\ast}\mathcal{F}(V)
\end{align*}
第一个等式是因为 $u^s = u^p$。第二个等式是米田引理。第三个等式由
$g^{-1}$ 与 $g_*$ 的伴随性得出。第四个等式由引理
\ref{sites-lemma-special-equivalence} (4) 得出。第五个等式由 $f^{-1}$ 与 $f_*$ 的
伴随性得出。第六个等式由米田引理得出。因此
$u^s g_*\mathcal{F} = f_*\mathcal{F}$，从而完成引理的证明。
% P01-ERRATA: P01-E0115 -- source gives several equality explanations as sentence fragments ("The first equality because", "The third equality by", etc.).
\end{proof}

\begin{remark}
\label{sites-remark-morphism-topoi-comes-from-morphism-sites}
记号和假设如引理 \ref{sites-lemma-morphism-topoi-comes-from-morphism-sites}。若位点
$\mathcal{D}$ 有终对象和纤维积，则函子
$u : \mathcal{D} \to \mathcal{C}'$ 满足命题
\ref{sites-proposition-get-morphism} 的所有假设。事实上，除引理中所述的性质外，
$u$ 还把 $\mathcal{D}$ 的终对象变为 $\mathcal{C}'$ 的终对象。这由 $u$ 的
构造显然可知。因此，若先把引理 \ref{sites-lemma-topos-good-site} 应用于
$\mathcal{D}$，再把引理
\ref{sites-lemma-morphism-topoi-comes-from-morphism-sites} 应用于所得拓扑斯态射
$\Sh(\mathcal{C}) \to \Sh(\mathcal{D}')$，便得到下列陈述：
% P01-ERRATA: P01-E0116 -- source says plural "Lemmas" before a single reference to lemma-topos-good-site.
任意拓扑斯态射 $f : \Sh(\mathcal{C}) \to \Sh(\mathcal{D})$ 都嵌入交换图
$$
\xymatrix{
\Sh(\mathcal{C}) \ar[d]_g \ar[r]_f &
\Sh(\mathcal{D}) \ar[d]^e \\
\Sh(\mathcal{C}') \ar[r]^{f'} &
\Sh(\mathcal{D}')
}
$$
并满足下列性质：
\begin{enumerate}
\item 态射 $e$ 与 $g$ 是由特殊余连续函子
$\mathcal{C} \to \mathcal{C}'$ 和 $\mathcal{D} \to \mathcal{D}'$ 给出的等价；
\item 位点 $\mathcal{C}'$ 与 $\mathcal{D}'$ 有纤维积、终对象，并具有
次典范拓扑；
\item 态射 $f' : \mathcal{C}' \to \mathcal{D}'$ 来自一个位点态射，它对应于
命题 \ref{sites-proposition-get-morphism} 适用的函子
$u : \mathcal{D}' \to \mathcal{C}'$；且
\item 给定 $\mathcal{C}$（相应地 $\mathcal{D}$）上的任意层集
$\mathcal{F}_i$（相应地 $\mathcal{G}_j$），可以假设其中每个层在
$\mathcal{C}'$（相应地 $\mathcal{D}'$）上都是可表层。
\end{enumerate}
把 $\mathcal{C}$ 和 $\mathcal{D}$ 替换为 $\mathcal{C}'$ 和
$\mathcal{D}'$ 往往很有用。
\end{remark}

\begin{remark}
\label{sites-remark-equivalence-topoi-comes-from-morphism-sites}
记号和假设如引理 \ref{sites-lemma-morphism-topoi-comes-from-morphism-sites}。进一步
假设原拓扑斯态射 $\Sh(\mathcal{C}) \to \Sh(\mathcal{D})$ 是等价。则引理
\ref{sites-lemma-morphism-topoi-comes-from-morphism-sites} 的证明中的构造给出两个
函子
$$
\mathcal{C} \rightarrow \mathcal{C}' \leftarrow \mathcal{D}
$$
它们都是特殊余连续函子。因此，在这种情况下，实际上可以把拓扑斯态射分解为
复合
$$
\Sh(\mathcal{C}) \rightarrow
\Sh(\mathcal{C}') =
\Sh(\mathcal{D}') \leftarrow
\Sh(\mathcal{D})
$$
如备注 \ref{sites-remark-morphism-topoi-comes-from-morphism-sites}，但中间态射为恒等
态射。
\end{remark}









\section{拓扑斯的局部化}
\label{sites-section-localize-topoi}

\noindent
我们把局部化的部分材料重复到表面上更一般的拓扑斯情形。实际上，这并不更
一般，因为总可以扩张底层位点，从而假设在位点的对象处作局部化。

\begin{lemma}
\label{sites-lemma-localize-topos}
设 $\mathcal{C}$ 为位点，$\mathcal{F}$ 是 $\mathcal{C}$ 上的层。则范畴
$\Sh(\mathcal{C})/\mathcal{F}$ 是拓扑斯。存在典范拓扑斯态射
$$
j_\mathcal{F} :
\Sh(\mathcal{C})/\mathcal{F}
\longrightarrow
\Sh(\mathcal{C})
$$
它是第 \ref{sites-section-localize} 节意义下的局部化，并且
\begin{enumerate}
\item 函子 $j_\mathcal{F}^{-1}$ 是函子
$\mathcal{H} \mapsto \mathcal{H} \times \mathcal{F}/\mathcal{F}$；且
\item 函子 $j_{\mathcal{F}!}$ 是遗忘函子
$\mathcal{G}/\mathcal{F} \mapsto \mathcal{G}$。
\end{enumerate}
\end{lemma}

\begin{proof}
应用引理 \ref{sites-lemma-topos-good-site}。这意味着可以假设 $\mathcal{C}$ 是
具有次典范拓扑的位点，并且对某个 $U \in \Ob(\mathcal{C})$ 有
$\mathcal{F} = h_U = h_U^\#$。因此，第 \ref{sites-section-localize} 节的材料适用。
特别地，存在等价
$\Sh(\mathcal{C}/U) = \Sh(\mathcal{C})/h_U^\#$，使复合
$$
\Sh(\mathcal{C}/U)
\to
\Sh(\mathcal{C})/h_U^\#
\to \Sh(\mathcal{C})
$$
等于 $j_{U!}$，见引理 \ref{sites-lemma-essential-image-j-shriek}。以
$a : \Sh(\mathcal{C})/h_U^\# \to \Sh(\mathcal{C}/U)$ 表示逆函子，于是
$j_{\mathcal{F}!} = j_{U!} \circ a$、
$j_\mathcal{F}^{-1} = a^{-1} \circ j_U^{-1}$，以及
$j_{\mathcal{F}, *} = j_{U, *} \circ a$。
% P01-ERRATA: P01-E0117 -- source "Denote a ... the inverse functor" omits "by" or otherwise needs recasting.
对 $j_{\mathcal{F}!}$ 的描述由上述内容得出。对
$j_\mathcal{F}^{-1}$ 的描述由引理
\ref{sites-lemma-compute-j-shriek-restrict} 得出。
\end{proof}

\begin{lemma}
\label{sites-lemma-localize-pushforward}
在引理 \ref{sites-lemma-localize-topos} 的情形中，函子
$j_{\mathcal{F}, *}$ 把 $\varphi : \mathcal{G} \to \mathcal{F}$ 送到层
$$
U
\longmapsto
\{\alpha : \mathcal{F}|_U \to \mathcal{G}|_U
\text{ 使得 } \alpha \text{ 是下列映射的右逆： }\varphi|_U \}.
$$
% P01-ERRATA: P01-E0118 -- source says "is the one associates to" rather than "is the one that associates to".
\end{lemma}

\begin{proof}
对任意 $\varphi : \mathcal{G} \to \mathcal{F}$，用记号
$\mathcal{G}/\mathcal{F}$ 表示 $\Sh(\mathcal{C})/\mathcal{F}$ 的相应对象。有
$$
(j_{\mathcal{F}, *}(\mathcal{G}/\mathcal{F}))(U) =
\Mor_{\Sh(\mathcal{C})}(h_U^\#, j_{\mathcal{F}, *}(\mathcal{G}/\mathcal{F})) =
\Mor_{\Sh(\mathcal{C})/\mathcal{F}}(j_{\mathcal{F}}^{-1}h_U^{\#},
(\mathcal{G}/\mathcal{F})).
$$
由引理 \ref{sites-lemma-localize-topos}，该集合是在集合映射
$$
\Mor_{\Sh(\mathcal{C})}(h_U^\# \times \mathcal{F}, \mathcal{G})
\xrightarrow{\varphi \circ}
\Mor_{\Sh(\mathcal{C})}(h_U^\# \times \mathcal{F}, \mathcal{F}).
$$
下元素 $h_U^\# \times \mathcal{F} \to \mathcal{F}$ 上的纤维。由引理
\ref{sites-lemma-internal-hom-sheaf} 中的伴随，有
\begin{align*}
\Mor_{\Sh(\mathcal{C})}(h_U^{\#}\times\mathcal{F}, \mathcal{G})
& =
\Mor_{\Sh(\mathcal{C})}(h_U^{\#},\SheafHom(\mathcal{F}, \mathcal{G})) \\
& =
\Mor_{\Sh(\mathcal{C}/U)}(\mathcal{F}|_{\mathcal{C}/U},
\mathcal{G}|_{\mathcal{C}/U}), \\
\Mor_{\Sh(\mathcal{C})}(h_U^{\#} \times \mathcal{F}, \mathcal{F})
& =
\Mor_{\Sh(\mathcal{C})}(h_U^{\#},\SheafHom(\mathcal{F},\mathcal{F})) \\
& =
\Mor_{\Sh(\mathcal{C}/U)}(\mathcal{F}|_{\mathcal{C}/U},
\mathcal{F}|_{\mathcal{C}/U}),
\end{align*}
在此伴随下，映射 $\varphi\circ$ 变为映射
$$
\Mor_{\Sh(\mathcal{C}/U)}(\mathcal{F}|_{\mathcal{C}/U},
\mathcal{G}|_{\mathcal{C}/U})
\longrightarrow
\Mor_{\Sh(\mathcal{C}/U)}(\mathcal{F}|_{\mathcal{C}/U},
\mathcal{F}|_{\mathcal{C}/U}),\quad
\psi \longmapsto \varphi|_{\mathcal{C}/U} \circ \psi,
$$
而元素 $h_U^\# \times \mathcal{F} \to \mathcal{F}$ 变为
$\text{id}_{\mathcal{F}|_{\mathcal{C}/U}}$。因此
$(j_{\mathcal{F}, *}\mathcal{G}/\mathcal{F})(U)$ 同构于映射
$$
\Mor_{\Sh(\mathcal{C}/U)}(\mathcal{F}|_{\mathcal{C}/U},
\mathcal{G}|_{\mathcal{C}/U})
\xrightarrow{\varphi|_{\mathcal{C}/U}\circ}
\Mor_{\Sh(\mathcal{C}/U)}(\mathcal{F}|_{\mathcal{C}/U},
\mathcal{F}|_{\mathcal{C}/U}),
$$
在 $\text{id}_{\mathcal{F}|_{\mathcal{C}/U}}$ 上的纤维，即
$\{\alpha : \mathcal{F}|_U \to \mathcal{G}|_U
\text{ 使得 } \alpha \text{ 是下列映射的右逆： }\varphi|_U \}$，
如所需。
\end{proof}




\begin{lemma}
\label{sites-lemma-localize-topos-site}
设 $\mathcal{C}$ 为位点，$\mathcal{F}$ 是 $\mathcal{C}$ 上的层。令
$\mathcal{C}/\mathcal{F}$ 为序对 $(U, s)$ 所成的范畴，其中
$U \in \Ob(\mathcal{C})$ 且 $s \in \mathcal{F}(U)$。规定
$\mathcal{C}/\mathcal{F}$ 的覆盖是一族
$\{(U_i, s_i) \to (U, s)\}$，使 $\{U_i \to U\}$ 是
$\mathcal{C}$ 的覆盖。则
$j : \mathcal{C}/\mathcal{F} \to \mathcal{C}$ 是位点的连续且余连续函子，
它诱导拓扑斯态射
$j : \Sh(\mathcal{C}/\mathcal{F}) \to \Sh(\mathcal{C})$。事实上，存在
等价 $\Sh(\mathcal{C}/\mathcal{F}) =
\Sh(\mathcal{C})/\mathcal{F}$，它把 $j$ 变为 $j_\mathcal{F}$。
\end{lemma}

\begin{proof}
我们略去验证 $\mathcal{C}/\mathcal{F}$ 是位点以及 $j$ 连续且余连续。由引理
\ref{sites-lemma-when-shriek}，存在所述拓扑斯态射 $j$，满足
$j^{-1}\mathcal{G}(U, s) = \mathcal{G}(U)$，并且 $j^{-1}$ 有左伴随
$j_!$。$\mathcal{C}/\mathcal{F}$ 上的态射
$\varphi : * \to j^{-1}\mathcal{G}$，等同于一条规则：给每个序对
$(U, s)$ 指定截面 $\varphi(s) \in \mathcal{G}(U)$，并与限制映射相容。
因此，这又等同于 $\mathcal{C}$ 上的态射
$\varphi : \mathcal{F} \to \mathcal{G}$。得出 $j_!* = \mathcal{F}$。特别地，
对每个 $\mathcal{H} \in \Sh(\mathcal{C}/\mathcal{F})$，都有典范映射
$$
j_!\mathcal{H} \to j_!* = \mathcal{F}
$$
也就是说，得到函子
$j'_! : \Sh(\mathcal{C}/\mathcal{F}) \to
\Sh(\mathcal{C})/\mathcal{F}$。该函子的逆，是把
$\Sh(\mathcal{C})/\mathcal{F}$ 的对象
$\varphi : \mathcal{G} \to \mathcal{F}$ 送到层
$$
a(\mathcal{G}/\mathcal{F}) :
(U, s) \longmapsto \{t \in \mathcal{G}(U) \mid \varphi(t) = s\}
$$
的规则。我们略去验证 $a(\mathcal{G}/\mathcal{F})$ 是层以及 $a$ 是
$j'_!$ 的逆。
\end{proof}

\begin{definition}
\label{sites-definition-localize-topos}
设 $\mathcal{C}$ 为位点，$\mathcal{F}$ 是 $\mathcal{C}$ 上的层。
\begin{enumerate}
\item 拓扑斯 $\Sh(\mathcal{C})/\mathcal{F}$ 称为
{\it 拓扑斯 $\Sh(\mathcal{C})$ 在 $\mathcal{F}$ 处的局部化}。
\item 引理 \ref{sites-lemma-localize-topos} 的拓扑斯态射
$j_\mathcal{F} :
\Sh(\mathcal{C})/\mathcal{F}
\to
\Sh(\mathcal{C})$ 称为{\it 局部化态射}。
\end{enumerate}
\end{definition}

\noindent
我们将证明：每当层 $\mathcal{F}$ 对位点的某个对象 $U$ 等于
$h_U^\#$ 时，拓扑斯的局部化就等于位点在 $U$ 处的局部化上的层范畴。此外，
还将验证所有函子性都与这一等同相容。

\begin{lemma}
\label{sites-lemma-localize-compare}
设 $\mathcal{C}$ 为位点，对 $\mathcal{C}$ 的某个对象 $U$，令
$\mathcal{F} = h_U^\#$。则引理 \ref{sites-lemma-localize-topos} 构造的
$j_\mathcal{F} : \Sh(\mathcal{C})/\mathcal{F}
\to \Sh(\mathcal{C})$ 与第 \ref{sites-section-localize} 节构造的拓扑斯态射
$j_U : \Sh(\mathcal{C}/U) \to \Sh(\mathcal{C})$ 相符；这里使用引理
\ref{sites-lemma-essential-image-j-shriek} 的等同
$\Sh(\mathcal{C}/U) = \Sh(\mathcal{C})/h_U^\#$。
\end{lemma}

\begin{proof}
我们已在引理 \ref{sites-lemma-essential-image-j-shriek} 中看到，复合
$\Sh(\mathcal{C}/U) \to \Sh(\mathcal{C})/h_U^\#
\to \Sh(\mathcal{C})$ 是 $j_{U!}$。由引理
\ref{sites-lemma-localize-topos}，函子
$\Sh(\mathcal{C})/h_U^\# \to \Sh(\mathcal{C})$ 是
$j_{\mathcal{F}!}$。因此，在这一等同下
$j_{\mathcal{F}!} = j_{U!}$。所以
$j_\mathcal{F}^{-1} = j_U^{-1}$（由伴随性），进而
$j_{\mathcal{F}, *} = j_{U, *}$（再次由伴随性）。
\end{proof}

\begin{lemma}
\label{sites-lemma-relocalize-topos}
设 $\mathcal{C}$ 为位点。若
$s : \mathcal{G} \to \mathcal{F}$ 是 $\mathcal{C}$ 上的层态射，则存在自然交换的
拓扑斯态射图式
$$
\xymatrix{
\Sh(\mathcal{C})/\mathcal{G} \ar[rd]_{j_\mathcal{G}} \ar[rr]_j & &
\Sh(\mathcal{C})/\mathcal{F} \ar[ld]^{j_\mathcal{F}} \\
& \Sh(\mathcal{C}) &
}
$$
其中 $j = j_{\mathcal{G}/\mathcal{F}}$ 是拓扑斯
$\Sh(\mathcal{C})/\mathcal{F}$ 在对象 $\mathcal{G}/\mathcal{F}$ 处的局部化。
特别地，有
$$
j^{-1}(\mathcal{H} \to \mathcal{F}) =
(\mathcal{H} \times_\mathcal{F} \mathcal{G} \to \mathcal{G})
$$
以及
$$
j_!(\mathcal{E} \xrightarrow{e} \mathcal{G}) =
(\mathcal{E} \xrightarrow{s \circ e} \mathcal{F}).
$$
\end{lemma}

\begin{proof}
对 $j^{-1}$ 与 $j_!$ 的描述来自引理
\ref{sites-lemma-localize-topos} 中对这些函子的描述。由这些函子作为遗忘函子的描述，
函子等式 $j_{\mathcal{G}!} = j_{\mathcal{F}!} \circ j_!$ 是清楚的。由伴随性还得到等式
$j_\mathcal{G}^{-1} = j^{-1} \circ j_\mathcal{F}^{-1}$ 以及
$j_{\mathcal{G}, *} = j_{\mathcal{F}, *} \circ j_*$。
\end{proof}

\begin{lemma}
\label{sites-lemma-relocalize-compare}
设 $\mathcal{C}$ 与 $s : \mathcal{G} \to \mathcal{F}$ 如引理
\ref{sites-lemma-relocalize-topos} 中所述。若
$\mathcal{G} = h_V^\#$、$\mathcal{F} = h_U^\#$，且
$s : \mathcal{G} \to \mathcal{F}$ 来自 $\mathcal{C}$ 的态射 $V \to U$，
则引理 \ref{sites-lemma-relocalize-topos} 中的图式等同于图式
(\ref{sites-equation-relocalize})；这里使用等同
$\Sh(\mathcal{C}/V) = \Sh(\mathcal{C})/h_V^\#$
与
$\Sh(\mathcal{C}/U) = \Sh(\mathcal{C})/h_U^\#$，
它们来自引理 \ref{sites-lemma-essential-image-j-shriek}。
\end{lemma}

\begin{proof}
这是因为对 $j^{-1}$ 的描述相符。见引理
\ref{sites-lemma-relocalize-explicit} 与引理
\ref{sites-lemma-relocalize-topos}。
\end{proof}






\section{局部化与拓扑斯态射}
\label{sites-section-localize-morphisms-topoi}

\noindent
本节是第 \ref{sites-section-localize-morphisms} 节对拓扑斯态射的类似版本。

\begin{lemma}
\label{sites-lemma-localize-morphism-topoi}
设 $f : \Sh(\mathcal{C}) \to \Sh(\mathcal{D})$
为拓扑斯态射，$\mathcal{G}$ 是 $\mathcal{D}$ 上的层。令
$\mathcal{F} = f^{-1}\mathcal{G}$。则存在拓扑斯的交换图式
$$
\xymatrix{
\Sh(\mathcal{C})/\mathcal{F} \ar[r]_{j_\mathcal{F}} \ar[d]_{f'} &
\Sh(\mathcal{C}) \ar[d]^f \\
\Sh(\mathcal{D})/\mathcal{G} \ar[r]^{j_\mathcal{G}} &
\Sh(\mathcal{D}).
}
$$
态射 $f'$ 由性质
$$
(f')^{-1}(\mathcal{H} \xrightarrow{\varphi} \mathcal{G})
=
(f^{-1}\mathcal{H} \xrightarrow{f^{-1}\varphi} \mathcal{F})
$$
刻画，并且有 $f'_*j_\mathcal{F}^{-1} = j_\mathcal{G}^{-1}f_*$。
\end{lemma}

\begin{proof}
由于命题是关于拓扑斯的，并不涉及底层位点，我们可以任意更换位点。因此，
由评注 \ref{sites-remark-morphism-topoi-comes-from-morphism-sites} 中的讨论，可以假设
$f$ 由连续函子 $u : \mathcal{D} \to \mathcal{C}$ 给出；该函子满足命题
\ref{sites-proposition-get-morphism} 的假设，两个位点都具有所有有限极限与次典范拓扑，
并且对 $\mathcal{D}$ 的某个对象 $V$ 有 $\mathcal{G} = h_V$。于是由引理
\ref{sites-lemma-pullback-representable-sheaf}，
$\mathcal{F} = f^{-1}h_V = h_{u(V)}$。由引理
\ref{sites-lemma-localize-morphism}，得到拓扑斯态射的交换图式
% P01-ERRATA: P01-E0119 -- the proof uses U below without defining U; it should set U = u(V) or use u(V) explicitly.
$$
\xymatrix{
\Sh(\mathcal{C}/U) \ar[r]_{j_U} \ar[d]_{f'} &
\Sh(\mathcal{C}) \ar[d]^f \\
\Sh(\mathcal{D}/V) \ar[r]^{j_V} &
\Sh(\mathcal{D}),
}
$$
并有 $f'_*j_U^{-1} = j_V^{-1}f_*$。由引理
\ref{sites-lemma-localize-compare}，可以把 $j_\mathcal{F}$ 与 $j_U$ 等同，
把 $j_\mathcal{G}$ 与 $j_V$ 等同。对 $(f')^{-1}$ 的描述见引理
\ref{sites-lemma-localize-morphism}。
\end{proof}

\begin{lemma}
\label{sites-lemma-localize-morphism-compare}
设 $f : \mathcal{C} \to \mathcal{D}$ 是由连续函子
$u : \mathcal{D} \to \mathcal{C}$ 给出的位点态射。设 $V$ 是
$\mathcal{D}$ 的对象，令 $U = u(V)$。令
$\mathcal{G} = h_V^\#$，并令
$\mathcal{F} = h_U^\# = f^{-1}h_V^\#$（见引理
\ref{sites-lemma-pullback-representable-sheaf}）。则引理
\ref{sites-lemma-localize-morphism-topoi} 的拓扑斯态射图式与引理
\ref{sites-lemma-localize-morphism} 的拓扑斯态射图式相符；这里使用引理
\ref{sites-lemma-localize-compare} 的等同
$j_\mathcal{F}= j_U$ 与 $j_\mathcal{G} = j_V$。
\end{lemma}

\begin{proof}
这并非全然平凡，因为在引理
\ref{sites-lemma-localize-morphism-topoi} 的证明中，为给出 $f$ 而选取的位点态射，
可能不同于本引理给定的位点态射。但在两种情形下，函子 $(f')^{-1}$ 都由同一规则
描述。因此它们相符，所关联的拓扑斯态射也相同。略去一些细节。
\end{proof}

\begin{lemma}
\label{sites-lemma-relocalize-morphism-topoi}
设 $f : \Sh(\mathcal{C}) \to \Sh(\mathcal{D})$
为拓扑斯态射。设 $\mathcal{G} \in \Sh(\mathcal{D})$、
$\mathcal{F} \in \Sh(\mathcal{C})$，且
$s : \mathcal{F} \to f^{-1}\mathcal{G}$ 是层态射。存在拓扑斯的交换图式
$$
\xymatrix{
\Sh(\mathcal{C})/\mathcal{F} \ar[r]_{j_\mathcal{F}} \ar[d]_{f_s} &
\Sh(\mathcal{C}) \ar[d]^f \\
\Sh(\mathcal{D})/\mathcal{G} \ar[r]^{j_\mathcal{G}} &
\Sh(\mathcal{D}).
}
$$
有 $f_s = f' \circ j_{\mathcal{F}/f^{-1}\mathcal{G}}$，其中
$f' :
\Sh(\mathcal{C})/f^{-1}\mathcal{G}
\to
\Sh(\mathcal{D})/\mathcal{F}$
% P01-ERRATA: P01-E0120 -- the codomain above should be Sh(D)/G; F is a sheaf on C and is ill-typed here.
如引理 \ref{sites-lemma-localize-morphism-topoi} 中所述，而
$j_{\mathcal{F}/f^{-1}\mathcal{G}} :
\Sh(\mathcal{C})/\mathcal{F}
\to
\Sh(\mathcal{C})/f^{-1}\mathcal{G}$
如引理 \ref{sites-lemma-relocalize-topos} 中所述。函子 $(f_s)^{-1}$ 由规则
$$
(f_s)^{-1}(\mathcal{H} \xrightarrow{\varphi} \mathcal{G})
=
(f^{-1}\mathcal{H} \times_{f^{-1}\varphi, f^{-1}\mathcal{G}, s} \mathcal{F}
\rightarrow \mathcal{F}).
$$
刻画。最后，给定任意态射 $b : \mathcal{G}' \to \mathcal{G}$、
$a : \mathcal{F}' \to \mathcal{F}$ 与
$s' : \mathcal{F}' \to f^{-1}\mathcal{G}'$，使得
$$
\xymatrix{
\mathcal{F}' \ar[r]_-{s'} \ar[d]_a & f^{-1}\mathcal{G}' \ar[d]^{f^{-1}b} \\
\mathcal{F} \ar[r]^-s & f^{-1}\mathcal{G}
}
$$
交换，则图式
$$
\xymatrix{
\Sh(\mathcal{C})/\mathcal{F}'
\ar[r]_{j_{\mathcal{F}'/\mathcal{F}}} \ar[d]_{f_{s'}} &
\Sh(\mathcal{C})/\mathcal{F} \ar[d]^{f_s} \\
\Sh(\mathcal{D})/\mathcal{G}' \ar[r]^{j_{\mathcal{G}'/\mathcal{G}}} &
\Sh(\mathcal{D})/\mathcal{G}.
}
$$
交换。
\end{lemma}

\begin{proof}
第一个方块的交换性由引理
\ref{sites-lemma-relocalize-topos} 中图式的交换性与引理
\ref{sites-lemma-localize-morphism-topoi} 中图式的交换性得出。对
$f_s^{-1}$ 的描述，则由引理
\ref{sites-lemma-localize-morphism-topoi} 中对 $(f')^{-1}$ 的描述与引理
\ref{sites-lemma-relocalize-topos} 中对
$(j_{\mathcal{F}/f^{-1}\mathcal{G}})^{-1}$ 的描述合并得出。最后一个方块的
交换性于是由等式
$$
f^{-1}\mathcal{H} \times_{f^{-1}\mathcal{G}, s} \mathcal{F}
\times_\mathcal{F} \mathcal{F}'
=
f^{-1}(\mathcal{H} \times_\mathcal{G} \mathcal{G}')
\times_{f^{-1}\mathcal{G}', s'} \mathcal{F}'
$$
得出，该等式是形式的。
\end{proof}

\begin{lemma}
\label{sites-lemma-relocalize-morphism-compare}
设 $f : \mathcal{C} \to \mathcal{D}$ 是由连续函子
$u : \mathcal{D} \to \mathcal{C}$ 给出的位点态射。设 $V$ 是
$\mathcal{D}$ 的对象，$c : U \to u(V)$ 是态射。令
$\mathcal{G} = h_V^\#$ 且 $\mathcal{F} = h_U^\# = f^{-1}h_V^\#$。
% P01-ERRATA: P01-E0121 -- h_U# need not equal f^{-1}h_V# = h_{u(V)}#; the statement should only set G = h_V# and F = h_U#.
令 $s : \mathcal{F} \to f^{-1}\mathcal{G}$ 为 $c$ 诱导的映射。则引理
\ref{sites-lemma-relocalize-morphism} 的拓扑斯态射图式与引理
\ref{sites-lemma-relocalize-morphism-topoi} 的拓扑斯态射图式相符；这里使用引理
\ref{sites-lemma-localize-compare} 的等同
$j_\mathcal{F} = j_U$ 与 $j_\mathcal{G} = j_V$。
\end{lemma}

\begin{proof}
这由合并引理 \ref{sites-lemma-relocalize-compare} 与
\ref{sites-lemma-localize-morphism-compare} 得出。
\end{proof}



















\section{点}
\label{sites-section-points}

\begin{definition}
\label{sites-definition-point-topos}
设 $\mathcal{C}$ 为位点。{it 拓扑斯 $\Sh(\mathcal{C})$ 的点}，是从
$\Sh(pt)$ 到 $\Sh(\mathcal{C})$ 的拓扑斯态射 $p$。
\end{definition}

\noindent
我们将用函子 $u : \mathcal{C} \to \textit{Sets}$ 定义位点的点。稍后会看到，
$u$ 将定义一个位点态射，它诱导与 $\mathcal{C}$ 关联的拓扑斯的一个点，见引理
\ref{sites-lemma-site-point-morphism}。

\medskip\noindent
设 $\mathcal{C}$ 为位点。令 $p = u$ 为函子
$u : \mathcal{C} \to \textit{Sets}$。引入这种看似奇特的说法，是因为谈论
函子的邻域似乎有些古怪；因此，我们把“点” $p$ 看作由一个范畴数据——即函子
$u$——给出的几何对象。$p$ 实际上等于 $u$ 这一事实并不重要。$p$ 的一个
{it 邻域}是序对 $(U, x)$，其中
$U \in \Ob(\mathcal{C})$ 且 $x \in u(U)$。一个{it 邻域态射}
$(V, y) \to (U, x)$，由 $\mathcal{C}$ 的态射
$\alpha :V \to U$ 给出，并满足 $u(\alpha)(y) = x$。注意，邻域范畴不是
一个“大”范畴。

\medskip\noindent
我们把预层 $\mathcal{F}$ 在 $p$ 处的{it 茎}定义为
\begin{equation}
\label{sites-equation-stalk}
\mathcal{F}_p = \colim_{\{(U, x)\}^{opp}} \mathcal{F}(U).
\end{equation}
该余极限取遍 $p$ 的邻域范畴的反范畴。换言之，$\mathcal{F}_p$ 的一个元素由
三元组 $(U, x, s)$ 给出，其中 $(U, x)$ 是 $p$ 的邻域，且
$s \in \mathcal{F}(U)$。三元组的相等关系，是由
$(U, x, s) \sim (V, y, \alpha^*s)$ 生成的等价关系，其中 $\alpha$ 如上。

\medskip\noindent
注意，若 $\varphi : \mathcal{F} \to \mathcal{G}$ 是集合预层的态射，
则得到典范茎映射
$\varphi_p : \mathcal{F}_p \to \mathcal{G}_p$。因而得到一个{it 茎函子}
$$
\textit{PSh}(\mathcal{C}) \longrightarrow \textit{Sets}, \quad
\mathcal{F} \longmapsto \mathcal{F}_p.
$$
我们用任意函子 $p = u : \mathcal{C} \to \textit{Sets}$ 定义了茎函子；
此定义无需任何条件便能成立\footnote{应当尽量避免对所有 $U$ 都有
$u(U) = \emptyset$ 的情形。}。另一方面，除非 $p$ 确实是一个点（见下述定义），
否则最好不要使用这一概念，因为一般来说茎函子不具有良好性质。

\begin{definition}
\label{sites-definition-point}
设 $\mathcal{C}$ 为位点。位点 $\mathcal{C}$ 的一个{it 点 $p$}，由函子
$u : \mathcal{C} \to \textit{Sets}$ 给出，并且满足
\begin{enumerate}
\item 对 $\mathcal{C}$ 的每个覆盖 $\{U_i \to U\}$，映射
$\coprod u(U_i) \to u(U)$ 是满射。
\item 对 $\mathcal{C}$ 的每个覆盖 $\{U_i \to U\}$ 以及每个态射
$V \to U$，映射
$u(U_i \times_U V) \to u(U_i) \times_{u(U)} u(V)$ 都是双射。
\item 茎函子 $\Sh(\mathcal{C}) \to \textit{Sets}$，
$\mathcal{F} \mapsto \mathcal{F}_p$，是左正合的。
\end{enumerate}
\end{definition}

\noindent
这些条件应当很熟悉，因为它们仿照定义
\ref{sites-definition-continuous} 与
\ref{sites-definition-morphism-sites} 中的条件。注意，(3) 蕴含
$*_p = \{*\}$，见例 \ref{sites-example-singleton-sheaf}。因此至少对某个 $U$ 有
$u(U) \not = \emptyset$（因为空余极限给出空集）。我们将在下文证明（引理
\ref{sites-lemma-point-site-topos}），这确实会诱导拓扑斯
$\Sh(\mathcal{C})$ 的一个点。在此之前，先对一般函子 $u$ 证明若干引理。

\begin{lemma}
\label{sites-lemma-points-recover}
设 $\mathcal{C}$ 为位点，$p = u : \mathcal{C} \to \textit{Sets}$ 是函子。
对 $U \in \Ob(\mathcal{C})$，存在函子性的同构
$(h_U)_p = u(U)$。
\end{lemma}

\begin{proof}
$(h_U)_p$ 的一个元素由三元组 $(V, y, f)$ 给出，其中
$V \in \Ob(\mathcal{C})$、$y\in u(V)$ 且
$f \in h_U(V) = \Mor_\mathcal{C}(V, U)$。若存在态射
$\phi : V \to V'$，满足 $u(\phi)(y) = y'$ 与
$f' \circ \phi = f$，则两个三元组 $(V, y, f)$、$(V', y', f')$
确定同一元素；一般还必须取这类等同的链，才能得到正确的等价关系。无论如何，
每个 $(V, y, f)$ 都等价于元素
$(U, u(f)(y), \text{id}_U)$。若存在上述 $\phi$，则三元组
$(V, y, f)$、$(V', y', f')$ 确定同一个元素
$(U, u(f)(y), \text{id}_U) = (U, u(f')(y'), \text{id}_U)$。
% P01-ERRATA: P01-E0122 -- the source calls these equivalent stalk representatives an "object" and later the "same triple"; both should be "element".
这证明了映射
$u(U) \to (h_U)_p$，$x \mapsto \text{下列对象的类： }(U, x, \text{id}_U)$
是双射。
\end{proof}

\noindent
设 $\mathcal{C}$ 为位点，$p = u : \mathcal{C} \to \textit{Sets}$ 是函子。
类比第 \ref{sites-section-functoriality-PSh} 节中的构造，对给定集合 $E$，按规则
\begin{equation}
\label{sites-equation-skyscraper}
U
\longmapsto
u^pE(U) = \Mor_{\textit{Sets}}(u(U), E) = \text{Map}(u(U), E).
\end{equation}
定义预层 $u^pE$。这定义了函子
$u^p : \textit{Sets} \to \textit{PSh}(\mathcal{C})$，$E \mapsto u^pE$。

\begin{lemma}
\label{sites-lemma-adjoint-point-push-stalk}
对任意函子 $u : \mathcal{C} \to \textit{Sets}$，函子 $u^p$ 是预层上的
茎函子的右伴随。
% P01-ERRATA: P01-E0123 -- the source's first sentence ends after "For any functor ...", leaving a sentence fragment.
\end{lemma}

\begin{proof}
设 $\mathcal{F}$ 是 $\mathcal{C}$ 上的预层，$E$ 是集合。态射
$\mathcal{F} \to u^pE$ 由一族相容映射
$\mathcal{F}(U) \to \text{Map}(u(U), E)$ 给出，也就是由一族相容映射
$\mathcal{F}(U) \times u(U) \to E$ 给出。而按 $\mathcal{F}_p$ 的定义，映射
$\mathcal{F}_p \to E$ 由一条规则给出：它给每个三元组 $(U, x, \sigma)$
指定 $E$ 中的一个元素，并使等价三元组映到同一元素；见围绕方程
(\ref{sites-equation-stalk}) 的讨论。这也正是一族相容映射
$\mathcal{F}(U) \times u(U) \to E$。
\end{proof}

\noindent
类比第 \ref{sites-section-continuous-functors} 节，有如下引理。

\begin{lemma}
\label{sites-lemma-point-pushforward-sheaf}
设 $\mathcal{C}$ 为位点，$p = u : \mathcal{C} \to \textit{Sets}$ 是函子。
假设对 $\mathcal{C}$ 的每个覆盖 $\{U_i \to U\}$，
\begin{enumerate}
\item 映射 $\coprod u(U_i) \to u(U)$ 是满射；且
\item 映射
$u(U_i \times_U U_j) \to u(U_i) \times_{u(U)} u(U_j)$ 是满射。
\end{enumerate}
则有
\begin{enumerate}
\item 对所有集合 $E$，预层 $u^pE$ 都是层，将它记为 $u^sE$；
\item 茎函子 $\Sh(\mathcal{C}) \to \textit{Sets}$ 与函子
$u^s: \textit{Sets} \to \Sh(\mathcal{C})$ 互为伴随；且
\item 对每个集合预层 $\mathcal{F}$，都有
$\mathcal{F}_p = \mathcal{F}^\#_p$。
\end{enumerate}
\end{lemma}

\begin{proof}
第一条断言直接来自层的定义、假设 (1) 和 (2) 以及 $u^pE$ 的定义。第二条是
$u^p$ 与茎函子之间的伴随性（引理
\ref{sites-lemma-adjoint-point-push-stalk}）限制到层之后的重述。第三条断言是因为，
对任意集合 $E$，由层化的伴随性质有
$$
\text{Map}(\mathcal{F}_p, E) =
\Mor_{\textit{PSh}(\mathcal{C})}(\mathcal{F}, u^pE) =
\Mor_{\Sh(\mathcal{C})}(\mathcal{F}^\#, u^sE) =
\text{Map}(\mathcal{F}^\#_p, E)
$$
\end{proof}

\noindent
特别地，当 $p = u$ 是点时，引理
\ref{sites-lemma-point-pushforward-sheaf} 成立。在这种情形下，我们把层 $u^sE$
看作在 $p$ 处取值 $E$ 的“摩天楼”层。

\begin{definition}
\label{sites-definition-pushforward-point}
设 $p$ 是由函子 $u$ 给出的位点 $\mathcal{C}$ 的点。对集合 $E$，定义
$p_*E = u^sE$，即上述引理 \ref{sites-lemma-point-pushforward-sheaf} 描述的层。
我们有时称它为{it 摩天楼层}。
% P01-ERRATA: P01-E0124 -- the source needs punctuation or "to be" in "define p_*E = u^sE the sheaf".
\end{definition}

\noindent
特别地，摩天楼层与茎具有如下伴随性质：
$$
\Mor_{\Sh(\mathcal{C})}(\mathcal{F}, p_*E)
=
\text{Map}(\mathcal{F}_p, E)
$$
这说明记号 $p^{-1}\mathcal{F} = \mathcal{F}_p$ 是自然的，我们有时会使用它。

\begin{lemma}
\label{sites-lemma-point-site-topos}
设 $\mathcal{C}$ 为位点。
\begin{enumerate}
\item 设 $p$ 是位点 $\mathcal{C}$ 的点。则上面引入的函子对
$(p_*, p^{-1})$ 定义拓扑斯态射
$\Sh(pt) \to \Sh(\mathcal{C})$。
\item 设 $p = (p_*, p^{-1})$ 是拓扑斯 $\Sh(\mathcal{C})$ 的点。
则函子 $u : U \mapsto p^{-1}(h_U^\#)$ 诱导位点 $\mathcal{C}$ 的点
$p'$，其关联的拓扑斯态射 $(p'_*, (p')^{-1})$ 等于 $p$。
\end{enumerate}
\end{lemma}

\begin{proof}
(1) 的证明。由上述结果，函子 $p_*$ 与 $p^{-1}$ 互为伴随。根据定义
\ref{sites-definition-point}，函子 $p^{-1}$ 必须是正合的。因此定义
\ref{sites-definition-topos} 中施加的条件全都满足，故 (1) 成立。

\medskip\noindent
(2) 的证明。设 $\{U_i \to U\}$ 是 $\mathcal{C}$ 的覆盖。则
$\coprod (h_{U_i})^\# \to h_U^\#$ 是满射，见引理
\ref{sites-lemma-covering-surjective-after-sheafification}。由于 $p^{-1}$ 是正合的
（由拓扑斯态射的定义），得出 $\coprod u(U_i) \to u(U)$ 是满射。这证明了定义
\ref{sites-definition-point} 的第 (1) 条。层化是正合的，见引理
\ref{sites-lemma-sheafification-exact}。因此，若 $U \times_V W$ 在
$\mathcal{C}$ 中存在，则
$$
h_{U \times_V W}^\# =  h_U^\# \times_{h_V^\#} h_W^\#
$$
又因 $p^{-1}$ 正合，看到
$u(U \times_V W) = u(U) \times_{u(V)} u(W)$。这证明了定义
\ref{sites-definition-point} 的第 (2) 条。令 $p' = u$，并令
$\mathcal{F}_{p'}$ 为用 $u$ 按方程
(\ref{sites-equation-stalk}) 定义的茎函子。存在典范比较映射
$c : \mathcal{F}_{p'} \to \mathcal{F}_p = p^{-1}\mathcal{F}$。具体而言，
给定表示 $\mathcal{F}_{p'}$ 的元素 $\xi$ 的三元组 $(U, x, \sigma)$，
把 $\sigma$ 看作映射 $\sigma : h_U^\# \to \mathcal{F}$；由于
$x \in u(U) = p^{-1}(h_U^\#)$，可以令
$c(\xi) = p^{-1}(\sigma)(x)$。由引理
\ref{sites-lemma-points-recover}，有 $(h_U)_{p'} = u(U)$。因为定义
\ref{sites-definition-point} 的条件 (1) 与 (2) 对 $p'$ 成立，由引理
\ref{sites-lemma-point-pushforward-sheaf}，还有
$(h_U^\#)_{p'} = (h_U)_{p'}$。因此
$$
(h_U^\#)_{p'} = (h_U)_{p'} = u(U) = p^{-1}(h_U^\#)
$$
我们断言这个双射等于比较映射
$c : (h_U^\#)_{p'} \to p^{-1}(h_U^\#)$（略去验证）。$\mathcal{C}$ 上任意层
都是形如 $h_U^\#$ 的层的余积之间映射的余等化子，见引理
\ref{sites-lemma-sheaf-coequalizer-representable}。茎函子
$\mathcal{F} \mapsto \mathcal{F}_{p'}$ 与函子 $p^{-1}$ 都与任意余极限交换
（因为二者都是左伴随）。得出对每个层 $\mathcal{F}$，$c$ 都是同构。因此茎函子
$\mathcal{F} \mapsto \mathcal{F}_{p'}$ 同构于 $p^{-1}$，特别地，它是正合的。
这证明了定义 \ref{sites-definition-point} 的条件 (3) 成立，且 $p'$ 是点。最后一条断言
已在上面证明，因为我们已经看到 $p^{-1} = (p')^{-1}$。
\end{proof}

\noindent
事实上，一个点总对应于某个位点态射，如下一个引理所示。

\begin{lemma}
\label{sites-lemma-site-point-morphism}
设 $\mathcal{C}$ 为位点，$p$ 是由
$u : \mathcal{C} \to \textit{Sets}$ 给出的 $\mathcal{C}$ 的点。设 $S_0$
为无限集，使得对所有 $U \in \Ob(\mathcal{C})$ 都有
$u(U) \subset S_0$。令 $\mathcal{S}$ 为评注
\ref{sites-remark-pt-topos} 中由幂集 $S = \mathcal{P}(S_0)$ 构造的位点。则
\begin{enumerate}
\item 存在等价 $i : \Sh(pt) \to \Sh(\mathcal{S})$；
\item 函子 $u : \mathcal{C} \to \mathcal{S}$ 诱导位点态射
$f : \mathcal{S} \to \mathcal{C}$；且
\item 复合
$$
\Sh(pt) \to
\Sh(\mathcal{S}) \to
\Sh(\mathcal{C})
$$
是引理 \ref{sites-lemma-point-site-topos} 的拓扑斯态射 $(p_*, p^{-1})$。
\end{enumerate}
\end{lemma}

\begin{proof}
第 (1) 条已在评注
\ref{sites-remark-pt-topos} 中看到。此外，回忆该等价把集合 $E$ 关联到
$\mathcal{S}$ 上的层 $i_*E$，它由规则
$V \mapsto \Mor_{\textit{Sets}}(V, E)$ 定义。第 (2) 条由
$\mathcal{C}$ 的点的定义（定义 \ref{sites-definition-point}）与位点态射的定义
（定义 \ref{sites-definition-morphism-sites}）直接得出。最后，考虑 $f_*i_*E$。
按构造有
$$
f_*i_*E(U) = i_*E(u(U)) = \Mor_{\textit{Sets}}(u(U), E)
$$
这等于 $p_*E(U)$，见方程
(\ref{sites-equation-skyscraper})。这证明了 (3)。
\end{proof}

\noindent
与拓扑情形不同，取值 $E$ 的摩天楼层的茎不总是 $E$。一般成立的是如下结论。

\begin{lemma}
\label{sites-lemma-stalk-skyscraper}
设 $\mathcal{C}$ 为位点，
$p : \Sh(pt) \to \Sh(\mathcal{C})$ 是与 $\mathcal{C}$ 关联的拓扑斯的点。
对任意集合 $E$，存在典范映射
$$
E \longrightarrow (p_*E)_p \longrightarrow E
$$
其复合为 $\text{id}_E$。
\end{lemma}

\begin{proof}
总有伴随映射 $(p_*E)_p = p^{-1}p_*E \to E$。当
$E = \{*\}$ 时，这一映射是同构，因为 $p_*$ 与 $p^{-1}$ 都是左正合的，
因而把终对象变为终对象。因此，给定 $e \in E$，可以考虑映射
$i_e : \{*\} \to E$，它给出
$$
\xymatrix{
p^{-1}p_*\{*\} \ar[rr]_{p^{-1}p_*i_e} \ar[d]_{\cong} & & p^{-1}p_*E \ar[d] \\
\{*\} \ar[rr]^{i_e} & & E
}
$$
由此得到所需映射 $E \to (p_*E)_p = p^{-1}p_*E$。
\end{proof}

\begin{lemma}
\label{sites-lemma-skyscraper-functor-exact}
设 $\mathcal{C}$ 为位点，
$p : \Sh(pt) \to \Sh(\mathcal{C})$ 是与 $\mathcal{C}$ 关联的拓扑斯的点。
函子 $p_* : \textit{Sets} \to \Sh(\mathcal{C})$ 具有如下性质：它与任意极限
交换；它是左正合的；它是忠实的；它把满射变为满射；它与余等化子交换；它反映
单射、满射和同构。
\end{lemma}

\begin{proof}
因为 $p_*$ 是右伴随，所以它与任意极限交换且左正合。$p^{-1}p_*E \to E$
是典范分裂满射这一事实，蕴含 $p_*$ 是忠实的，并反映单射、满射与同构。由引理
\ref{sites-lemma-point-site-topos}，可以假设 $p$ 来自底层位点 $\mathcal{C}$ 的一个点
$u : \mathcal{C} \to \textit{Sets}$。在此情形下，层 $p_*E$ 由
$$
p_*E(U) = \Mor_{\textit{Sets}}(u(U), E)
$$
给出，见方程 \ref{sites-equation-skyscraper} 与定义
\ref{sites-definition-pushforward-point}。从该公式立即得出，$p_*$ 把满射变为满射，
并把余等化子变为余等化子。
\end{proof}






\section{构造点}
\label{sites-section-construct-points}

\noindent
本节给出判据，判断从位点到集合范畴的函子何时定义该位点的点。

\begin{lemma}
\label{sites-lemma-neighbourhoods-cofiltered}
设 $\mathcal{C}$ 为位点，$p = u : \mathcal{C} \to \textit{Sets}$ 是函子。
若 $p$ 的邻域范畴是余滤过的，则茎函子
(\ref{sites-equation-stalk}) 是左正合的。
\end{lemma}

\begin{proof}
设 $\mathcal{I} \to \Sh(\mathcal{C})$，
$i \mapsto \mathcal{F}_i$ 是层的有限图式。需要证明该系统极限的茎与茎的极限
相符。令 $\mathcal{F}$ 为该系统作为{it 预层}的极限。根据引理
\ref{sites-lemma-limit-sheaf}，它是层，也是层范畴中的极限。因此只须证明
$\mathcal{F}_p = \lim_\mathcal{I} \mathcal{F}_{i, p}$。还要回忆
$\mathcal{F}$ 具有简单描述，见第
\ref{sites-section-limits-colimits-PSh} 节。故只须证明
$$
\lim_i \colim_{\{(U, x)\}^{opp}} \mathcal{F}_i(U)
=
\colim_{\{(U, x)\}^{opp}} \lim_i \mathcal{F}_i(U).
$$
这由 Categories, Lemma
\ref{categories-lemma-directed-commutes} 成立，因为按假设，邻域范畴的反范畴
是滤过的。
\end{proof}

\begin{lemma}
\label{sites-lemma-neighbourhoods-directed}
设 $\mathcal{C}$ 为位点。假设 $\mathcal{C}$ 有终对象 $X$ 和纤维积。设
$p = u : \mathcal{C} \to \textit{Sets}$ 是函子，满足
\begin{enumerate}
\item $u(X)$ 是单点集；且
\item 对每一对具有同一余域的态射 $U \to W$ 与 $V \to W$，映射
$u(U \times_W V) \to u(U) \times_{u(W)} u(V)$ 是双射。
\end{enumerate}
则 $p$ 的邻域范畴是余滤过的，因而茎函子
$\Sh(\mathcal{C}) \to \textit{Sets}$，
$\mathcal{F} \to \mathcal{F}_p$，与有限极限交换。
% P01-ERRATA: P01-E0125 -- the source has the duplicated phrase "Then the the category".
\end{lemma}

\begin{proof}
请注意它与引理 \ref{sites-lemma-directed} 的相似性。关于 $\mathcal{C}$ 的假设蕴含
$\mathcal{C}$ 有有限极限。见 Categories, Lemma
\ref{categories-lemma-finite-limits-exist}。假设 (1) 蕴含邻域范畴非空。
设 $(U, x)$ 与 $(V, y)$ 是邻域。则由 (2)，
$u(U \times V) = u(U \times_X V) =
u(U) \times_{u(X)} u(V) = u(U) \times u(V)$。因此存在同时映到
$(U, x)$ 与 $(V, y)$ 的邻域 $(U \times_X V, z)$。
设 $a, b : (V, y) \to (U, x)$ 是邻域范畴中的两个态射。令 $W$ 为
$a, b : V \to U$ 的等化子。与 Categories, Lemma
\ref{categories-lemma-finite-limits-exist} 的证明相同，可以用纤维积写出 $W$：
$$
W = (V \times_{a, U, b} V) \times_{(pr_1, pr_2), V \times V, \Delta} V
$$
(2) 中的双射性保证存在 $z \in u(W)$ 映到 $((y, y), y)$。于是
$(W, z) \to (V, y)$ 如所需地等化 $a, b$。“因而”部分就是引理
\ref{sites-lemma-neighbourhoods-cofiltered}。
\end{proof}

\begin{proposition}
\label{sites-proposition-point-limits}
设 $\mathcal{C}$ 为位点，并假设 $\mathcal{C}$ 中存在有限极限（即
$\mathcal{C}$ 有纤维积和终对象）。这种位点 $\mathcal{C}$ 的一个点 $p$，由函子
$u : \mathcal{C} \to \textit{Sets}$ 给出，并满足
\begin{enumerate}
\item $u$ 与有限极限交换；且
\item 若 $\{U_i \to U\}$ 是覆盖，则
$\coprod_i u(U_i) \to u(U)$ 是满射。
\end{enumerate}
\end{proposition}

\begin{proof}
先设 $p$ 是由函子 $u$ 给出的点（定义
\ref{sites-definition-point}）。条件 (2) 直接由点的定义满足。由引理
\ref{sites-lemma-points-recover}，有 $(h_U)_p = u(U)$。由引理
\ref{sites-lemma-point-pushforward-sheaf}，有
$(h_U^\#)_p = (h_U)_p$。因此看到 $u$ 等于函子的复合
$$
\mathcal{C} \xrightarrow{h}
\textit{PSh}(\mathcal{C}) \xrightarrow{{}^\#}
\Sh(\mathcal{C}) \xrightarrow{()_p}
\textit{Sets}
$$
这些函子中的每一个都是左正合的，故 $u$ 满足 (1)。

\medskip\noindent
反之，设 $u$ 满足 (1) 和 (2)。此时立即看出，$u$ 满足定义
\ref{sites-definition-point} 的前两个条件。它的茎函子是正合的，因为由引理
\ref{sites-lemma-point-pushforward-sheaf} 它是左伴随，而由引理
\ref{sites-lemma-neighbourhoods-directed} 它与有限极限交换。
\end{proof}

\begin{remark}
\label{sites-remark-improve-proposition-points-limits}
事实上，设 $\mathcal{C}$ 为位点，假设 $\mathcal{C}$ 有终对象 $X$ 与纤维积。
设 $p = u: \mathcal{C} \to \textit{Sets}$ 为函子，满足
\begin{enumerate}
\item $u(X) = \{*\}$ 是单点集；
% P01-ERRATA: P01-E0126 -- the source omits "is" in "$u(X)=\{*\}$ a singleton".
\item 对每一对具有同一余域的态射 $U \to W$ 与 $V \to W$，映射
$u(U \times_W V) \to u(U) \times_{u(W)} u(V)$ 是满射。
\item 对每个覆盖 $\{U_i \to U\}$，映射
$\coprod u(U_i) \to u(U)$ 是满射。
\end{enumerate}
一般来说，$p$ {\bf 不是} $\mathcal{C}$ 的点。一个例子是具有两个对象
$\{U, X\}$ 的范畴 $\mathcal{C}$，其中恰有一个非恒等箭头，即
$U \to X$。在 $\mathcal{C}$ 上赋予平凡拓扑，即仅有的覆盖是
$\{U \to U\}$ 与 $\{X \to X\}$。层 $\mathcal{F}$ 与预层是一回事，并由
三元组 $(A, B, A \to B)$ 组成：即 $A = \mathcal{F}(X)$、
$B = \mathcal{F}(U)$，而 $A \to B$ 是对应于 $U \to X$ 的限制映射。
注意 $U \times_X U = U$，故纤维积存在。考虑函子 $u = p$，其中
$u(X) = \{*\}$ 且 $u(U) = \{*_1, *_2\}$。它满足 (1)、(2) 和 (3)，
但相应的茎函子 (\ref{sites-equation-stalk}) 是函子
$$
(A, B, A \to B) \longmapsto B \amalg_A B
$$
它不是正合的。具体来说，考虑层的单射
$(\emptyset, \{1\}, \emptyset \to \{1\}) \to (\{1\}, \{1\}, \{1\} \to \{1\})$，
茎函子把它变为集合间的非单射
$$
\{1\} \amalg \{1\} \longrightarrow \{1\} \amalg_{\{1\}} \{1\}
$$
\end{remark}

\begin{example}
\label{sites-example-point-topological}
设 $X$ 为拓扑空间，$X_{Zar}$ 是例
\ref{sites-example-site-topological} 的位点。设 $x \in X$ 是一点。考虑函子
$$
u : X_{Zar} \longrightarrow \textit{Sets}, \quad
U \mapsto
\left\{
\begin{matrix}
\emptyset & \text{若} & x \not \in U \\
\{*\} & \text{若} & x \in U
\end{matrix}
\right.
$$
该函子与积和纤维积交换，并把覆盖变为满射族。因此得到位点
$X_{Zar}$ 的一个点 $p$。立即可验证，茎函子与 Sheaves, Section
\ref{sheaves-section-stalks} 中定义的 $x$ 处的茎相符。
\end{example}

\begin{example}
\label{sites-example-point-topology}
设 $X$ 为拓扑空间。拓扑斯 $\Sh(X)$ 的点是什么？为看清这一点，令
$X_{Zar}$ 为例 \ref{sites-example-site-topological} 的位点。由引理
\ref{sites-lemma-point-site-topos}，$\Sh(X)$ 的点对应于该位点的点。设 $p$ 是
由函子 $u : X_{Zar} \to \textit{Sets}$ 给出的位点 $X_{Zar}$ 的点。
我们将使用命题 \ref{sites-proposition-point-limits} 中对此类 $u$ 的刻画。它立即蕴含
$u(\emptyset) = \emptyset$ 与 $u(U \cap V) = u(U) \times u(V)$。特别地，
通过对角映射有 $u(U) = u(U) \times u(U)$，这蕴含 $u(U)$ 或为单点集，
或为空集。此外，若 $U = \bigcup U_i$ 是开覆盖，则
$$
u(U) = \emptyset \Rightarrow \forall i, \ u(U_i) = \emptyset
\quad\text{且}\quad
u(U) \not = \emptyset \Rightarrow \exists i, \ u(U_i) \not = \emptyset.
$$
由此得出，存在唯一的最大开集 $W \subset X$ 满足 $u(W) = \emptyset$，即所有满足
$u(U) = \emptyset$ 的开集 $U$ 的并。令 $Z = X \setminus W$。若
$Z = Z_1 \cup Z_2$，其中 $Z_i \subset Z$ 闭，则
$W = (X \setminus Z_1) \cap (X \setminus Z_2)$，故
$\emptyset = u(W) = u(X \setminus Z_1) \times u(X \setminus Z_2)$；于是
$u(X \setminus Z_1) = \emptyset$ 或
$u(X \setminus Z_2) = \emptyset$。这意味着
$X \setminus Z_1 = W$ 或 $X \setminus Z_2 = W$。换言之，$Z$ 不可约。
现在看到 $u$ 由规则
$$
u : X_{Zar} \longrightarrow \textit{Sets}, \quad
U \mapsto
\left\{
\begin{matrix}
\emptyset & \text{若} & Z \cap U = \emptyset \\
\{*\} & \text{若} & Z \cap U \not = \emptyset
\end{matrix}
\right.
$$
描述。注意，对任意不可约闭子集 $Z \subset X$，该函子都满足命题
\ref{sites-proposition-point-limits} 的假设 (1)、(2)，因而定义一个点。换言之，
位点 $X_{Zar}$ 的点与 $X$ 的不可约闭子集一一对应。特别地，若 $X$ 是清醒拓扑
空间，则 $X_{Zar}$ 的点与 $X$ 的点一一对应，见例
\ref{sites-example-point-topological}。
\end{example}

\begin{example}
\label{sites-example-point-G-sets}
考虑例 \ref{sites-example-site-on-group} 与第
\ref{sites-section-example-sheaf-G-sets} 节描述的位点 $\mathcal{T}_G$。遗忘函子
$u : \mathcal{T}_G \to \textit{Sets}$ 与积和纤维积交换，并把覆盖变为满射族。
因此它定义 $\mathcal{T}_G$ 的一个点。我们把
$\Sh(\mathcal{T}_G)$ 与 $G\textit{-Sets}$ 等同。茎函子
$$
p^{-1} :
\Sh(\mathcal{T}_G) = G\textit{-Sets}
\longrightarrow
\textit{Sets}
$$
就是遗忘函子。正向像 $p_*$ 是函子
$$
\textit{Sets}
\longrightarrow
\Sh(\mathcal{T}_G) = G\textit{-Sets}
$$
它把集合 $S$ 映到 $G$-集合 $\text{Map}(G, S)$，作用为
$g \cdot \psi = \psi \circ R_g$，其中 $R_g$ 是右乘。特别地，作为集合有
$p^{-1}p_*S = \text{Map}(G, S)$，而引理
\ref{sites-lemma-stalk-skyscraper} 的映射
$S \to \text{Map}(G, S) \to S$ 就是显然的映射。
\end{example}

\begin{example}
\label{sites-example-indiscrete-points}
设 $\mathcal{C}$ 是赋予混沌拓扑的范畴（例
\ref{sites-example-indiscrete}）。对 $\mathcal{C}$ 的每个对象 $U_0$，函子
$u : U \mapsto \Mor_\mathcal{C}(U_0, U)$ 都定义 $\mathcal{C}$ 的一个点
$p$。事实上，定义 \ref{sites-definition-point} 的条件 (1) 与 (2) 是直接的，
因为仅有的覆盖由恒等映射给出。条件 (3) 成立，因为
$\mathcal{F}_p = \mathcal{F}(U_0)$，并且由于该拓扑是混沌拓扑，在 $U_0$
上取截面是正合函子。
% P01-ERRATA: P01-E0127 -- the source repeats "Condition (2)" where exactness proves condition (3).
% P01-ERRATA: P01-E0128 -- the source calls the topology "discrete" although Example example-indiscrete names it chaotic/indiscrete.
\end{example}





\section{点与拓扑斯态射}
\label{sites-section-functorial-points}

\noindent
本节对点与拓扑斯态射作若干评注。

\begin{lemma}
\label{sites-lemma-point-functor}
设 $u : \mathcal{C} \to \mathcal{D}$ 是函子，
$v : \mathcal{D} \to \textit{Sets}$ 是函子，并令
$w = v \circ u$。分别以 $q$、$p$ 表示与 $v$、$w$ 关联的茎函子
(\ref{sites-equation-stalk})。则对 $\mathcal{C}$ 上的预层 $\mathcal{F}$，函子性地有
$(u_p\mathcal{F})_q = \mathcal{F}_p$。
\end{lemma}

\begin{proof}
这是一个简单的范畴事实。我们有
\begin{align*}
(u_p\mathcal{F})_q
& =
\colim_{(V, y)} \colim_{U, \phi : V \to u(U)} \mathcal{F}(U) \\
& = \colim_{(V, y, U, \phi : V \to u(U))} \mathcal{F}(U) \\
& = \colim_{(U, x)} \mathcal{F}(U) \\
& = \mathcal{F}_p
\end{align*}
第一个等式由 $u_p$ 的定义与茎函子的定义成立。注意 $y \in v(V)$。在第二个
等式中，只是合并了余极限。为说明第三个等式，把 Categories, Lemma
\ref{categories-lemma-colimit-constant-connected-fibers} 应用于按规则
$$
F((V, y, U, \phi : V \to u(U))) = (U, v(\phi)(y)).
$$
定义的图式范畴之间的函子 $F$。这是有意义的，因为
$w(U) = v(u(U))$。来验证 Categories, Lemma
\ref{categories-lemma-colimit-constant-connected-fibers} 的假设。注意 $F$ 有右逆
$(U, x) \mapsto (u(U), x, U, \text{id} : u(U) \to u(U))$。这同样有意义，
因为 $x \in w(U) = v(u(U))$。另一方面，总有 $(U, x)$ 上的纤维范畴中的态射
$$
(V, y, U, \phi : V \to u(U))
\longrightarrow
(u(U), v(\phi)(y), U, \text{id} : u(U) \to u(U))
$$
这表明各纤维范畴连通。第四个也是最后一个等式是清楚的。
\end{proof}

\begin{lemma}
\label{sites-lemma-point-morphism-sites}
\begin{slogan}
位点的映射定义点上的映射，而且拉回保持茎。
\end{slogan}
设 $f : \mathcal{D} \to \mathcal{C}$ 是由连续函子
$u : \mathcal{C} \to \mathcal{D}$ 给出的位点态射。设 $q$ 是由函子
$v : \mathcal{D} \to \textit{Sets}$ 给出的 $\mathcal{D}$ 的点，见定义
\ref{sites-definition-point}。则函子
$v \circ u : \mathcal{C} \to \textit{Sets}$ 定义 $\mathcal{C}$ 的一个点
$p$；此外，对 $\mathcal{C}$ 上的任意层 $\mathcal{F}$，存在典范等同
$$
(f^{-1}\mathcal{F})_q = \mathcal{F}_p
$$
\end{lemma}

\begin{proof}[引理 \ref{sites-lemma-point-morphism-sites} 的第一种证明]
注意，由于 $u$ 连续且 $v$ 定义一个点，$v \circ u$ 立即满足定义
\ref{sites-definition-point} 的条件 (1) 与 (2)。来证明所显示的等式。设
$\mathcal{F}$ 是 $\mathcal{C}$ 上的层。则
$$
(f^{-1}\mathcal{F})_q = (u_s\mathcal{F})_q =
(u_p \mathcal{F})_q = \mathcal{F}_p
$$
第一个等式因为 $f^{-1} = u_s$，第二个等式由引理
\ref{sites-lemma-point-pushforward-sheaf}，第三个等式由引理
\ref{sites-lemma-point-functor}。现在又看到，$p$ 也满足定义
\ref{sites-definition-point} 的条件 (3)，因为它是正合函子的复合。证明完毕。
\end{proof}

\begin{proof}[引理 \ref{sites-lemma-point-morphism-sites} 的第二种证明]
由引理
\ref{sites-lemma-site-point-morphism}，可把 $(q_*, q^{-1})$ 分解为
$$
\Sh(pt) \xrightarrow{i} \Sh(\mathcal{S})
\xrightarrow{h} \Sh(\mathcal{D})
$$
其中第二个拓扑斯态射来自由函子
$v : \mathcal{D} \to \mathcal{S}$ 诱导的位点态射
$h : \mathcal{S} \to \mathcal{D}$（这是有意义的，因为
$\mathcal{S} \subset \textit{Sets}$ 是包含 $v$ 的像中每个对象的满子范畴）。由引理
\ref{sites-lemma-composition-morphisms-sites}，复合
$v \circ u : \mathcal{C} \to \mathcal{S}$ 定义位点态射
$g : \mathcal{S} \to \mathcal{C}$。特别地，函子
$v \circ u : \mathcal{C} \to \mathcal{S}$ 连续；根据
$\mathcal{S}$ 中覆盖的定义（见评注
\ref{sites-remark-pt-topos}），这意味着 $v \circ u$ 满足定义
\ref{sites-definition-point} 的条件 (1) 与 (2)。另一方面，看到
$$
g_*i_*E(U) = i_*E(v(u(U)) = \Mor_{\textit{Sets}}(v(u(U)), E)
$$
% P01-ERRATA: P01-E0129 -- the frozen display is missing a closing parenthesis after i_*E(v(u(U))).
这是由评注 \ref{sites-remark-pt-topos} 中对 $i$ 的构造得出的。注意，这与如下函子的
公式相同，而它等于 $(v \circ u)^pE$，见方程
(\ref{sites-equation-skyscraper})。
% P01-ERRATA: P01-E0130 -- the source says "the same as the formula for which is equal to" and omits the object of "for".
由引理
\ref{sites-lemma-point-pushforward-sheaf}，函子
$g_*i_* = (v \circ u)^p = (v \circ u)^s$ 是茎函子
$\mathcal{F} \mapsto \mathcal{F}_q$ 的右伴随。
% P01-ERRATA: P01-E0131 -- this stalk functor is p^{-1} on Sh(C), so the subscript should be p, not q.
因此看到茎函子 $p^{-1}$ 典范同构于 $i^{-1} \circ g^{-1}$。所以它是正合的，
并得出 $p$ 是点。最后，由构造有 $g = f \circ h$，故
$p^{-1} = i^{-1} \circ h^{-1} \circ f^{-1} = q^{-1} \circ f^{-1}$；
也就是说，得到本引理所显示的公式。
\end{proof}

\begin{lemma}
\label{sites-lemma-point-morphism-topoi}
设 $f : \Sh(\mathcal{D}) \to \Sh(\mathcal{C})$ 是拓扑斯态射，
$q : \Sh(pt) \to \Sh(\mathcal{D})$ 是一个点。则
$p = f \circ q$ 是拓扑斯 $\Sh(\mathcal{C})$ 的点，而且对
$\mathcal{C}$ 上的任意层 $\mathcal{F}$，存在典范等同
$$
(f^{-1}\mathcal{F})_q = \mathcal{F}_p
$$
\end{lemma}

\begin{proof}
这由定义以及拓扑斯态射可以复合这一事实立即得出。
\end{proof}





\section{局部化与点}
\label{sites-section-localize-points}

\noindent
本节证明，局部化 $\mathcal{C}/U$ 的点可以用简单方式从 $\mathcal{C}$ 的点构造。

\begin{lemma}
\label{sites-lemma-point-localize}
设 $\mathcal{C}$ 为位点，$p$ 是由
$u : \mathcal{C} \to \textit{Sets}$ 给出的 $\mathcal{C}$ 的点。设 $U$ 是
$\mathcal{C}$ 的对象，且 $x \in u(U)$。函子
$$
v : \mathcal{C}/U \longrightarrow \textit{Sets}, \quad
(\varphi : V \to U) \longmapsto \{y \in u(V) \mid u(\varphi)(y) = x\}
$$
定义位点 $\mathcal{C}/U$ 的一个点 $q$，使图式
$$
\xymatrix{
& \Sh(pt) \ar[d]^p \ar[ld]_q \\
\Sh(\mathcal{C}/U) \ar[r]^{j_U} &
\Sh(\mathcal{C})
}
$$
交换。换言之，对 $\mathcal{C}$ 上的任意层有
$\mathcal{F}_p = (j_U^{-1}\mathcal{F})_q$。
% P01-ERRATA: P01-E0132 -- the source says "for any sheaf on C" but omits the sheaf symbol F.
\end{lemma}

\begin{proof}
如引理 \ref{sites-lemma-site-point-morphism} 中那样选取 $S$ 与 $\mathcal{S}$。
可以如该引理中一样等同 $\Sh(pt) = \Sh(\mathcal{S})$，并把由 $u$ 诱导的
拓扑斯态射写成
$p = f : \Sh(\mathcal{S}) \to \Sh(\mathcal{C})$。由引理
\ref{sites-lemma-localize-morphism}，得到拓扑斯的交换图式
$$
\xymatrix{
\Sh(\mathcal{S}/u(U)) \ar[r]_-{j_{u(U)}} \ar[d]_{p'} &
\Sh(\mathcal{S}) \ar[d]^p \\
\Sh(\mathcal{C}/U) \ar[r]^{j_U} &
\Sh(\mathcal{C}),
}
$$
其中 $p'$ 由函子 $u' : \mathcal{C}/U \to \mathcal{S}/u(U)$ 给出，
$V/U \mapsto u(V)/u(U)$。考虑函子
$j_x : \mathcal{S} \cong \mathcal{S}/x$，它把集合 $E$ 关联到集合 $E$，
并赋予取常值 $x$ 的常值映射 $E \to u(U)$。则 $j_x$ 是全忠实的余连续函子，
并有连续右伴随
$v_x : (\psi : E \to u(U)) \mapsto \psi^{-1}(\{x\})$。注意
$j_{u(U)} \circ j_x = \text{id}_\mathcal{S}$ 且
$v_x \circ u' = v$。这些观察蕴含如下拓扑斯交换图式
$$
\xymatrix{
\Sh(\mathcal{S}) \ar[rd]^a \ar[rdd]_q
\ar `r[rrr] `d[dd]^p [rrdd] & & & \\
& \Sh(\mathcal{S}/u(U)) \ar[r]_-{j_{u(U)}} \ar[d]^{p'} &
\Sh(\mathcal{S}) \ar[d]^p & \\
& \Sh(\mathcal{C}/U) \ar[r]^{j_U} &
\Sh(\mathcal{C}) &
}
$$
具体如下：
\begin{enumerate}
\item 态射
$a : \Sh(\mathcal{S}) \to \Sh(\mathcal{S}/u(U))$ 是与余连续函子
$j_x$ 关联的拓扑斯态射；它等于与连续函子 $v_x$ 关联的态射，见引理
\ref{sites-lemma-cocontinuous-morphism-topoi} 与第
\ref{sites-section-cocontinuous-adjoint} 节。
\item 复合 $p \circ j_{u(U)} \circ a = p$，因为
$j_{u(U)} \circ j_x = \text{id}_\mathcal{S}$。
\item 复合 $p' \circ a$ 给出拓扑斯态射。此外，它是与连续函子
$v_x \circ u' = v$ 关联的拓扑斯态射。因此 $v$ 确实定义
$\mathcal{C}/U$ 的一个点 $q$，并且由构造它嵌入上述图式。
\end{enumerate}
本引理的证明结束。
\end{proof}

\begin{lemma}
\label{sites-lemma-points-above-point}
设 $\mathcal{C}$、$p$、$u$、$U$ 如引理
\ref{sites-lemma-point-localize} 中所述。引理
\ref{sites-lemma-point-localize} 的构造，在位于 $p$ 上方的
$\mathcal{C}/U$ 的点 $q$ 与 $u(U)$ 的元素 $x$ 之间给出一一对应。
\end{lemma}

\begin{proof}
设 $q$ 是由函子 $v : \mathcal{C}/U \to \textit{Sets}$ 给出的
$\mathcal{C}/U$ 的点，并且作为拓扑斯态射有 $j_U \circ q = p$。回忆对
$\mathcal{C}$ 的任意对象 $V$，有
$u(V) = p^{-1}(h_V^\#)$，见引理
\ref{sites-lemma-point-site-topos}。类似地，对 $\mathcal{C}/U$ 的任意对象
$V/U$，有 $v(V/U) = q^{-1}(h_{V/U}^\#)$。考虑如下两个图式
$$
\vcenter{
\xymatrix{
\Mor_{\mathcal{C}/U}(W/U, V/U) \ar[r] \ar[d] &
\Mor_\mathcal{C}(W, V) \ar[d] \\
\Mor_{\mathcal{C}/U}(W/U, U/U) \ar[r] &
\Mor_\mathcal{C}(W, U)
}
}
\quad
\vcenter{
\xymatrix{
h_{V/U}^\# \ar[r] \ar[d] & j_U^{-1}(h_V^\#) \ar[d] \\
h_{U/U}^\# \ar[r] & j_U^{-1}(h_U^\#)
}
}
$$
右侧图式是 $\mathcal{C}/U$ 上预层图式的层化，该预层图式把 $W/U$ 映到
左侧的集合图式。（这里有一个小技术点，即
$(j_U^{-1}h_V)^\# = j_U^{-1}(h_V^\#)$，对 $h_U$ 亦然；见引理
\ref{sites-lemma-technical-pu}。）注意左侧的集合图式是笛卡尔的。由于层化正合
（引理 \ref{sites-lemma-sheafification-exact}），得出右侧图式也是笛卡尔的。

\medskip\noindent
把正合函子 $q^{-1}$ 应用于右侧图式，得到集合的笛卡尔图式
$$
\xymatrix{
v(V/U) \ar[r] \ar[d] & u(V) \ar[d] \\
v(U/U) \ar[r] & u(U)
}
$$
这里使用了 $q^{-1} \circ j^{-1} = p^{-1}$。
% P01-ERRATA: P01-E0133 -- the source uses undefined j^{-1}; the localization functor here is j_U^{-1}.
由于 $U/U$ 是 $\mathcal{C}/U$ 的终对象，$v(U/U)$ 是单点集。因此
$v(U/U)$ 在 $u(U)$ 中的像是一个元素 $x$，而顶部水平映射给出双射
$v(V/U) \to  \{y \in u(V) \mid y \mapsto x\text{ 于 }u(U)\}$，如所需。
\end{proof}

\begin{lemma}
\label{sites-lemma-stalk-j-shriek}
设 $\mathcal{C}$ 为位点，$p$ 是由
$u : \mathcal{C} \to \textit{Sets}$ 给出的 $\mathcal{C}$ 的点，$U$ 是
$\mathcal{C}$ 的对象。对 $\mathcal{C}/U$ 上任意层 $\mathcal{G}$，有
$$
(j_{U!}\mathcal{G})_p =
\coprod\nolimits_q \mathcal{G}_q
$$
其中余积取遍按引理
\ref{sites-lemma-point-localize} 与元素 $x \in u(U)$ 关联的 $\mathcal{C}/U$ 的点 $q$。
\end{lemma}

\begin{proof}
使用引理 \ref{sites-lemma-describe-j-shriek} 中的描述：
$j_{U!}\mathcal{G}$ 是预层
$V \mapsto
\coprod\nolimits_{\varphi \in \Mor_\mathcal{C}(V, U)}
\mathcal{G}(V/_\varphi U)$
所关联的层。此外，由引理
\ref{sites-lemma-point-pushforward-sheaf}，$j_{U!}\mathcal{G}$ 在 $p$ 处的茎等于
该预层的茎。因此
$$
(j_{U!}\mathcal{G})_p =
\colim_{(V, y)} \coprod\nolimits_{\varphi : V \to U} \mathcal{G}(V/_\varphi U)
$$
对该余极限的每个元素 $(V, y, \varphi, s)$，可以指定
$x = u(\varphi)(y) \in u(U)$。因此得到
$$
(j_{U!}\mathcal{G})_p =
\coprod\nolimits_{x \in u(U)}
\colim_{(\varphi : V \to U, y), \ u(\varphi)(y) = x} \mathcal{G}(V/_\varphi U).
$$
根据我们对点 $q$ 的构造，这等于本引理中的表达式。
\end{proof}

\begin{remark}
\label{sites-remark-not-pushforward}
警告：引理 \ref{sites-lemma-stalk-j-shriek} 的结果对 $j_{U, *}$ 没有类似结论。
\end{remark}





\section{拓扑斯的 2-态射}
\label{sites-section-2-category}

\noindent
本节简要讨论拓扑斯的 $2$-态射这一概念。

\begin{definition}
\label{sites-definition-2-morphism-topoi}
设 $f, g : \Sh(\mathcal{C}) \to \Sh(\mathcal{D})$
是两个拓扑斯态射。一个{it 从 $f$ 到 $g$ 的 2-态射}由函子间的变换
$t : f_* \to g_*$ 给出。
\end{definition}

\noindent
在图示中，我们有时把 $t$ 表示如下：
$$
\xymatrix{
\Sh(\mathcal{C})
\rrtwocell^f_g{t}
&
&
\Sh(\mathcal{D})
}
$$
注意，由于 $f^{-1}$ 伴随于 $f_*$，$g^{-1}$ 伴随于 $g_*$，$t$ 还诱导函子间的
变换 $t : g^{-1} \to f^{-1}$（通常用同一符号表示）；它由下列图式交换这一条件
唯一刻画：
$$
\xymatrix{
\Mor_{\Sh(\mathcal{D})}(\mathcal{G}, f_*\mathcal{F})
\ar[d]_{t \circ -} \ar@{=}[r] &
\Mor_{\Sh(\mathcal{C})}(f^{-1}\mathcal{G}, \mathcal{F})
\ar[d]^{- \circ t}
\\
\Mor_{\Sh(\mathcal{D})}(\mathcal{G}, g_*\mathcal{F})
\ar@{=}[r] &
\Mor_{\Sh(\mathcal{C})}(g^{-1}\mathcal{G}, \mathcal{F})
}
$$
由于集合论上的困难（见评注
\ref{sites-remark-morphism-topoi-big}），我们得不到拓扑斯的 2-范畴。不过仍可定义
水平复合与垂直复合，并证明 Categories, Section
\ref{categories-section-2-categories} 所列的严格 2-范畴公理成立。具体来说，
2-态射的垂直复合是清楚的（只须复合函子间的变换）；1-态射的复合已在定义
\ref{sites-definition-topos} 中定义；而
$$
\xymatrix{
\Sh(\mathcal{C})
\rtwocell^f_g{t}
&
\Sh(\mathcal{D})
\rtwocell^{f'}_{g'}{s}
&
\Sh(\mathcal{E})
}
$$
的水平复合，由 Categories, Definition
\ref{categories-definition-horizontal-composition} 中引入的函子间变换
$s \star t$ 定义。明确地说，$s \star t$ 由
$$
\xymatrix{
f'_*f_*\mathcal{F}
\ar[r]^{f'_*t} &
f'_*g_*\mathcal{F}
\ar[r]^s &
g'_*g_*\mathcal{F}
}
\quad\text{或}\quad
\xymatrix{
f'_*f_*\mathcal{F}
\ar[r]^s &
g'_*f_*\mathcal{F}
\ar[r]^{g'_*t} &
g'_*g_*\mathcal{F}
}
$$
给出（这些映射相等）。由于这些定义与 Categories, Section
\ref{categories-section-formal-cat-cat} 中的定义相符，由 Categories, Lemma
\ref{categories-lemma-properties-2-cat-cats} 得出：在这些定义下，严格 2-范畴的
公理成立。




\section{点之间的态射}
\label{sites-section-morphisms-points}

\begin{lemma}
\label{sites-lemma-maps-u-points}
设 $\mathcal{C}$ 为位点，
$u, u' : \mathcal{C} \to \textit{Sets}$ 是两个函子，而
$t : u' \to u$ 是函子间的变换。则得到典范茎函子变换
$t_{stalk} : \mathcal{F}_{p'} \to \mathcal{F}_p$；通过引理
\ref{sites-lemma-points-recover} 中的等同，它与 $t$ 相符。
\end{lemma}

\begin{proof}
略。
\end{proof}

\begin{definition}
\label{sites-definition-morphism-points}
设 $\mathcal{C}$ 为位点。设 $p, p'$ 是由函子
$u, u' : \mathcal{C} \to \textit{Sets}$ 给出的 $\mathcal{C}$ 的点。一个
{it 态射 $f : p \to p'$}，由函子间的变换
$$
f_u : u' \to u.
$$
给出。
\end{definition}

\noindent
注意函子间的变换沿相反方向。正如稍后会看到的，这可以通过把态射 $f$ 看作
一种 $2$-箭头来理解；图示如下：
$$
\xymatrix{
\textit{Sets}
=
\Sh(pt)
\rrtwocell^p_{p'}{f}
&
&
\Sh(\mathcal{C})
}
$$
具体来说，稍后会看到，$f_u$ 在摩天楼层构造之间诱导典范函子变换
$p_* \to p'_*$。

\medskip\noindent
这是一个相当重要的概念，值得在此作更完整的处理。期望事项如下：
\begin{enumerate}
\item 描述例 \ref{sites-example-point-G-sets} 中所述 $\mathcal{T}_G$ 的点的自同构。
\item 用 $\Mor(p_*, p'_*)$ 描述 $\Mor(p, p')$。
\item 拓扑空间中点的特殊化。证明：若在拓扑空间 $X$ 中
$x' \in \overline{\{x\}}$，则存在态射 $p \to p'$，其中 $p$（相应地，$p'$）
是 $X_{Zar}$ 中与 $x$（相应地，$x'$）关联的点。
\end{enumerate}




\section{有足够多点的位点}
\label{sites-section-sites-enough-points}

\begin{definition}
\label{sites-definition-enough-points}
设 $\mathcal{C}$ 为位点。
\begin{enumerate}
\item 若每个层映射 $\phi : \mathcal{F} \to \mathcal{G}$，只要在所有茎
$\mathcal{F}_{p_i} \to \mathcal{G}_{p_i}$ 上都是同构，就必为同构，则称点族
$\{p_i\}_{i\in I}$ 是{it 保守的}。
% P01-ERRATA: P01-E0134 -- the source calls F_{p_i}->G_{p_i} "fibres" although they are stalks.
\item 若存在一个保守点族，则称 $\mathcal{C}$ {it 有足够多的点}。
\end{enumerate}
\end{definition}

\noindent
事实证明，这时可以在茎上检验“正合性”。

\begin{lemma}
\label{sites-lemma-exactness-stalks}
设 $\mathcal{C}$ 为位点，$\{p_i\}_{i\in I}$ 是保守点族。则
\begin{enumerate}
\item 给定任意层映射 $\varphi : \mathcal{F} \to \mathcal{G}$，
$\forall i, \varphi_{p_i}$ 为单射蕴含 $\varphi$ 为单射。
\item 给定任意层映射 $\varphi : \mathcal{F} \to \mathcal{G}$，
$\forall i, \varphi_{p_i}$ 为满射蕴含 $\varphi$ 为满射。
\item 给定任意一对层映射
$\varphi_1, \varphi_2 : \mathcal{F} \to \mathcal{G}$，
$\forall i, \varphi_{1, p_i} = \varphi_{2, p_i}$ 蕴含
$\varphi_1 = \varphi_2$。
\item 给定有限图式 $\mathcal{G} : \mathcal{J}
\to \Sh(\mathcal{C})$、层 $\mathcal{F}$ 以及态射
$q_j : \mathcal{F} \to \mathcal{G}_j$，则 $(\mathcal{F}, q_j)$ 是该图式的
极限，当且仅当对每个 $i$，茎
$(\mathcal{F}_{p_i}, (q_j)_{p_i})$ 是相应极限。
\item 给定有限图式 $\mathcal{F} : \mathcal{J}
\to \Sh(\mathcal{C})$、层 $\mathcal{G}$ 以及态射
$e_j : \mathcal{F}_j \to \mathcal{G}$，则 $(\mathcal{G}, e_j)$ 是该图式的
余极限，当且仅当对每个 $i$，茎
$(\mathcal{G}_{p_i}, (e_j)_{p_i})$ 是相应余极限。
\end{enumerate}
\end{lemma}

\begin{proof}
将反复使用如下事实：所有茎函子都与任意有限极限和余极限交换，因而与积、纤维积
等交换。还将使用单射就是单态射、满射就是满态射这一事实。层映射
$\varphi : \mathcal{F} \to \mathcal{G}$ 为单射，当且仅当
$\mathcal{F} \to \mathcal{F} \times_\mathcal{G} \mathcal{F}$
是同构。因此 (1) 成立。类似地，
$\varphi : \mathcal{F} \to \mathcal{G}$ 为满射，当且仅当
$\mathcal{G} \amalg_\mathcal{F} \mathcal{G} \to \mathcal{G}$
是同构。因此 (2) 成立。映射 $a, b : \mathcal{F} \to \mathcal{G}$ 相等，
当且仅当
$\text{Equalizer}(a, b : \mathcal{F} \to \mathcal{G}) \to \mathcal{F}$
是同构\footnote{等化子可以用积与纤维积表示；见 Categories, Lemma
\ref{categories-lemma-finite-limits-exist} 的证明。}。因此 (3) 成立。
断言 (4) 与 (5) 直接由定义及本证明开头的评注得出。
\end{proof}

\begin{lemma}
\label{sites-lemma-enough}
设 $\mathcal{C}$ 为位点，$\{(p_i, u_i)\}_{i\in I}$ 是点族。该族保守，当且仅当
对每个层 $\mathcal{F}$、每个 $U\in \Ob(\mathcal{C})$ 以及每一对互异截面
$s, s' \in \mathcal{F}(U)$，$s \not = s'$，都存在 $i$ 与
$x\in u_i(U)$，使三元组 $(U, x, s)$ 与 $(U, x, s')$ 定义
$\mathcal{F}_{p_i}$ 的不同元素。
\end{lemma}

\begin{proof}
设该族保守，且 $\mathcal{F}$、$U$ 与 $s, s'$ 如引理中所述。截面 $s$、$s'$
定义两个不同映射 $a, a' : (h_U)^\# \to \mathcal{F}$。因此由引理
\ref{sites-lemma-exactness-stalks}，存在 $i$ 使
$a_{p_i} \not = a'_{p_i}$。回忆由引理
\ref{sites-lemma-points-recover} 与
\ref{sites-lemma-point-pushforward-sheaf}，有
$(h_U)^\#_{p_i} = u_i(U)$。因此存在 $x \in u_i(U)$，使
$\mathcal{F}_{p_i}$ 中 $a_{p_i}(x) \not = a'_{p_i}(x)$。展开定义便看到
$(U, x, s)$ 与 $(U, x, s')$ 如本引理的陈述所述。

\medskip\noindent
为证明逆命题，假设本引理关于点的存在性的条件。设
$\phi : \mathcal{F} \to \mathcal{G}$ 是在所有茎上都为同构的层映射。需要证明
$\phi$ 既单又满，见引理
\ref{sites-lemma-mono-epi-sheaves}。单射性是该假设的直接推论。为证明满射性，需要证明
$\mathcal{G} \amalg_\mathcal{F} \mathcal{G} \to \mathcal{G}$ 是同构
（Categories, Lemma
\ref{categories-lemma-characterize-mono-epi}）。由于该映射显然满，只须检验其单射性；
这又由假设中
$\mathcal{G} \amalg_\mathcal{F} \mathcal{G} \to \mathcal{G}$
在所有茎上都是单射得出。
\end{proof}

\noindent
在下一个引理中，根据引理
\ref{sites-lemma-points-above-point}，点 $q_{i, x}$ 恰好是位于点 $p_i$ 上方的
$\mathcal{C}/U$ 的所有点。

\begin{lemma}
\label{sites-lemma-localize-enough}
设 $\mathcal{C}$ 为位点，$U$ 是 $\mathcal{C}$ 的对象。设
$\{(p_i, u_i)\}_{i\in I}$ 是 $\mathcal{C}$ 的点族。对
$x \in u_i(U)$，令 $q_{i, x}$ 为引理
\ref{sites-lemma-point-localize} 中构造的 $\mathcal{C}/U$ 的点。若
$\{p_i\}$ 是保守点族，则
$\{q_{i, x}\}_{i \in I, x \in u_i(U)}$ 是 $\mathcal{C}/U$ 的保守点族。
特别地，若 $\mathcal{C}$ 有足够多的点，则每个局部化 $\mathcal{C}/U$ 也如此。
% P01-ERRATA: P01-E0135 -- the source begins the third sentence with lowercase "let" after a period.
\end{lemma}

\begin{proof}
我们知道 $j_{U!}$ 诱导等价
$j_{U!} : \Sh(\mathcal{C}/U) \to \Sh(\mathcal{C})/h_U^\#$，见引理
\ref{sites-lemma-essential-image-j-shriek}。此外，还知道
$(j_{U!}\mathcal{G})_{p_i} = \coprod_x \mathcal{G}_{q_{i, x}}$，见引理
\ref{sites-lemma-stalk-j-shriek}。因此结论形式地成立。
\end{proof}

\noindent
下一个引理告诉我们，可以在位点上局部地检验点的存在性。

\begin{lemma}
\label{sites-lemma-enough-points-local}
设 $\mathcal{C}$ 为位点，$\{U_i\}_{i \in I}$ 是 $\mathcal{C}$ 的对象族。假设
\begin{enumerate}
\item $\coprod h_{U_i}^\# \to *$ 是层的满射；且
\item 每个局部化 $\mathcal{C}/U_i$ 都有足够多的点。
\end{enumerate}
则 $\mathcal{C}$ 有足够多的点。
\end{lemma}

\begin{proof}
对每个 $i \in I$，设 $\{p_j\}_{j \in J_i}$ 是 $\mathcal{C}/U_i$ 的保守点族。
对 $j \in J_i$，以
$q_j : \Sh(pt) \to \Sh(\mathcal{C})$ 表示 $p_j$ 与局部化态射
$\Sh(\mathcal{C}/U_i) \to \Sh(\mathcal{C})$ 的复合。
% P01-ERRATA: P01-E0136 -- the source's "denote q_j ... the composition" omits "by" or requires recasting.
则 $q_j$ 是点，见引理
\ref{sites-lemma-point-morphism-topoi}。我们断言点族
$\{q_j\}_{j \in \coprod J_i}$ 保守。事实上，设
$\mathcal{F} \to \mathcal{G}$ 是 $\mathcal{C}$ 上的层映射，使得对所有
$j \in \coprod J_i$，映射
$\mathcal{F}_{q_j} \to \mathcal{G}_{q_j}$ 都是同构。设 $W$ 是
$\mathcal{C}$ 的对象。由假设 (1)，存在覆盖 $\{W_a \to W\}$ 与态射
$W_a \to U_{i(a)}$。由于由引理
\ref{sites-lemma-point-morphism-topoi} 有
$(\mathcal{F}|_{\mathcal{C}/U_{i(a)}})_{p_j} = \mathcal{F}_{q_j}$ 与
$(\mathcal{G}|_{\mathcal{C}/U_{i(a)}})_{p_j} = \mathcal{G}_{q_j}$，
而点族 $\{p_j\}_{j \in J_{i(a)}}$ 保守，所以
$\mathcal{F}|_{U_{i(a)}} \to \mathcal{G}|_{U_{i(a)}}$ 是同构。因此对每个 $a$，
$\mathcal{F}(W_a) \to \mathcal{G}(W_a)$ 是双射。类似地，对每个 $a, b$，
$\mathcal{F}(W_a \times_W W_b) \to \mathcal{G}(W_a \times_W W_b)$
是双射。由层条件，这表明
$\mathcal{F}(W) \to \mathcal{G}(W)$ 是双射，即
$\mathcal{F} \to \mathcal{G}$ 是同构。
\end{proof}

\begin{lemma}
\label{sites-lemma-check-morphism-sites}
设 $u : \mathcal{C} \to \mathcal{D}$ 是位点的连续函子。设
$\{(q_i, v_i)\}_{i\in I}$ 是 $\mathcal{D}$ 的保守点族。若每个函子
$u_i = v_i \circ u$ 都定义 $\mathcal{C}$ 的一个点，则 $u$ 定义位点态射
$f : \mathcal{D} \to \mathcal{C}$。
\end{lemma}

\begin{proof}
以 $p_i$ 表示 $\textit{PSh}(\mathcal{C})$ 上与函子 $u_i$ 对应的茎函子
(\ref{sites-equation-stalk})。有
$$
(f^{-1}\mathcal{F})_{q_i} =
(u_s\mathcal{F})_{q_i} =
(u_p\mathcal{F})_{q_i} =
\mathcal{F}_{p_i}
$$
% P01-ERRATA: P01-E0137 -- f is used before the proof establishes that u defines a morphism of sites; u_s should be used directly until exactness is proved.
第一个等式因为 $f^{-1} = u_s$，第二个等式由引理
\ref{sites-lemma-point-pushforward-sheaf}，第三个等式由引理
\ref{sites-lemma-point-functor}。因此，若 $p_i$ 是点，则先沿 $f$ 拉回、再在
$q_i$ 处取茎是正合函子。由于点族 $\{q_i\}$ 保守，这蕴含 $f^{-1}$ 是正合函子；
于是由定义 \ref{sites-definition-morphism-sites}，看到 $f$ 是位点态射。
\end{proof}





\section{点的存在性判据}
\label{sites-section-criterion-points}

\noindent
本节对应于 Deligne 为 \cite[Expos\'e VI]{SGA4} 所写的附录。事实上，
两者几乎逐字相同。

\medskip\noindent
设 $\mathcal{C}$ 为位点。假设 $(I, \geq)$ 是有向集，并且
$(U_i, f_{ii'})$ 是 $I$ 上的逆系统，见 Categories, Definition
\ref{categories-definition-system-over-poset}。给定数据
$(I, \geq, U_i, f_{ii'})$，定义
$$
u : \mathcal{C} \longrightarrow \textit{Sets}, \quad
u(V) = \colim_i \Mor_\mathcal{C}(U_i , V)
$$
令 $\mathcal{F} \mapsto \mathcal{F}_p$ 为第
\ref{sites-section-points} 节中与 $u$ 关联的茎函子。由定义直接可见，在这种特殊情形下，
实际上有
$$
\mathcal{F}_p = \colim_i \mathcal{F}(U_i)
$$
注意，$u$ 与所有有限极限（指 $\mathcal{C}$ 中可表示的那些）交换，因为每个函子
$V \mapsto \Mor_\mathcal{C}(U_i , V)$ 都如此；见 Categories, Lemma
\ref{categories-lemma-directed-commutes}。
% P01-ERRATA: P01-E0138 -- the source uses plural "do" after singular "each of the functors"; "does" is required.

\medskip\noindent
若 $J \subset I$，$J$ 上的序由 $I$ 上的序诱导，并且
$V_j = U_j$、$g_{jj'} = f_{jj'}$（换言之，$J$ 上的逆系统由 $I$ 上的逆系统
诱导），则称系统 $(I, \geq, U_i, f_{ii'})$ 是
$(J, \geq, V_j, g_{jj'})$ 的一个{it 加细}。
% P01-ERRATA: P01-E0139 -- the source omits "is" in "the ordering on J induced from that of I".
令 $u$ 为与 $(I, \geq, U_i, f_{ii'})$ 关联的函子，$u'$ 为与
$(J, \geq, V_j, g_{jj'})$ 关联的函子。这诱导函子间的变换
$$
u' \longrightarrow u
$$
原因仅仅是 $u'$ 的余极限取遍 $u$ 的余极限中各系统的一个子系统。特别地，
得到关联的茎函子变换 $\mathcal{F}_{p'} \to \mathcal{F}_p$，见引理
\ref{sites-lemma-maps-u-points}。

\begin{lemma}
\label{sites-lemma-refine}
设 $\mathcal{C}$ 为位点。设 $(J, \geq, V_j, g_{jj'})$ 是上述系统，关联的函子对
为 $(u', p')$。设 $\mathcal{F}$ 是 $\mathcal{C}$ 上的层，
$s, s' \in \mathcal{F}_{p'}$ 是不同元素。设
$\{W_k \to W\}$ 是 $\mathcal{C}$ 的有限覆盖，且 $f \in u'(W)$。
则存在 $(J, \geq, V_j, g_{jj'})$ 的一个加细
$(I, \geq, U_i, f_{ii'})$，使 $s, s'$ 映到
$\mathcal{F}_p$ 的不同元素，而且 $f$ 在 $u(W)$ 中的像落在某个
$u(W_k)$ 的像中。
\end{lemma}

\begin{proof}
存在 $j_0 \in J$，使 $f$ 由 $f' : V_{j_0} \to W$ 定义。对
$j \geq j_0$，令
$V_{j, k} = V_j \times_{f'\circ f_{j j_0}, W} W_k$。
% P01-ERRATA: P01-E0140 -- the transition map in the J-system is g_{j j_0}; f_{j j_0} is undefined here.
则 $\{V_{j, k} \to V_j\}$ 是位点 $\mathcal{C}$ 中的有限覆盖。因此
$\mathcal{F}(V_j) \subset \prod_k \mathcal{F}(V_{j, k})$。再次由
Categories, Lemma \ref{categories-lemma-directed-commutes}，看到
$$
\mathcal{F}_{p'} =
\colim_j \mathcal{F}(V_j)
\longrightarrow
\prod\nolimits_k \colim_j \mathcal{F}(V_{j, k})
$$
是单射。因此存在 $k$，使 $s$ 与 $s'$ 在
$\colim_j \mathcal{F}(V_{j, k})$ 中的像不同。令
$J_0 = \{j \in J, j \geq j_0\}$ 且 $I = J \amalg J_0$。在 $I$ 上规定序，
使第二个并项的任何元素都不小于第一个并项的任何元素，而其余次序使用 $J$ 上的序。
若 $j \in I$ 位于第一个并项，就使用 $V_j$；若 $j \in I$ 位于第二个并项，
就使用 $V_{j, k}$。略去该逆系统的转移映射的定义。由上述结果，$s, s'$ 在
$\mathcal{F}_p$ 中的像不同。此外，由构造，$f'$ 在 $V_{j, k}$ 上的限制分解经过
$W_k$。
\end{proof}

\begin{lemma}
\label{sites-lemma-refine-all-at-once}
设 $\mathcal{C}$ 为位点。设 $(J, \geq, V_j, g_{jj'})$ 是上述系统，关联的函子对
为 $(u', p')$。设 $\mathcal{F}$ 是 $\mathcal{C}$ 上的层，
$s, s' \in \mathcal{F}_{p'}$ 是不同元素。则存在
$(J, \geq, V_j, g_{jj'})$ 的一个加细
$(I, \geq, U_i, f_{ii'})$，使 $s, s'$ 映到
$\mathcal{F}_p$ 的不同元素，并且对位点 $\mathcal{C}$ 的每个有限覆盖
$\{W_k \to W\}$ 以及任意 $f \in u'(W)$，$f$ 在 $u(W)$ 中的像都落在某个
$u(W_k)$ 的像中。
\end{lemma}

\begin{proof}
令 $E$ 为序对 $(\{W_k \to W\}, f\in u'(W))$ 所成的集合。考虑序对
$(E' \subset E, (I, \geq, U_i, f_{ii'}))$，满足
\begin{enumerate}
\item $(I, \geq, U_i, g_{ii'})$ 是
$(J, \geq, V_j, g_{jj'})$ 的加细；
% P01-ERRATA: P01-E0141 -- the tuple is defined with transition maps f_{ii'}, but item (1) switches to undefined g_{ii'}.
\item $s, s'$ 映到 $\mathcal{F}_p$ 的不同元素；且
\item 对每个序对 $(\{W_k \to W\}, f\in u'(W)) \in E'$，$f$ 在
$u(W)$ 中的像都落在某个 $u(W_k)$ 的像中。
\end{enumerate}
在第一因子上按包含关系、第二因子上按加细关系对这些序对排序。以
$\mathcal{S}$ 表示所有上述序对
$(E' \subset E, (I, \geq, U_i, f_{ii'}))$ 所成的类。
% P01-ERRATA: P01-E0142 -- the source's "Denote S the class" omits "by" or requires recasting.
我们断言 $\mathcal{S}$ 满足 Zorn 引理的假设。事实上，设
$(E'_a, (I_a, \geq, U_i, f_{ii'}))_{a \in A}$ 是 $\mathcal{S}$ 的全序子集。
则可以定义 $E' = \bigcup_{a \in A} E'_a$，并令
$I = \bigcup_{a \in A} I_a$。我们断言相应序对
$(E' , (I, \geq, U_i, f_{ii'}))$ 是 $\mathcal{S}$ 的元素。条件 (1) 与 (3)
是清楚的。对条件 (2)，注意
$$
u = \colim_{a \in A} u_a
\text{ 并相应地 }
\mathcal{F}_p = \colim_{a \in A} \mathcal{F}_{p_a}
$$
$s, s'$ 的像在该茎中互异，由集合的有向余极限的描述得出；见
Categories, Section \ref{categories-section-directed-colimits}。在这种情形下，
我们将简写为
$(E', (I, \ldots)) = \bigcup_{a \in A}(E'_a, (I_a, \ldots))$。

\medskip\noindent
好；若 $\mathcal{S}$ 是集合，则 Zorn 引理适用，并与上面的引理
\ref{sites-lemma-refine} 结合而容易证明本引理。但它并不适用，因为
$\mathcal{S}$ 是类。为绕开这一点，在 $E$ 上选取一个良序。对 $e \in E$，令
$E'_e = \{e' \in E \mid e' \leq e\}$。用超限递归构造序对
$(E'_e, (I_e, \ldots)) \in \mathcal{S}$，使得
$e_1 \leq e_2 \Rightarrow (E'_{e_1}, (I_{e_1}, \ldots))
\leq (E'_{e_2}, (I_{e_2}, \ldots))$。设 $e \in E$，比如
$e = (\{W_k \to W\}, f\in u'(W))$。若 $e$ 有前驱 $e - 1$，就按引理
\ref{sites-lemma-refine} 对系统
$e = (\{W_k \to W\}, f\in u'(W))$ 取
$(I_{e - 1}, \ldots)$ 的一个加细 $(I_e, \ldots)$。若 $e$ 没有前驱，就对系统
$e = (\{W_k \to W\}, f\in u'(W))$ 取
$\bigcup_{e' < e} (I_{e'}, \ldots)$ 的一个加细 $(I_e, \ldots)$。最后，
并 $\bigcup_{e \in E} I_e$ 就是本引理所提出问题的一个解。
\end{proof}

\begin{proposition}
\label{sites-proposition-criterion-points}
\begin{reference}
\cite[Expos\'e VI, Appendix by Deligne, Proposition 9.0]{SGA4}
\end{reference}
设 $\mathcal{C}$ 为位点。假设
\begin{enumerate}
\item $\mathcal{C}$ 中存在有限极限；且
\item 每个覆盖 $\{U_i \to U\}_{i \in I}$ 都有一个由
$\mathcal{C}$ 的有限覆盖给出的加细。
\end{enumerate}
则 $\mathcal{C}$ 有足够多的点。
\end{proposition}

\begin{proof}
需要证明：给定 $\mathcal{C}$ 上任意层 $\mathcal{F}$、任意
$U \in \Ob(\mathcal{C})$ 以及任意两个不同截面
$s, s' \in \mathcal{F}(U)$，存在一点 $p$，使 $s, s'$ 在
$\mathcal{F}_p$ 中的像不同。见引理
\ref{sites-lemma-enough}。考虑系统 $(J, \geq, V_j, g_{jj'})$，其中
$J = \{1\}$、$V_1 = U$、$g_{11} = \text{id}_U$。应用引理
\ref{sites-lemma-refine-all-at-once}。由该引理的结果，得到加细原系统的系统
$(I, \geq, U_i, f_{ii'})$，使
$s_p \not = s'_p$，而且对位点 $\mathcal{C}$ 的每个有限覆盖
$\{W_k \to W\}$，映射 $\coprod_k u(W_k) \to u(W)$ 都是满射。由于
$\mathcal{C}$ 的每个覆盖都可以由有限覆盖加细，得出对位点
$\mathcal{C}$ 的{\it 任意}覆盖 $\{W_k \to W\}$，
$\coprod_k u(W_k) \to u(W)$ 都是满射。这蕴含 $u = p$ 是点，见命题
\ref{sites-proposition-point-limits}（以及本节开头保证 $u$ 与有限极限交换的讨论）。
\end{proof}

\begin{lemma}
\label{sites-lemma-criterion-points}
设 $\mathcal{C}$ 为位点。设 $I$ 是集合，并且对每个 $i \in I$，设
$U_i$ 是 $\mathcal{C}$ 的对象，满足
\begin{enumerate}
\item $\coprod h_{U_i}$ 满射到 $\Sh(\mathcal{C})$ 的终对象；且
% P01-ERRATA: P01-E0143 -- h_{U_i} is generally only a presheaf; the object of Sh(C) should be h_{U_i}^# as in lemma-enough-points-local.
\item $\mathcal{C}/U_i$ 满足命题
\ref{sites-proposition-criterion-points} 的假设。
\end{enumerate}
则 $\mathcal{C}$ 有足够多的点。
\end{lemma}

\begin{proof}
由假设 (2) 与上述命题，$\mathcal{C}/U_i$ 有足够多的点。通过引理
\ref{sites-lemma-point-morphism-sites} 的过程，$\mathcal{C}/U_i$ 的点给出
$\mathcal{C}$ 的点。因此只须证明：若
$\phi : \mathcal{F} \to \mathcal{G}$ 是 $\mathcal{C}$ 上的层映射，且对所有
$i$，$\phi|_{\mathcal{C}/U_i}$ 都是同构，则 $\phi$ 是同构。由假设 (1)，
对 $\mathcal{C}$ 的每个对象 $W$，存在覆盖
$\{W_j \to W\}_{j \in J}$，使对每个 $j \in J$，存在 $i \in I$ 与态射
$f_j : W_j \to U_i$。于是映射
$\mathcal{F}(W_j) \to \mathcal{G}(W_j)$ 都是双射；类似地，
$\mathcal{F}(W_j \times_W W_{j'}) \to \mathcal{G}(W_j \times_W W_{j'})$
也都是双射。层条件告诉我们，$\mathcal{F}(W) \to \mathcal{G}(W)$ 是双射，
如所需。
\end{proof}



\section{弱可缩对象}
\label{sites-section-w-contractible}

\noindent
位点的一个{\it 弱可缩}对象，是指满足下述引理中各等价条件的对象。

\begin{lemma}
\label{sites-lemma-w-contractible}
设 $\mathcal{C}$ 为位点，$U$ 是 $\mathcal{C}$ 的对象。下列条件等价：
\begin{enumerate}
\item 对每个覆盖 $\{U_i \to U\}$，存在层映射
$h_U^\# \to \coprod h_{U_i}^\#$，它是
$\coprod h_{U_i} \to h_U$ 的层化的右逆。
\item 对每个集合层的满射 $\mathcal{F} \to \mathcal{G}$，映射
$\mathcal{F}(U) \to \mathcal{G}(U)$ 都是满射。
\end{enumerate}
\end{lemma}

\begin{proof}
假设 (1)，并设 $\mathcal{F} \to \mathcal{G}$ 是集合层的满射。对
$s \in \mathcal{G}(U)$，存在覆盖 $\{U_i \to U\}$ 以及
$t_i \in \mathcal{F}(U_i)$，其像为 $s|_{U_i}$；见定义
\ref{sites-definition-sheaves-injective-surjective}。通过
(\ref{sites-equation-map-representable-into-sheaf})，把 $t_i$ 看成映射
$t_i : h_{U_i}^\# \to \mathcal{F}$。将
$\coprod t_i : \coprod h_{U_i}^\# \to \mathcal{F}$ 与 (1) 给出的映射
$h_U^\# \to \coprod h_{U_i}^\#$ 复合，就得到映到 $s$ 的截面
$t \in \mathcal{F}(U)$。因此 (2) 成立。
% P01-ERRATA: P01-E0144 -- the source duplicates the predicate in "we get from (1) we obtain".

\medskip\noindent
假设 (2) 成立。设 $\{U_i \to U\}$ 为覆盖。则
$\coprod h_{U_i}^\# \to h_U^\#$ 是满射（引理
\ref{sites-lemma-covering-surjective-after-sheafification}）。因此由 (2)，存在
$\coprod h_{U_i}^\#$ 的截面 $s$，映到 $h_U^\#$ 的截面
$\text{id}_U$。该截面对应于映射
$h_U^\# \to \coprod h_{U_i}^\#$，它是
$\coprod h_{U_i} \to h_U$ 的层化的右逆，从而证明 (1)。
\end{proof}

\begin{definition}
\label{sites-definition-w-contractible}
设 $\mathcal{C}$ 为位点。
\begin{enumerate}
\item 若引理 \ref{sites-lemma-w-contractible} 的等价条件成立，则称
$\mathcal{C}$ 的对象 $U$ 是{\it 弱可缩的}。
\item 若 $\mathcal{C}$ 的每个对象 $U$ 都有一个覆盖
$\{U_i \to U\}$，且对所有 $i$，$U_i$ 都是弱可缩的，则称该位点有
{\it 足够多的弱可缩对象}。
\item 更一般地，若 $P$ 是 $\mathcal{C}$ 的对象的一种性质，并且
$\mathcal{C}$ 的每个对象 $U$ 都有覆盖 $\{U_i \to U\}$，使所有
对所有 $i$，$U_i$ 都具有 $P$，则称 $\mathcal{C}$ 有{\it 足够多的 $P$ 对象}。
\end{enumerate}
\end{definition}

\noindent
$\mathbf{A}^1_\mathbf{C}$ 的小 \'etale 位点没有任何弱可缩对象。另一方面，
任意概形的小 pro-\'etale 位点都有足够多的可缩对象。







\section{正像的正合性}
\label{sites-section-pushforward}

\noindent
设 $f$ 为拓扑斯态射。一般而言，函子 $f_*$ 仅左正合。可以对此函子附加许多条件，
以挑出特定类别的拓扑斯态射。我们在此汇集这些条件，并指出其中若干逻辑依赖关系。
下述引理有些部分纯属范畴论（也就是说，它们不依赖于拓扑斯态射；有一对伴随函子
就足够了）。

\begin{lemma}
\label{sites-lemma-exactness-properties}
设 $f : \Sh(\mathcal{C}) \to \Sh(\mathcal{D})$ 是拓扑斯态射。考虑下列性质
（针对集合层）：
\begin{enumerate}
\item $f_*$ 是忠实函子；
\item $f_*$ 是全忠实函子；
\item 对 $\Sh(\mathcal{C})$ 中的所有 $\mathcal{F}$，
$f^{-1}f_*\mathcal{F} \to \mathcal{F}$ 都是满射；
\item $f_*$ 将满射映为满射；
\item $f_*$ 与余等化子交换；
\item $f_*$ 与推出交换；
\item 对 $\Sh(\mathcal{C})$ 中的所有 $\mathcal{F}$，
$f^{-1}f_*\mathcal{F} \to \mathcal{F}$ 都是同构；
\item $f_*$ 反映单射；
\item $f_*$ 反映满射；
\item $f_*$ 反映双射；且
\item 对任意满射 $\mathcal{F} \to f^{-1}\mathcal{G}$，存在满射
$\mathcal{G}' \to \mathcal{G}$，使
$f^{-1}\mathcal{G}' \to f^{-1}\mathcal{G}$ 通过
$\mathcal{F} \to f^{-1}\mathcal{G}$ 分解。
\end{enumerate}
则有下列蕴含关系：
\begin{enumerate}
\item[(a)] (2) $\Rightarrow$ (1),
\item[(b)] (3) $\Rightarrow$ (1),
\item[(c)] (7) $\Rightarrow$ (1), (2), (3), (8), (9), (10).
\item[(d)] (3) $\Leftrightarrow$ (9),
\item[(e)] (6) $\Rightarrow$ (4) and (5) $\Rightarrow$ (4),
\item[(f)] (4) $\Leftrightarrow$ (11),
\item[(g)] (9) $\Rightarrow$ (8), (10), and
\item[(h)] (2) $\Leftrightarrow$ (7).
\end{enumerate}
图示如下：
$$
\xymatrix{
(6) \ar@{=>}[rd] & & & & & (9) \ar@{=>}[r] \ar@{=>}[rd] & (8) \\
& (4) \ar@{<=>}[r] & (11) &
(2) \ar@{<=>}[r] &
(7) \ar@{=>}[ru] \ar@{=>}[rd] & & (10) \\
(5) \ar@{=>}[ur] & & & & & (3) \ar@{=>}[r] & (1)
}
$$
\end{lemma}

\begin{proof}
(a) 的证明：由定义立即得出。

\medskip\noindent
(b) 的证明。设 $a, b : \mathcal{F} \to \mathcal{F}'$ 是
$\mathcal{C}$ 上的层映射。若 $f_*a = f_*b$，则
$f^{-1}f_*a = f^{-1}f_*b$。考虑交换图
$$
\xymatrix{
\mathcal{F} \ar@<-1ex>[r] \ar@<1ex>[r] & \mathcal{F}' \\
f^{-1}f_*\mathcal{F} \ar@<-1ex>[r] \ar@<1ex>[r] \ar[u] &
f^{-1}f_*\mathcal{F}' \ar[u]
}
$$
若下面两条箭头相等而竖直箭头为满射，则上面两条箭头相等。因此 (b) 得证。

\medskip\noindent
(c) 的证明。设 $a : \mathcal{F} \to \mathcal{F}'$ 是
$\mathcal{C}$ 上的层映射。考虑交换图
$$
\xymatrix{
\mathcal{F} \ar[r] & \mathcal{F}' \\
f^{-1}f_*\mathcal{F} \ar[r] \ar[u] &
f^{-1}f_*\mathcal{F}' \ar[u]
}
$$
若 (7) 成立，则竖直箭头均为同构。因此，若 $f_*a$ 为单射（相应地，满射或双射），
则下面的箭头为单射（相应地，满射或双射），从而上面的箭头也为单射（相应地，满射或
双射）。于是 (7) 蕴含 (8)、(9)、(10)。显然 (7) 蕴含 (3)。蕴含关系
(7) $\Rightarrow$ (2), (1) 由下面将证明的 (a) 和 (h) 得出。

\medskip\noindent
(d) 的证明。假设 (3)。设 $a : \mathcal{F} \to \mathcal{F}'$ 是
$\mathcal{C}$ 上的层映射，并且 $f_*a$ 是满射。由于 $f^{-1}$ 正合，这蕴含
$f^{-1}f_*a : f^{-1}f_*\mathcal{F} \to f^{-1}f_*\mathcal{F}'$
为满射。结合 (3)，得出 $a$ 是满射。这就是说 (9) 成立。
假设 (9)。设 $\mathcal{F}$ 是 $\mathcal{C}$ 上的层。需要证明映射
$f^{-1}f_*\mathcal{F} \to \mathcal{F}$ 是满射。只须证明
$f_*f^{-1}f_*\mathcal{F} \to f_*\mathcal{F}$ 是满射。而这是成立的，因为存在典范映射
$f_*\mathcal{F} \to f_*f^{-1}f_*\mathcal{F}$，它是该映射的一侧逆。

\medskip\noindent
(e) 的证明。以下不再说明地使用 Categories, Lemma
\ref{categories-lemma-characterize-mono-epi}。若
$\mathcal{F} \to \mathcal{F}'$ 是满射，则
$\mathcal{F}' \amalg_\mathcal{F} \mathcal{F}' \to \mathcal{F}'$
是同构。因此 (6) 蕴含
$$
f_*\mathcal{F}' \amalg_{f_*\mathcal{F}} f_*\mathcal{F}' =
f_*(\mathcal{F}' \amalg_\mathcal{F} \mathcal{F}')
\longrightarrow
f_*\mathcal{F}'
$$
也是同构。进而这蕴含 $f_*\mathcal{F} \to f_*\mathcal{F}'$ 是满射。因此 (6)
蕴含 (4)。若 $\mathcal{F} \to \mathcal{F}'$ 是满射，则由引理
\ref{sites-lemma-coequalizer-surjection}，$\mathcal{F}'$ 是两个投影
$\mathcal{F} \times_{\mathcal{F}'} \mathcal{F} \to \mathcal{F}$
的余等化子。因此若 (5) 成立，则 $f_*\mathcal{F}'$ 是两个投影
$$
f_*(\mathcal{F} \times_{\mathcal{F}'} \mathcal{F}) =
f_*\mathcal{F} \times_{f_*\mathcal{F}'} f_*\mathcal{F}
\longrightarrow
f_*\mathcal{F}
$$
的余等化子，而这显然意味着 $f_*\mathcal{F} \to f_*\mathcal{F}'$ 是满射。因此 (5)
也蕴含 (4)。

\medskip\noindent
(f) 的证明。假设 (4)。设 $\mathcal{F} \to f^{-1}\mathcal{G}$ 是
$\mathcal{C}$ 上的层的满射。由 (4)，
$f_*\mathcal{F} \to f_*f^{-1}\mathcal{G}$ 是满射。令
$\mathcal{G}'$ 为纤维积
$$
\xymatrix{
f_*\mathcal{F} \ar[r] & f_*f^{-1}\mathcal{G} \\
\mathcal{G}' \ar[u] \ar[r] & \mathcal{G} \ar[u]
}
$$
于是 $\mathcal{G}' \to \mathcal{G}$ 也是满射。考虑交换图
$$
\xymatrix{
\mathcal{F} \ar[r] & f^{-1}\mathcal{G} \\
f^{-1}f_*\mathcal{F} \ar[r] \ar[u] & f^{-1}f_*f^{-1}\mathcal{G} \ar[u] \\
f^{-1}\mathcal{G}' \ar[u] \ar[r] & f^{-1}\mathcal{G} \ar[u]
}
$$
由此得到所需结论。反过来，假设 (11)。设
$a : \mathcal{F} \to \mathcal{F}'$ 是 $\mathcal{C}$ 上的层的满射。
% P01-ERRATA: P01-E0145 -- the source omits the article in "be surjective map of sheaves".
考虑纤维积图
$$
\xymatrix{
\mathcal{F} \ar[r] & \mathcal{F}' \\
\mathcal{F}'' \ar[u] \ar[r] & f^{-1}f_*\mathcal{F}' \ar[u]
}
$$
由于下面的水平箭头是满射，由 (11) 可找到满射
$\gamma : \mathcal{G}' \to f_*\mathcal{F}'$，使 $f^{-1}\gamma$ 通过
$\mathcal{F}'' \to f^{-1}f_*\mathcal{F}'$ 分解：
$$
\xymatrix{
& \mathcal{F} \ar[r] & \mathcal{F}' \\
f^{-1}\mathcal{G}' \ar[r] & \mathcal{F}'' \ar[u] \ar[r] &
f^{-1}f_*\mathcal{F}' \ar[u]
}
$$
用 $f_*$ 将其向下推，得到交换图
$$
\xymatrix{
& f_*\mathcal{F} \ar[r] & f_*\mathcal{F}' \\
f_*f^{-1}\mathcal{G}' \ar[r] & f_*\mathcal{F}'' \ar[u] \ar[r] &
f_*f^{-1}f_*\mathcal{F}' \ar[u] \\
\mathcal{G}' \ar[u] \ar[rr] & & f_*\mathcal{F}' \ar[u]
}
$$
这证明了 (4)。

\medskip\noindent
(g) 的证明。假设 (9)。以下不再说明地使用 Categories, Lemma
\ref{categories-lemma-characterize-mono-epi}。设
$a : \mathcal{F} \to \mathcal{F}'$ 是 $\mathcal{C}$ 上的层映射，且
$f_*a$ 是单射。这意味着
$f_*\mathcal{F} \to f_*\mathcal{F} \times_{f_*\mathcal{F}'} f_*\mathcal{F} =
f_*(\mathcal{F} \times_{\mathcal{F}'} \mathcal{F})$ 是同构。于是由 (9)，
$\mathcal{F} \to \mathcal{F} \times_{\mathcal{F}'} \mathcal{F}$ 是满射，
即为同构。因此 $a$ 是单射，也就是说 (8) 成立。由于 (10) 显然等价于
(8) $+$ (9)，(g) 得证。

\medskip\noindent
(h) 的证明。这就是 Categories, Lemma
\ref{categories-lemma-adjoint-fully-faithful}。
\end{proof}

\noindent
下面给出位点态射的一个条件，它保证函子 $f_*$ 将满射映为满射。

\begin{lemma}
\label{sites-lemma-weaker}
设 $f : \mathcal{D} \to \mathcal{C}$ 是由连续函子
$u : \mathcal{C} \to \mathcal{D}$ 所关联的位点态射。假设对
$\mathcal{C}$ 的任意对象 $U$ 以及 $\mathcal{D}$ 中的任意覆盖
$\{V_j \to u(U)\}$，都存在 $\mathcal{C}$ 中的覆盖 $\{U_i \to U\}$，使层映射
$$
\coprod h_{u(U_i)}^\# \to h_{u(U)}^\#
$$
通过层映射
$$
\coprod h_{V_j}^\# \to h_{u(U)}^\#.
$$
分解。则 $f_*$ 将层的满射映为层的满射。
\end{lemma}

\begin{proof}
设 $a : \mathcal{F} \to \mathcal{G}$ 是 $\mathcal{D}$ 上层的满射。设 $U$ 是
$\mathcal{C}$ 的对象，并设
$s \in f_*\mathcal{G}(U) = \mathcal{G}(u(U))$。由假设，存在覆盖
$\{V_j \to u(U)\}$ 以及截面 $s_j \in \mathcal{F}(V_j)$，使
$a(s_j) = s|_{V_j}$。现在可以把截面 $s$、$s_j$ 与 $a$ 看成给出下列层映射交换图：
$$
\xymatrix{
\coprod h_{V_j}^\# \ar[r]_-{\coprod s_j} \ar[d] & \mathcal{F} \ar[d]^a \\
h_{u(U)}^\# \ar[r]^s & \mathcal{G}
}
$$
由假设，存在覆盖 $\{U_i \to U\}$，使上面的交换图可以扩充如下：
$$
\xymatrix{
& \coprod h_{V_j}^\# \ar[r]_-{\coprod s_j} \ar[d] & \mathcal{F} \ar[d]^a \\
\coprod h_{u(U_i)}^\# \ar[r] \ar[ur] &
h_{u(U)}^\# \ar[r]^s & \mathcal{G}
}
$$
因为 $\mathcal{F}$ 是层，从左下角到右上角的映射对应于一族截面
$s_i \in \mathcal{F}(u(U_i))$，即截面 $s_i \in f_*\mathcal{F}(U_i)$。
图的交换性蕴含 $a(s_i)$ 等于 $s$ 在 $U_i$ 上的限制。换言之，我们已经证明
$f_*a$ 是层的满射。
\end{proof}

\begin{example}
\label{sites-example-weaker-not-coequalizer}
假设 $f : \mathcal{D} \to \mathcal{C}$ 满足引理 \ref{sites-lemma-weaker}
的假设。一般而言，$f_*$ 并不与余等化子或推出交换。事实上，设 $f$ 是拓扑空间态射
$X = \{1, 2\} \to Y = \{*\}$ 所关联的位点态射（见例
\ref{sites-example-continuous-map}），其中 $Y$ 是单点空间，而
$X = \{1, 2\}$ 是具有两个点的离散空间。$X$ 上的层 $\mathcal{F}$ 由一对集合
$(A_1, A_2)$ 给出。于是 $f_*\mathcal{F}$ 对应于集合 $A_1 \times A_2$。
因此，若
$a = (a_1, a_2), b = (b_1, b_2) : (A_1, A_2) \to (B_1, B_2)$
是 $X$ 上的层映射，则 $a, b$ 的余等化子为 $(C_1, C_2)$，其中 $C_i$ 是
$a_i, b_i$ 的余等化子；而 $f_*a, f_*b$ 的余等化子是
$$
a_1 \times a_2, b_1 \times b_2 :
A_1 \times A_2 \longrightarrow B_1 \times B_2
$$
的余等化子，一般不同于 $C_1 \times C_2$。事实上，若
$A_2 = \emptyset$，则 $A_1 \times A_2 = \emptyset$，故所示两箭头的余等化子为
$B_1 \times B_2$；但一般而言 $C_1 \not = B_1$。对推出可作类似例子。
\end{example}

\noindent
下述引理给出一个判据，用以判断何时位点态射的函子 $f_*$ 反映单射与满射。
注意，这也蕴含 $f_*$ 是忠实函子，并且映射
$f^{-1}f_*\mathcal{F} \to \mathcal{F}$ 总是满射。

\begin{lemma}
\label{sites-lemma-cover-from-below}
设 $f : \mathcal{D} \to \mathcal{C}$ 是由函子
$u : \mathcal{C} \to \mathcal{D}$ 给出的位点态射。假设对
$\mathcal{D}$ 的每个对象 $V$，都存在 $\mathcal{C}$ 的对象 $U_i$ 以及态射
$u(U_i) \to V$，使 $\{u(U_i) \to V\}$ 是 $\mathcal{D}$ 的覆盖。在这种情形下，
函子 $f_* : \Sh(\mathcal{D}) \to \Sh(\mathcal{C})$ 反映单射与满射。
\end{lemma}

\begin{proof}
设 $\alpha : \mathcal{F} \to \mathcal{G}$ 是 $\mathcal{D}$ 上的层映射。
% P01-ERRATA: P01-E0146 -- the source says "be maps of sheaves" for the single map alpha.
由假设，对 $\mathcal{D}$ 的每个对象 $V$，层条件给出
$\mathcal{F}(V) \subset \prod \mathcal{F}(u(U_i)) =
\prod f_*\mathcal{F}(U_i)$，其中 $U_i \in \Ob(\mathcal{C})$ 适当；对
$\mathcal{G}$ 亦然。因此，若 $f_*\alpha$ 是单射，则显然 $\alpha$ 是单射。
换言之，$f_*$ 反映单射。

\medskip\noindent
假设 $f_*\alpha$ 是满射。对上述 $V, U_i, u(U_i) \to V$ 以及截面
$s \in \mathcal{G}(V)$，存在覆盖 $\{U_{ij} \to U_i\}$，使
$s|_{u(U_{ij})}$ 落在 $\mathcal{F}(u(U_{ij}))$ 的像中。由于
$\{u(U_{ij}) \to V\}$ 是覆盖（因为 $u$ 连续，并由位点公理），我们得出
$s$ 局部落在像中。因此 $\alpha$ 是满射。换言之，$f_*$ 反映满射。
\end{proof}

\begin{example}
\label{sites-example-cover-from-below}
我们构造一个满足引理 \ref{sites-lemma-cover-from-below} 假设的态射
$f : \mathcal{D} \to \mathcal{C}$。具体地，设
$\varphi : G \to H$ 是有限群的态射。考虑可数 $G$-集合与 $H$-集合的位点
$\mathcal{D} = \mathcal{T}_G$ 与 $\mathcal{C} = \mathcal{T}_H$，其覆盖为联合满射
的可数态射族（例 \ref{sites-example-site-on-group}）。设
$u : \mathcal{T}_H \to \mathcal{T}_G$ 是节
\ref{sites-section-G-sets-morphisms} 中描述的函子，而
$f : \mathcal{T}_G \to \mathcal{T}_H$ 是相应的位点态射。若 $\varphi$ 是单射，
则每个可数 $G$-集合，作为 $G$-集合，都是某个可数 $H$-集合的商（若
$\varphi$ 不是单射，这一点不成立）。因此 $f$ 满足引理
\ref{sites-lemma-cover-from-below} 的假设。
若 $\mathcal{T}_G$ 上的层 $\mathcal{F}$ 对应于 $G$-集合 $S$，则典范映射
$$
f^{-1}f_*\mathcal{F} \longrightarrow \mathcal{F}
$$
对应于映射
$$
\text{Map}_G(H, S) \longrightarrow S,\quad
a \longmapsto a(1_H)
$$
若 $\varphi$ 单射但不满射，则该映射是满射（正如引理
\ref{sites-lemma-cover-from-below} 所应有的结论），但一般不是单射（例如取
$G = \{1\}$、$H = \{1, \sigma\}$ 以及 $S = \{1, 2\}$）。此外，函子
$f_*$ 不与余等化子或推出交换（取 $G = \{1\}$ 且
$H = \{1, \sigma\}$）。
\end{example}

\section{几乎余连续的函子}
\label{sites-section-almost-cocontinuous}

\noindent
设 $\mathcal{C}$ 为位点。范畴 $\textit{PSh}(\mathcal{C})$ 有始对象，即把空集
赋给 $\mathcal{C}$ 的每个对象的预层。将该预层记为 $\emptyset$。由层化的性质，
$\emptyset$ 的层化 $\emptyset^\#$ 是 $\mathcal{C}$ 上层的范畴
$\Sh(\mathcal{C})$ 的始对象。

\begin{definition}
\label{sites-definition-empty}
设 $\mathcal{C}$ 为位点。若 $\emptyset^\# \to h_U^\#$ 是层的同构，则称
$\mathcal{C}$ 的对象 $U$ 在{\it 层论意义下为空}。
\end{definition}

\noindent
下述引理把这个概念表述得更明确。

\begin{lemma}
\label{sites-lemma-characterize-empty}
设 $\mathcal{C}$ 为位点，$U$ 是 $\mathcal{C}$ 的对象。下列条件等价：
\begin{enumerate}
\item $U$ 在层论意义下为空；
\item 对每个层 $\mathcal{F}$，$\mathcal{F}(U)$ 都是单点集；
\item $\emptyset^\#(U)$ 是单点集；
\item $\emptyset^\#(U)$ 非空；且
\item 空族是 $\mathcal{C}$ 中 $U$ 的一个覆盖。
\end{enumerate}
此外，若 $U$ 在层论意义下为空，则对 $\mathcal{C}$ 的任意态射
$U' \to U$，对象 $U'$ 也在层论意义下为空。
\end{lemma}

\begin{proof}
对任意层 $\mathcal{F}$，有
$\mathcal{F}(U) = \Mor_{\Sh(\mathcal{C})}(h_U^\#, \mathcal{F})$。
因此 (1) 与 (2) 等价。显然 (2) 蕴含 (3)，而 (3) 蕴含 (4)。
若 $U$ 的每个覆盖都由非空族给出，则由加号构造的定义，
$\emptyset^+(U)$ 为空。注意 $\emptyset^+ = \emptyset^\#$，因为
$\emptyset$ 是分离预层；见定理 \ref{sites-theorem-plus}。所以 (4) 蕴含 (5)。
若 (5) 成立，则由 $\mathcal{F}$ 的层条件，对每个层 $\mathcal{F}$，
$\mathcal{F}(U)$ 都是单点集；见注 \ref{sites-remark-sheaf-condition-empty-covering}。
因此 (5) 蕴含 (2)，从而 (1) -- (5) 等价。引理的最后一个断言由定义
\ref{sites-definition-site} 的公理 (3) 对 $U$ 的空覆盖应用而得。
% P01-ERRATA: P01-E0147 -- the source omits "to" in "applied the empty covering of U".
\end{proof}

\begin{definition}
\label{sites-definition-almost-cocontinuous}
设 $\mathcal{C}$、$\mathcal{D}$ 为位点，并设
$u : \mathcal{C} \to \mathcal{D}$ 为函子。若对 $\mathcal{C}$ 的每个对象
$U$ 以及每个覆盖 $\{V_j \to u(U)\}_{j \in J}$，都存在 $\mathcal{C}$ 中的覆盖
$\{U_i \to U\}_{i \in I}$，使对 $I$ 中每个 $i$，下列两个条件至少有一个成立：
\begin{enumerate}
\item $u(U_i)$ 在层论意义下为空；或
\item 对某个 $j \in J$，态射 $u(U_i) \to u(U)$ 通过 $V_j$ 分解，
\end{enumerate}
则称 $u$ 是{\it 几乎余连续的}。
\end{definition}

\noindent
该定义的动机来自拓扑空间的闭浸入 $i : Z \to X$。如例
\ref{sites-example-closed-map-cocontinuous-false} 所讨论的，连续函子
$X_{Zar} \to Z_{Zar}$，$U \mapsto Z \cap U$ 并不余连续，但它在上述意义下
几乎余连续。我们知道，$i_*$ 对集合层虽不正合，对阿贝尔群层却正合；见
Sheaves, Remark \ref{sheaves-remark-i-star-not-exact}。
% P01-ERRATA: P01-E0148 -- the source needs commas around the parenthetical "while not exact on sheaves of sets".
这对连续且几乎余连续的函子普遍成立。

\begin{lemma}
\label{sites-lemma-almost-cocontinuous-sheafification}
设 $\mathcal{C}$、$\mathcal{D}$ 为位点，并设
$u : \mathcal{C} \to \mathcal{D}$ 为函子。假设 $u$ 连续且几乎余连续。设
$\mathcal{G}$ 是 $\mathcal{D}$ 上的预层，并且每当 $V$ 在层论意义下为空时，
$\mathcal{G}(V)$ 都是单点集。则
$(u^p\mathcal{G})^\# = u^p(\mathcal{G}^\#)$。
\end{lemma}

\begin{proof}
设 $U \in \Ob(\mathcal{C})$。需要证明
$(u^p\mathcal{G})^\#(U) = u^p(\mathcal{G}^\#)(U)$。只须证明
$(u^p\mathcal{G})^+(U) = u^p(\mathcal{G}^+)(U)$，因为
$\mathcal{G}^+$ 仍是满足本引理假设的预层。我们有
$$
u^p(\mathcal{G}^+)(U) =
\mathcal{G}^+(u(U)) =
\colim_\mathcal{V} \check H^0(\mathcal{V}, \mathcal{G})
$$
其中余极限遍历 $\mathcal{D}$ 中 $u(U)$ 的所有覆盖 $\mathcal{V}$。另一方面，
$$
u^p(\mathcal{G})^+(U) =
\colim_\mathcal{U} \check H^0(u(\mathcal{U}), \mathcal{G})
$$
其中余极限遍历 $\mathcal{C}$ 中 $U$ 的覆盖
$\mathcal{U} = \{U_i \to U\}_{i \in I}$ 所成的范畴，并且
$u(\mathcal{U}) = \{u(U_i) \to u(U)\}_{i \in I}$。$u$ 连续这一条件意味着每个
$u(\mathcal{U})$ 都是覆盖。写成 $I = I_1 \amalg I_2$，其中
$$
I_2 = \{i \in I \mid u(U_i)\text{ 在层论意义下为空}\}
$$
则 $u(\mathcal{U})' = \{u(U_i) \to u(U)\}_{i \in I_1}$ 仍是该对象的覆盖，
% P01-ERRATA: P01-E0149 -- the source says "still a covering of because" and omits the covered object u(U).
因为其余每一块都可由空族覆盖，所以可以由定义 \ref{sites-definition-site} 的公理 (2)
删去。此外，由我们对 $\mathcal{G}$ 的假设，
$\check H^0(u(\mathcal{U}), \mathcal{G}) =
\check H^0(u(\mathcal{U})', \mathcal{G})$。最后，$u$ 几乎余连续这一条件蕴含：
对 $u(U)$ 的每个覆盖 $\mathcal{V}$，都存在 $U$ 的覆盖 $\mathcal{U}$，使
$u(\mathcal{U})'$ 加细 $\mathcal{V}$。所以，上面两个余极限具有相同的值，
如所需。
\end{proof}

\begin{lemma}
\label{sites-lemma-continuous-almost-cocontinuous}
设 $\mathcal{C}$、$\mathcal{D}$ 为位点，并设
$u : \mathcal{C} \to \mathcal{D}$ 为函子。假设 $u$ 连续且几乎余连续。则
$u^s = u^p : \Sh(\mathcal{D}) \to \Sh(\mathcal{C})$ 与推出和余等化子交换
（更一般地，与有限连通余极限交换）。
\end{lemma}

\begin{proof}
设 $\mathcal{I}$ 为有限连通指标范畴。设
$\mathcal{I} \to \Sh(\mathcal{D})$，$i \mapsto \mathcal{G}_i$ 为一个图。
% P01-ERRATA: P01-E0150 -- the source says "Let ... by a diagram" instead of "be a diagram".
我们知道，该图的余极限是预层范畴中余极限的层化；见引理
\ref{sites-lemma-colimit-sheaves}。以 $\colim^{Psh}$ 表示预层范畴中的余极限。
由于 $\mathcal{I}$ 有限且连通，预层
$\colim^{Psh}_i \mathcal{G}_i$ 满足引理
\ref{sites-lemma-almost-cocontinuous-sheafification} 的假设（因为单点集的有限连通余极限
是单点集）。所以该引理给出
\begin{align*}
u^s(\colim_i \mathcal{G}_i) & =
u^s((\colim^{Psh}_i \mathcal{G}_i)^\#) \\
& = (u^p(\colim^{Psh}_i \mathcal{G}_i))^\# \\
& = (\colim^{PSh}_i u^p(\mathcal{G}_i))^\# \\
& = \colim_i u^s(\mathcal{G}_i)
\end{align*}
如所需。
% P01-ERRATA: P01-E0151 -- the source inconsistently writes Psh and PSh for the same presheaf colimit superscript.
\end{proof}

\begin{lemma}
\label{sites-lemma-morphism-of-sites-almost-cocontinuous}
设 $f : \mathcal{D} \to \mathcal{C}$ 是与连续函子
$u : \mathcal{C} \to \mathcal{D}$ 关联的位点态射。若 $u$ 几乎余连续，则
$f_*$ 与推出和余等化子交换（更一般地，与有限连通余极限交换）。
\end{lemma}

\begin{proof}
这是引理 \ref{sites-lemma-continuous-almost-cocontinuous} 的特例。
\end{proof}

\section{子拓扑斯}
\label{sites-section-subtopoi}

\noindent
定义如下。

\begin{definition}
\label{sites-definition-embedding}
设 $\mathcal{C}$ 和 $\mathcal{D}$ 为位点。若 $f_*$ 全忠实，则称拓扑斯态射
$f : \Sh(\mathcal{D}) \to \Sh(\mathcal{C})$ 为一个{\it 嵌入}。
\end{definition}

\noindent
由引理 \ref{sites-lemma-exactness-properties}，这等价于要求：对
$\mathcal{D}$ 上的每个层 $\mathcal{F}$，伴随映射
$f^{-1}f_*\mathcal{F} \to \mathcal{F}$ 都是同构。

\begin{definition}
\label{sites-definition-subtopos}
设 $\mathcal{C}$ 为位点。若存在拓扑斯嵌入
$f : \Sh(\mathcal{D}) \to \Sh(\mathcal{C})$，使严格满子范畴
$E \subset \Sh(\mathcal{C})$ 等于函子 $f_*$ 的本质像，则称 $E$ 为一个
{\it 子拓扑斯}。
\end{definition}

\noindent
下述引理构造的子拓扑斯将在后面的定义中称为“开”的。

\begin{lemma}
\label{sites-lemma-open-subtopos}
设 $\mathcal{C}$ 为位点，$\mathcal{F}$ 是 $\mathcal{C}$ 上的层。下列条件等价：
\begin{enumerate}
\item $\mathcal{F}$ 是 $\Sh(\mathcal{C})$ 的终对象的子对象；且
\item 拓扑斯 $\Sh(\mathcal{C})/\mathcal{F}$ 是
$\Sh(\mathcal{C})$ 的子拓扑斯。
\end{enumerate}
\end{lemma}

\begin{proof}
我们已在引理 \ref{sites-lemma-localize-topos} 中看到，
$\Sh(\mathcal{C})/\mathcal{F}$ 是拓扑斯。这里回顾其证明。先应用引理
\ref{sites-lemma-topos-good-site}，可假设 $\mathcal{C}$ 是具有次典范拓扑、纤维积、
终对象 $X$ 以及满足 $\mathcal{F} = h_U$ 的对象 $U$ 的位点。引理
\ref{sites-lemma-localize-topos} 的证明表明，拓扑斯态射
$j_\mathcal{F} : \Sh(\mathcal{C})/\mathcal{F} \to \Sh(\mathcal{C})$
（在某些识别下）等于局部化态射
$j_U : \Sh(\mathcal{C}/U) \to \Sh(\mathcal{C})$。

\medskip\noindent
假设 (2)。这意味着对 $\mathcal{C}/U$ 上的所有层 $\mathcal{G}$，
$j_U^{-1}j_{U, *}\mathcal{G} \to \mathcal{G}$ 都是同构。对
$\mathcal{C}/U$ 的任意对象 $Z/U$，由引理
\ref{sites-lemma-localize-given-products}，有
$$
(j_{U, *}h_{Z/U})(U) = \Mor_{\mathcal{C}/U}(U \times_X U/U, Z/U)
$$
在上面的等式中令 $\mathcal{G} = h_{Z/U}$，便对所有 $Z/U$ 得到
$$
\Mor_{\mathcal{C}/U}(U \times_X U/U, Z/U) = \Mor_{\mathcal{C}/U}(U, Z/U)
$$
由 Yoneda 引理（Categories, Lemma \ref{categories-lemma-yoneda}），这蕴含
$U \times_X U = U$。由 Categories, Lemma
\ref{categories-lemma-characterize-mono-epi}，$U \to X$ 是单态射；换言之，
(1) 成立。

\medskip\noindent
假设 (1)。则由引理 \ref{sites-lemma-restrict-back}，
$j_U^{-1} j_{U, *} = \text{id}$。
\end{proof}

\begin{definition}
\label{sites-definition-open-subtopos}
设 $\mathcal{C}$ 为位点。若存在 $\Sh(\mathcal{C})$ 的终对象的子层
$\mathcal{F}$，使严格满子范畴 $E \subset \Sh(\mathcal{C})$ 是引理
\ref{sites-lemma-open-subtopos} 所述的子拓扑斯
$\Sh(\mathcal{C})/\mathcal{F}$，则称 $E$ 为一个{\it 开子拓扑斯}。
\end{definition}

\noindent
这意味着 $\Sh(\mathcal{C})$ 的开子拓扑斯之集，与
$\Sh(\mathcal{C})$ 的终对象的子对象之集之间存在双射。给定一个开子拓扑斯，
存在一个“闭”补。

\begin{lemma}
\label{sites-lemma-closed-subtopos}
设 $\mathcal{C}$ 为位点，$\mathcal{F}$ 是 $\Sh(\mathcal{C})$ 的终对象
$*$ 的子层。由所有满足
$\mathcal{F} \times \mathcal{G} \to \mathcal{F}$ 为同构的层
$\mathcal{G}$ 组成的满子范畴，是 $\Sh(\mathcal{C})$ 的子拓扑斯。
\end{lemma}

\begin{proof}
应用引理 \ref{sites-lemma-topos-good-site}，可假设 $\mathcal{C}$ 是具有该引理所列性质
的位点。特别地，$\mathcal{C}$ 有终对象 $X$（因而 $* = h_X$），也有满足
$\mathcal{F} = h_U$ 的对象 $U$。

\medskip\noindent
令 $\mathcal{D} = \mathcal{C}$ 作为范畴，但规定一族态射
$\{V_j \to V\}$ 是覆盖，当且仅当
$\{V_i \to V\} \cup \{U \times_X V \to V\}$ 是覆盖。
% P01-ERRATA: P01-E0152 -- the source switches the covering-family index from j to i without declaration.
由我们对 $\mathcal{C}$ 的选择，这恰好意味着
$$
h_{U \times_X V} \amalg \coprod h_{V_i} \longrightarrow h_V
$$
是满射。我们断言 $\mathcal{D}$ 是位点，即覆盖满足定义
\ref{sites-definition-site} 的条件 (1)、(2)、(3)。条件 (1) 成立。对条件 (2)，设
$\{V_i \to V\}$ 与 $\{V_{ij} \to V_i\}$ 是 $\mathcal{D}$ 的覆盖。则复合
$$
\coprod \left(
h_{U \times_X V_i} \amalg \coprod h_{V_{ij}}
\right) \longrightarrow
h_{U \times_X V} \amalg \coprod h_{V_i} \longrightarrow h_V
$$
是满射。由于每个态射 $U \times_X V_i \to V$ 都通过
$U \times_X V$ 分解，得出
$$
h_{U \times_X V} \amalg \coprod h_{V_{ij}} \longrightarrow h_V
$$
是满射，即 $\{V_{ij} \to V\}$ 是 $\mathcal{D}$ 中 $V$ 的覆盖。条件 (3)
类似地得出，因为层的满射在基变换下仍是满射。

\medskip\noindent
注意，（恒等）函子 $u : \mathcal{C} \to \mathcal{D}$ 连续，并与纤维积及终对象
交换。因此得到位点态射 $f : \mathcal{D} \to \mathcal{C}$（命题
\ref{sites-proposition-get-morphism}）。注意 $f_*$ 在底层预层上是恒等函子，因而全忠实。
为完成证明，需要说明 $f_*$ 的本质像正是本引理指定的满子范畴
$E \subset \Sh(\mathcal{C})$。为此，注意
$\mathcal{G} \in \Ob(\Sh(\mathcal{C}))$ 属于 $E$，当且仅当对
$\mathcal{C}$ 的所有对象 $V$，$\mathcal{G}(U \times_X V)$ 都是单点集。因此，
这样的层对 $\mathcal{D}$ 的所有覆盖都满足层性质（论证略）。反过来，若
$\mathcal{G}$ 对 $\mathcal{D}$ 的所有覆盖都满足层性质，则
$\mathcal{G}(U \times_X V)$ 是单点集，因为在 $\mathcal{D}$ 中，对象
$U \times_X V$ 由空覆盖覆盖。
\end{proof}

\begin{definition}
\label{sites-definition-closed-subtopos}
设 $\mathcal{C}$ 为位点。若存在 $\Sh(\mathcal{C})$ 的终对象的子层
$\mathcal{F}$，使严格满子范畴 $E \subset \Sh(\mathcal{C})$ 是引理
\ref{sites-lemma-closed-subtopos} 所述的子拓扑斯，则称 $E$ 为一个
{\it 闭子拓扑斯}。
\end{definition}

\noindent
现在可以定义拓扑斯的闭浸入与开浸入。

\begin{definition}
\label{sites-definition-immersion-topoi}
设 $f : \Sh(\mathcal{D}) \to \Sh(\mathcal{C})$ 为拓扑斯态射。
\begin{enumerate}
\item 若 $f$ 是嵌入，且 $f_*$ 的本质像是开子拓扑斯，则称 $f$ 是
{\it 开浸入}。
\item 若 $f$ 是嵌入，且 $f_*$ 的本质像是闭子拓扑斯，则称 $f$ 是
{\it 闭浸入}。
\end{enumerate}
\end{definition}

\begin{lemma}
\label{sites-lemma-closed-immersion}
设 $i : \Sh(\mathcal{D}) \to \Sh(\mathcal{C})$ 为拓扑斯的闭浸入。则
$i_*$ 全忠实，将满射映为满射，与余等化子交换，与推出交换，并反映单射、满射与双射。
\end{lemma}

\begin{proof}
设 $\mathcal{F}$ 是 $\Sh(\mathcal{C})$ 的终对象 $*$ 的子层，并设
$E \subset \Sh(\mathcal{C})$ 是由所有满足
$\mathcal{F} \times \mathcal{G} \to \mathcal{F}$ 为同构的
$\mathcal{G}$ 组成的满子范畴。由引理 \ref{sites-lemma-closed-subtopos}，函子
$i_*$ 同构于包含函子 $\iota : E \to \Sh(\mathcal{C})$。

\medskip\noindent
设 $j_{\mathcal{F}} : \Sh(\mathcal{C})/\mathcal{F} \to \Sh(\mathcal{C})$
为局部化函子（引理 \ref{sites-lemma-localize-topos}）。注意，$E$ 也可以描述为所有满足
$j_\mathcal{F}^{-1}\mathcal{G} = *$ 的层 $\mathcal{G}$ 之集。

\medskip\noindent
设 $a, b : \mathcal{G}_1 \to \mathcal{G}_2$ 是 $E$ 的两个态射。
% P01-ERRATA: P01-E0153 -- the source says "two morphism" instead of "two morphisms".
要证明 $\iota$ 与余等化子交换，只须证明 $a$、$b$ 在
$\Sh(\mathcal{C})$ 中的余等化子属于 $E$。这是清楚的，因为两个态射
$* \to *$ 的余等化子是 $*$，并且 $j_\mathcal{F}^{-1}$ 正合。推出的情形类似。

\medskip\noindent
因此 $i_*$ 满足引理 \ref{sites-lemma-exactness-properties} 的性质 (5)、(6)、(7)，
从而态射 $i$ 满足该引理中提到的所有性质，特别是本引理提到的那些性质。
\end{proof}

\section{代数结构层}
\label{sites-section-sheaves-algebraic-structures}

\noindent
在 Sheaves, Section \ref{sheaves-section-algebraic-structures} 中，我们把一种
代数结构类型定义为序对 $(\mathcal{A}, s)$，其中 $\mathcal{A}$ 是范畴，而
$s : \mathcal{A} \to \textit{Sets}$ 是满足下列条件的函子：
\begin{enumerate}
\item $s$ 忠实；
\item $\mathcal{A}$ 有极限，且 $s$ 与极限交换；
\item $\mathcal{A}$ 有滤过余极限，且 $s$ 与它们交换；且
\item $s$ 反映同构。
\end{enumerate}
对这样的代数结构类型，我们已经看到：空间 $X$ 上取值于 $\mathcal{A}$ 的预层
$\mathcal{F}$ 是层，当且仅当与它关联的集合预层是层。此外，我们阐明了茎的概念；
并且给定连续映射 $f : X \to Y$，定义了代数结构层上的伴随函子正像与逆像，
它们与底层集合层上的正像与逆像一致。再者，把基上的代数结构层扩张到空间的所有
开集，也如预期那样运作。
% P01-ERRATA: P01-E0154 -- the source's compound subject "pushforward and pullback" takes "agree", not "agrees".

\medskip\noindent
这部分材料在位点与层的框架中仍然成立。设 $(\mathcal{A}, s)$ 为一种代数结构类型，
$\mathcal{C}$ 为位点。分别以
$\textit{PSh}(\mathcal{C}, \mathcal{A})$ 与
$\Sh(\mathcal{C}, \mathcal{A})$ 表示 $\mathcal{C}$ 上取值于
$\mathcal{A}$ 的预层范畴与层范畴。
\begin{itemize}
\item[] ($\alpha$) 取值于 $\mathcal{A}$ 的预层是层，当且仅当其底层集合预层是层。
见 Sheaves, Lemma \ref{sheaves-lemma-sheaves-structure} 的证明。
\item[] ($\beta$) 给定取值于 $\mathcal{A}$ 的预层 $\mathcal{F}$，预层
${\mathcal{F}}^\# = (\mathcal{F}^+)^+$ 是层。这是因为层化过程中出现的余极限
是滤过的，甚至是遍历有向集的余极限（见节 \ref{sites-section-sheafification}，特别是
引理 \ref{sites-lemma-sheafification-exact} 的证明），并且 $s$ 与滤过余极限交换。
\item[] ($\gamma$) 得到交换图
$$
\xymatrix{
\Sh(\mathcal{C}, \mathcal{A}) \ar@<1ex>[r] \ar[d]_s &
\textit{PSh}(\mathcal{C}, \mathcal{A}) \ar@<1ex>[l]^{{}^\#} \ar[d]^s\\
\Sh(\mathcal{C}) \ar@<1ex>[r] &
\textit{PSh}(\mathcal{C}) \ar@<1ex>[l]
}
$$
\item[] ($\delta$) 当且仅当 $\mathcal{F}$ 是代数结构层时，
$\mathcal{F} = \mathcal{F}^\#$。
\item[] ($\epsilon$) 函子 ${}^\#$ 左伴随于包含函子：
$$
\Mor_{\textit{PSh}(\mathcal{C}, \mathcal{A})}(\mathcal{G}, \mathcal{F})
=
\Mor_{\Sh(\mathcal{C}, \mathcal{A})}(\mathcal{G}^\#, \mathcal{F})
$$
其证明与命题 \ref{sites-proposition-sheafification-adjoint} 的证明相同。
\item[] ($\zeta$) 函子 $\mathcal{F} \mapsto \mathcal{F}^\#$ 左正合。
其证明与引理 \ref{sites-lemma-sheafification-exact} 的证明相同。
\end{itemize}

\begin{definition}
\label{sites-definition-pushforward-algebraic-structures}
设 $f : \mathcal{D} \to \mathcal{C}$ 是由函子
$u : \mathcal{C} \to \mathcal{D}$ 给出的位点态射。对代数结构预层，按规则
$u^p\mathcal{F}(U) = \mathcal{F}(uU)$ 定义{\it 正像}函子；对代数结构层，
以同一规则定义，即 $f_*\mathcal{F}(U) = \mathcal{F}(uU)$。
\end{definition}

\noindent
问题出现在试图定义逆像时。原因是节 \ref{sites-section-functoriality-PSh} 中定义函子
$u_p$ 的余极限未必滤过。因此，一般而言，上面的公理不足以定义代数结构（预）层的
逆像。尽管如此，几乎在所有情形中，下述引理都足以定义代数结构（预）层的正像与逆像。
% P01-ERRATA: P01-E0155 -- the source says "trying the define" instead of "trying to define".

\begin{lemma}
\label{sites-lemma-push-pull-good-case}
假设函子 $u : \mathcal{C} \to \mathcal{D}$ 满足命题
\ref{sites-proposition-get-morphism} 的假设，因而产生位点态射
$f : \mathcal{D} \to \mathcal{C}$。在这种情形下，逆像函子
$f^{-1}$（相应地，$u_p$）与正像函子 $f_*$（相应地，$u^p$）扩张为代数结构层
（相应地，预层）范畴上的一对伴随函子。此外，这些函子与取底层集合层
（相应地，预层）的操作交换。
\end{lemma}

\begin{proof}
上面已经定义了 $f_* = u^p$。在命题 \ref{sites-proposition-get-morphism} 的证明过程中，
我们看到，在该命题的假设下，定义 $u_p$ 所用的所有余极限都是滤过的。因此由代数结构
类型的定义，可以用完全相同的余极限在代数结构预层上定义 $u_p$。$u_p$ 与 $u^p$
的伴随性，其证明与引理 \ref{sites-lemma-adjoints-u} 的证明完全相同。上面对代数结构预层
之层化的讨论，于是蕴含可以定义
$f^{-1}(\mathcal{F}) = (u_p\mathcal{F})^\#$。
\end{proof}

\noindent
下面简要讨论如何处理一般位点态射、更一般地任意拓扑斯态射的逆像与正像。

\medskip\noindent
设 $\mathcal{C}$ 为位点。在 $\mathcal{A} = \textit{Ab}$ 的情形，可以把
$\mathcal{C}$ 上的阿贝尔（预）层看成四元组
$(\mathcal{F}, +, 0, i)$。其中数据为
\begin{enumerate}
\item[(D1)] $\mathcal{F}$ 是集合层；
\item[(D2)] $+ : \mathcal{F} \times \mathcal{F} \to \mathcal{F}$ 是集合层的态射；
\item[(D3)] $0 : * \to \mathcal{F}$ 是从单点层（见例
\ref{sites-example-singleton-sheaf}）到 $\mathcal{F}$ 的态射；且
\item[(D4)] $i : \mathcal{F} \to \mathcal{F}$ 是集合层的态射。
\end{enumerate}
这些数据必须满足下列公理：
\begin{enumerate}
\item[(A1)] $+$ 结合且交换；
\item[(A2)] $0$ 是 $+$ 的单位元；且
\item[(A3)] $+ \circ (1, i) = 0 \circ (\mathcal{F} \to *)$。
\end{enumerate}
对照 Sheaves, Lemma \ref{sheaves-lemma-abelian-presheaves}。设
$f : \mathcal{D} \to \mathcal{C}$ 为位点态射。注意，由于 $f^{-1}$ 正合，有
$f^{-1}* = *$ 以及
$f^{-1}(\mathcal{F} \times \mathcal{F}) =
f^{-1}\mathcal{F} \times f^{-1}\mathcal{F}$。因此可以直接把
$f^{-1}\mathcal{F}$ 定义为四元组
$(f^{-1}\mathcal{F}, f^{-1}+, f^{-1}0, f^{-1}i)$。由于 $f^{-1}$ 是与有限极限
交换的函子，这些公理得以保持。最后，不难验证 $f_*$ 与 $f^{-1}$ 如通常那样伴随。

\medskip\noindent
在 \cite{SGA4} 中使用了这种方法。他们引入所谓
``{\it esp\`ece the structure alg\'ebrique $\ll$d\'efinie
par limites projectives finie$\gg$}''。对这样的 esp\`ece，可以用上述方法定义一对伴随
函子 $f^{-1}$ 与 $f_*$。这显然适用于实践中遇到的大多数代数结构。我们不把该构造
形式化，而只列出这种方法适用的代数结构（逐一验证）。事实上，该方法适用于任意
拓扑斯态射。
% P01-ERRATA: P01-E0156 -- the quoted French has "the" for "de" and singular "finie" where "limites" requires plural "finies".

\begin{proposition}
\label{sites-proposition-functoriality-algebraic-structures-topoi}
\begin{slogan}
拓扑斯态射保持代数结构。
\end{slogan}
设 $\mathcal{C}$、$\mathcal{D}$ 为位点。设 $f = (f^{-1}, f_*)$ 为从
$\Sh(\mathcal{D}) \to \Sh(\mathcal{C})$ 的拓扑斯态射。上述方法对下列代数结构
类型给出代数结构层上的一对伴随函子 $(f^{-1}, f_*)$，且与取底层集合层相容：
\begin{enumerate}
\item 带基点集合；
\item 阿贝尔群；
\item 群；
\item 幺半群；
\item 环；
\item 固定环上的模；且
\item 固定域上的李代数。
\end{enumerate}
此外，在每种情形中，上面标为 ($\alpha$)、($\beta$)、($\gamma$)、($\delta$)、
($\epsilon$)、($\zeta$) 的结果都成立。
\end{proposition}

\begin{proof}
命题的最后一个断言成立，是因为每个所列范畴配备显然的遗忘函子后，确实都是本节
开头所解释意义下的一种代数结构类型。见 Sheaves, Lemma
\ref{sheaves-lemma-list-algebraic-structures}。

\medskip\noindent
(2) 的证明。按照命题前的讨论，把阿贝尔群层看成四元组
$(\mathcal{F}, +, 0, i)$。若 $(\mathcal{F}, +, 0, i)$ 定义在
$\mathcal{C}$ 上，则其逆像定义为
$(f^{-1}\mathcal{F}, f^{-1}+, f^{-1}0, f^{-1}i)$。若
$(\mathcal{G}, +, 0, i)$ 定义在 $\mathcal{D}$ 上，则其正像定义为
$(f_*\mathcal{G}, f_*+, f_*0, f_*i)$。这是可行的，因为
$f_*(\mathcal{G} \times \mathcal{G}) = f_*\mathcal{G} \times f_*\mathcal{G}$。
伴随性由集合层函子的伴随性得出，因为
$$
\Mor_{\textit{Ab}(\mathcal{C})}
((\mathcal{F}_1, +, 0, i), (\mathcal{F}_2, +, 0, i))
=
\left\{
\begin{matrix}
\varphi \in \Mor_{\Sh(\mathcal{C})}
(\mathcal{F}_1, \mathcal{F}_2) \\
\varphi \text{ 与下列映射相容： }+, 0, i
\end{matrix}
\right\}
$$
细节留给读者。

\medskip\noindent
这种方法也适用于环层：把一个（有单位元的）环层看成满足初等代数教材中所列公理的
六元组 $(\mathcal{O}, + , 0, i, \cdot, 1)$ 即可。

\medskip\noindent
带基点集合层是序对 $(\mathcal{F}, p)$，其中 $\mathcal{F}$ 是集合层，而
$p : * \to \mathcal{F}$ 是集合层的映射。

\medskip\noindent
群层由带有适当公理的四元组 $(\mathcal{F}, \cdot, 1, i)$ 给出。

\medskip\noindent
幺半群层由带有适当公理的序对 $(\mathcal{F}, \cdot)$ 给出。
% P01-ERRATA: P01-E0157 -- the source has singular "axiom" after "suitable" although monoid structure requires multiple laws.

\medskip\noindent
设 $R$ 为环。$R$-模层由五元组
$(\mathcal{F}, +, 0, i, \{\lambda_r\}_{r \in R})$ 给出，其中四元组
$(\mathcal{F}, +, 0, i)$ 是上述阿贝尔群层，而
$\lambda_r : \mathcal{F} \to \mathcal{F}$ 是满足下列条件的一族集合层态射：
% P01-ERRATA: P01-E0158 -- the source says "An sheaf" instead of "A sheaf".
$\lambda_r \circ 0 = 0$，
$\lambda_r \circ + = + \circ (\lambda_r, \lambda_r)$，
$\lambda_{r + r'} =
+ \circ \lambda_r \times \lambda_{r'} \circ (\text{id}, \text{id})$，
$\lambda_{rr'} = \lambda_r \circ \lambda_{r'}$，
$\lambda_1 = \text{id}$，$\lambda_0 = 0 \circ (\mathcal{F} \to *)$。
% P01-ERRATA: P01-E0159 -- the distributivity formula needs parentheses around (lambda_r times lambda_r') to make the composition typed and unambiguous.
\end{proof}

\noindent
我们将在 Modules on Sites, Section
\ref{sites-modules-section-sheaves-modules} 中讨论环层上的模层范畴。

\begin{remark}
\label{sites-remark-no-pullback-presheaves}
设 $\mathcal{C}$、$\mathcal{D}$ 为位点。设
$u : \mathcal{D} \to \mathcal{C}$ 是连续函子，它产生位点态射
$\mathcal{C} \to \mathcal{D}$。注意，即使在阿贝尔群的情形，我们也尚未定义阿贝尔群
预层的逆像函子。由于阿贝尔群范畴中所有余极限均可表出，当然可以用定义集合预层上
$u_p$ 时所用的相同余极限，在阿贝尔预层上定义函子 $u_p^{ab}$。在阿贝尔预层范畴上，
$u_p^{ab}$ 也将左伴随于 $u^p$。然而，在底层集合预层上，$u_p^{ab}$ 未必总与
$u_p$ 一致。
\end{remark}

\section{拉回映射}
\label{sites-section-pullback}

\noindent
有时位点 $\mathcal{C}$ 没有终对象。在这种情形下，按如下方式定义整体截面函子。

\begin{definition}
\label{sites-definition-global-sections}
位点 $\mathcal{C}$ 上的集合预层 $\mathcal{F}$ 的{\it 整体截面}是集合
$$
\Gamma(\mathcal{C}, \mathcal{F}) =
\Mor_{\textit{PSh}(\mathcal{C})}(*, \mathcal{F})
$$
其中 $*$ 是 $\mathcal{C}$ 上预层范畴的终对象，即把单点集赋给每个对象的预层。
% P01-ERRATA: P01-E0160 -- the source's plural subject "global sections" is followed by singular "is".
\end{definition}

\noindent
当然，同一定义也适用于层。下面给出一种计算整体截面的方式。

\begin{lemma}
\label{sites-lemma-compute-global-sections}
设 $\mathcal{C}$ 为位点。设 $a, b : V \to U$ 是 $\mathcal{C}$ 的态射，使
$$
\xymatrix{
h_V^\# \ar@<1ex>[r] \ar@<-1ex>[r] & h_U^\# \ar[r] & {*}
}
$$
是 $\Sh(\mathcal{C})$ 中的余等化子。则
$\Gamma(\mathcal{C}, \mathcal{F})$ 是
$a^*, b^* : \mathcal{F}(U) \to \mathcal{F}(V)$ 的等化子。
% P01-ERRATA: P01-E0161 -- the source calls a,b:V->U "objects" of C; they are morphisms.
\end{lemma}

\begin{proof}
因为 $\Mor_{\Sh(\mathcal{C})}(h_U^\#, \mathcal{F}) = \mathcal{F}(U)$，
这由定义即得。
\end{proof}

\noindent
现在设 $f : \Sh(\mathcal{D}) \to \Sh(\mathcal{C})$ 为拓扑斯态射。则对
$\mathcal{C}$ 上的任意层 $\mathcal{F}$，存在拉回映射
$$
f^{-1} :
\Gamma(\mathcal{C}, \mathcal{F})
\longrightarrow
\Gamma(\mathcal{D}, f^{-1}\mathcal{F})
$$
事实上，由于 $f^{-1}$ 正合，它把 $*$ 映为 $*$。所以 $\mathcal{F}$ 在
$\mathcal{C}$ 上的整体截面 $s$，即层映射 $s : * \to \mathcal{F}$，可以拉回为
$f^{-1}s : * = f^{-1}* \to f^{-1}\mathcal{F}$。

\medskip\noindent
可以稍作推广：考虑分别定义在 $\mathcal{C}$、$\mathcal{D}$ 上的一对层
$\mathcal{F}$、$\mathcal{G}$，连同一个映射
$f^{-1}\mathcal{F} \to \mathcal{G}$。将上述构造与显然的映射
$\Gamma(\mathcal{D}, f^{-1}\mathcal{F}) \to \Gamma(\mathcal{D}, \mathcal{G})$
复合，便得到映射
$$
\Gamma(\mathcal{C}, \mathcal{F})
\longrightarrow
\Gamma(\mathcal{D}, \mathcal{G})
$$
该映射有时也称为拉回映射。

\medskip\noindent
自然出现的一个稍微更一般的构造如下。假设有拓扑斯态射的交换图
$$
\xymatrix{
\Sh(\mathcal{D}) \ar[rd]_h \ar[rr]_f & &
\Sh(\mathcal{C}) \ar[ld]^g \\
& \Sh(\mathcal{B})
}
$$
再假设 $\mathcal{F}$ 是 $\mathcal{C}$ 上的层。则存在{\it 拉回映射}
$$
f^{-1} : g_*\mathcal{F} \longrightarrow h_*f^{-1}\mathcal{F}
$$
它正是由识别
$g_*f_*f^{-1}\mathcal{F} = h_*f^{-1}\mathcal{F}$，以及对典范映射
$\mathcal{F} \to f_*f^{-1}\mathcal{F}$ 应用 $g_*$ 所得到的映射。若 $g$
是恒等态射，则该映射在整体截面上与上面的拉回映射一致。

\medskip\noindent
在前一段的情形下，再设有分别定义在 $\mathcal{C}$、$\mathcal{D}$ 上的一对层
$\mathcal{F}$、$\mathcal{G}$，连同一个映射
$f^{-1}\mathcal{F} \to \mathcal{G}$。把上述拉回映射与对
$f^{-1}\mathcal{F} \to \mathcal{G}$ 应用 $h_*$ 所得的映射复合，得到
$$
g_*\mathcal{F} \longrightarrow h_*\mathcal{G}
$$
限制到 $\mathcal{B}$ 的某个对象上的截面，就恢复上面讨论的整体截面“拉回映射”
（选取适当位点）。

\medskip\noindent
更一般地，设有拓扑斯的交换图
$$
\xymatrix{
\Sh(\mathcal{D}) \ar[d]_h \ar[r]_f &
\Sh(\mathcal{C}) \ar[d]^g \\
\Sh(\mathcal{B}) \ar[r]^e & \Sh(\mathcal{A})
}
$$
以及 $\mathcal{C}$ 上的层 $\mathcal{G}$。则存在{\it 基变换映射}
$$
e^{-1}g_*\mathcal{G} \longrightarrow h_*f^{-1}\mathcal{G}.
$$
该映射伴随于映射
$g_*\mathcal{G} \to e_*h_*f^{-1}\mathcal{G} = (e \circ h)_*f^{-1}\mathcal{G}$，
后者就是刚才描述的拉回映射。

\begin{remark}
\label{sites-remark-compose-base-change}
考虑交换图
$$
\xymatrix{
\Sh(\mathcal{B}') \ar[r]_k \ar[d]_{f'} &
\Sh(\mathcal{B}) \ar[d]^f \\
\Sh(\mathcal{C}') \ar[r]^l \ar[d]_{g'} &
\Sh(\mathcal{C}) \ar[d]^g \\
\Sh(\mathcal{D}') \ar[r]^m &
\Sh(\mathcal{D})
}
$$
则两个方块的基变换映射复合后给出外部矩形的基变换映射。更精确地，复合
\begin{align*}
m^{-1} \circ (g \circ f)_*
& =
m^{-1} \circ g_* \circ f_* \\
& \to g'_* \circ l^{-1} \circ f_* \\
& \to g'_* \circ f'_* \circ k^{-1} \\
& = (g' \circ f')_* \circ k^{-1}
\end{align*}
是该矩形的基变换映射。验证略去。
\end{remark}

\begin{remark}
\label{sites-remark-compose-base-change-horizontal}
考虑交换图
$$
\xymatrix{
\Sh(\mathcal{C}'') \ar[r]_{g'} \ar[d]_{f''} &
\Sh(\mathcal{C}') \ar[r]_g \ar[d]_{f'} &
\Sh(\mathcal{C}) \ar[d]^f \\
\Sh(\mathcal{D}'') \ar[r]^{h'} &
\Sh(\mathcal{D}') \ar[r]^h &
\Sh(\mathcal{D})
}
$$
的环化拓扑斯。则两个方块的基变换映射复合后给出外部矩形的基变换映射。更精确地，
复合
\begin{align*}
(h \circ h')^{-1} \circ f_*
& =
(h')^{-1} \circ h^{-1} \circ f_* \\
& \to (h')^{-1} \circ f'_* \circ g^{-1} \\
& \to f''_* \circ (g')^{-1} \circ g^{-1} \\
& = f''_* \circ (g \circ g')^{-1}
\end{align*}
是该矩形的基变换映射。验证略去。
% P01-ERRATA: P01-E0162 -- the source calls this a diagram of ringed topoi although only bare sheaf topoi are shown and no ring data enter the construction.
\end{remark}

\section{与 SGA4 的比较}
\label{sites-section-notation-comparison}

\noindent
节 \ref{sites-section-functoriality-PSh} 中函子 $u^p$、$u_p$ 的记号，以及节
\ref{sites-section-continuous-functors} 中 $u^s$、$u_s$ 的记号，取自
\cite[pages 14 and 42]{ArtinTopologies}。作出这些选择之后，节
\ref{sites-section-more-functoriality-PSh} 中函子 ${}_pu$ 与节
\ref{sites-section-cocontinuous-functors} 中函子 ${}_su$ 的记号看起来也很合理。
本节将我们的记号与 SGA4 的记号比较。

\medskip\noindent
{\bf 预层：}设 $u : \mathcal{C} \to \mathcal{D}$ 为范畴间的函子。
在 \cite[Exposee I, Section 5]{SGA4} 中，函子 $u^p$ 记为 $u^*$。
在 \cite[Exposee I, Proposition 5.1]{SGA4} 中，函子 $u_p$ 记为 $u_!$；
在 \cite[Exposee I, Proposition 5.1]{SGA4} 中，函子 ${}_pu$ 记为 $u_*$。
换言之，
$$
u_p, u^p, {}_pu\quad(SP)
\quad\text{与}\quad
u_!, u^*, u_*\quad(SGA4)
$$
需要提醒读者，SGA4 的不同部分对这些函子使用不同记号。

\medskip\noindent
{\bf 层与连续函子：}假设 $\mathcal{C}$ 与 $\mathcal{D}$ 为位点，而
$u : \mathcal{C} \to \mathcal{D}$ 是连续函子（定义
\ref{sites-definition-continuous}）。在 \cite[Exposee III, 1.11]{SGA4} 中，
函子 $u^s$ 记为 $u_s$；在 \cite[Exposee III, Proposition 1.2]{SGA4} 中，
函子 $u_s$ 记为 $u^s$。换言之，
$$
u_s, u^s\quad(SP)
\quad\text{与}\quad
u^s, u_s\quad(SGA4)
$$
当 $u$ 定义位点态射 $f : \mathcal{D} \to \mathcal{C}$（定义
\ref{sites-definition-morphism-sites}）时，相应的拓扑斯态射（引理
\ref{sites-lemma-morphism-sites-topoi}）与 \cite[Exposee IV, (4.9.1.1)]{SGA4}
中的相同。

\medskip\noindent
{\bf 层与余连续函子：}假设 $\mathcal{C}$ 与 $\mathcal{D}$ 为位点，而
$u : \mathcal{C} \to \mathcal{D}$ 是余连续函子（定义
\ref{sites-definition-cocontinuous}）。在
\cite[Exposee III, Proposition 2.3]{SGA4} 中，函子 ${}_su$（引理
\ref{sites-lemma-pu-sheaf}）记为 $u_*$；在
\cite[Exposee III, Proposition 2.3]{SGA4} 中，函子 $(u^p\ )^\#$ 记为 $u^*$。
换言之，
$$
(u^p\ )^\#, {}_su\quad(SP)
\quad\text{与}\quad
u^*, u_*\quad(SGA4)
$$
因此，引理 \ref{sites-lemma-cocontinuous-morphism-topoi} 中与 $u$ 关联的拓扑斯态射，
与 \cite[Exposee IV, 4.7]{SGA4} 中的相同。

\medskip\noindent
{\bf 拓扑斯态射：}若 $f$ 是由函子 $(f^{-1}, f_*)$ 给出的拓扑斯态射，则在
\cite[Exposee IV, Definition 3.1]{SGA4} 中，函子 $f^{-1}$ 记为 $f^*$。
我们将用 $f^{-1}$ 表示集合层、或更一般地代数结构层的逆像（节
\ref{sites-section-sheaves-algebraic-structures}）。对于环化拓扑斯的态射，我们将用
$f^*$ 表示模层的逆像（Modules on Sites, Definition
\ref{sites-modules-definition-pushforward}）。

\section{拓扑}
\label{sites-section-topologies}

\noindent
本节按照 \cite{SGA4} 定义范畴上的拓扑。可以用这种语言发展层与拓扑斯的全部理论。
关于这些材料的现代论述见 \cite{KS}。不过，对代数几何最重要的情形，是位点在其
底层范畴上定义的拓扑。因此，我们强烈建议初次阅读的读者
{\bf 跳过本节以及本章后面的所有其他章节}！

\begin{definition}
\label{sites-definition-sieve}
设 $\mathcal{C}$ 为范畴，$U \in \Ob(\mathcal{C})$。$U$ 上的一个
{\it 筛 $S$} 是子预层 $S \subset h_U$。
\end{definition}

\noindent
换言之，$U$ 上的筛对每个对象 $T \in \Ob(\mathcal{C})$，从所有态射
$T \to U$ 所成的集合中选出一个子集 $S(T)$。事实上，对子集族
$S(T) \subset h_U(T) = \Mor_\mathcal{C}(T, U)$ 的唯一要求是下列规则：
\begin{equation}
\label{sites-equation-property-sieve}
\left.
\begin{matrix}
(\alpha : T \to U) \in S(T) \\
g : T' \to T
\end{matrix}
\right\} \Rightarrow
(\alpha \circ g : T' \to U) \in S(T')
\end{equation}
一个有用的直观图景，是把映射 $S \to h_U$ 看成“从 $S$ 到 $U$ 的态射”。

\begin{lemma}
\label{sites-lemma-sieves-set}
设 $\mathcal{C}$ 为范畴，$U \in \Ob(\mathcal{C})$。
\begin{enumerate}
\item $U$ 上所有筛构成集合。
\item 包含关系在该集合上定义偏序。
\item 筛的并与交仍是筛。
\item
\label{sites-item-sieve-generated}
给定 $\mathcal{C}$ 中以 $U$ 为目标的态射族
$\{U_i \to U\}_{i\in I}$，存在唯一最小的 $U$ 上的筛 $S$，使每个
$U_i \to U$ 都属于 $S(U_i)$。
\item 筛 $S = h_U$ 是最大筛。
\item 空子预层是最小筛。
\end{enumerate}
\end{lemma}

\begin{proof}
由子预层的定义，预层 $\mathcal{F}$ 的所有子预层所成的集族，是
$\prod_{U \in \Ob(\mathcal{C})} \mathcal{P}(\mathcal{F}(U))$ 的子集。
后者是集合。（这里 $\mathcal{P}(A)$ 表示 $A$ 的幂集。）因此 $U$ 上所有筛构成集合。

\medskip\noindent
偏序定义如下：$S \leq S'$ 当且仅当对所有 $T \to U$，有
$S(T) \subset S'(T)$。记作 $S \subset S'$。

\medskip\noindent
给定 $U$ 上的一族筛 $S_i$（$i \in I$），可以把 $\bigcup S_i$ 定义为满足
$(\bigcup S_i)(T) = \bigcup S_i(T)$（对所有 $T \in \Ob(\mathcal{C})$）的筛。
交 $\bigcap S_i$ 以同样方式定义。

\medskip\noindent
给定命题中的 $\{U_i \to U\}_{i\in I}$，考虑预层态射
$h_{U_i} \to h_U$。直接把 $S$ 定义为这些预层映射之像（定义
\ref{sites-definition-image}）的并。

\medskip\noindent
引理的最后两个断言显然。
\end{proof}

\begin{definition}
\label{sites-definition-sieve-generated}
设 $\mathcal{C}$ 为范畴。给定 $\mathcal{C}$ 中以 $U$ 为目标的态射族
$\{f_i : U_i \to U\}_{i\in I}$，称引理 \ref{sites-lemma-sieves-set} 第
(\ref{sites-item-sieve-generated}) 部分所述的 $U$ 上的筛 $S$ 为
{\it 由态射 $f_i$ 生成的 $U$ 上的筛}。
\end{definition}

\begin{definition}
\label{sites-definition-pullback-sieve}
设 $\mathcal{C}$ 为范畴，$f : V \to U$ 是 $\mathcal{C}$ 的态射，而
$S \subset h_U$ 是筛。定义 $S$ 沿 $f$ 的{\it 拉回}为 $V$ 上的筛
$S \times_U V$，规则为
$$
(\alpha : T \to V) \in (S \times_U V)(T)
\Leftrightarrow
(f \circ \alpha : T \to U) \in S(T)
$$
\end{definition}

\noindent
留给读者验证这确实是筛（提示：使用等式 \ref{sites-equation-property-sieve}）。
有时也称 $S \times_U V$ 为 $S$ 沿 $f : V \to U$ 的{\it 基变换}。

\begin{lemma}
\label{sites-lemma-pullback-sieve-section}
设 $\mathcal{C}$ 为范畴，$U \in \Ob(\mathcal{C})$，而 $S$ 是 $U$ 上的筛。
若 $f : V \to U$ 属于 $S$，则 $S \times_U V = h_V$ 是最大筛。
\end{lemma}

\begin{proof}
由定义显然。
\end{proof}

\begin{definition}
\label{sites-definition-topology}
设 $\mathcal{C}$ 为范畴。$\mathcal{C}$ 上的一个{\it 拓扑}，是把 $U$ 上所有筛
所成集合的一个子集 $J(U)$ 赋给每个 $U \in \Ob(\mathcal{C})$ 的规则，并满足：
\begin{enumerate}
\item 对 $\mathcal{C}$ 中每个态射 $f : V \to U$ 以及每个
$S \in J(U)$，拉回 $S \times_U V$ 属于 $J(V)$。
\item 若 $S$、$S'$ 是 $U \in \Ob(\mathcal{C})$ 上的筛，$S \in J(U)$，
并且对所有 $f \in S(V)$，拉回 $S' \times_U V$ 都属于 $J(V)$，则
$S'$ 属于 $J(U)$。
\item 对每个 $U \in \Ob(\mathcal{C})$，最大筛 $S = h_U$ 属于 $J(U)$。
\end{enumerate}
\end{definition}

\noindent
在这种情形下，属于 $J(U)$ 的筛称为{\it 覆盖筛}。

\begin{lemma}
\label{sites-lemma-topology-basic}
设 $\mathcal{C}$ 为范畴，$J$ 是 $\mathcal{C}$ 上的拓扑，而
$U \in \Ob(\mathcal{C})$。
\begin{enumerate}
\item $J(U)$ 中元素的有限交仍属于 $J(U)$。
\item 若 $S \in J(U)$ 且 $S' \supset S$，则 $S' \in J(U)$。
\end{enumerate}
\end{lemma}

\begin{proof}
(1) 的证明。空族之交的情形由定义 \ref{sites-definition-topology} 的条件 (3) 得出。
设 $S, S' \in J(U)$，并考虑 $S'' = S \cap S'$。对 $S(V)$ 中的每个
$V \to U$，有
$$
S' \times_U V = S'' \times_U V
$$
这是因为 $V \to U$ 已经属于 $S$。所以由定义的第二条公理，
$S'' \in J(U)$。

\medskip\noindent
(2) 的证明。设 $S \in J(U)$ 且 $S' \supset S$。对 $S(V)$ 中的每个
$V \to U$，由引理 \ref{sites-lemma-pullback-sieve-section}，有
$S' \times_U V = h_V$。所以由第三条公理，$S' \times_U V \in J(V)$；
再由第二条公理，$S' \in J(U)$。
\end{proof}

\begin{definition}
\label{sites-definition-finer}
设 $\mathcal{C}$ 为范畴，$J$、$J'$ 是 $\mathcal{C}$ 上的两个拓扑。若对
$\mathcal{C}$ 的每个对象 $U$ 都有 $J'(U) \subset J(U)$，则称 $J$ 比
$J'$ {\it 更细}或{\it 更强}。在这种情形下，也称 $J'$ 比 $J$
{\it 更粗}或{\it 更弱}。
\end{definition}

\noindent
换言之，$J'$ 的任意覆盖筛都是 $J$ 的覆盖筛。$\mathcal{C}$ 上存在最细拓扑，
即每个筛都是覆盖筛的拓扑，称为 $\mathcal{C}$ 的{\it 离散拓扑}。也存在最粗拓扑，
即对所有对象 $U$ 都有 $J(U) = \{h_U\}$ 的拓扑，称为{\it 混沌拓扑}或
{\it 平凡拓扑}。

\begin{lemma}
\label{sites-lemma-play-with-topologies}
设 $\mathcal{C}$ 为范畴，而 $\{J_i\}_{i\in I}$ 是一族拓扑。
\begin{enumerate}
\item 规则 $J(U) = \bigcap J_i(U)$ 定义 $\mathcal{C}$ 上的拓扑。
\item 存在比所有拓扑 $J_i$ 都更细的最粗拓扑。
\end{enumerate}
\end{lemma}

\begin{proof}
第一部分由定义直接得出。第二部分由取所有比各 $J_i$ 都更细的拓扑之交得出。
\end{proof}

\noindent
至此，可以不加任何动机地定义层。

\begin{definition}
\label{sites-definition-sheaf-sets-topology}
设范畴 $\mathcal{C}$ 配备拓扑 $J$，而 $\mathcal{F}$ 是 $\mathcal{C}$ 上的集合预层。
若对每个 $U \in \Ob(\mathcal{C})$ 以及 $U$ 的每个覆盖筛 $S$，典范映射
$$
\Mor_{\textit{PSh}(\mathcal{C})}(h_U, \mathcal{F})
\longrightarrow
\Mor_{\textit{PSh}(\mathcal{C})}(S, \mathcal{F})
$$
都是双射，则称 $\mathcal{F}$ 是 $\mathcal{C}$ 上的{\it 层}。
\end{definition}

\noindent
回忆所示公式的左端等于 $\mathcal{F}(U)$。换言之，$\mathcal{F}$ 是层，当且仅当
$\mathcal{F}$ 在 $U$ 上的一个截面，等同于由
$(\alpha : T \to U) \in S(T)$ 参数化的一族相容截面
$s_{T, \alpha} \in \mathcal{F}(T)$；并且这对 $U$ 上每个覆盖筛 $S$ 都成立。

\begin{lemma}
\label{sites-lemma-topology-presheaves-sheaves}
设 $\mathcal{C}$ 为范畴，$\{ \mathcal{F}_i \}_{i\in I}$ 是
$\mathcal{C}$ 上的一族集合预层。对每个 $U \in \Ob(\mathcal{C})$，以
$J(U)$ 表示 $U$ 上满足下列性质的筛 $S$ 所成的集合：对每个态射
$V \to U$，映射
% P01-ERRATA: P01-E0163 -- the source's "denote J(U) the set" omits "by" or needs recasting.
$$
\Mor_{\textit{PSh}(\mathcal{C})}(h_V, \mathcal{F}_i)
\longrightarrow
\Mor_{\textit{PSh}(\mathcal{C})}(S \times_U V, \mathcal{F}_i)
$$
对所有 $i \in I$ 都是双射。则 $J$ 定义 $\mathcal{C}$ 上的拓扑。它是使所有
$\mathcal{F}_i$ 都成为层的最细拓扑。
\end{lemma}

\begin{proof}
首先注意，若 $J$ 是拓扑，则引理的最后一个断言立即成立。

\medskip\noindent
对每个 $i \in I$ 与 $U \in \Ob(\mathcal{C})$，令 $J_i(U)$ 为 $U$ 上满足
下列性质的筛 $S$ 所成的集合：对每个态射 $V \to U$，映射
$$
\Mor_{\textit{PSh}(\mathcal{C})}(h_V, \mathcal{F}_i)
\longrightarrow
\Mor_{\textit{PSh}(\mathcal{C})}(S \times_U V, \mathcal{F}_i)
$$
是双射。注意 $J(U) = \bigcap J_i(U)$。因此，若证明 $J_i$ 定义拓扑，就由引理
\ref{sites-lemma-play-with-topologies} 得到结论。于是归约到下一段讨论的情形。

\medskip\noindent
证明当预层族仅含一个预层 $\mathcal{F}$ 时，$J$ 是拓扑。拓扑的第一条和第三条公理
立即得到验证。设有对象 $U$ 以及 $U$ 上的筛 $S, S'$，使
$S \in J(U)$，并且对 $S(V)$ 中的所有 $V \to U$，都有
$S' \times_U V \in J(V)$。需要证明 $S' \in J(U)$。换言之，需要证明对任意
$f : W \to U$，映射
$$
\mathcal{F}(W) =
\Mor_{\textit{PSh}(\mathcal{C})}(h_W, \mathcal{F})
\longrightarrow
\Mor_{\textit{PSh}(\mathcal{C})}(S' \times_U W, \mathcal{F})
$$
都是双射。取元素
$\varphi \in \Mor_{\textit{PSh}(\mathcal{C})}(S' \times_U W, \mathcal{F})$。
我们将构造映到 $\varphi$ 的截面 $s \in \mathcal{F}(W)$。

\medskip\noindent
设 $\alpha : V \to W$ 是 $S \times_U W$ 的一个元素。由拉回的定义，复合
$f \circ \alpha : V \to W  \to U$ 属于 $S$。因此，由对筛对 $S, S'$ 的假设，
$S' \times_U V$ 属于 $J(V)$。现在有预层的交换图
$$
\xymatrix{
S' \times_U V \ar[r] \ar[d] & h_V \ar[d] \\
S' \times_U W \ar[r] & h_W
}
$$
$\varphi$ 在 $S' \times_U V$ 上的限制对应于元素
$s_{V, \alpha} \in \mathcal{F}(V)$。这是由 $J$ 的定义以及
$S' \times_U V$ 属于 $J(V)$ 得出的。留给读者验证：规则
$(V, \alpha) \mapsto s_{V, \alpha}$ 定义元素
$\psi \in \Mor_{\textit{PSh}(\mathcal{C})}(S \times_U W, \mathcal{F})$。
由于 $S \in J(U)$，从 $J$ 的定义立即看出，$\psi$ 对应于
$\mathcal{F}(W)$ 的一个元素 $s$。

\medskip\noindent
留给读者验证：构造 $\varphi \mapsto s$ 是上面所示自然映射的逆。
\end{proof}

\begin{definition}
\label{sites-definition-canonical-topology}
设 $\mathcal{C}$ 为范畴。使所有可表预层均为层的 $\mathcal{C}$ 上最细拓扑（见引理
\ref{sites-lemma-topology-presheaves-sheaves}），称为 $\mathcal{C}$ 的
{\it 典范拓扑}。
\end{definition}

\section{位点定义的拓扑}
\label{sites-section-topology-site}

\noindent
假设 $\mathcal{C}$ 是范畴，而 $\text{Cov}_1(\mathcal{C})$ 与
$\text{Cov}_2(\mathcal{C})$ 是在 $\mathcal{C}$ 上定义位点结构的两组覆盖。
在这种情形下，$\text{Cov}_1(\mathcal{C})$ 与
$\text{Cov}_2(\mathcal{C})$ 的（集合）层范畴可能相同；例如见引理
\ref{sites-lemma-refine-same-topology}。

\medskip\noindent
事实上，$\mathcal{C}$ 关于某个拓扑 $J$ 的层范畴与 $J$ 彼此决定。
这是一个非平凡的断言，稍后将在定理 \ref{sites-theorem-topology-and-topos}
中讨论。

\medskip\noindent
暂且接受这一点，研究位点所决定的拓扑就是合理的。

\begin{lemma}
\label{sites-lemma-site-gives-topology}
设 $\mathcal{C}$ 是覆盖为 $\text{Cov}(\mathcal{C})$ 的位点。对
$\mathcal{C}$ 的每个对象 $U$，令 $J(U)$ 为 $U$ 上具有下述性质的筛 $S$
所成的集合：存在覆盖
$\{f_i : U_i \to U\}_{i\in I} \in \text{Cov}(\mathcal{C})$，使由 $f_i$
生成的筛 $S'$（见定义 \ref{sites-definition-sieve-generated}）包含于 $S$。
\begin{enumerate}
\item 该 $J$ 是 $\mathcal{C}$ 上的拓扑。
\item 预层 $\mathcal{F}$ 是该拓扑的层（见定义
\ref{sites-definition-sheaf-sets-topology}），当且仅当它是该位点上的层（见定义
\ref{sites-definition-sheaf-sets}）。
\end{enumerate}
\end{lemma}

\begin{proof}
为证明第一个断言，只需注意位点定义（定义 \ref{sites-definition-site}）的公理
(1)、(2)、(3) 分别直接蕴含拓扑定义（定义 \ref{sites-definition-topology}）的公理
(3)、(2)、(1)。作为例子，证明 $J$ 具有性质 (2)。设 $U$ 是
$\mathcal{C}$ 的对象，而 $S, S'$ 是 $U$ 上的筛，满足 $S \in J(U)$，并且对
$S(V)$ 中的每个 $V \to U$ 都有 $S' \times_U V \in J(V)$。由 $J(U)$
的定义，可找到该位点的覆盖 $\{f_i : U_i \to U\}$，使 $S$，即
$h_{U_i} \to h_U$ 的像所成的筛，包含于 $S$。
% P01-ERRATA: P01-E0164 -- the source reads "such that S the image ... is contained in S", duplicating S and omitting the grammatical construction.
由于每个 $S'\times_U U_i$ 属于 $J(U_i)$，可见存在该位点的覆盖
$\{U_{ij} \to U_i\}$，使 $h_{U_{ij}} \to h_{U_i}$ 包含于
$S' \times_U U_i$。由基变换的定义，这意味着
$h_{U_{ij}} \to h_U$ 包含于子预层 $S' \subset h_U$。由位点的公理 (2)，
$\{U_{ij} \to U\}$ 是 $U$ 的覆盖；再由 $J$ 的定义，得出
$S' \in J(U)$。

\medskip\noindent
设 $\mathcal{F}$ 为预层。假设 $\mathcal{F}$ 是拓扑 $J$ 的层。我们将证明
$\mathcal{F}$ 也是该位点上的层。设
$\{f_i : U_i \to U\}_{i\in I}$ 是该位点的覆盖，而
$s_i \in \mathcal{F}(U_i)$ 是满足
$s_i|_{U_i \times_U U_j} = s_j|_{U_i \times_U U_j}$（对所有 $i, j$）的一族截面。
需要证明存在唯一的截面 $s \in \mathcal{F}(U)$，其在各 $U_i$ 上的限制为
$s_i$。设 $S \subset h_U$ 为由 $f_i$ 生成的筛。按定义，$S \in J(U)$。
与其直接构造 $s$，由拓扑中的层条件，只须构造元素
% P01-ERRATA: P01-E0165 -- the source spells "Instead" as two words, "In stead".
$$
\varphi \in \Mor_{\textit{PSh}(\mathcal{C})}(S, \mathcal{F}).
$$
对某个对象 $T \in \mathcal{U}$，取 $\alpha \in S(T)$。
% P01-ERRATA: P01-E0166 -- the category U is undefined here; T should be an object of C.
这恰好意味着 $\alpha : T \to U$ 是一个通过某个 $i\in I$ 的 $f_i$ 分解的态射
（可能通过多于 $1$ 个）。选取这样的指标 $i$ 与分解
$\alpha = f_i \circ \alpha_i$，并定义 $\varphi(\alpha) = \alpha_i^* s_i$。
若 $i'$、$\alpha = f_i \circ \alpha_{i'}'$ 是第二种选择，则
% P01-ERRATA: P01-E0167 -- the second factorization uses f_i again; it should use f_{i'}.
$\alpha_i^* s_i = (\alpha_{i'}')^* s_{i'}$，这恰由对所有 $i, j$ 的条件
$s_i|_{U_i \times_U U_j} = s_j|_{U_i \times_U U_j}$ 得出。因此
$\varphi(\alpha)$ 良定义。留给读者验证：$\varphi$（它进而决定 $s$）是正确的，
即 $s$ 的相应限制为 $s_i$。

\medskip\noindent
设 $\mathcal{F}$ 为预层。假设 $\mathcal{F}$ 是位点
$(\mathcal{C}, \text{Cov}(\mathcal{C}))$ 上的层。我们将证明
$\mathcal{F}$ 也是拓扑 $J$ 的层。设 $U$ 是 $\mathcal{C}$ 的对象，而 $S$ 是
$U$ 上关于拓扑 $J$ 的覆盖筛。设
$$
\varphi \in \Mor_{\textit{PSh}(\mathcal{C})}(S, \mathcal{F}).
$$
需要证明 $\mathcal{F}(U) =
\Mor_{\textit{PSh}(\mathcal{C})}(h_U, \mathcal{F})$ 中存在唯一元素，限制为
$\varphi$。由定义，存在覆盖
$\{f_i : U_i \to U\}_{i\in I} \in \text{Cov}(\mathcal{C})$，使
$f_i : U_i \to U$ 属于 $S(U_i)$。因此可以令
$s_i = \varphi(f_i) \in \mathcal{F}(U_i)$。容易验证，对所有 $i, j$，有
$s_i|_{U_i \times_U U_j} = s_j|_{U_i \times_U U_j}$。于是由
$\mathcal{F}$ 在位点 $(\mathcal{C}, \text{Cov}(\mathcal{C}))$ 上的层条件，
得到所需截面 $s$。细节留给读者。
\end{proof}

\begin{definition}
\label{sites-definition-topology-associated-site}
设 $\mathcal{C}$ 是覆盖为 $\text{Cov}(\mathcal{C})$ 的位点。上面引理
\ref{sites-lemma-site-gives-topology} 中构造的拓扑 $J$，称为
{\it 与 $\mathcal{C}$ 关联的拓扑}。
\end{definition}

\noindent
设 $\mathcal{C}$ 为范畴。设 $\text{Cov}_1(\mathcal{C})$ 与
$\text{Cov}_2(\mathcal{C})$ 是在 $\mathcal{C}$ 上定义位点结构的两组覆盖。
它们定义的拓扑完全可能相同。若如此，则称
$\text{Cov}_1(\mathcal{C})$ 与 $\text{Cov}_2(\mathcal{C})$
在 $\mathcal{C}$ 上{\it 定义相同的拓扑}。并且在这种情形下，由引理
\ref{sites-lemma-site-gives-topology}，层范畴相同。
% P01-ERRATA: P01-E0168 -- the source calls Cov_1 and Cov_2 "two coverings" although each denotes a set/class of coverings defining a site.

\medskip\noindent
通常，我们只关心一组覆盖所定义的拓扑，并把选择不同覆盖集的可能性看作研究该拓扑
的工具。

\begin{remark}
\label{sites-remark-enlarge-coverings}
扩大覆盖类。显然，若 $\text{Cov}(\mathcal{C})$ 在 $\mathcal{C}$ 上定义位点结构，
则可以向 $\mathcal{C}$ 加入任意一组具有固定目标的态射族，只要它们与
$\text{Cov}(\mathcal{C})$ 的元素恒真等价（见定义
\ref{sites-definition-combinatorial-tautological}），而不改变拓扑。
% P01-ERRATA: P01-E0169 -- the source says to add covering families "to C"; they should be added to Cov(C), not to the category itself.
\end{remark}

\begin{remark}
\label{sites-remark-shrink-coverings}
缩小覆盖类。设 $\mathcal{C}$ 为位点。考虑集合
$$
\mathcal{S} = P(\text{Arrows}(\mathcal{C})) \times \Ob(\mathcal{C})
$$
其中 $P(\text{Arrows}(\mathcal{C}))$ 是态射集的幂集，即由所有态射集组成的集合。
令 $\mathcal{S}_\tau \subset \mathcal{S}$ 为所有满足下列条件的
$(T, U) \in \mathcal{S}$ 所成的子集：
(a) 所有 $\varphi \in T$ 的目标都是 $U$；
(b) 族 $\{\varphi\}_{\varphi \in T}$ 与 $\text{Cov}(\mathcal{C})$
中的某个覆盖恒真等价（见定义 \ref{sites-definition-combinatorial-tautological}）。
显然，若把 $\mathcal{S}_\tau$ 的元素看作覆盖，并不会恰好得到定义
\ref{sites-definition-site} 中所定义的位点概念。所得结构
$(\mathcal{C}, \mathcal{S}_\tau)$ 满足略经修改的条件。修改后的条件为：
\begin{enumerate}
\item[(0')] $\text{Cov}(\mathcal{C}) \subset
P(\text{Arrows}(\mathcal{C})) \times \Ob(\mathcal{C})$,
\item[(1')] 若 $V \to U$ 是同构，则
$(\{V \to U\}, U) \in \text{Cov}(\mathcal{C})$。
\item[(2')] 若 $(T, U) \in \text{Cov}(\mathcal{C})$，并且对
$T$ 中的 $f : U' \to U$ 给定
$(T_f, U') \in \text{Cov}(\mathcal{C})$，则令
$T' = \{f \circ f' \mid f \in T,\ f' \in T_f\}$，有
$(T', U) \in \text{Cov}(\mathcal{C})$。
\item[(3')] 若 $(T, U) \in \text{Cov}(\mathcal{C})$，且
$g : V \to U$ 是 $\mathcal{C}$ 的态射，则
\begin{enumerate}
\item 对 $T$ 中的 $f : U' \to U$，$U' \times_{f, U, g} V$ 存在；且
\item 对纤维积的某种选择令
$T' = \{\text{pr}_2 : U' \times_{f, U, g} V \to V \mid f : U' \to U \in T\}$，
则 $(T', V) \in \text{Cov}(\mathcal{C})$。
\end{enumerate}
\end{enumerate}
容易验证：给定满足上述 (0') -- (3') 的结构，适当扩大
$\text{Cov}(\mathcal{C})$ 后（对照 Sets, Section
\ref{sets-section-coverings-site}），得到一个位点。至少从拓扑角度看，该概念与实际的
位点概念显然差别不大。它有两个好处：由上面的条件 (0')，覆盖自动构成集合；同样由
(0')，所有这类结构的总体也构成集合。代价是必须到处写“恒真等价”。
\end{remark}

\section{拓扑中的层化}
\label{sites-section-sheafification-topology}

\noindent
本节说明拓扑中与层化构造对应的构造。

\medskip\noindent
设 $\mathcal{C}$ 为范畴，$J$ 是 $\mathcal{C}$ 上的拓扑，而
$\mathcal{F}$ 是集合预层。对每个 $U \in \Ob(\mathcal{C})$，把
$$
L\mathcal{F}(U)
=
\colim_{S \in J(U)^{opp}}
\Mor_{\textit{PSh}(\mathcal{C})}(S, \mathcal{F})
$$
定义为一个余极限。这里把 $J(U)$ 看成按包含关系排序的偏序集；见引理
\ref{sites-lemma-sieves-set}。该系统的转移映射定义如下。若
$S \subset S'$ 属于 $J(U)$，则 $S \to S'$ 是预层态射，因而存在自然的限制映射
$$
\Mor_{\textit{PSh}(\mathcal{C})}(S', \mathcal{F})
\longrightarrow
\Mor_{\textit{PSh}(\mathcal{C})}(S, \mathcal{F}).
$$
所以，只要反转 $J(U)$ 上的次序（这正是上标 ${}^{opp}$ 所表示的意思），
$S \mapsto \Mor_{\textit{PSh}(\mathcal{C})}(S, \mathcal{F})$
就是 Categories, Definition \ref{categories-definition-system-over-poset}
意义下的有向系统。特别地，因为 $h_U \in J(U)$，存在来自识别
$\mathcal{F}(U) = \Mor_{\textit{PSh}(\mathcal{C})}(h_U, \mathcal{F})$
的典范映射
$$
\ell : \mathcal{F}(U) \longrightarrow L\mathcal{F}(U)
$$
此外，定义 $L\mathcal{F}(U)$ 的余极限是有向的，因为对 $U$ 上任意一对覆盖筛
$S, S'$，筛 $S \cap S'$ 仍是覆盖筛；见引理 \ref{sites-lemma-sieves-set}。
% P01-ERRATA: P01-E0170 -- closure of covering sieves under finite intersections is Lemma lemma-topology-basic, not Lemma lemma-sieves-set.

\medskip\noindent
设 $f : V \to U$ 是 $\mathcal{C}$ 中的态射，并设 $S \in J(U)$。有交换图
$$
\xymatrix{
S \times_U V \ar[r] \ar[d] & h_V \ar[d] \\
S \ar[r] & h_U
}
$$
可以用左侧竖直映射得到典范限制映射
$$
\Mor_{\textit{PSh}(\mathcal{C})}(S, \mathcal{F})
\to \Mor_{\textit{PSh}(\mathcal{C})}(S \times_U V, \mathcal{F}).
$$
基变换 $S \mapsto S \times_U V$ 诱导保序映射 $J(U) \to J(V)$，而限制映射
定义 Categories, Lemma {categories-lemma-functorial-colimit} 所述的函子变换。
% P01-ERRATA: P01-E0171 -- the source omits the \ref command around categories-lemma-functorial-colimit.
所以得到自然的限制映射
$$
L\mathcal{F}(U) \longrightarrow L\mathcal{F}(V).
$$

\begin{lemma}
\label{sites-lemma-L-presheaf}
在上述情形中：
% P01-ERRATA: P01-E0172 -- the source's "In the situation above." is a sentence fragment before the list.
\begin{enumerate}
\item 赋值 $U \mapsto L\mathcal{F}(U)$ 连同上面定义的限制映射构成预层。
\item 各映射 $\ell$ 粘合为预层态射
$\ell : \mathcal{F} \to L\mathcal{F}$。
\item 规则 $\mathcal{F} \mapsto (\mathcal{F} \xrightarrow{\ell}
L\mathcal{F})$ 是函子。
\item 若 $\mathcal{F}$ 是 $\mathcal{G}$ 的子预层，则
$L\mathcal{F}$ 是 $L\mathcal{G}$ 的子预层。
\item 映射 $\ell : \mathcal{F} \to L\mathcal{F}$ 具有下列性质：对每个截面
$s \in L\mathcal{F}(U)$，存在 $U$ 上的覆盖筛 $S$ 以及元素
$\varphi \in \Mor_{\textit{PSh}(\mathcal{C})}(S, \mathcal{F})$，使
$\ell(\varphi)$ 等于 $s$ 在 $S$ 上的限制。
\end{enumerate}
\end{lemma}

\begin{proof}
略。
\end{proof}

\begin{definition}
\label{sites-definition-presheaf-separated-topology}
设 $\mathcal{C}$ 为范畴，而 $J$ 是 $\mathcal{C}$ 上的拓扑。若对每个对象 $U$
以及 $U$ 上的每个覆盖筛 $S$，典范映射
$\mathcal{F}(U) \to \Mor_{\textit{PSh}(\mathcal{C})}(S, \mathcal{F})$
都是单射，则称集合预层 $\mathcal{F}$ 是{\it 分离的}。
\end{definition}

\begin{theorem}
\label{sites-theorem-L-topology}
设 $\mathcal{C}$ 为范畴，$J$ 是 $\mathcal{C}$ 上的拓扑，而
$\mathcal{F}$ 是集合预层。
\begin{enumerate}
\item 预层 $L\mathcal{F}$ 是分离的。
\item 若 $\mathcal{F}$ 分离，则 $L\mathcal{F}$ 是层，并且预层映射
$\mathcal{F} \to L\mathcal{F}$ 是单射。
\item 若 $\mathcal{F}$ 是层，则 $\mathcal{F} \to L\mathcal{F}$ 是同构。
\item 预层 $LL\mathcal{F}$ 总是层。
\end{enumerate}
\end{theorem}

\begin{proof}
(3) 由 $L$ 的定义与层的定义（定义 \ref{sites-definition-sheaf-sets-topology}）
立即得出。(4) 形式地由其余部分得出。

\medskip\noindent
概述 (1) 的证明。设 $S$ 是对象 $U$ 的覆盖筛。设
$\varphi_i \in L\mathcal{F}(U)$（$i = 1, 2$）在
$\Mor_{\textit{PSh}(\mathcal{C})}(S, L\mathcal{F})$ 中映到同一元素。
可以找到 $U$ 上的一个覆盖筛 $S'$，使两个 $\varphi_i$ 都由元素
$\varphi_i \in \Mor_{\textit{PSh}(\mathcal{C})}(S', \mathcal{F})$ 表示。
把 $S$ 与 $S'$ 都替换为仍是覆盖筛的 $S' \cap S$（见引理
\ref{sites-lemma-sieves-set}），可假设 $S' = S$。设
$V\in \Ob(\mathcal{C})$，而 $\alpha : V \to U$ 属于 $S(V)$。则
$S \times_U V = h_V$；见引理 \ref{sites-lemma-pullback-sieve-section}。所以
$\varphi_i$ 沿 $V \to U$ 的限制对应于 $\mathcal{F}$ 在 $V$ 上的截面
$s_{i, V, \alpha}$。假设意味着存在 $V$ 的覆盖筛 $S_{V, \alpha}$，使
$s_{i, V, \alpha}$ 在
$\Mor_{\textit{PSh}(\mathcal{C})}(S_{V, \alpha}, \mathcal{F})$
中限制为同一元素。考虑 $U$ 上由下列规则定义的筛 $S''$：
\begin{eqnarray}
\label{sites-equation-S-prime-prime}
(f : T \to U) \in S''(T)
& \Leftrightarrow &
\exists\ V , \ \alpha : V \to U, \ \alpha \in S(V), \nonumber \\
& &
\exists\ g : T \to V, \ g \in S_{V, \alpha}(T), \\
& &
f = \alpha \circ g \nonumber
\end{eqnarray}
由拓扑的公理 (2)，$S''$ 是 $U$ 上的覆盖筛。由构造，
$\varphi_1$ 与 $\varphi_2$ 在
$\Mor_{\textit{PSh}(\mathcal{C})}(S'', L\mathcal{F})$ 中限制为同一元素，
如所需。

\medskip\noindent
概述 (2) 的证明。假设 $\mathcal{F}$ 是 $\mathcal{C}$ 上关于拓扑 $J$
分离的集合预层。设 $S$ 是 $\mathcal{C}$ 的对象 $U$ 的覆盖筛。假设
$\varphi \in \Mor_\mathcal{C}(S, L\mathcal{F})$。需要找到限制为 $\varphi$
的元素 $s \in L\mathcal{F}(U)$。
% P01-ERRATA: P01-E0173 -- S and LF are presheaves, so the Hom must be in PSh(C), not in C.
设 $V\in \Ob(\mathcal{C})$，而 $\alpha : V \to U$ 属于 $S(V)$。值
$\varphi(\alpha) \in L\mathcal{F}(V)$ 由 $V$ 的一个覆盖筛
$S_{V, \alpha}$ 以及预层态射
$\varphi_{V, \alpha} : S_{V, \alpha} \to \mathcal{F}$ 给出。与上面的证明一样，
通过等式 (\ref{sites-equation-S-prime-prime}) 定义 $U$ 上的覆盖筛 $S''$。用下列简单
规则定义
$$
\varphi'' : S'' \longrightarrow \mathcal{F}
$$
：对每个 $f : T \to U$、$f \in S''(T)$，按照等式
(\ref{sites-equation-S-prime-prime}) 选取 $V, \alpha, g$，然后令
$$
\varphi''(f) = \varphi_{V, \alpha}(g).
$$
我们断言这与 $V, \alpha, g$ 的选择无关。考虑另一个这样的选择$ V', \alpha', g'$。
% P01-ERRATA: P01-E0174 -- the source omits a space before the math span in "choice$ V'".
$\varphi_{V, \alpha}$ 与 $\varphi_{V', \alpha'}$ 在 $T$ 上下列两个覆盖筛之交上的
限制相同：
$$
(S_{V, \alpha} \times_{V, g} T) \cap (S_{V', \alpha'} \times_{V', g'} T)
$$
事实上，这两个限制都对应于 $\varphi$ 通过 $f$ 在 $T$ 上的限制，而所需等式由
$\mathcal{F}$ 分离得出。把公共限制记为 $\psi$。
% P01-ERRATA: P01-E0175 -- the source's "Denote the common restriction psi" omits "by".
选择无关性由
$\varphi_{V, \alpha}(g) = \psi(\text{id}_T) =
\varphi_{V', \alpha'}(g')$ 得出。于是 $\varphi''$ 给出元素
$s \in L\mathcal{F}(U)$。留给读者检查 $s$ 限制为 $\varphi$。
\end{proof}

\begin{definition}
\label{sites-definition-associated-sheaf-topology}
设范畴 $\mathcal{C}$ 配备拓扑 $J$，而 $\mathcal{F}$ 是
$\mathcal{C}$ 上的集合预层。层 $\mathcal{F}^\# := LL\mathcal{F}$ 连同典范映射
$\mathcal{F} \to \mathcal{F}^\#$，称为{\it 与 $\mathcal{F}$ 关联的层}。
\end{definition}

\begin{proposition}
\label{sites-proposition-sheafification-adjoint-topology}
设范畴 $\mathcal{C}$ 配备一个拓扑，而 $\mathcal{F}$ 是
$\mathcal{C}$ 上的集合预层。典范映射 $\mathcal{F} \to \mathcal{F}^\#$ 具有下列
泛性质：对任意映射 $\mathcal{F} \to \mathcal{G}$，其中
$\mathcal{G}$ 是集合层，存在唯一映射 $\mathcal{F}^\# \to \mathcal{G}$，使
$\mathcal{F} \to \mathcal{F}^\# \to \mathcal{G}$ 等于给定映射。
\end{proposition}

\begin{proof}
与命题 \ref{sites-proposition-sheafification-adjoint} 的证明相同。
\end{proof}

\section{拓扑与层}
\label{sites-section-topology-and-sheaves}

\begin{lemma}
\label{sites-lemma-sieve-sheafification}
设范畴 $\mathcal{C}$ 配备拓扑 $J$，$U$ 是 $\mathcal{C}$ 的对象，而
$S$ 是 $U$ 上的筛。下列条件等价：
\begin{enumerate}
\item 筛 $S$ 是覆盖筛。
\item 映射 $S \to h_U$ 的层化 $S^\# \to h_U^\#$ 是同构。
\end{enumerate}
\end{lemma}

\begin{proof}
先作两点一般说明。将使用 $S^\# = LLS$ 与 $h_U^\# = LLh_U$。特别地，
由引理 \ref{sites-lemma-L-presheaf}，$S^\# \to h_U^\#$ 是单射。注意
$\text{id}_U \in h_U(U)$。因此，它产生 $Lh_U$ 以及
$h_U^\# = LLh_U$ 在 $U$ 上的截面，我们也将这些截面记为 $\text{id}_U$。

\medskip\noindent
假设 $S$ 是覆盖筛。显然，只须找到态射 $h_U \to S^\#$，使其与前述单射的复合
$h_U \to h_U^\#$ 是典范映射。为找到这种映射，
只须找到限制为 $\text{id}_U$ 的截面 $s \in S^\#(U)$。但由于 $S$ 是覆盖筛，
元素 $\text{id}_S \in \Mor_{\textit{PSh}(\mathcal{C})}(S, S)$
产生 $LS$ 在 $U$ 上的一个截面，它在 $Lh_U$ 中限制为 $\text{id}_U$。
所以结论成立。

\medskip\noindent
假设 $S^\# \to h_U^\#$ 是同构。设 $1 \in S^\#(U)$ 是与
$h_U^\#(U)$ 中的 $\text{id}_U$ 对应的元素。由于 $S^\# = LLS$，存在
$U$ 上的覆盖筛 $S'$，使 $1$ 来自
$$
\varphi \in \Mor_{\textit{PSh}(\mathcal{C})}(S', LS).
$$
这又意味着，对每个 $\alpha : V \to U$、$\alpha\in S'(V)$，存在
$V$ 上的覆盖筛 $S_{V, \alpha}$，使 $\varphi(\alpha)$ 对应于预层态射
$S_{V, \alpha} \to S$。换言之，$S_{V, \alpha}$ 包含于
$S \times_U V$。由拓扑的第二条公理，$S$ 是覆盖筛。
\end{proof}

\begin{theorem}
\label{sites-theorem-topology-and-topos}
设 $\mathcal{C}$ 为范畴，而 $J$、$J'$ 是 $\mathcal{C}$ 上的拓扑。下列条件等价：
\begin{enumerate}
\item $J = J'$；
\item 拓扑 $J$ 的层与拓扑 $J'$ 的层相同。
\end{enumerate}
\end{theorem}

\begin{proof}
若 $J = J'$，则层的概念相同，这是恒真式。反过来，引理
\ref{sites-lemma-sieve-sheafification} 用层化函子刻画覆盖筛。而层化函子
$\textit{PSh}(\mathcal{C}) \to \Sh(\mathcal{C}, J)$ 左伴随于包含函子
$\Sh(\mathcal{C}, J) \to \textit{PSh}(\mathcal{C})$。因此，若子范畴
$\Sh(\mathcal{C}, J)$ 与 $\Sh(\mathcal{C}, J')$ 相同，则层化函子相同，
因而覆盖筛族相同。
\end{proof}

\begin{lemma}
\label{sites-lemma-finer-topology}
假设与记号如定理 \ref{sites-theorem-topology-and-topos}。则
$J \subset J'$ 当且仅当拓扑 $J'$ 的每个层都是拓扑 $J$ 的层。
% P01-ERRATA: P01-E0176 -- the source's fragment "Assumption and notation as in" needs a finite verb and plural "Assumptions".
\end{lemma}

\begin{proof}
一个方向是清楚的。对另一个方向，假设
$\Sh(\mathcal{C}, J') \subset \Sh(\mathcal{C}, J)$。由形式论证，这蕴含：
若 $\mathcal{F}$ 是集合预层，而
$\mathcal{F} \to \mathcal{F}^\#$、相应地
$\mathcal{F} \to \mathcal{F}^{\#, \prime}$ 是关于 $J$、相应地 $J'$ 的层化，
则存在典范映射 $\mathcal{F}^\# \to \mathcal{F}^{\#, \prime}$，使
$\mathcal{F} \to \mathcal{F}^\# \to \mathcal{F}^{\#, \prime}$
等于典范映射 $\mathcal{F} \to \mathcal{F}^{\#, \prime}$。当然，
$\mathcal{F}^\# \to \mathcal{F}^{\#, \prime}$ 把第二个层识别为第一个层关于拓扑
$J'$ 的层化。把它应用于引理 \ref{sites-lemma-sieve-sheafification} 的映射
$S \to h_U$，得到交换图
$$
\xymatrix{
S \ar[r] \ar[d] &
S^\# \ar[r] \ar[d] &
S^{\#, \prime} \ar[d] \\
h_U \ar[r] &
h_U^\# \ar[r] &
h_U^{\#, \prime}
}
$$
显然，若 $S$ 是拓扑 $J$ 的覆盖筛，则中间的竖直映射是同构（由该引理），
从而右侧竖直映射也是同构，因为它是中间竖直映射关于 $J'$ 的层化。再次由该引理，
$S$ 也是 $J'$ 的覆盖筛。
\end{proof}

\section{拓扑与连续函子}
\label{sites-section-topologies-continuous-functors}

\noindent
说明连续函子如何在层上给出一对伴随函子。





\section{点与拓扑}
\label{sites-section-points-topologies}

\noindent
回忆节 \ref{sites-section-points}：给定函子
$p = u : \mathcal{C} \to \textit{Sets}$，可以定义茎函子
$$
\textit{PSh}(\mathcal{C}) \longrightarrow \textit{Sets},
\mathcal{F} \longmapsto \mathcal{F}_p.
$$

\begin{definition}
\label{sites-definition-point-topology}
设 $\mathcal{C}$ 为范畴，而 $J$ 是 $\mathcal{C}$ 上的拓扑。该拓扑的一个
{\it 点 $p$} 由满足下列条件的函子
$u : \mathcal{C} \to \textit{Sets}$ 给出：
\begin{enumerate}
\item 对 $U$ 上的每个覆盖筛 $S$，映射
$S_p \to (h_U)_p$ 是满射。
\item 茎函子 $\Sh(\mathcal{C}) \to \textit{Sets}$，
$\mathcal{F} \to \mathcal{F}_p$，是正合的。
\end{enumerate}
% P01-ERRATA: P01-E0177 -- the source uses \to rather than \mapsto for the action F to F_p, unlike the same stalk functor displayed immediately above.
\end{definition}
